\documentclass[11pt]{article}

\usepackage[margin=1.02in]{geometry}
\usepackage{amsmath,amssymb,amsthm,mathtools,mathrsfs}
\usepackage{bm}
\usepackage{enumitem}
\usepackage{booktabs,array,tabularx,longtable}
\usepackage{xcolor}
\usepackage{microtype}
\usepackage[colorlinks=true,linkcolor=blue,citecolor=blue,urlcolor=blue]{hyperref}
\usepackage{aliascnt}
\usepackage[nameinlink,capitalise]{cleveref}
\usepackage{tikz-cd}
\setlength{\emergencystretch}{4em}
\raggedbottom
\hbadness=10000

% -----------------------------------------------------------------------------
% theorem environments
% -----------------------------------------------------------------------------
\newtheorem{theorem}{Theorem}[section]

\newaliascnt{proposition}{theorem}
\newtheorem{proposition}[proposition]{Proposition}
\aliascntresetthe{proposition}

\newaliascnt{lemma}{theorem}
\newtheorem{lemma}[lemma]{Lemma}
\aliascntresetthe{lemma}

\newaliascnt{corollary}{theorem}
\newtheorem{corollary}[corollary]{Corollary}
\aliascntresetthe{corollary}

\theoremstyle{definition}
\newaliascnt{definition}{theorem}
\newtheorem{definition}[definition]{Definition}
\aliascntresetthe{definition}

\newaliascnt{construction}{theorem}
\newtheorem{construction}[construction]{Construction}
\aliascntresetthe{construction}

\newaliascnt{assumption}{theorem}
\newtheorem{assumption}[assumption]{Assumption}
\aliascntresetthe{assumption}

\theoremstyle{remark}
\newaliascnt{remark}{theorem}
\newtheorem{remark}[remark]{Remark}
\aliascntresetthe{remark}

\crefname{theorem}{Theorem}{Theorems}
\Crefname{theorem}{Theorem}{Theorems}
\crefname{proposition}{Proposition}{Propositions}
\Crefname{proposition}{Proposition}{Propositions}
\crefname{lemma}{Lemma}{Lemmas}
\Crefname{lemma}{Lemma}{Lemmas}
\crefname{corollary}{Corollary}{Corollaries}
\Crefname{corollary}{Corollary}{Corollaries}
\crefname{definition}{Definition}{Definitions}
\Crefname{definition}{Definition}{Definitions}
\crefname{construction}{Construction}{Constructions}
\Crefname{construction}{Construction}{Constructions}
\crefname{assumption}{Assumption}{Assumptions}
\Crefname{assumption}{Assumption}{Assumptions}
\crefname{remark}{Remark}{Remarks}
\Crefname{remark}{Remark}{Remarks}

% -----------------------------------------------------------------------------
% notation
% -----------------------------------------------------------------------------
\newcommand{\N}{\mathbb N}
\newcommand{\Z}{\mathbb Z}
\newcommand{\R}{\mathbb R}
\newcommand{\C}{\mathbb C}
\newcommand{\one}{\mathbf 1}
\newcommand{\id}{\mathrm{id}}
\newcommand{\dd}{\mathrm d}
\newcommand{\Tr}{\operatorname{Tr}}
\newcommand{\tr}{\operatorname{tr}}
\newcommand{\rank}{\operatorname{rank}}
\newcommand{\Ker}{\operatorname{Ker}}
\newcommand{\Ran}{\operatorname{Ran}}
\newcommand{\Span}{\operatorname{span}}
\newcommand{\Hom}{\operatorname{Hom}}
\newcommand{\supp}{\operatorname{supp}}
\newcommand{\spec}{\operatorname{spec}}
\newcommand{\diag}{\operatorname{diag}}
\newcommand{\norm}[1]{\lVert#1\rVert}
\newcommand{\abs}[1]{\lvert#1\rvert}
\newcommand{\ip}[2]{\left\langle #1,#2\right\rangle}
\newcommand{\Grand}{\mathrm{Grand}}
\newcommand{\GT}{\boldsymbol{\mathbb T}_{\Grand}}
\newcommand{\GTX}{\boldsymbol{\mathbb T}_{\Grand,X}}
\newcommand{\Hist}{\mathcal A^{\mathrm{hist}}}
\newcommand{\Sym}{\operatorname{Sym}}
\newcommand{\Alt}{\operatorname{Alt}}
\newcommand{\Cl}{\mathrm{Cl}}
\newcommand{\HS}{\mathrm{HS}}
\newcommand{\op}{\mathrm{op}}
\newcommand{\CP}{\mathrm{CP}}
\newcommand{\UCP}{\mathrm{UCP}}
\newcommand{\vol}{\mathrm{vol}}
\newcommand{\Ric}{\mathrm{Ric}}
\newcommand{\Ein}{\mathrm{Ein}}
\newcommand{\Krein}{\mathcal K}
\newcolumntype{L}[1]{>{\raggedright\arraybackslash}p{#1}}


\title{\textbf{Renewal Geometry and the Emergence of Lorentzian Spacetime}}
\author{%
  Aur\'elien P\'elissier$^{1,\dagger}$\\
  {\small $^{1}$\'Ecole normale sup\'erieure Paris-Saclay, 91190 Gif-sur-Yvette, France}\\[-0.25em]
  {\small $^{\dagger}$Correspondence: \href{mailto:aurelien.pelissier.38@gmail.com}{aurelien.pelissier.38@gmail.com}}%
}
\date{}

\begin{document}
\maketitle

\begin{abstract}
We construct a finite predictive framework in which Lorentzian spacetime and vacuum Einstein dynamics emerge from operational renewal data rather than being assumed at the outset. Histories indistinguishable by all possible futures are identified, yielding a canonical predictive structure from which relational geometry is reconstructed. Two independent finite criteria select three spatial dimensions, while one additional protected signed direction fixes a $(1,3)$ Lorentzian signature. Calibrated renewal statistics then reconstruct the spatial metric, lapse, and shift.
%
A common variational structure is reconstructed directly from the finite dynamics. On certified noncollapsed continuum branches, it converges to Palatini--Holst gravity and yields the vacuum Einstein equation, while failed conditions produce explicit obstructions. The same framework also supports a cutoff-uniform nonsymmetric \(3+1\) Einstein population with arbitrary three-dimensional dependence, stable across an open family of renewal laws sharing the same continuum limit.
\end{abstract}




\medskip
\noindent\textbf{Keywords.} Emergent spacetime; Lorentzian geometry; predictive dynamics; renewal processes; variational dynamics; Palatini--Holst gravity; ADM decomposition.
\setcounter{tocdepth}{2}
%\tableofcontents


% =============================================================================
\section{Introduction}
\label{sec:introduction}
% =============================================================================

One of the oldest ambitions of fundamental physics is to explain why spacetime
has the structure that it does, rather than to take that structure as the stage
on which all other physics occurs.  General relativity transformed gravity
from a force propagating through spacetime into the dynamics of spacetime
itself, but it still begins with a differentiable manifold carrying a
Lorentzian metric.  Most subsequent approaches to quantum gravity likewise
replace some part of the classical geometric description while retaining
another part as primitive.  Causal-set theory begins with a locally finite
causal order \cite{bombelli-sorkin,surya}; graph, simplicial, and group-field
approaches begin with combinatorial incidence data whose collective phases are
expected to approximate geometry
\cite{konopkamarkopouloseverini2008,ambjornjurkiewiczloll2009,oriti2009};
and noncommutative geometry replaces the underlying space by spectral and
operator-algebraic data \cite{Connes1994}.  Other programmes seek gravitational
dynamics from thermodynamics, entropy, or quantum information
\cite{jacobson1995,verlinde2011,vanraamsdonk2010}.  These approaches differ
profoundly in their microscopic assumptions, yet they share a central
difficulty: the mathematical language used to reconstruct spacetime often
already contains some analogue of geometry, locality, causal order, dimension,
or temporal structure.

The question pursued here is more primitive.  We ask whether the structures
normally used to define spacetime can instead arise from the organization of
\emph{predictive histories}.  The microscopic object is not a point of a
manifold, a simplex, a causal element, or a prescribed spectral triple.  It is
an operational renewal event: an intervention occurs, a finite outcome is
registered, some distinctions remain capable of influencing later operations,
and the process may subsequently renew.  At this level a history is simply a
sequence of operational events.  It has no assigned position, distance, metric,
connection, curvature, or spacetime dimension, and its ordering is not yet
interpreted as motion through a pre-existing geometry.

\paragraph{Prediction before geometry.}
This choice of starting point is motivated by a simple observation.  Physical
description is never sensitive to an arbitrary microscopic history in its full
detail.  What matters operationally is which distinctions in that history can
still affect a future experiment.  Two past histories that produce exactly the
same statistics under every admissible continuation are therefore physically
indistinguishable for all subsequent purposes.  Identifying such histories is
a predictive rather than a geometric reduction.  It is closely related in
spirit to causal-state constructions in computational mechanics, where pasts
are quotiented by equality of their conditional futures
\cite{shalizicrutchfield,TraversCrutchfield2014}, and to continuous-time
predictive representations of stochastic processes
\cite{MarzenCrutchfield2017,JurgensCrutchfield2021}.  The difference is that
the present construction retains the complete operational intervention data
and the positive process structure required for the subsequent algebraic and
geometric reconstruction.

Because future equivalence is preserved under admissible continuation, the
quotient is not merely a compressed description of the past.  It carries a
canonical successor structure and therefore becomes the minimum persistent
predictive state of the process.  Persistence is thus reconstructed from the
organization of future distinguishability rather than supplied as an
independent state space, graph adjacency, or background coordinate system.

\paragraph{Renewal as a microscopic principle.}
This perspective leads naturally to renewal theory.  Renewal processes are
usually introduced as stochastic models of systems that repeatedly reset,
regenerate, or lose part of their past.  They underlie classical results in
probability and statistical mechanics \cite{feller}, while non-Markovian
renewal processes include heavy-tailed and infinite-mean regimes with
qualitatively different long-time behavior
\cite{garsialamperti1962,erickson1970,caravennadoney2019}.  In conventional
applications, however, renewal occurs \emph{inside} an already given time and
state space.  Waiting times are measured by an external clock, and the process
evolves on a previously specified collection of states.

Here we reverse that logical order.  Renewal is not placed in spacetime;
renewal belongs to the microscopic organization from which spacetime is
reconstructed.  The microscopic law may retain history through waiting-time,
record, or return statistics and need not be Markovian.  A local Markovian
evolution can nevertheless arise after predictive compression and coarse
graining, once unresolved renewal memory has collapsed onto a sufficient
future state.  In this sense the familiar locality of continuum evolution is
naturally interpreted as a \emph{Markovian scaling regime of a more general
renewal dynamics}, rather than as the fundamental form of microscopic
evolution.

There is a useful conceptual analogy with statistical mechanics and
renormalization.  Thermodynamic variables emerge after most microscopic
information has been discarded, while renormalization retains degrees of
freedom that remain relevant at longer scales.  Here the reduction is instead
organized by \emph{future distinguishability}: microscopic information is
discarded precisely when no admissible future can read it.  What survives is
therefore not an arbitrary coarse graining but a canonical predictive quotient.
The geometry developed below is built on that quotient.  Loosely speaking,
spacetime is the stable organization of those relational distinctions that
prediction cannot forget.

\paragraph{Main result.}
The central result is a finite-to-continuum reconstruction and certification
theorem.  Complete conditional future statistics first determine a canonical
minimum predictive carrier and a positive history algebra, whose intrinsic
word-Gram hierarchy together form the \emph{Renewal Grand Tensor}.  The finite relational dimension is selected in two independent ways.  Exact
endpoint--cycle source-rank balance among connected simple regular relations
selects the tetrahedral cell $K_4$ without an alternating-form hypothesis;
within the stronger pair-homogeneous simplex category, minimal faithful
alternating completion selects the same cell.  Its standard module is the
three-dimensional $A_3$ root space.  Calibrated root-race statistics then
reconstruct the ten ADM variables.  The positive Grand Tensor cannot itself
supply a negative direction: an independently occurring protected binary
orientation is an irreducible signed datum, and adjoining one signed direction
to the positive three-dimensional spatial form fixes Lorentzian inertia
$(1,3)$ before a Clifford realization is chosen.

The dynamical part is reconstructed from the complete accepted physical
cylinder rather than imposed as a continuum action.  A marked renewal pressure
recovers the full action slope, finite curl and period tests determine whether
the increments integrate to one scalar common action, and a positive
Fisher--Bregman gap, with its physical test-lift and cutoff constants retained,
certifies stationarity of the accepted field dynamics for that independently
reconstructed action.  Finite Gibbs/Kullback--Leibler identities separately
compare the accepted probability rows with their reconstructed costs.  After resolved memory is removed,
Lorentz covariance and the finite bivector commutant classify the first-order
gravitational score as Palatini--Holst plus determinant volume.  A positive
spatial cut margin supplies regulator-uniform Sobolev, Poincar\'e, counting,
and compact-screen bounds.  Continuous and fully discrete propagation estimates
give a spacetime $L^2$ curvature bound from initial energy and explicit
growth, transfer and differentiated-source budgets; a sharp action-current
comparison can supply the source budget.  Connection convergence and
interface-complete Cartan comparisons identify the resulting curvature.  On the branch where these finite and cofinal
certificates pass, the weak--strong Palatini limit gives
\begin{equation}
\boxed{G_{\mu\nu}+\Lambda g_{\mu\nu}=0}
\end{equation}
distributionally on the determining test core.  If a required condition fails,
the construction instead returns the corresponding finite or cofinal
obstruction.  The result is therefore an Einstein-or-obstruction theorem,
rather than a claim that every microscopic renewal process flows to general
relativity.  A failed sufficient certificate does not exclude an Einstein
subsequence obtainable by a different argument.

The positive branch is nonempty and is not confined to symmetry-reduced
kinematics.  In addition to the homogeneous and Gowdy sectors described
below, a full three-dimensional local harmonic renewal writer admits a
cutoff-independent nonlinear Cauchy theorem in a fixed Sobolev ball for all
ten metric components and all spatial lattice modes.  More strongly, this is
not an isolated engineered update: the same estimates and Einstein limit hold
uniformly on a fixed open $55$-dimensional family of local differential laws,
with no convergence of the law parameter required along the cutoff sequence.
Around finite realizations of these laws, conditional physical residual
moments and a profile-transport metric define relatively open neighborhoods
of history-dependent transition rows whose successful histories have the same
Einstein limit in probability and, along the declared odd refinement, eventually
almost surely.  The compatible initial records can also be produced at finite
cutoff by solving the original reconstructed action's complete lapse--shift
constraint rows from unconstrained small free tensor records; no continuum
vacuum metric is required as an input to that preparation.

These population results and the general certification theorem have different
logical roles.  \Cref{thm:main-open-3plus1,thm:main-open-law-basin} construct
nonsymmetric $3+1$ classes and reach the Einstein equation by cutoff-uniform
evolution, consistency and finite preparation estimates, whereas
\Cref{thm:main-renewal-einstein} is a recognition theorem for arbitrary
complete gravitational renewal cylinders.  Its five certificates are not the
definition of the constructive population.  Nor is the stable law family
silently identified with the raw exact finite Euler flow of the reconstructed
action.  That stronger route is analyzed independently in
\Cref{thm:main-exact-action-response,prop:main-exact-action-tangency} and
Appendix~\ref{supp:exact-action-audit}.  The exact-action analysis now contains
an all-cutoff regular transverse source block, an exact scalar--trace response
decomposition, source-neutral finite preparations, controlled literal initial
curvature on a nontrivial branch, and a hierarchy of temporal compatibility
tests.  A natural exactly prepared branch nevertheless develops a higher
tangency defect under its own flow; this defect cannot be removed by changing
only the leading multiplier tilt at fixed first jet, but it changes
nontrivially under admissible canonical first-jet deformations.  The remaining
direct exact-action problem is therefore a nonlinear cofinal tangency and
propagation problem, not merely a finite determinant test.

The homogeneous isotropic sector gives a complementary exact variational
realization: arbitrary positive renewal variables yield the full lapse-varied
vacuum Friedmann equations, so the flat and de~Sitter rate profiles arise as
stationary solutions rather than being prescribed in advance.  A finite
positive accept/reject process reconstructs the same reduced action directly
from measured trial probabilities and converges quantitatively to the
stationary branch.  Exact conservative time blocking of this quadratic class
preserves the scale-free action--law defect, so it does not generate an open
basin toward the matched surface.  Beyond homogeneity, a nonlinear unpolarized
Gowdy $T^3$ regulator preserves an exact finite momentum invariant and, on its
finite offset alphabet, has pathwise $O(h)$ convergence of the full Riemann
tensor on smooth bounded slabs.  These constructions provide complementary
open, homogeneous, and nonlinear symmetry-reduced populations without
asserting generic microscopic attraction to the certified branch.

The precise operational entrance is intentionally narrower than a claim that
all geometric information is generated from positive probabilities alone.
No predictive state space or persistent relational carrier is postulated
independently: histories are identified exactly when every admissible
continuation has the same conditional law, and the resulting
future-equivalence quotient is the unique minimum persistent predictive
carrier through which all future distinctions factor.  The geometric branch
also retains two kinds of information that this positive predictive quotient
cannot manufacture.  A protected binary orientation is an actual two-valued
operational distinction, but no interpretation as past and future,
positive and negative norm, or Lorentzian signature is assumed microscopically.
Likewise, normalized race probabilities determine only relative rates, so
absolute waiting-time and measure calibrations are required when those rates
are interpreted geometrically.

A further distinction is essential.  The future-equivalence class is derived
as an abstract predictive state, but spacetime reconstruction must physically
read that class at a cut and correlate it on the same history with the
occurring root event, scheduler-separated fast path, orientation, calibration,
and other declared marks.  These data form a minimal faithful readable relation
packet.  The packet does not introduce a second predictive state: its state
coordinate is the already derived future-equivalence class.  Its additional
content is physical occurrence and same-cut readability.  The enrichment is
conservative: forgetting the added marks recovers the original event cylinders
exactly, while future minimization removes every unread refinement.  Thus
predictive persistence is derived, the physical information needed for the
geometric reconstruction must actually occur together, and metric, connection,
curvature, and spacetime dimension remain downstream outputs.

\paragraph{The Renewal Grand Tensor and Lorentzian reconstruction.}
The first major structural output is the object we call the
\emph{Renewal Grand Tensor}.  The name does not refer to a tensor field on a
pre-existing manifold.  Predictively indistinguishable histories are first
identified, producing the unique minimum persistent predictive carrier.
Tomographically complete intervention statistics then reconstruct the finite
positive process tensors, while support-minimal completely positive
realization removes unobservable environmental degrees of freedom
\cite{Stinespring1955,Choi1975}.  The surviving mixed histories generate a
finite $C^*$-algebra, and their intrinsic positive overlaps form a word-Gram
hierarchy.  Together these structures record which distinctions survive into
the future, how renewal operations compose, and how the resulting histories
overlap.  The Grand Tensor is therefore not geometry imposed on a renewal
process; it is the canonical algebraic and positive structure retained by the
process after every future-invisible distinction has been removed.

This places the construction near noncommutative and spectral geometry, but
with a different starting point.  In ordinary spectral geometry an algebra,
Hilbert-space representation, and Dirac operator are supplied as geometric
data \cite{Connes1994}.  Lorentzian variants enrich this structure with an
indefinite inner product, fundamental symmetry, temporal datum, or twist
\cite{strohmaier,franco,franco-eckstein,
vandendungenpaschkerennie2013,vandendungen2016,
DevastatoFarnsworthLizziMartinetti2018,
DevastatoFilaciMartinettiSingh2020}.  Recent work has also examined the
emergence of temporal and Lorentzian structures within twisted
almost-commutative settings \cite{nieuviarts2026a,nieuviarts2026b}.  The
present construction operates one level upstream.  Neither spectral geometry
nor a Krein structure is taken as the initial object.  The positive history
algebra is reconstructed from operational renewal statistics, while the
protected binary orientation becomes an indefinite/Krein structure only after
it acts on the predictively reduced relational carrier.  Its Clifford
realization then distinguishes one timelike direction.

Predictive collapse does more than compress a finite history.  Because future
equivalence is preserved under admissible continuation, the minimum predictive
state propagates canonically from one event to the next and therefore supplies
a persistent relational carrier.  The local dimension is then tested by two
logically distinct finite selectors.  The first uses no alternating structure:
for a connected simple regular relation, exact equality of the faithful
endpoint-source and cycle-source ranks has the unique nontrivial solution
$K_4$.  The second works in the more restrictive pair-homogeneous simplex
category and asks for a nondegenerate alternating temporal--spatial coefficient
pairing together with minimum faithful source innovation; it independently
selects the same $K_4$.  Thus the alternating form is not the sole route to
three spatial dimensions.  It remains a branch axiom of the second selector
and supplies only an antisymmetric coefficient pairing, not a symplectic phase
space, Hamiltonian flow, metric, or Lorentzian sign.  Neither selector is a
universality theorem over arbitrary higher-arity microscopic systems.

The metric appears one layer later.  The relational carrier supports a finite
family of elementary root transitions.  Their occurrence rates and waiting
times provide observable first- and second-order statistics.  The second-order
statistics determine the spatial kinetic form, the directed first moment
determines the shift, and a physical speed density fixes the remaining scale.
These quantities invert uniquely to the ADM spatial metric, lapse, and shift.
The metric coefficients are therefore not microscopic labels but collective
statistical quantities reconstructed from the renewal process.

The geometric clock is likewise not identified with the slower renewal clock.
A nondegenerate continuum bracket requires a scheduler-separated root race
with total rate of order $h^{-2}$.  Its winner and waiting-time statistics
supply the absolute kinetic scale, whereas root-direction probabilities alone
leave the common rate scale undetermined.  The protected binary orientation then acquires its geometric meaning.  The
positive Grand-Tensor and root-race forms remain positive and therefore cannot
produce an indefinite direction.  The orientation Read supplies one
independent signed line; leaving the reconstructed three-dimensional spatial
form positive, the minimal nondegenerate signed extension has inertia
$(-,+,+,+)$, up to overall sign convention.  Only after this inertia is fixed
do we choose a minimal Clifford representation.  The gamma matrices therefore
realize a signature already determined by the one-sign extension; they do not
select the signature.  Thus the Lorentzian sign is neither obtained by analytic
continuation of a positive diffusion tensor nor inserted through a Dirac-matrix
ansatz.

\paragraph{From kinematics to Einstein dynamics.}
Producing Lorentzian kinematics is not yet a derivation of gravitational
dynamics.  In particular, prescribing a renewal-rate history whose
reconstructed metric is de~Sitter would establish only realizability of one
Einstein solution.  The dynamical part therefore begins instead from the
complete accepted physical cylinder.  Future refinement reconstructs the
minimum determining-field quotient and the accepted transition kernel acting
on it.  On the same resolved history, a protected vector reward together with
physical duration determines an analytic renewal pressure.  Its first
derivative restores the affine action anchor lost by a centered probability
score, while its higher jets reconstruct the mixed action Gram.  Finite
face-curl and period tests then determine whether these increments integrate
to a single scalar common action and reconstruct that action when they do.

The actually accepted transition is next compared with the update generated
by this independently reconstructed action.  On the interior variational
branch, the Fisher--Bregman gap controls the squared represented first
variation, and explicit physical test-lift bounds retain the cutoff dependence
needed in the continuum passage.  Exact Gibbs/Kullback--Leibler identities
provide a separate comparison of finite probability rows.  The field-action
comparison, rather than a cost fitted after observing the transition law,
supplies the stationarity test.

The gravitational handoff is likewise formulated as a reconstruction and
certification problem.  The protected cycle reward entering this step is
initially a generic measured face reward: no Palatini, Holst,
Einstein--Hilbert, or cosmological Lagrangian is named in the entrance data.
After resolved cut memory is removed, binary renewal subdivision and stable
additivity produce a first-order face functional without assuming a derivative
at a zero face.  Lorentz covariance and the residual bivector commutant then
leave only the identity and internal Hodge star, so only at this stage is the
local first-order gravitational score classified as Palatini--Holst plus
determinant volume.

Connection stationarity yields zero torsion on the spinless branch.  Global
spatial compactness is not introduced as a separate geometric axiom: the
physical endpoint/capacity record determines an exact cut constant, and a
uniform positive margin gives Sobolev, Poincar\'e, low-energy-counting, and
finite-rank compact-screen bounds.  Curvature control is supplied by a separate
positive propagation estimate.  A bounded initial curvature energy together
with regulator-uniform observer/transport and differentiated-source budgets
yields a spacetime $L^2$ curvature bound, also for a fully discrete curvature
record.  Strong coframe convergence supplies the nonlinear product passage;
strong connection convergence and interface-complete consistency separately
identify its curvature.  These are the two analytic inputs to the
weak--strong Palatini passage.  On the positive certificate branch,
the limiting metric variation is
\begin{equation}
\chi\int_K\sqrt{-g}\,
(G_{\mu\nu}+\Lambda g_{\mu\nu})k^{\mu\nu},
\end{equation}
and determining stationarity therefore gives the vacuum Einstein equation
distributionally.

These finite reconstructions assemble into a certification-or-obstruction
alternative.  For a complete gravitational renewal cylinder, the
determining-field quotient, accepted kernel, renewal pressure, common-action
exactness, action gap, spatial cut constant, face linearization, memory short,
connection residual, curvature-propagation budgets, and Cartan compatibility
are finite or cofinally testable.  If the corresponding positive margins
persist, the Einstein continuum follows.  If they do not, the construction
returns the failed object itself: a nonzero common-action curl or period, a
positive Bregman/KL row, a macroscopic cut/potential witness, a
face-linearity or memory residual, a torsion/connection defect, an unbounded
curvature-propagation budget, a Cartan mismatch, or a cutoff compactness
defect.

The homogeneous isotropic sector addresses a narrower question not settled by
the abstract Einstein handoff.  Instead of prescribing a de~Sitter rate
profile, we begin with arbitrary positive root-symmetric variables
$\kappa(t)$ and $\varrho(t)$.  The reconstructed ADM map converts them
bijectively to the scale factor and lapse.  Rewriting the classified vacuum
action in $(\kappa,\varrho)$ and varying both variables before fixing the lapse
yields the Hamiltonian constraint and scale equation, equivalently the full
spatially flat vacuum Friedmann system.  The exponential de~Sitter profile
therefore appears only after the stationary equations have been solved.

The same distinction persists at finite cutoff.  A midpoint variational action
retains independent lapse-weighted intervals, and its stationary recurrence
converges at second order to the continuum flat/de~Sitter branch.  A finite
positive accept/reject law reconstructs this action directly from its proposal
distribution, accepted row, and total acceptance probability.  A measured rank
defect determines whether the local low-noise update and the cosmological
parameter of the reconstructed action agree.  On the compatible branch, finite
stochastic cylinders have vanishing Euler residual and converge quantitatively
to the stationary de~Sitter solution.

The finite evolution mechanism is not confined to homogeneous trajectories.
In the unpolarized Gowdy $T^3$ vacuum class both wave polarizations and a
spatially varying nonlinear frame are retained.  A
source/characteristic/source splitting preserves an exact staggered momentum
invariant, while the finite offset support and shared Hermite readout give
pathwise second-jet control and hence $O(h)$ convergence of the full Riemann
tensor on fixed smooth slabs.  The independently varied unreduced action has
the same vanishing continuum first variations.  This symmetry-reduced theorem
does not imply generic attraction to de~Sitter or unrestricted four-dimensional
regularity, but it provides a nonlinear inhomogeneous class in which full
curvature compactness is derived rather than assumed.

\paragraph{Scope and organization.}
The emergence claim is relative to a precise operational entrance.  Actual
occurrence, complete contextual prediction, protected orientation, and
absolute physical calibration are irreducible rows of the declared finite
language.  For gravitational dynamics one additionally needs the complete
accepted field/action cylinder to occur physically.  Once these records are
present, however, the predictive carrier, determining field, accepted kernel,
action slope, common scalar action, ADM geometry, compact spatial screen, and
stationarity residual are reconstructed or certified rather than supplied as
independent continuum geometry.  Vanishing of the action gap and persistence
of the positive noncollapse and curvature-propagation margins remain genuine
physical conditions that can fail.

The dimensional claim has the same controlled scope.  The $K_4/A_3$ result
does not state that every possible microscopic theory must have three spatial
dimensions.  Rather, two independent finite classifications select the same
cell: endpoint--cycle source-rank balance does so without an alternating-form
assumption, while the alternating/minimal-source simplex selector does so in a
stronger category.  Their agreement removes the alternating form from the role
of sole dimensional selector, while still falling short of universality over
arbitrary higher-arity systems.  Likewise the protected binary orientation is
not a microscopic Lorentzian metric.  The positive operational rows cannot
reconstruct it, as shown by explicit separating models; once the $A_3$ spatial
form is positive, this single independent sign fixes the minimal indefinite
extension to Lorentzian inertia $(1,3)$.  The subsequent Clifford matrices are
a representation of that quadratic form, not the origin of its signature.

The conceptual proposal is therefore that renewal may be more fundamental
than its conventional role in probability theory suggests.  From the usual
point of view, renewal describes memory and resetting \emph{within} time.  From
the present point of view, retention and forgetting in operational history
determine which distinctions survive at all, while predictive equivalence
reconstructs the persistent state on which relational geometry is built.
Effective Markovian locality then appears not as the microscopic starting
principle but as a limiting organization of predictive memory.  In this sense,
the usual hierarchy
\begin{equation}
\boxed{\text{spacetime}\longrightarrow\text{stochastic dynamics}}
\end{equation}
is replaced by
\begin{equation}
\boxed{
\text{renewal dynamics}
\longrightarrow
\text{predictive structure}
\longrightarrow
\text{relational geometry}
\longrightarrow
\text{Lorentzian spacetime}.}
\end{equation}
The claim is not that renewal theory in its classical form already contains
general relativity.  It is that when renewal is treated operationally,
predictively, and with a minimal protected orientation, persistence and
relational structure can be reconstructed before ordinary spacetime geometry
is defined.

Section~\ref{sec:grand-tensor-main} constructs the minimum predictive carrier
and the Renewal Grand Tensor from finite operational data.
Section~\ref{sec:adm-main} identifies the minimal faithful relational branch
and reconstructs the ADM and Lorentzian structures from the calibrated root
process.  Section~\ref{sec:einstein-dynamics} reconstructs the
determining-field kernel and common action, introduces the accepted/action
stationarity criterion, derives compact spatial screens from a finite cut
margin, propagates curvature control from finite initial-energy and scaled
source budgets in continuous or discrete time, and identifies connection
curvature including interfaces.  It then gives a cutoff-uniform nonsymmetric
$3+1$ population theorem, enlarges it to an open family of differential laws
and relatively open finite transition-row neighborhoods with exact
finite-action initial preparation, develops the exact-action response and
dynamical-tangency hierarchy, and only afterward states the maximally general
Palatini--Holst certification and Einstein-or-obstruction alternative.
Section~\ref{sec:vacuum-main} studies complementary explicit vacuum sectors:
it derives the full homogeneous Friedmann system and
its positive operational realization, proves that conservative homogeneous
blocking preserves the relative matching defect and characterizes matching
by optimizer composition, develops a nonlinear
inhomogeneous unpolarized Gowdy regulator with exact discrete momentum
propagation and pathwise full-curvature convergence, and recovers the flat and
de~Sitter regulators as stationary solutions rather than prescribed rate
profiles.
Section~\ref{sec:discussion} discusses the exact scope of the result and the
remaining selection, exact-action tangency/propagation, and universality questions.
Technical reconstruction, variational, representation-theoretic,
operational-dynamical, and continuum estimates are collected in
Appendices~\ref{supp:grand-reconstruction}--\ref{supp:vacuum}.





% =============================================================================
\section{Predictive reconstruction and the Renewal Grand Tensor}
\label{sec:grand-tensor-main}
% =============================================================================

\subsection{Operational histories and future equivalence}

At a finite cutoff $X$, the operational entrance is a finite typed event
protocol with finite admissible histories, finite intervention and outcome
alphabets, and normalized conditional probabilities for every reachable
history.  The physical intervention letters include tomographically complete
frames on each exposed finite operator port; their multitime probabilities are
affine under randomization, additive under coarse graining, causal under
deterministic discard, and compatible on common prefixes.  A distinguished
binary orientation is retained as an actual protected operational distinction.
No predictive state space, persistent relational carrier, metric, scheduler,
connection, or spacetime dimension is part of this entrance.

For a reachable history $h$ ending at cut type $x$, let $\mathsf F_{X,x}$ be
the set of all finite admissible future words starting at $x$.  Write
$p_X(w\mid h)$ for the corresponding conditional cylinder probability.  Define
\begin{equation}
 h\sim_X h'
 \quad\Longleftrightarrow\quad
 p_X(w\mid h)=p_X(w\mid h')
 \quad\text{for every }w\in\mathsf F_{X,x}.
 \label{eq:future-equivalence-main}
\end{equation}
Histories of different cut type are never identified.  If additional actual
terminal Reads or bounded readout queries are present, they are included in
the future word alphabet; thus \eqref{eq:future-equivalence-main} means equality
of the complete declared future signature.

Future equivalence is a typed right congruence.  Indeed, if $h\sim_Xh'$ and an
event $a:x\to y$ is admissible, then the one-letter future gives
$p_X(a\mid h)=p_X(a\mid h')$.  On a positive-probability branch, the chain rule
then gives
\[
 p_X(v\mid ha)
 =\frac{p_X(av\mid h)}{p_X(a\mid h)}
 =\frac{p_X(av\mid h')}{p_X(a\mid h')}
 =p_X(v\mid h'a)
\]
for every continuation $v$.  Hence $ha\sim_Xh'a$.

\subsection{The minimum persistent predictive carrier}

Define
\begin{equation}
 \boxed{
 Z_{X,x}^{\min}:=\mathsf H_{X,x}/\!\sim_X,
 \qquad q_{X,x}:\mathsf H_{X,x}\twoheadrightarrow Z_{X,x}^{\min}.}
 \label{eq:predictive-state-main}
\end{equation}
For an admissible event $a:x\to y$, the right-congruence property defines
\begin{equation}
 \boxed{
 T_{X,a}([h])=[ha],
 \qquad
 \kappa_{X,a}([h])=p_X(a\mid h).}
 \label{eq:predictive-update-main}
\end{equation}
Thus the event law itself reconstructs a deterministic predictive state
machine.  If $A_1,A_2,\ldots$ is the realized event sequence,
$H_n=A_1\cdots A_n$, and $Z_n=q_X(H_n)$, then
\begin{equation}
 \boxed{
 Z_{n+1}=T_{X,A_{n+1}}(Z_n),
 \qquad
 \mathfrak p_{X,n}^{\min}
   =(Z_n,A_{n+1},Z_{n+1}).}
 \label{eq:derived-relation-packet-main}
\end{equation}
The persistent relation is therefore the transport of a future-equivalence
class under event composition, not an independently supplied state variable.

The quotient is universal.  Every reachable deterministic predictive
realization $S_{X,x}$ of the same conditional event law admits a unique
surjection
\begin{equation}
 \pi_{X,x}:S_{X,x}\twoheadrightarrow Z_{X,x}^{\min}
 \label{eq:predictive-universal-main}
\end{equation}
intertwining branch probabilities and updates.  Consequently every declared
downstream readout depending only on the complete future signature factors
uniquely through $Z_X^{\min}$.  Generated ledgers, hidden-state labels, and
other deterministic predictive presentations are therefore representations of
the canonical quotient rather than additional foundational coordinates.

The physical geometric branch considered below requires more than this
predictive reduction.  In particular, same-cut readability of the predictive
class, the protected orientation, and the absolute winner--waiting calibration
of the fast root race are actual physical rows.  They are not determined by
the positive predictive quotient.  The purpose of
\eqref{eq:predictive-state-main}--\eqref{eq:derived-relation-packet-main} is to
remove independently postulated persistence and every future-invisible state
refinement.

\subsection{Conservative readable relational completion}

On a spacetime-capable branch, let $O_{X,n}$ denote the occurring root-bearing
event at the $n$th slow cut, let $\Gamma_{X,n}$ denote the scheduler-separated
fast root path between slow cuts, and let $\Xi_{X,n}$ collect the remaining
same-history physical marks used by the geometric reconstruction, including the
protected orientation and absolute calibrations.  A raw readable relation
packet is
\begin{equation}
 \boxed{
 \mathfrak p_{X,n}^{\rm rel}
 =\bigl(
 Z_{X,n}^{\min},\Gamma_{X,n},O_{X,n},\Xi_{X,n},Z_{X,n+1}^{\min}
 \bigr).}
 \label{eq:readable-relation-packet-main}
\end{equation}
The state coordinates in \eqref{eq:readable-relation-packet-main} are not new
hidden variables: they are same-cut physical realizations of the predictive
classes already reconstructed in \eqref{eq:predictive-state-main}.  Two raw
packets are identified when every admitted continuation, readout query, and
terminal Read has the same conditional law.  Denote the resulting finite
future-minimal quotient by $\mathfrak P_X^{\rm rel,min}$.

\begin{proposition}[Conservative readable completion]
\label{prop:readable-relational-completion}
Suppose the additional readable marks in
\eqref{eq:readable-relation-packet-main} are attached to the underlying finite
operational process by normalized same-history conditional kernels.  Then:
\begin{enumerate}[label=\textup{(\roman*)},leftmargin=2.8em]
\item forgetting $\Gamma$, $\Xi$, and the explicit same-cut realization of the
      predictive class recovers every cylinder probability and branch map of
      the underlying operational process exactly;
\item future-signature refinement terminates at
      $\mathfrak P_X^{\rm rel,min}$, and every other finite readable packet
      sufficient for the same geometric reconstruction has a unique surjection onto
      it;
\item unread refinements add neither predictive content nor source dimension;
      metric, coframe, connection, curvature, and Einstein data remain
      downstream readout outputs and are not coordinates of the relation
      packet.
\end{enumerate}
Thus physical readability is a conservative enrichment of a derived
predictive carrier, not a second postulated state space.
\end{proposition}

The proof is given in Appendix~\ref{supp:readable-relational-completion}.  This
proposition is the interface between predictive reconstruction and the
geometric branch: the first determines which state distinctions exist, while
the second records which of those distinctions and physical calibrations
actually occur together on one history.

\subsection{The Grand Tensor}

The operational multitime table on the physical intervention frames
reconstructs a positive causal process tensor.  Its support-minimal square-root
realization supplies canonical sequential memories, unique up to
record-preserving memory unitaries.  The represented minimal-state projections,
mixed forward maps, and their Hilbert--Schmidt reverses generate a finite
history algebra $\mathcal A_X^{\mathrm{hist}}$.  If $[w]$ denotes the
represented history word and $\widehat\tau_{\mathrm{reg}}$ is the normalized
regular trace, define
\begin{equation}
 \mathbb K_{X,r}(v,w)
 =\widehat\tau_{\mathrm{reg}}\bigl([v]^*[w]\bigr),
 \qquad |v|,|w|\le r.
 \label{eq:word-gram-main}
\end{equation}
The positive structural Renewal Grand Tensor is
\begin{equation}
 \boxed{
 \GTX
 =\left(
 \mathcal A_X^{\mathrm{hist}},
 \{\mathbb K_{X,r}\}_{r\ge0},
 Z_X^{\min}
 \right).}
 \label{eq:grand-tensor-main}
\end{equation}
It is a tensor in the operational sense of a complete positive moment object:
the word-Gram hierarchy records all represented mixed moments, while the
history algebra records the multiplication and involution those moments
represent.  Finite dimensionality forces the represented word hierarchy to
become flat after a finite depth, at which point one further word layer
reconstructs the complete represented algebra.

For the spacetime branch, \eqref{eq:grand-tensor-main} is used together with
the minimal faithful readable packet $\mathfrak P_X^{\rm rel,min}$ of
\Cref{prop:readable-relational-completion}.  The distinction is deliberate:
$\GTX$ records the canonical positive predictive-history structure, whereas
the readable packet records the actual same-history occurrence of orientation,
calibration, scheduler, and root-path data required to interpret that
structure geometrically.  Forgetting the latter returns the positive
structural tensor exactly.

\begin{theorem}[Canonical predictive Grand-Tensor reconstruction]
\label{mt:grand-tensor}
Let a finite typed operational renewal process have a finite event grammar,
normalized prefix-compatible conditional probabilities, tomographically
complete physical intervention frames, finite operational positivity, causal
trace recursion, and a protected binary orientation as described above.  Then:
\begin{enumerate}[label=\textup{(\roman*)},leftmargin=2.8em]
\item future equivalence \eqref{eq:future-equivalence-main} is a typed right
      congruence, and $Z_X^{\min}$ is the unique coarsest reachable
      deterministic predictive realization of the event law;
\item the successor maps \eqref{eq:predictive-update-main} are canonical and
      the persistent event packet \eqref{eq:derived-relation-packet-main} is
      derived from composition rather than independently postulated;
\item the state-conditioned Hankel quotient is the unique reachable,
      future-separated minimum linear predictor and determines the executable
      forward algebra;
\item the physical probability table reconstructs unique positive causal
      prefix combs whose square-root supports determine a canonical
      support-minimal sequential $\CP$ class, unique up to
      minimal-state-preserving unitaries;
\item the contextual history quotient, normalized regular trace, and positive
      word-Gram hierarchy are canonical, and finite flatness reconstructs the
      complete represented monoidal $*$-algebra;
\item every deterministic predictive presentation and every future-operational
      compiler factors uniquely through $Z_X^{\min}$; consequently
      \eqref{eq:grand-tensor-main} is canonical without an independently
      supplied ledger, hidden state, persistent state space, preferred Kraus
      frame, primitive metric, or preselected history algebra;
\item any finite same-history readable enrichment by normalized
      conditional kernels admits the minimal faithful quotient
      $\mathfrak P_X^{\rm rel,min}$ of
      \eqref{eq:readable-relation-packet-main}; its forgetful image is the
      underlying operational law, so same-cut readability does not alter the
      predictive reconstruction or introduce new geometric coordinates.
\end{enumerate}
\end{theorem}

The detailed proof is given in Appendix~\ref{supp:grand-reconstruction}.  Three minimizations
remain logically distinct: predictive state minimization, linear Hankel
minimization, and contextual $C^*$-history separation.  Conflating them either
retains future-invisible spectator structure or erases distinctions visible to
an actual operational continuation.

The foundational status of the retained operational rows is sharpened in
\Cref{thm:supp-relative-primitive-floor}.  Relative to the fixed typed audit
language used here, actual occurrence, complete contextual prediction,
protected orientation, and absolute physical calibration form a minimal known
independent operational basis for the paper-level reconstruction.  The
remaining calibration freedom is itself classified in
\Cref{thm:supp-calibration-fibre}: physically established log-affine incidence
relations leave an exact kernel quotient, and one scalar anchor per independent
physical quotient direction is necessary and sufficient.  These are relative
minimality statements for the declared language, not claims of absolute
metaphysical minimality.

\begin{remark}[What the Grand Tensor represents]
The Grand Tensor is not an array attached to an already existing geometry.  It is the finite positive history object from
which the geometry is subsequently read.  Its predictive component says which
past distinctions remain effective and how they persist under continuation;
its algebra says how represented operations compose; and its Gram hierarchy
says how those histories overlap in the canonical intrinsic positive form.
Geometry becomes possible only after these questions have been answered on one
common, future-minimal carrier.
\end{remark}

% =============================================================================
\section{Minimal relational geometry and Lorentzian ADM reconstruction}
\label{sec:adm-main}
% =============================================================================

\subsection{Two independent finite selectors for the tetrahedral cell}

The local dimension is not inferred from a single structural axiom.  We use
instead two finite selectors with different hypotheses.  The first is the
logically weaker one and requires no alternating form: exact balance between
faithful endpoint and cycle source ranks inside connected simple regular
relations uniquely gives $K_4$.  The second works in the stronger
pair-homogeneous simplex category and asks for a nondegenerate alternating
coefficient pairing together with minimum faithful source innovation; it gives
the same cell.  Their agreement is the dimensional statement used below.

\begin{proposition}[Endpoint--cycle balance selects $K_4$ without an alternating form]
\label{prop:main-regular-graph-k4}
Let $G$ be a connected simple $k$-regular relation graph on $v>1$ endpoints.
If its faithful endpoint source and cycle source have equal rank, then
\begin{equation}
 v-1=\frac{kv}{2}-v+1,
 \qquad\hbox{so}\qquad
 (4-k)v=4.
 \label{eq:main-regular-graph-balance}
\end{equation}
The unique nontrivial connected simple regular solution is
\begin{equation}
 \boxed{v=4,\qquad k=3,\qquad G\cong K_4.}
 \label{eq:main-regular-graph-k4}
\end{equation}
Thus the associated faithful endpoint module has dimension $v-1=3$.  This
selector contains no temporal--spatial alternating-form hypothesis.
\end{proposition}

The proof is included in Appendix~\ref{supp:dimension-adm}, where the same
balance is derived from the cut--cycle complex before the stronger simplex
selector is introduced.

For the simplex selector, consider a local relation event on the minimum
persistent predictive carrier of \Cref{mt:grand-tensor}, with one
future-readable source and one future-readable target.  Physical reversal
turns it into a signed edge.  Future minimization removes parallel labels that
no future can distinguish, while pair-homogeneous recurrence completes the
orbit of one nonzero relation.  The relation-colour-free skeleton is therefore
$K_N$.  Writing $\R^N=\R u_0\oplus W_N$, with
$u_0=N^{-1/2}\one_N$, the signed edge space decomposes as
\begin{equation}
 C_1(K_N)=u_0\wedge W_N\oplus\Lambda^2W_N.
 \label{eq:cut-cycle-main}
\end{equation}
The first summand is the endpoint or cut source and the second the complete
cycle or interference source, of dimensions $N-1$ and
$\binom{N-1}{2}$ respectively.

\begin{proposition}[Alternating/minimal-source simplex selector]
\label{prop:main-alternating-k4}
Assume in addition that the temporal--spatial coefficient space
$\R t\oplus W_N$ carries a nondegenerate alternating form, that the endpoint
and complete signed-edge sources are faithful, that the cycle source is
nonzero, and that the future-readable source innovation is minimal among
conservative completions satisfying these queries.  Then
\begin{equation}
 \boxed{N=4,\qquad G_{\mathrm{rel}}=K_4,
 \qquad W_{\mathrm{sp}}=W_4\cong\R^3.}
 \label{eq:k4-main}
\end{equation}
Moreover, after complexification,
\begin{equation}
 \Hom_{S_N}(\Lambda^2W_N,W_N\otimes\operatorname{sgn})
 =\begin{cases}\C,&N=4,\\0,&N\ne4,\end{cases}
 \label{eq:main-oriented-copy}
\end{equation}
so only at $N=4$ is the complete cycle representation one oriented copy of the
spatial standard module rather than an additional inequivalent source type.
\end{proposition}

The proof is given in Appendix~\ref{supp:dimension-adm}.  The first conclusion
uses only parity, nonzero interference, and source minimality; the
representation-theoretic statement is independent evidence for the exceptional
economy of $N=4$ inside the simplex family.

\begin{remark}[Exact scope of the dimensional claim]
\label{rem:main-dimension-scope}
The endpoint--cycle balance theorem and the alternating/minimal-source theorem
are independent finite classifications with different hypotheses.  The first
shows that the alternating form is not the sole origin of the three-dimensional
spatial module.  The second does not become a hidden symplectic or Hamiltonian
assumption: its alternating form supplies only an antisymmetric coefficient
pairing and no metric, lapse, shift, connection, curvature, or Lorentzian sign.
Neither theorem is a classification of arbitrary graphs, hypergraphs, or
higher-arity microscopic relation systems.
\end{remark}

\subsection{The \texorpdfstring{$A_3$}{A3} root race}

Choose an anchored basis of $W_4$ in which the six unoriented $A_3$ roots are
\begin{equation}
 r_1=e_1,\quad r_2=e_2,\quad r_3=e_3,\quad
 r_4=e_2-e_1,\quad r_5=e_3-e_1,\quad r_6=e_3-e_2.
 \label{eq:a3-roots-main}
\end{equation}
Let $k_a^\pm$ be the two directed rates associated with $\pm r_a$, and let $h$
be the mesh scale.  The predictable second and first moments are
\begin{equation}
 \mathcal B_h
 =h^2\sum_{a=1}^{6}(k_a^++k_a^-)r_ar_a^{\mathsf T},
 \qquad
 \beta_h
 =h\sum_{a=1}^{6}(k_a^+-k_a^-)r_a.
 \label{eq:bracket-drift-main}
\end{equation}
The six dyadics $r_ar_a^{\mathsf T}$ form a basis of $\Sym_3$, so the rate sums
reconstruct all six spatial inverse-metric coefficients.  The rate differences
map with rank three onto $\beta_h$; their remaining three coordinates are
circulating root currents.  Winner probabilities alone lose the common rate
scale.  Observing both the winning root and its waiting time gives a faithful
score frame for all directed rates.

\begin{remark}[The missing scale is an exact calibration fibre]
\label{rem:main-calibration-fibre}
More generally, write all positive scales and additive origins in log-affine
coordinates $x\in\mathbb R^N$.  If the physically established relative
incidences are $Lx=b$, the remaining calibration class is an affine translate
of $\Ker L$; after quotienting conventional changes of base units it is
$\Ker L/\operatorname{Im}U$, where $LU=0$.  One scalar anchor per independent
quotient direction is necessary and sufficient.  The winner-only root race is
the simplest instance: common rescaling of all rates is invisible to winner
probabilities, and the waiting-time observation supplies the missing rate
anchor.  The speed density $\varrho_h$ below is a separate physical measure
calibration.  The complete finite statement is
\Cref{thm:supp-calibration-fibre}.
\end{remark}

Let $\varrho_h>0$ be the physical speed density.  The spatial metric and lapse
are reconstructed by
\begin{equation}
 \boxed{
 g_h=\varrho_h^{2/3}(\det\mathcal B_h)^{1/3}\mathcal B_h^{-1},
 \qquad
 N_h=\varrho_h^{1/3}(\det\mathcal B_h)^{1/6}.}
 \label{eq:adm-inversion-main}
\end{equation}
Equivalently,
\begin{equation}
 \sqrt{\det g_h}=\varrho_h,
 \qquad
 \mathcal B_h=N_h^2g_h^{-1}.
 \label{eq:adm-forward-main}
\end{equation}
Together with the shift $\beta_h$, these are exactly the ten ADM coefficients.
The formula is an analytic bijection on the positive cone.  A nondegenerate
bracket requires total root rate of order $h^{-2}$; identifying the geometric
scheduler with a slower $O(h^{-1})$ renewal clock would make
$\mathcal B_h\to0$ and collapse the metric.

\subsection{From positive predictive geometry to one irreducible signed direction}

The Grand Tensor reconstructed in the previous section is positive: its
canonical word-Gram hierarchy is a Hilbert-space Gram structure, and the
root-race bracket is a positive quadratic form on the three-dimensional
spatial module.  Neither can contain a negative direction.  The additional
physical datum is a protected two-valued operational Read
\begin{equation}
 \epsilon_X\in\{+1,-1\},
 \label{eq:orientation-read-main}
\end{equation}
with both values separated by the declared Read policy.  In the support-minimal
Hilbert realization it has orthogonal readout projections $P_+$ and $P_-$,
$P_++P_-=I$, and therefore
\begin{equation}
 \boxed{J:=P_+-P_-,\qquad J=J^*=J^{-1},}
 \qquad [u,v]_J:=\langle u,Jv\rangle.
 \label{eq:krein-main}
\end{equation}
This is the operator representation of an already occurring $\mathbb Z_2$
mark.  The microscopic record itself carries no norm, cone, lapse, temporal
coordinate, or indefinite metric.

The sign is not reconstructible from the positive packet.  The separating
models of \Cref{thm:supp-relative-primitive-floor} keep the complete positive
operational data fixed while reversing the protected orientation, and the
obstruction survives common conservative refinement.  The signed row is
therefore irreducible relative to the declared operational language rather than
a reparametrization of the positive Grand Tensor.

\begin{proposition}[Minimal one-sign extension has Lorentzian inertia]
\label{prop:main-signature-minimality}
Let $(W,g)$ be the positive three-dimensional spatial quadratic space
reconstructed from the $K_4/A_3$ root race.  Adjoin one independent real line
$\R t$ whose protected orientation supplies the opposite sign and leave
$g|_W$ unchanged.  Then every minimal nondegenerate orthogonal signed extension
is, up to multiplication of the whole quadratic form by $-1$,
\begin{equation}
 \boxed{q(\tau t+w)=-N^{-2}\tau^2+g(w,w),
 \qquad \operatorname{inertia}(q)=(1,3).}
 \label{eq:main-signature-inertia}
\end{equation}
No choice of Clifford representation enters this conclusion.
\end{proposition}

The proof and the corresponding independence statement are given in
Appendix~\ref{supp:dimension-adm}.  The point is that the inertia conclusion
precedes the Clifford realization.

Thus the conceptual chain is
\begin{equation}
\boxed{
\begin{aligned}
\text{predictive collapse}
&\longrightarrow \text{positive Grand-Tensor geometry}
\longrightarrow K_4/A_3\text{ spatial geometry}\\
&\longrightarrow \text{one irreducible signed direction}
\longrightarrow \text{Lorentzian inertia }(1,3)\\
&\longrightarrow \text{Clifford realization}.
\end{aligned}}
\label{eq:predictive-krein-chain}
\end{equation}
The Krein/Clifford layer therefore realizes an already determined signed
quadratic structure; it is not the device that selects its inertia.

\subsection{Clifford realization of the already determined Lorentzian form}

The matrix representation is now introduced only after
\Cref{prop:main-signature-minimality}.  A minimal depth-two realization uses
binary factors $A,B$ to represent the four Clifford axes.  The physical sign
content remains the single Read \eqref{eq:orientation-read-main}; the second
matrix factor is representation-theoretic multiplicity and is not a second
operational time-orientation choice.  Set
\begin{equation}
 \gamma^0=-iX_A\otimes I_B,
 \quad
 \gamma^1=Z_A\otimes Z_B,
 \quad
 \gamma^2=Z_A\otimes Y_B,
 \quad
 \gamma^3=Y_A\otimes I_B.
 \label{eq:clifford-main}
\end{equation}
They satisfy
$(\gamma^0)^2=-I$, $(\gamma^a)^2=I$, and
$\gamma^\mu\gamma^\nu+\gamma^\nu\gamma^\mu=0$ for $\mu\ne\nu$; together they
generate $M_4(\C)\cong\Cl_{1,3}(\C)$.  If $E_hE_h^{\mathsf T}=g_h^{-1}$, the
principal symbol
\begin{equation}
 \sigma_h(\xi)
 =\gamma^0N_h^{-1}(\xi_0-\beta_h\!\cdot\!\vec\xi)
  +\gamma^a(E_h)_a{}^i\xi_i
 \label{eq:dirac-symbol-main}
\end{equation}
has square
\begin{equation}
 \sigma_h(\xi)^2
 =\left[-N_h^{-2}(\xi_0-\beta_h\!\cdot\!\vec\xi)^2
        +g_h^{ij}\xi_i\xi_j\right]I.
 \label{eq:dirac-square-main}
\end{equation}
The Clifford cell therefore represents the Lorentzian form fixed in
\eqref{eq:main-signature-inertia}; it does not create its one-minus-three-plus
signature.

\begin{theorem}[Minimal faithful $3+1$ ADM reconstruction from either finite selector]
\label{mt:adm}
Let the minimum persistent predictive carrier of \Cref{mt:grand-tensor} admit a
geometric branch with one source and one target, physical reversal, faithful
endpoint data, and a local relation selected as $K_4$ by either
\Cref{prop:main-regular-graph-k4} or the alternating/minimal-source simplex
criterion of \Cref{prop:main-alternating-k4}.  Assume an occurring
scheduler-separated root-race calibration supplies directed winners, their
waiting times, and a positive physical speed density.  Assume also that the
protected orientation packet supplies the Read
\eqref{eq:orientation-read-main}, independent of the positive packet, with one
nontrivial signed direction.  Then:
\begin{enumerate}[label=\textup{(\roman*)},leftmargin=2.8em]
\item the selected local relation cell is $K_4$ and its spatial module is the
      $A_3$ root space $W_4\cong\R^3$;
\item the winner--waiting statistics of the scheduler-separated root race form
      a faithful coordinate system for all directed rates;
\item the moments \eqref{eq:bracket-drift-main} and the speed density invert
      uniquely to the ADM spatial metric, lapse, and shift through
      \eqref{eq:adm-inversion-main};
\item a nondegenerate continuum bracket requires an $O(h^{-2})$ fast scheduler;
\item positivity of the reconstructed spatial form together with exactly one
      independent signed direction fixes Lorentzian inertia $(1,3)$ up to
      overall sign, before any Clifford representation is chosen;
\item the Clifford cell \eqref{eq:clifford-main} faithfully realizes that
      already determined quadratic form, and the principal symbol has square
      \eqref{eq:dirac-square-main};
\item under summably correctable mixed-Gram and refinement defects, the
      dimension, source ranks, pair gap, ADM packet, and signed inertia persist
      on the canonical cofinal limit quotient.
\end{enumerate}
No persistent predictive state, metric, coframe, connection, curvature tensor,
or spacetime dimension is independently supplied.  Persistence is reconstructed
by \Cref{mt:grand-tensor}; the protected orientation and absolute
winner--waiting calibration remain irreducible physical input rows on this
branch.  The alternating form is required only on the second dimensional
selector and is not a prerequisite of the first.
\end{theorem}

The residual status of the last two rows is not merely a limitation of the
construction: \Cref{thm:supp-relative-primitive-floor} gives explicit
separating models showing that neither the protected orientation nor the
absolute calibration is functorially reconstructed from the remaining
positive operational data, even after common cofinal refinement.

The detailed proof is given in Appendix~\ref{supp:dimension-adm}.  The local algebraic audit is
finite: on the free $A_3$ accumulator, all equal-endpoint word relations are
generated by inverse digons, root triangles, and commutator squares of depth at
most four.  Global topology enters separately through a finite period bank.

% =============================================================================
\section{Renewal action reconstruction and the Einstein equation}
\label{sec:einstein-dynamics}
% =============================================================================

The ADM reconstruction of the preceding section is kinematical: it identifies
which Lorentzian geometry is encoded by a calibrated renewal packet.  We now
show that, once the corresponding physical history records actually occur, the
dynamical state, action, stationarity test, and spatial compactness screen can
be reconstructed or certified from finite renewal data.  The only residual
hypotheses in the final handoff are therefore genuine positive-branch and
continuum-control conditions rather than an independently supplied Einstein
action.

\subsection{The complete accepted cylinder determines the dynamical state}

At cutoff $X$, let $\widetilde\Theta_X$ denote the finite same-history record
carrying the reconstructed geometric variables used by the gravitational
sector, generic protected reward and duration coordinates, and the resolved
accepted-branch labels.  At this stage the reward coordinates are not named as
Palatini, Holst, Einstein--Hilbert, or cosmological terms; they are simply actual
finite Reads (or fixed functions of actual protected record coordinates) whose
variation content will be reconstructed and classified below.  Two values
$x,y\in\widetilde\Theta_X$ are equivalent
when every admitted future word and every declared determining coordinate has
the same value from $x$ and $y$, including branch-domain indicators.  Put
\begin{equation}
 \boxed{
 \Theta_{{\rm fld},X}^{\min}
 :=\widetilde\Theta_X/\!\sim_{{\rm fld},X}.}
 \label{eq:main-determining-field}
\end{equation}
For resolved accepted subbranches $a$, let $p_{a,X}(\theta)$ be their physical
probabilities and $u_a^{\min}$ their descended updates.  On the positive
acceptance domain define
\begin{equation}
 \boxed{
 K_X^{\rm acc}(\theta'\mid\theta)
 =\frac{1}{p_{{\rm acc},X}(\theta)}
   \sum_{\substack{a:\,u_a^{\min}(\theta)=\theta'}}p_{a,X}(\theta).}
 \label{eq:main-accepted-kernel}
\end{equation}

\begin{theorem}[Determining-field and accepted-kernel reconstruction]
\label{thm:main-determining-kernel}
The quotient \eqref{eq:main-determining-field} is the unique coarsest finite
record through which every declared determining-field future and every
primitive update factors.  Finite partition refinement terminates, and
\eqref{eq:main-accepted-kernel} is the unique table-native accepted transition
reproducing all selected post-accepted field futures.
\end{theorem}

Thus the accepted field dynamics is not selected after the fact once the
complete accepted cylinder is present.  What cannot be reconstructed from the
marginal accepted/not-accepted bit is the complete field cylinder itself; this
is a physical occurrence requirement rather than an implicit state-space
assumption.  The proof and the corresponding non-identifiability boundary are
given in Appendix~\ref{supp:dynamics}.

\subsection{Renewal pressure and reconstruction of a common action}

Let $(\Omega_X,p_X)$ be a finite faithful marked first-return law, let
$\tau_X:\Omega_X\to(0,\infty)$ be the protected physical duration, and let
$E_X$ be the finite variation space.  A protected cycle-reward synthesis is a
linear map
\begin{equation}
 \mathscr G_X:E_X\longrightarrow L^2(\Omega_X,p_X),
 \qquad v\longmapsto g_{v,X},
 \label{eq:main-reward-synthesis}
\end{equation}
whose values are actual Reads or fixed functions of protected record
coordinates.  Define
\begin{equation}
 \mathcal L_X(q,r)
 =\mathbb E_X\!\left[e^{-g_{q,X}-r\tau_X}\right].
 \label{eq:main-reward-laplace}
\end{equation}
Since $\partial_r\mathcal L_X(0,0)=-\mathbb E_X\tau_X<0$, there is a unique
analytic pressure $\mathcal P_X$ near the origin with
\begin{equation}
 \boxed{\mathcal P_X(0)=0,\qquad
 \mathcal L_X(q,\mathcal P_X(q))=1.}
 \label{eq:main-renewal-pressure}
\end{equation}
Let $S_Xv=D\log p_{X,q}|_{q=0}[v]$, let
$\pi_X=D\mathcal P_X(0)$, and let $T_Xc=c\tau_X$.

\begin{theorem}[Renewal-pressure action reconstruction]
\label{thm:main-action-reconstruction}
The marked renewal table reconstructs $\mathcal P_X$ and all of its finite
jets.  If $\bar\tau_X=\mathbb E_X\tau_X$, then
\begin{equation}
 \boxed{
 \pi_X(v)=-\frac{\mathbb E_X g_{v,X}}{\bar\tau_X},
 \qquad
 \mathscr G_X=-S_X-T_X\pi_X.}
 \label{eq:main-action-slope-reconstruction}
\end{equation}
Hence the centered probability score together with physical duration recovers
the complete uncentered action slope, including its affine anchor.  The pressure
Hessian and duration cross-jets reconstruct the full marginal and mixed action
Gram whenever all reward sectors occur on the same marked cycle record.
\end{theorem}

The remaining question is whether the reconstructed increments belong to one
scalar action.  Let $\mathscr V_X$ be the finite protected variation complex and
let $\alpha_X\in C^1(\mathscr V_X;\R)$ be the reconstructed action increment.

\begin{theorem}[Finite common-action exactness]
\label{thm:main-common-action}
There is a scalar common action $\mathscr S_X$ satisfying
\begin{equation}
 \boxed{\alpha_X=d_0\mathscr S_X}
 \label{eq:main-common-action}
\end{equation}
if and only if $d_1\alpha_X=0$ on every declared $2$-cell and the periods of
$\alpha_X$ vanish on one integral basis of $H_1(\mathscr V_X;\Z)$.  After one
additive anchor is fixed, $\mathscr S_X$ is unique and is reconstructed by path
integration.  Moreover, if $P_{\rm har}$ is the harmonic projection and
$\sigma_{\rm coex}>0$ is the least singular value of $d_1$ on the coexact
range, then
\begin{equation}
 \boxed{
 \operatorname{dist}(\alpha_X,\Ran d_0)^2
 \le \|P_{\rm har}\alpha_X\|^2
 +\sigma_{\rm coex}^{-2}\|d_1\alpha_X\|^2.}
 \label{eq:main-common-action-obstruction}
\end{equation}
Thus failure of a common action is itself a finite positive obstruction, not an
unseen modeling choice.
\end{theorem}

\begin{remark}[Non-circularity of the gravitational reward]
\label{rem:main-gravitational-reward-noncircular}
The synthesis $\mathscr G_X$ is fixed by the protected cycle record before it is
compared with the accepted kernel $K_X^{\rm acc}$.  Neither the transition
probabilities nor a desired continuum equation are used to fit a gravitational
cost after the fact.  Moreover, no Palatini or Holst decomposition is assumed at
this stage.  That decomposition appears only after derivative-free face
linearization, complete resolved-memory subtraction, and the Lorentz-equivariant
commutant classification below.  The subsequent Bregman/KL test is therefore an
independent question: it asks whether the actually accepted dynamics is
stationary for the previously reconstructed action.
\end{remark}

\subsection{Accepted transitions and common-action stationarity}

On a compact convex chart $\Theta\subset\R^d$ of the reconstructed determining
field, let $\Psi$ and the reconstructed common action $\mathcal A$ be $C^2$ on
an open neighbourhood of $\Theta$.  The operational potential is strictly
convex, with Fisher metric
\begin{equation}
 G(\theta)=\nabla^2\Psi(\theta),
 \qquad mI\preceq G(\theta)\preceq MI,
 \label{eq:main-fisher-metric}
\end{equation}
and, for $\lambda,L\ge0$ and $\eta>0$,
\begin{equation}
 -\lambda G(\theta)\preceq D^2\mathcal A(\theta)
 \preceq L G(\theta),\qquad \eta\lambda<1.
 \label{eq:main-action-hessian}
\end{equation}
For the Bregman divergence~\cite{Bregman1967}
\begin{equation}
 D_\Psi(y,\theta)=\Psi(y)-\Psi(\theta)
       -\langle\nabla\Psi(\theta),y-\theta\rangle,
 \label{eq:main-bregman}
\end{equation}
put
\begin{equation}
 \Phi_\theta(y)=\mathcal A(y)+\eta^{-1}D_\Psi(y,\theta),
 \qquad y_\theta^*=\arg\min_{y\in\Theta}\Phi_\theta(y).
 \label{eq:main-proximal-action}
\end{equation}
The minimum is unique.  On the branch used for unrestricted first variations,
$y_\theta^*$ lies in the interior of $\Theta$ for every supported source
$\theta$.  This interiority is important: a constrained minimum at a chart
boundary need not have zero ordinary action gradient.

For an actually accepted transition $\theta\to y$, define
\begin{equation}
 \boxed{\Delta_\Psi(\theta,y)
       =\Phi_\theta(y)-\Phi_\theta(y_\theta^*)\ge0.}
 \label{eq:main-action-gap}
\end{equation}
If $K=(1-\vartheta)I+\vartheta K^{\rm acc}$, with state-independent
$0<\vartheta\le1$, has stationary law $\nu$, set
\begin{equation}
 \begin{aligned}
 \overline\Delta_\Psi
  &=\int\nu(\dd\theta)K^{\rm acc}(\dd y\mid\theta)
                          \Delta_\Psi(\theta,y),\\
 \mathfrak S(\theta)
  &=\nabla\mathcal A(\theta)^*G(\theta)^{-1}\nabla\mathcal A(\theta).
 \end{aligned}
 \label{eq:main-stationarity-gap}
\end{equation}

\begin{theorem}[Accepted action gap controls interior stationarity]
\label{thm:main-action-stationarity}
Under the preceding hypotheses, including interiority of the supported
proximal minimizers, an explicit finite constant
$C_{\Psi,\mathcal A,\eta}$ satisfies
\begin{equation}
 \boxed{\int\mathfrak S(\theta)\,\nu(\dd\theta)
       \le C_{\Psi,\mathcal A,\eta}\,\overline\Delta_\Psi.}
 \label{eq:main-stationarity-bound}
\end{equation}
In particular, zero gap forces the represented first variation of the common
action to vanish $\nu$-almost surely.
\end{theorem}

The proof in Appendix~\ref{supp:dynamics} keeps the convexity and chart
constants explicit.  To pass to physical metric tests at varying cutoffs, let
$v_{X,K}(k;\theta)$ be the actual represented lift of a fixed determining test
$k$ on a compact cylinder $K$, and put
\begin{equation}
 B_{X,K}(k)^2:=\sup_{\theta\in\supp\nu_X}
 \|v_{X,K}(k;\theta)\|_{G_X(\theta)}^2.
 \label{eq:main-test-lift-budget}
\end{equation}
Duality then gives the quantitative physical-test estimate
\begin{equation}
 \boxed{\mathbb E_{\nu_X}
       |\delta\mathcal A_X[v_{X,K}(k)]|^2
       \le B_{X,K}(k)^2C_X\overline\Delta_{\Psi,X},}
 \qquad C_X=C_{\Psi_X,\mathcal A_X,\eta_X}.
 \label{eq:main-physical-stationarity}
\end{equation}
Thus the right-hand side, rather than the unweighted gap alone, is the
quantity required to vanish.  On a countable determining core this gives
convergence of the first variations in mean square and an almost-sure
vanishing subsequence under any common coupling; the deterministic continuum
handoff is applied to the resulting realized sequences.  The construction and
proof are given in \Cref{prop:supp-physical-test-stationarity}.

For a finite probability row with reference law $p_\theta(y)>0$ and
same-history action cost $c_\theta(y)$, the entropy-proximal row is
\begin{equation}
 q_\theta^*(y)=\frac{p_\theta(y)e^{-\eta c_\theta(y)}}
                    {\sum_zp_\theta(z)e^{-\eta c_\theta(z)}},
 \label{eq:main-gibbs-row}
\end{equation}
and its row variational gap is exactly
$\eta^{-1}D_{\rm KL}(K^{\rm acc}(\cdot\mid\theta)\Vert q_\theta^*)$.
This is an independent likelihood test against the reconstructed cost.
Equality of two finite-temperature probability rows does not by itself make
every sampled field configuration a critical point of $\mathcal A$; the
field-action estimate and its physical test lifts remain the stationarity
criterion used in the Einstein passage.

\subsection{Global noncollapse and compact spatial screens}

Let the physical endpoint quotient of a global regulator carry masses
$m_{v,X}$, symmetric spatial conductances $c_{e,X}$, physical cell scale
$h_X$, and provenance-safe efficient capacities $u_{e,X}$.  The Dirichlet
form and its nonnegative self-adjoint Laplacian are
\begin{equation}
 \mathcal E_X^{\rm sp}(f)
 =\sum_{\{u,v\}\in E_X}c_{uv,X}|f(u)-f(v)|^2,
 \qquad
 \mathcal E_X^{\rm sp}(f)=\langle f,\Delta_X^{\rm sp}f\rangle_{L^2(\mu_X)}.
 \label{eq:main-spatial-dirichlet}
\end{equation}
The endpoint/action packet also contains a finite capacity--conductance
comparison.  On the positive comparison branch there are cutoff-independent
constants $0<a_-\le a_+<\infty$ such that
\begin{equation}
 \boxed{
 a_-h_Xc_{e,X}\le u_{e,X}\le a_+h_Xc_{e,X}
 \qquad(e\in E_X).}
 \label{eq:main-capacity-conductance}
\end{equation}
This is an explicit same-history comparison between two reconstructed physical
rows; it is not an identification made after taking the continuum limit.
Define the efficient and analytic cut constants
\begin{equation}
 \boxed{
 I_X^{\rm eff}:=
 \min_{0<\mu_X(A)\le\mu_X(V_X)/2}
 \frac{u_X(\partial A)}{\mu_X(A)^{2/3}},
 \qquad
 I_X^{\rm sp}:=
 \min_{0<\mu_X(A)\le\mu_X(V_X)/2}
 \frac{h_Xc_X(\partial A)}{\mu_X(A)^{2/3}}.}
 \label{eq:main-cut-constant}
\end{equation}
Then
\begin{equation}
 a_- I_X^{\rm sp}\le I_X^{\rm eff}\le a_+I_X^{\rm sp}.
 \label{eq:main-cut-comparison}
\end{equation}
Thus a positive efficient-capacity margin is exactly the operational input
needed to obtain a positive spatial isoperimetric margin on this branch.

For a nonempty half-volume cut $A$, write $A^c=V_X\setminus A$ and define
the mean-zero demand
\begin{equation}
 b_{A,X}(v)=\mu_X(A)^{2/3}
 \left[\frac{m_{v,X}\one_A(v)}{\mu_X(A)}
       -\frac{m_{v,X}\one_{A^c}(v)}{\mu_X(A^c)}\right].
 \label{eq:main-cut-demand}
\end{equation}
This fixes the normalization of the dual cut witness below.

\begin{theorem}[Cut margin and compact-screen alternative]
\label{thm:main-spatial-screen}
Assume the total physical volume is uniformly bounded,
$\mu_X(V_X)\le V_*$, the weighted degree satisfies
$\sum_uc_{uv,X}\le D_*h_X^{-2}m_{v,X}$, and
\eqref{eq:main-capacity-conductance} holds.  If
$\inf_X I_X^{\rm eff}=I_*^{\rm eff}>0$, then
$I_X^{\rm sp}\ge I_*^{\rm eff}/a_+=:I_*^{\rm sp}>0$ and every mean-zero $f$
obeys regulator-uniform Sobolev and Poincar\'e bounds
\begin{equation}
 \|f\|_{L^6(\mu_X)}^2\le C_{\rm S}\mathcal E_X^{\rm sp}(f),
 \qquad
 \|f\|_{L^2(\mu_X)}^2\le C_{\rm P}\mathcal E_X^{\rm sp}(f),
 \label{eq:main-sobolev-poincare}
\end{equation}
with
\begin{equation}
 C_{\rm S}=\frac{128D_*}{(I_*^{\rm sp})^2},
 \qquad
 C_{\rm P}=V_*^{2/3}C_{\rm S}.
 \label{eq:main-sobolev-constants}
\end{equation}
The low-energy counting function satisfies
\begin{equation}
 \boxed{
 N_X(R)\le1+4e^{3/2}V_*(C_{\rm S}R)^{3/2}}
 \qquad(R>0),
 \label{eq:main-weyl-count}
\end{equation}
and for $P_{X,R}=\one_{[0,R]}(\Delta_X^{\rm sp})$,
\begin{equation}
 \boxed{
 \|(I-P_{X,R})f\|_{L^2(\mu_X)}^2
 \le R^{-1}\mathcal E_X^{\rm sp}(f).}
 \label{eq:main-screen-tail}
\end{equation}
Hence the spectral projections form uniformly exhausted finite-rank compact
screens.  If instead $I_{X_j}^{\rm eff}>0$ tends to zero, the minimizing
cuts $A_j$ give normalized potentials $\phi_j$ with
\begin{equation}
 \sum_eu_{e,X_j}|\dd\phi_j(e)|=1,
 \qquad
 \left|\sum_v b_{A_j,X_j}(v)\phi_j(v)\right|
 =(I_{X_j}^{\rm eff})^{-1}\longrightarrow\infty,
 \label{eq:main-neck-obstruction}
\end{equation}
which is an explicit macroscopic-neck obstruction.  If a cut has zero
capacity already at finite cutoff, that disconnected-support cut is retained
directly, without dividing by its capacity.
\end{theorem}

The discrete Sobolev--Nash--heat-kernel mechanism behind this theorem is
standard in spirit; see, for example, Chung~\cite[Chs.~10--11]{Chung1997} and
Saloff-Coste~\cite{SaloffCoste2002}.  Because the regulator normalization is
part of the present result, Appendix~\ref{supp:dynamics} gives the complete
finite proof, including the constants, the $L^1\!\to L^\infty$ heat-semigroup
bound, and the derivation of the $R^{3/2}$ counting law.

Thus compact spatial control is derived on the positive cut branch and is
replaced by a concrete obstruction on the collapsed branch.  Uniform spatial
ranks and tails alone do not establish spacetime compactness.  The
common-cylinder Hodge/action packet also retains convergence of transported
finite screens and temporal control of the field coefficients.
Appendix~\ref{supp:dynamics} gives the compact-screen upgrade from this
packet to the strong convergence needed for quadratic Palatini insertions.

\subsection{Lorentz-natural classification of the gravitational score}

Let $e^I$ be the reconstructed coframe, $\omega^{IJ}$ a metric-compatible
reconstructed Cartan connection, $B^{IJ}=e^I\wedge e^J$, and $T^I=De^I$.
Before any continuum gravitational Lagrangian is named, the generic protected
face reward is subjected to two finite operations described in
Appendix~\ref{supp:dynamics}: binary renewal subdivision, which produces a unique
first-order linear functional on its exact or summably stable branch, and
subtraction of the complete resolved cut-memory log-partition.  Denote the
resulting generic first-order gravitational score by $A^{\rm grav}_{X,\rm sh}$.
For later reference, write the three Lorentz-natural first-order densities (up
to the overall normalization convention) as
\begin{align}
 L_{\rm Palatini}
 &=\frac12\epsilon_{IJKL}e^I\wedge e^J\wedge R^{KL},\\
 L_{\rm Holst}
 &=e^I\wedge e^J\wedge R_{IJ},\\
 L_{\rm vol}
 &=\frac1{4!}\epsilon_{IJKL}
 e^I\wedge e^J\wedge e^K\wedge e^L.
 \label{eq:main-phv-densities}
\end{align}

\begin{theorem}[Lorentz-natural classification of the generic gravitational score]
\label{thm:main-gravitational-classification}
Assume $A^{\rm grav}_{X,\rm sh}$ is real and Lorentz natural on the represented
coframe-bivector/curvature variables, its residual typed commutant on the
bivector carrier is scalar, and determinant coframe volume is represented.
Then, up to coefficient normalization, its only first-order Lorentz-natural
bulk terms are
\begin{equation}
 \boxed{
 A^{\rm grav}_{X,\rm sh}
 \in\Span_{\R}\{L_{\rm Holst},L_{\rm Palatini},L_{\rm vol}\}.}
 \label{eq:main-phv-classification}
\end{equation}
Equivalently, there are real coefficients $\alpha_X,\beta_X,\lambda_X$ such that
\begin{equation}
 A^{\rm grav}_{X,\rm sh}
 =\alpha_XL_{\rm Holst}+\beta_XL_{\rm Palatini}+\lambda_XL_{\rm vol}.
 \label{eq:main-phv-coefficients}
\end{equation}
No Palatini, Holst, or Einstein action is part of the hypothesis.
\end{theorem}

The classification follows because both the coframe bivector and curvature
transform in $\Lambda^2\R^{1,3}$ and
$\operatorname{End}_{SO^+(1,3)}(\Lambda^2\R^{1,3})
=\Span_{\R}\{I,\star\}$; the scalar residual commutant excludes any additional
internal tensor, and determinant volume is the only independent alternating
four-coframe scalar at this order.  The full proof is in
Appendix~\ref{supp:dynamics}.

\subsection{Initial-energy propagation and identified curvature}

The remaining analytic issue is curvature compactness.  Rather than assume a
spacetime-wide curvature bound as an independent hypothesis of the Einstein
handoff, we use a positive curvature-energy propagation certificate.  This is
the finite analogue of the standard Einstein--Bianchi energy method; compare
Anderson--Choquet-Bruhat--York~\cite{ACY1998} and the bounded $L^2$ curvature
framework of Klainerman--Rodnianski--Szeftel~\cite{KRS2015}.

On a compact foliated cylinder $K=[0,T_*]\times\Sigma$, let
$y_X(t)=(E_X(t),B_X(t))$ be the actual retained electric/magnetic curvature
coefficients in the reconstructed observer frame, with positive physical Gram
$M_X(t)=\operatorname{diag}(M_{E,X},M_{B,X})\succ0$.  Put
\begin{equation}
 z_X(t)^2:=y_X(t)^TM_X(t)y_X(t).
 \label{eq:main-curvature-energy}
\end{equation}
The finite Bianchi/transport propagation writer is required to occur in the
form
\begin{equation}
 \dot y_X=(K_X+L_X)y_X+f_X,
 \qquad M_XK_X+K_X^TM_X=0,
 \label{eq:main-curvature-propagation-writer}
\end{equation}
where $K_X$ is the positive-Gram skew principal curl block, $L_X$ contains
lower-order observer, lapse, connection, and transport terms, and $f_X$ is the
complete differentiated curvature source, including any incoming boundary-energy
contribution when the cylinder has a timelike spatial boundary.  Define
\begin{equation}
 a_X(t):=\frac12\max\left\{0,
 \lambda_{\max}\!\left[
 M_X^{-1/2}(\dot M_X+M_XL_X+L_X^TM_X)M_X^{-1/2}
 \right]\right\}.
 \label{eq:main-curvature-growth}
\end{equation}
On the gravitational shared-link branch, $f_X$ contains the physically
transferred, volume-normalized Palatini metric row together with the
differentiated Cartan and finite Bianchi defects; Appendix~\ref{supp:dynamics}
keeps these terms separate.

The principal cancellation is tested on the actual record.  For example, a
recorded curl $C_X:V_{B,X}\to V_{E,X}$ has positive-Gram adjoint
$C_X^\dagger=M_{B,X}^{-1}C_X^TM_{E,X}$, and the block
$\left(\begin{smallmatrix}0&C_X\\-C_X^\dagger&0\end{smallmatrix}\right)$
is exactly $M_X$-skew.  Defining this adjoint does not prove that an actual
geometric second curl equals $-C_X^\dagger$: their discrepancy is retained
in the source or symmetric-growth record.  Consequently no hidden
$h_X^{-1}$ principal norm is discarded in the estimate.

\begin{theorem}[Initial-energy curvature propagation certificate]
\label{thm:main-curvature-propagation}
Assume \eqref{eq:main-curvature-propagation-writer}.  Then
\begin{equation}
 \boxed{
 z_X(t)\le e^{\int_0^t a_X}
 \left[z_X(0)+\int_0^t e^{-\int_0^s a_X}
 \|f_X(s)\|_{M_X(s)}\,\dd s\right].}
 \label{eq:main-curvature-energy-propagation}
\end{equation}
Suppose moreover that the actual curvature field entering the Palatini test has
a uniformly stable interpolation
\begin{equation}
 R_X=\mathcal I_Xy_X+\zeta_X,
 \qquad
 \|\mathcal I_Xy_X(t)\|_{L^2(\Sigma_t)}\le C_Iz_X(t),
 \qquad
 \|\zeta_X\|_{L^2(K)}\le C_\zeta,
 \label{eq:main-curvature-incidence}
\end{equation}
and that the lapse is bounded by $N_*$.  If, uniformly in $X$,
\begin{equation}
 z_X(0)\le Z_*,\qquad
 \int_0^{T_*}a_X(t)\,\dd t\le A_*,\qquad
 \int_0^{T_*}\|f_X(t)\|_{M_X(t)}\,\dd t\le F_*,
 \label{eq:main-curvature-budgets}
\end{equation}
then
\begin{equation}
 \boxed{
 \|R_X\|_{L^2(K)}
 \le C_I\sqrt{N_*T_*}\,e^{A_*}(Z_*+F_*)+C_\zeta.}
 \label{eq:main-derived-l2-curvature}
\end{equation}
Thus the curvature-size requirement needed below follows with $p=2>3/2$.
\end{theorem}

\paragraph{Propagation on a finite time alphabet.}
The preceding estimate has a fully discrete counterpart, so a differentiable
curvature path is not needed merely to propagate its size.  At one spatial
cutoff, take a partition of the selected time interval with lengths $s_j>0$.
Let $Y_j=M_j^{1/2}y_j$ be the normalized actual curvature records and
let $T_j$ be their recorded physical transfer in these coordinates.  Suppose
\begin{equation}
 Y_{j+1}-T_jY_j
 =s_j(K_j+L_j)\frac{Y_{j+1}+T_jY_j}{2}+\varepsilon_j,
 \label{eq:main-discrete-curvature-update}
\end{equation}
where $K_j^T=-K_j$, $(L_j+L_j^T)/2\preceq a_jI$,
$\|T_j\|_{\op}\le e^{d_j}$, and $a_j,d_j\ge0$.  The residual
$\varepsilon_j$ includes the complete source, solver and coordinate-transfer
errors of the actual step.  In particular, nonlinear changes of field
coordinates are not assumed to commute with midpoint discretization.

\begin{theorem}[Fully discrete initial-energy curvature certificate]
\label{thm:main-discrete-curvature}
Assume $s_ja_j/2\le b<1$, and set
\begin{equation}
 D_X=\sum_jd_j,\quad A_X^{\rm d}=\sum_js_ja_j,\quad
 \mathcal D_X=\sum_j\frac{\|\varepsilon_j\|^2}{s_j},\quad
 T_X=\sum_js_j.
 \label{eq:main-discrete-curvature-budgets}
\end{equation}
Each recorded linear step is uniquely solvable, and
\begin{equation}
 \boxed{\max_j\|Y_j\|
 \le e^{D_X+A_X^{\rm d}/(1-b)}
       \left(\|Y_0\|+\frac{\sqrt{T_X\mathcal D_X}}{1-b}\right).}
 \label{eq:main-discrete-curvature-bound}
\end{equation}
There is no restriction involving $s_j\|K_j\|_{\op}$.  With a uniformly
stable interval interpolation of these actual curvature records, a uniform
lapse bound, and a bounded reconstruction error as in
\eqref{eq:main-curvature-incidence}, uniform bounds on the quantities in
\eqref{eq:main-discrete-curvature-bound} imply
$\sup_X\|R_X\|_{L^2(K)}<\infty$.
\end{theorem}

The statement concerns the represented linear step with its coefficients
fixed.  For an implicit nonlinear evolution, existence of that evolution is
not inferred from this matrix inverse; its actual coefficient and residual
records must satisfy \eqref{eq:main-discrete-curvature-update}.
The proof in Appendix~\ref{supp:dynamics} uses the skew cancellation before
taking norms and retains the physical transfer between successive Grams.

\paragraph{The source norm is determined by the action-current transfer.}
The forcing budget can itself be checked against the reconstructed action,
provided the actual differentiated current has the required incidence.  Write
$\mathfrak a_X=D\mathcal A_X$ for its action-gradient covector and suppose
\begin{equation}
 f_X=\mathsf T_X\mathfrak a_X+r_X^{\rm src}.
 \label{eq:main-current-incidence}
\end{equation}
The transfer $\mathsf T_X$ includes the physical metric-test lift, the nonzero
Palatini normalization, trace reversal, the differentiated Bianchi current,
and its frame comparison.  Torsion, Cartan, boundary and discretization
contributions not in this map remain in $r_X^{\rm src}$.

\begin{proposition}[Sharp differentiated action-current certificate]
\label{prop:main-action-current}
Let $H_X\succ0$ be the actual source Gram, transferred to the curvature-energy
space.  The least squared amplification from the action dual norm is
\begin{equation}
 \boxed{\Gamma_X
   =\|H_X^{1/2}\mathsf T_XG_X^{1/2}\|_{\op}^2,\qquad
 \|\mathsf T_X\mathfrak a\|_{H_X}^2
       \le\Gamma_X\|\mathfrak a\|_{G_X^{-1}}^2.}
 \label{eq:main-sharp-current-transfer}
\end{equation}
Consequently a pointwise certificate
$\|\mathfrak a_X(t)\|_{G_X^{-1}}^2\le C_X(t)\Delta_X(t)$ gives
\begin{equation}
 \int_0^{T_*}\|f_X\|_{H_X}\,\dd t
 \le\sqrt{2T_*\int_0^{T_*}
       \bigl(\Gamma_XC_X\Delta_X+
                         \|r_X^{\rm src}\|_{H_X}^2\bigr)\,\dd t}.
 \label{eq:main-action-source-budget}
\end{equation}
For the normalized discrete source
$\varepsilon_j/s_j=\mathsf T_j\mathfrak a_j+r_j^{\rm src}$, the
corresponding squared-energy estimate is
$\mathcal D_X\le2\sum_js_j(\Gamma_jC_j\Delta_j+
\|r_j^{\rm src}\|_{H_j}^2)$.
\end{proposition}

This replaces a separately supplied source-size bound by an explicitly scaled
action-side bound on the verified transfer branch.  It does not turn a small
undifferentiated Euler residual into a small Bianchi current: a single centered
grid derivative can have $\Gamma_X=h_X^{-2}$.  The sharpness example and the
expected-energy version of the certificate are proved in
Appendix~\ref{supp:dynamics}.  An estimate only in expectation yields a bound
in probability; almost-sure uniform curvature control additionally uses the
pathwise or summable source conditions stated there.

\paragraph{Identification with the actual connection curvature.}
Curvature size and curvature identification are separate questions.  Let
$\mathcal T_{X,K}$ be the finite determining curvature-test bank, and let
$\mathcal R_X(\omega_X)$ denote the finite Cartan writer on the common-cylinder
interpolation.  Define
\begin{equation}
 \boxed{\mathfrak C_X(K):=
 \sup_{\substack{\Phi\in\mathcal T_{X,K}\\\|\Phi\|_{L^2(K)}\le1}}
 \left|\int_K\!\langle R_X-\mathcal R_X(\omega_X),\Phi\rangle\right|.}
 \label{eq:main-cartan-curvature-residual}
\end{equation}
The banks exhaust a fixed countable determining smooth core, with compatible
test restrictions and vanishing test-interpolation errors.  Their geometric
consistency is recorded separately by
\begin{equation}
 \mathfrak I_X(\Phi):=
 \int_K\!\langle\mathcal R_X(\omega_X),\Phi\rangle
 -\langle\dd_{\rm dist}\omega_X+\omega_X\wedge\omega_X,\Phi\rangle.
 \label{eq:main-cartan-incidence}
\end{equation}
The distributional derivative includes every interface contribution of a
piecewise connection.  Thus matching two finite writers alone is not used to
identify the nonlinear curvature of a weakly convergent connection.

\begin{theorem}[Identified curvature compactness after propagation]
\label{thm:main-curvature-compactness}
Suppose either \Cref{thm:main-curvature-propagation} or
\Cref{thm:main-discrete-curvature} supplies the uniform curvature bound on
every compact $K$.  Assume also
\begin{equation}
 \omega_X\to\omega\ \hbox{strongly in }L^2_{\rm loc},\qquad
 \mathfrak C_X(K)\to0,\qquad
 \mathfrak I_X(\Phi)\to0
 \label{eq:main-identified-curvature-hypotheses}
\end{equation}
for the determining smooth core.  Then every cofinal sequence has a subsequence
such that
\begin{equation}
 \boxed{R_X\rightharpoonup R\ \hbox{in }L^2(K),\qquad
 R=\dd\omega+\omega\wedge\omega\ \hbox{distributionally}.}
 \label{eq:main-curvature-weak}
\end{equation}
A common subsequence works on a countable compact exhaustion.
\end{theorem}

For discontinuous connection reconstructions, the identification has a concrete
sufficient realization.  On a uniformly regular cubical chart put
$J_f=n_f^\flat\wedge[\omega_h]_f$ at each interior face and
\begin{equation}
 \mathcal J_h^2=\sum_fh^{-1}\|J_f\|_{L^2(f)}^2.
 \label{eq:main-interface-budget}
\end{equation}
A bounded-overlap volume lifting $\mathcal L_hJ$ of the face measures gives
\begin{equation}
 \begin{aligned}
 R_h^{\rm lift}
   &=\dd_b\omega_h+\omega_h\wedge\omega_h+\mathcal L_hJ,\\
 \|R_h^{\rm lift}\|_{L^2(K)}
   &\le\|\dd_b\omega_h\|_2+C\|\omega_h\|_4^2+C\mathcal J_h.
 \end{aligned}
 \label{eq:main-lifted-curvature}
\end{equation}
Its discrepancy from the full distributional curvature is at most
$C_Kh\mathcal J_h\|\nabla\Phi\|_\infty$ on a fixed smooth test.  The
construction and proof are given in \Cref{thm:supp-interface-curvature}.
These are explicit additional interface budgets, not consequences of a small
connection Euler residual.  The lifting is used only when the action's actual
curvature writer agrees with it on the determining tests, up to a vanishing
recorded discrepancy.

\begin{remark}[Scope of the propagated curvature estimate]
\label{rem:main-curvature-not-circular}
The curvature bounds are obtained before the limiting Einstein equation is
invoked.  They replace an all-time $L^p$ curvature assumption by initial
energy, finite propagation and transfer budgets, and a separately identified
connection limit.  The action-current certificate can discharge the source
budget when its scaled errors are controlled; it does not establish an
unrestricted nonlinear evolution bound for the remaining coefficients.
Flat and de~Sitter regulators satisfy the required bounds directly.  In the
nonlinear Gowdy class below, the actual full-curvature bound and its
identification instead follow pathwise from the second-jet estimate on the
finite history alphabet.
\end{remark}


\paragraph{Two theorem levels.}

The curvature and compactness machinery above is needed for the most general
renewal-to-Einstein handoff, but the Einstein branch is not established only
by postulating its eventual certification conditions.  We first give a
direct nonlinear population theorem.  For one explicitly specified local
renewal writer, cutoff-uniform energy, consistency, subsidiary propagation,
and finite-word realization are proved on a fixed Sobolev neighbourhood.
The Einstein limit in this subsection is therefore obtained without taking
\textup{(E1)--(E5)} below as independent hypotheses.

\subsection{A constructive open nonsymmetric \texorpdfstring{$3+1$}{3+1} Einstein population}
\label{subsec:main-open-3plus1}

The general certification theorem stated below deliberately separates reconstruction from
dynamical propagation.  We now exhibit a concrete population on which the
continuum conclusion can be proved directly from a cutoff-uniform nonlinear
energy, without assuming the five handoff certificates as independent
analytic hypotheses.  The construction uses the same reconstructed
Lorentzian coefficient packet but a declared local harmonic writer.  It is
therefore a positive population theorem, not a claim that every accepted
kernel or the raw finite Euler flow of the reconstructed action must choose
this writer.

Let $N=2m+1$, $h=N^{-1}$, and
$\Lambda_h=(h\mathbb Z/\mathbb Z)^3$.  With shifts $S_i$ set
\begin{equation}
 D_i^+=(S_i-I)/h,\qquad D_i^-=(I-S_i^{-1})/h,\qquad
 D_i^0=(D_i^++D_i^-)/2.
 \label{eq:main-open-differences}
\end{equation}
For $g=\eta+q$ in a small Lorentzian chart, write the harmonic normal row as
\begin{equation}
 a(g)q_{tt}+2b^i(g)\partial_iq_t-c^{ij}(g)\partial_i\partial_jq
             -F(g,\partial g)=0,
 \qquad a=-g^{00},\ b^i=-g^{0i},\ c^{ij}=g^{ij},
 \label{eq:main-harmonic-normal-row}
\end{equation}
where $F$ is the analytic first-jet source obtained from the reduced vacuum
Einstein tensor.  Put
\begin{equation}
 \mathsf K_b=\sum_i\{M_{b^i}D_i^0+D_i^0M_{b^i}\},
 \label{eq:main-open-skew-transport}
\end{equation}
and, with $Y_h=(v,D_1^+q,D_2^+q,D_3^+q)$,
\begin{equation}
 \mathsf G_h(q,v)=F(g,Y_h)
 -\sum_{i,j}Dc^{ij}(g)[D_i^+q]D_j^+q
 +\sum_iDb^i(g)[D_i^+q]v.
 \label{eq:main-open-compensator}
\end{equation}
The local evolution law is
\begin{equation}
 \boxed{
 q_t=v,\qquad
 a(g)v_t=\sum_{i,j}D_i^-\!\bigl(c^{ij}(g)D_j^+q\bigr)
           -\mathsf K_bv+\mathsf G_h(q,v).}
 \label{eq:main-open-writer}
\end{equation}
The divergence and skew placements are fixed by the positive quadratic
principal energy.  The compensator removes exactly the first-derivative terms
introduced by these placements, so the smooth continuum limit is
\eqref{eq:main-harmonic-normal-row}.  A parity obstruction of the bare
$A_3\cong D_3$ execution graph requires one additional odd-parity execution
incidence for a cutoff-uniform Cartesian Sobolev realization; this incidence
is used only to execute the writer and is not added to the root race entering
\eqref{eq:bracket-drift-main}.  The complete proof is in
Appendix~\ref{supp:open-3plus1}.  In the main theorem we use the discrete
Sobolev norms
\begin{equation}
 \|u\|_{H_h^r}^2=\sum_{|\alpha|\le r}\|D^\alpha u\|_h^2,
 \qquad
 \|X\|_{\mathcal X_h^r}^2=\|q\|_{H_h^{r+1}}^2+\|v\|_{H_h^r}^2,
 \label{eq:main-open-sobolev-norms}
\end{equation}
where $D^\alpha$ is formed from the commuting forward differences.

For an integer $s\ge11$ define the compatible continuum data class
\begin{equation}
 \begin{split}
 \mathcal D_s(\varepsilon)=\bigl\{(\gamma,K):\;&
 \gamma\in H^{s+1}(\mathbb T^3),\ K\in H^s(\mathbb T^3),\\
 &R(\gamma)+(\tr_\gamma K)^2-|K|_\gamma^2=0,\qquad
 \nabla_\gamma^iK_{ij}-\partial_j\tr_\gamma K=0,\\
 &\|\gamma-I\|_{H^{s+1}}+\|K\|_{H^s}<\varepsilon\bigr\}.
 \end{split}
 \label{eq:main-open-data-class}
\end{equation}
This is relative openness in the vacuum constraint set; it is not an ambient
open ball of fixed nonzero constraint violations.

\begin{theorem}[Open nonsymmetric $3+1$ renewal--Einstein population]
\label{thm:main-open-3plus1}
There are $\varepsilon_s,T_s,h_s,C_s>0$, independent of the cutoff, with the
following properties.
\begin{enumerate}[label=\textup{(O\arabic*)},leftmargin=3em]
\item For every $h<h_s$, every real ten-component initial record
      $X_h(0)=(q_h(0),v_h(0))$ in the fixed ball
      $\|q_h(0)\|_{H_h^{s+1}}+\|v_h(0)\|_{H_h^s}\le\varepsilon_s$
      has a unique solution of \eqref{eq:main-open-writer} on $[0,T_s]$, and
      \begin{equation}
       \sup_{t\le T_s}
       \bigl(\|q_h(t)\|_{H_h^{s+1}}+\|v_h(t)\|_{H_h^s}\bigr)
       \le C_s\|X_h(0)\|_{\mathcal X_h^s}.
       \label{eq:main-open-uniform-time}
      \end{equation}
      No symmetry, Fourier-support restriction, analyticity, or
      cutoff-shrinking amplitude is imposed.
\item Every sufficiently small $(\gamma,K)\in\mathcal D_s(\varepsilon_s)$
      has a harmonic finite preparation for which the whole cutoff sequence
      converges on $[0,T_s]\times\mathbb T^3$ to a unique metric $g$ with the
      prescribed initial data.  Writing
      $\varepsilon_0=\|\gamma-I\|_{H^{s+1}}+\|K\|_{H^s}$, one has
      \begin{equation}
       \boxed{G(g)=0,\qquad
       \sup_{t\le T_s}\left(
       \|g_h-g\|_{H^{s-1}}+\|\partial_tg_h-\partial_tg\|_{H^{s-2}}
       +\|\partial_t^2g_h-\partial_t^2g\|_{H^{s-3}}
       \right)\le C_sh\varepsilon_0.}
       \label{eq:main-open-einstein-rate}
      \end{equation}
      Here $g_h$ denotes the canonical interpolation of $\eta+q_h$.
\item The harmonic subsidiary field and Einstein residual obey, at the
      corresponding lower orders,
      \begin{equation}
       \sup_{t\le T_s}
       \bigl(\|c(g_h)\|_{H^{s-2}}+\|\partial_tc(g_h)\|_{H^{s-3}}
                  +\|G(g_h)\|_{H^{s-3}}\bigr)
       \le C_sh\varepsilon_s,
       \label{eq:main-open-subsidiary}
      \end{equation}
      while the positive $A_3$ rate lift remains in an interior cone and its
      spatial cut margin stays uniformly positive.  In particular the
      reconstructed metric is uniformly noncollapsed and its metric curvature
      is uniformly bounded on the common slab and converges to the curvature
      of $g$ at the lower Sobolev order supplied by
      \eqref{eq:main-open-einstein-rate}.
\item The differential writer has a genuinely finite renewal realization.
      With segment duration $\tau_h\asymp h^2$, fixed-radius local residual
      corrections and finite precision can produce committed successor words
      whose complete acceleration defect is $O(\varepsilon h^4)$ and whose
      differentiated defect is $O(\varepsilon h^2)$ on the common slab.
      Every such certified history inherits the bounds above.  This statement
      concerns finite-range information flow; it does not identify the
      arithmetic service time with the physical root clock.
\end{enumerate}
Consequently the positive Einstein population is not confined to homogeneous
or one-dimensional inhomogeneous reductions: it contains a relatively open
small-data class with arbitrary three-dimensional dependence.  The next
result shows that the differential writer itself is also not isolated.
\end{theorem}

\begin{remark}[Relation to the certification theorem and to exact-action flow]
\label{rem:main-open-3plus1-scope}
Theorem~\ref{thm:main-open-3plus1} removes the symmetry-reduced and
cutoff-shrinking-lifespan limitations for the declared writer, but it is not a
universal attraction theorem.  Its evolution law is an energy-compatible
local realization of the harmonic Palatini envelope; it is not asserted to
be the exact finite Euler flow of every reconstructed action cylinder.
Appendix~\ref{supp:exact-action-audit} studies that stronger identification
separately.  On a regular finite exact-action branch the constraint quotient
need not already have the local first-class structure of continuum general
relativity.  Nevertheless a large transverse multiplier block is exactly
regular at every odd cutoff, the complete signed response admits an exact
positive-reproduction plus scalar--trace decomposition, and selected original-
action preparations eliminate the leading weak target while retaining
nontrivial three-dimensional data.  On a further prepared branch the leading
active source is trace-free and uniformly rank-controlled, so the $2/3$ trace
resonance is absent there at leading order; the actual transverse geometric
response is correspondingly suppressed.  What remains is dynamical: the
prepared exact-action manifold must be tangent to its own flow, the resulting
higher signed system must remain regular under refinement, and the estimates
must persist on a common physical interval.  None of these stronger
exact-action conditions is assumed in \Cref{thm:main-open-3plus1}.
\end{remark}




\subsection{An open renewal-law basin and action-prepared finite seeds}
\label{subsec:main-open-law-basin}

The fixed writer of \Cref{thm:main-open-3plus1} is the centre of a larger
finite law class.  Let
\begin{equation}
 \Lambda_h=-\sum_{i=1}^3D_i^-D_i^+,
 \qquad
 \mathcal B=\{B=B^T\in\mathbb R^{10\times10}:\|B\|_{\op}<b_0\},
 \label{eq:main-law-ball}
\end{equation}
where $b_0>0$ is fixed below the spatial coercivity margin of the common small
metric chart.  Thus $\mathcal B$ is a $55$-dimensional open ball.  For
$B\in\mathcal B$ define
\begin{equation}
 \boxed{
 q_t=v,\qquad
 a(g)v_t=\sum_{i,j}D_i^-\!\bigl(c^{ij}(g)D_j^+q\bigr)
       -\mathsf K_bv+\mathsf G_h(q,v)-h^2B\Lambda_h^2q.}
 \label{eq:main-open-law-family}
\end{equation}
The added fourth-order stencil has fixed execution radius.  It need not be
small in the top ultraviolet vector-field norm; its incidence and sign are
instead controlled by the same positive shifted energy.

For a finite chronological realization of \eqref{eq:main-open-law-family},
let $I_j$ have duration $\tau_j\asymp h^2$.  A candidate physical letter $e$
at a live prefix $\pi$ carries a piecewise-Hermite segment $Q_e$ with actual
acceleration defect
\begin{equation}
 f_e=Q_e''-V_{B,h}(Q_e,Q_e'),
 \label{eq:main-law-residual}
\end{equation}
and, with unit cost assigned to a grammar or execution failure, define
\begin{equation}
 c_{h,j}(\pi,e)=\varepsilon^{-2}
 \int_{I_j}\bigl(\|f_e\|_{H_h^s}^2
       +\|\partial_tf_e\|_{H_h^{s-3}}^2\bigr)\,\dd t,
 \qquad C_h=\sum_jc_{h,j}.
 \label{eq:main-law-cost}
\end{equation}
The cost is determined by the physical letter and its prefix, not by its
probability.

Finally let $\mathscr H_h(\gamma,\pi,N,\beta)$ be the canonical Hamiltonian
obtained from the reconstructed phase-compatible action after stationary
connection and Legendre elimination, and set
\begin{equation}
 \boxed{\mathscr C_h(X)=
 \left(-\nabla_N\mathscr H_h,\ \nabla_\beta\mathscr H_h\right)
       \big|_{N=1,\,\beta=0}.}
 \label{eq:main-action-initial-constraints}
\end{equation}
This is the original four-row finite initial constraint map; it is not the
residual of the local evolution writer.

\begin{theorem}[Open renewal-law basin with action-prepared finite seeds]
\label{thm:main-open-law-basin}
After decreasing the common small-data radius if necessary, the following
statements hold with constants independent of the cutoff.
\begin{enumerate}[label=\textup{(L\arabic*)},leftmargin=3em]
\item For every $B\in\overline{\mathcal B}$ the law
      \eqref{eq:main-open-law-family} has the same fixed-radius nonlinear
      lifespan as \Cref{thm:main-open-3plus1}.  There is an energy
      $\mathscr E^B_{s,h}$ uniformly equivalent to
      $\|X\|_{\mathcal X_h^s}^2$ such that
      \begin{equation}
       (\mathscr E^B_{s,h})'
       \le C_s\mathscr E^B_{s,h}
          +C_s(\mathscr E^B_{s,h})^{3/2}.
       \label{eq:main-law-energy}
      \end{equation}
\item Let $B_h\in\overline{\mathcal B}$ be arbitrary, with no convergence of
      $B_h$ required.  For the compatible data of
      \eqref{eq:main-open-data-class}, the whole sequence of solutions of
      \eqref{eq:main-open-law-family} converges to the same harmonic vacuum
      development $g$ as the central writer, and
      \begin{equation}
       \sum_{j=0}^3
       \|q_{B_h,h}^{(j)}-q_{0,h}^{(j)}\|_{C_tH_h^{s-1-j}}
          \le C_sh\|B_h\|_{\op}\varepsilon.
       \label{eq:main-law-jet-comparison}
      \end{equation}
      Hence the metric, subsidiary, noncollapse and curvature conclusions of
      \Cref{thm:main-open-3plus1} are uniform on the entire coefficient ball.
\item Let $K_{h,j}(e\mid\pi)$ be arbitrary history-dependent finite transition
      rows satisfying, at every live prefix,
      \begin{equation}
       \boxed{\sum_eK_{h,j}(e\mid\pi)c_{h,j}(\pi,e)<M\tau_jh^4.}
       \label{eq:main-law-row-moment}
      \end{equation}
      No independence, invariant starting law, or per-letter small-defect
      condition is imposed.  Then
      \begin{equation}
       \boxed{\mathbb EC_h\le MTh^4,
       \qquad \mathbb P\{C_h>h^2\}\le MTh^2.}
       \label{eq:main-law-probability}
      \end{equation}
      On the complementary event the finite history shadows a member of the
      differential-law family closely enough to have the same Einstein limit.
      At each fixed cutoff and grammar the strict conditions
      \eqref{eq:main-law-row-moment} form nonempty relatively open convex
      subsets of the product of row simplexes, and may contain strictly
      positive rows.  Along the odd refinement sequence the exceptional
      probabilities are summable, so under any common coupling the Einstein
      conclusion holds eventually almost surely.
\item For every fixed sufficiently small nonzero preparation amplitude $t$,
      the exact map \eqref{eq:main-action-initial-constraints} has a finite,
      nonsymmetric analytic zero-set chart generated from arbitrary small free
      tensor records $W$ after removal of four fixed balancing modes.  Its
      prepared records $\mathfrak I_h(t,W)$ satisfy
      \begin{equation}
       \boxed{\mathscr C_h(\mathfrak I_h(t,W))=0,\qquad
       \|\mathfrak I_h(t,W)\|\le C_s|t|.}
       \label{eq:main-action-prepared-zero}
      \end{equation}
      The chart is relatively open in the complete finite zero set, has
      codimension $4N^3$, and the linearized constraint map has a right inverse
      obeying
      \begin{equation}
       \|\mathscr R_{h,t}f\|
       \le C_s\bigl(\|P_\perp f\|+|t|^{-1}|P_0f|\bigr).
       \label{eq:main-action-prepared-inverse}
      \end{equation}
      No CMC, conformal-flatness, spatial symmetry, finite Fourier support, or
      continuum-compatible metric is imposed on the free record $W$.
      Harmonic completion of these exact finite records feeds every sequence
      $B_h\in\overline{\mathcal B}$ into the common Einstein population, and
      the original four lapse--shift rows vanish identically at the initial
      cut independently of the subsequently chosen acceleration law.
\end{enumerate}
The law-space statement is deliberately quantitative.  It gives a fixed open
ball of differential laws and relatively open finite probability-row
neighborhoods in a physically scaled residual topology; it does not assert a
fixed unscaled total-variation ball of arbitrary microscopic kernels or
attraction from arbitrary leading-order forces.
\end{theorem}

\begin{remark}[The robustness class is nontrivial but not arbitrary]
\label{rem:main-law-basin-scope}
The family in \Cref{thm:main-open-law-basin} can differ from the central
writer by a nonvanishing amount at the highest lattice frequencies.  The law
parameter is, however, retained during a physical history.  Unrestricted
rapid switching between two individually stable coefficient marks can
parametrically amplify a resolved high-frequency mode.  Likewise, a row may
be close in ordinary total variation while accumulating order-one failure
probability over $O(h^{-2})$ stages.  These boundaries explain why the
neighborhood is formulated by energy-compatible coefficients and complete
physical residual profiles rather than by an unscaled norm on arbitrary
kernels.  The proofs are in Appendix~\ref{supp:open-law-basin}.
\end{remark}

\subsection{Exact-action response, preparation, and dynamical tangency}
\label{subsec:main-exact-action-audit}

The constructive laws above are not identified with the raw finite Euler flow
of the reconstructed phase-compatible action.  That stronger route can be
analyzed directly.  Let the finite canonical phase have dimension $12V$,
let $\lambda=(N,\beta)\in\mathbb R^{4V}$ denote lapse--shift multipliers,
and let $P_h=\nabla_\lambda H_h$ be the exact multiplier row of the retained
finite Hamiltonian.  Write
\begin{equation}
 F_h^{\rm can}=\mathbb J\nabla_XH_h,\qquad
 B_h=D_XP_hF_h^{\rm can},\qquad M_h=D_\lambda P_h.
 \label{eq:main-exact-action-basic}
\end{equation}
The exact constrained evolution satisfies
$X_t=F_h^{\rm can}$, $P_h=0$, and
$B_h+M_h\dot\lambda=0$.

For the flat Cartan tangent let $\mathsf P_{\rm tr}$ denote the orthogonal
projection onto scalar spatial matrices and put
\begin{equation}
 K_0=2I-6\mathsf P_{\rm tr},\qquad
 K_0^{-1}=\tfrac12I-\tfrac34\mathsf P_{\rm tr}.
 \label{eq:main-exact-action-K0}
\end{equation}
For any full-rank finite source inclusion $F$ and actual target $b$, define
\begin{equation}
 \Pi=F(F^*F)^{-1}F^*,\qquad
 w=F(F^*F)^{-1}b,
 \label{eq:main-exact-action-positive-reproduction}
\end{equation}
and let $E$ be the isometric inclusion of scalar matrices,
$E\tau=\tau I/\sqrt3$.  Set
\begin{equation}
 C=E^*\Pi E,\qquad H=2I-3C,\qquad \zeta=E^*w.
 \label{eq:main-exact-action-trace-data}
\end{equation}
These objects distinguish the positive source-reproduction problem from the
indefinite Cartan correction.

\begin{theorem}[Exact-action response and prepared-branch structure]
\label{thm:main-exact-action-response}
The exact finite-action system has the following properties.
\begin{enumerate}[label=\textup{(X\arabic*)},leftmargin=3em]
\item On a regular one-null branch the extended constraint manifold has
      restricted symplectic characteristic quotient dimension
      \begin{equation}
       \boxed{12V-2.}
       \label{eq:main-exact-action-quotient}
      \end{equation}
      Thus the finite branch is not regularly canonically equivalent to a
      $12V$-dimensional system with $4V$ independent local first-class
      constraints and no additional second-class constraints.  A continuum
      GR identification, if present, must involve a selected or singular
      cofinal mechanism rather than a finite-cutoff relabeling.
\item On the scalar-momentum branch of Appendix~\ref{supp:exact-action-audit},
      the transverse-shift source $A_h$ is injective on the
      $2(V-1)$-dimensional transverse space at every odd cutoff and obeys
      \begin{equation}
       c h^2\|\Omega_h\beta\|_h^2
       \le\|A_h\beta\|_h^2
       \le C h^2\|\Omega_h\beta\|_h^2,
       \qquad
       A_h^*K_0^{-1}A_h=\tfrac12A_h^*A_h.
       \label{eq:main-exact-action-transverse}
      \end{equation}
      Consequently these columns can be eliminated exactly from the signed
      response system without discarding any remaining source equation.
\item For every full-rank remaining source $F$, the signed Gram
      $F^*K_0^{-1}F$ is invertible if and only if $H$ in
      \eqref{eq:main-exact-action-trace-data} is invertible.  On that branch,
      with $\chi=H^{-1}\zeta$, the unchanged physical response is
      \begin{equation}
       \boxed{v=w-3(I-\Pi)E\chi,\qquad E^*v=-\chi,}
       \label{eq:main-exact-action-scalar-response}
      \end{equation}
      and
      \begin{equation}
       \boxed{\|v\|^2=\|w\|^2+9\langle\chi,(I-C)\chi\rangle.}
       \label{eq:main-exact-action-scalar-energy}
      \end{equation}
      Hence $v_h\to0$ is equivalent to the two separately measurable
      conditions $w_h\to0$ and $H_h^{-1}E^*w_h\to0$.  Diagonalizing $C$
      recovers the unique trace-angle resonance $\mu=2/3$ as a corollary,
      rather than as a separate physical postulate.
\item There are nontrivial exactly prepared original-action families with
      arbitrary sufficiently fine odd cutoff for which the complete weak
      quadratic target vanishes identically while the three spatial axial
      modes remain nonzero.  The exact finite constraints admit analytic
      preparations
      \begin{equation}
       \boxed{\mathscr C_h(X_h(a))=0,\qquad
       X_h(a)=aY_h+O(a^2).}
       \label{eq:main-exact-action-neutral-preparation}
      \end{equation}
      A distinct rationally detuned family has, at every sufficiently fine
      odd cutoff, complete leading active rank $3V-3$, trace-free active
      range, and
      \begin{equation}
       \boxed{(\mathcal A_h^{\rm act})^*K_0^{-1}\mathcal A_h^{\rm act}
       =\tfrac12(\mathcal A_h^{\rm act})^*\mathcal A_h^{\rm act}
       \succeq c h^{16}I.}
       \label{eq:main-exact-action-active-gram}
      \end{equation}
      Thus the leading $2/3$ signed trace resonance is absent on that active
      range.  At a verified regular source-neutral seed the exact normalized
      original-action continuation has literal temporal and spatial link
      curvature $O_h(|a|)$ at the initial cut.
\item On the selected weak-neutral active family treated in the supplement,
      the true transverse normal operator has a cutoff-independent spectral
      interval and the actual direct transverse solution satisfies
      \begin{equation}
       \boxed{\|\mathfrak E_h(\beta_{0,h})\|_{H_h^s}\le C_sh.}
       \label{eq:main-exact-action-direct-output}
      \end{equation}
      The remaining scalar lapse equation is exactly positive:
      \begin{equation}
       \boxed{\mathscr S_h
       =\tfrac12\mathscr D_h^*\mathscr D_h
        +\tfrac{h^2}{2}\sigma_h^*K_{Q,h}^{-1}\sigma_h,
       \qquad cI\preceq K_{Q,h}\preceq CI.}
       \label{eq:main-exact-action-scalar-flux}
      \end{equation}
      The retained flux lifts the bare diagonal soft sector; a large response
      predicted by the diagonal compression alone is therefore not a theorem
      about the original full scalar solve.
\item Static response suppression and metric-jet shadowing are distinct.  For
      a $C^r$ family of full-rank sources, differentiation of
      $H\chi=\zeta$ gives
      \begin{equation}
       \boxed{\chi^{(j)}=H^{-1}\left[\zeta^{(j)}
        +3\sum_{\ell=1}^j\binom j\ell C^{(\ell)}
          \chi^{(j-\ell)}\right].}
       \label{eq:main-exact-action-time-hierarchy}
      \end{equation}
      In particular a small instantaneous electric response does not control
      its time derivative.  In the active/weak amplitude grading, harmonic
      first-jet shadowing additionally requires the normalized lapse--shift
      rate to satisfy
      \begin{equation}
       a^{-1}\|(u_a)_A\|_h\to0,
       \qquad a^{-2}\|(u_a)_W\|_h\to0.
       \label{eq:main-exact-action-clock-test}
      \end{equation}
\end{enumerate}
The theorem gives exact finite response structure and nonempty prepared
branches, but it does not identify those branches with a cutoff-uniform
Einstein evolution.  Dynamical invariance of the prepared set is a separate
question.
\end{theorem}

\begin{proposition}[Exact-action tangency obstruction and upstream deformability]
\label{prop:main-exact-action-tangency}
There exists an exactly prepared nontrivial original-action branch at a fixed
finite cutoff for which the complete weak multiplier velocity vanishes at the
initial cut and the literal curvature is $O_h(a)$, but the derivative of an
explicit weak Poisson-null tangency functional is nonzero.  The corresponding
exact histories are controlled on every interval $0\le t\le a^3S$ for fixed
$S$, while no bound of the form
\begin{equation}
 \sup_{0\le t\le T}\|R_{0i0j}(g_h)(a,t)\|\le Ca
 \label{eq:main-exact-action-no-fixed-linear}
\end{equation}
can hold uniformly for all sufficiently small $a$ on one fixed positive
physical interval $T$ for that family.  Altering only the admissible leading
multiplier tilt at the same canonical first jet does not remove the leading
tangency defect.  However, inside the finite canonical first-jet chart there
is an exact lower-incidence tangent $(Z,m)$ for which
\begin{equation}
 \boxed{D\mathcal L[(Z,m)]=0,
 \qquad D\Delta_\diamond[(Z,m)]\ne0.}
 \label{eq:main-exact-action-upstream-tangent}
\end{equation}
Thus the obstruction is attached to the chosen first jet rather than an
invariant no-go property of the unchanged finite action.  What remains open is
an exact regular nonlinear root of the complete transported tangency system,
its cofinal continuation, higher signed-response propagation, differentiated
curvature control, and a cutoff-independent exact-action lifespan.
\end{proposition}

The proofs and constructions are collected in
Appendix~\ref{supp:exact-action-audit}.  The logical boundary is therefore
sharper than a comparison between two discretizations.  A stable open family
of renewal laws is already known to converge to Einstein evolution, while the
raw action has an independently understood finite response geometry and a
concrete hierarchy of compatibility conditions whose nonlinear cofinal
closure remains to be proved.

\subsection{General renewal-to-Einstein certification}

\begin{theorem}[General renewal-to-Einstein certification]
\label{thm:main-renewal-einstein}
The preceding constructive theorem treats one explicit nonlinear writer.
For a general complete gravitational renewal cylinder, consider a cofinal
sequence of Lorentzian relational regulators of \Cref{mt:adm}.  On every
compact spacetime cylinder $K$, define the following five certification
conditions on compatible restrictions of the same operational process.
\begin{enumerate}[label=\textup{(E\arabic*)},leftmargin=3em]
\item The complete accepted cylinder and generic protected reward-duration record
      reconstruct the determining kernel and scalar common action.  Its actual
      first variations tend to zero on a fixed determining metric-test core.
      The scaled bound \eqref{eq:main-physical-stationarity}, with the realized
      subsequence of \Cref{prop:supp-physical-test-stationarity}, is a sufficient
      operational route to this condition.  The physical coframe lifts of these
      tests have uniformly bounded, convergent coefficient tensors in the
      curvature and volume insertions.
\item Complete resolved-memory subtraction and exact or summably stable binary
      face subdivision give the first-order score classified by
      \Cref{thm:main-gravitational-classification}.
\item In a fixed physical action calibration the spacetime-constant
      coefficients converge,
      $(\alpha_X,\beta_X,\lambda_X)\to(\alpha,\beta,\lambda)$, with
      $\beta\ne0$ the coefficient of $L_{\rm Palatini}$.  The spinless
      connection Euler residual tends to zero and the Cartan map is uniformly
      nondegenerate on the retained coframe chart.
\item The endpoint packet satisfies the hypotheses of
      \Cref{thm:main-spatial-screen}, with
      $\liminf_X I_X^{\rm eff}>0$.  Its common-cylinder compact-screen packet,
      including transported-screen convergence and temporal control, gives a
      nondegenerate limit with $e_X\to e$ strongly in local $L^2$ and
      $\sup_X\|e_X\|_{L^6(K)}<\infty$.
\item The actual curvature records satisfy either the continuous propagation
      certificate of \Cref{thm:main-curvature-propagation} or the fully discrete
      certificate of \Cref{thm:main-discrete-curvature}, with uniform initial,
      growth, transfer and source budgets and stable curvature interpolation.
      The scaled action-current estimate
      \eqref{eq:main-action-source-budget} may supply the source budget.
      The connection and its actual Cartan writer satisfy
      \eqref{eq:main-identified-curvature-hypotheses}.
\end{enumerate}
If conditions \textup{(E1)--(E5)} hold on a cofinal tail, then a subsequence
has zero limiting torsion and, for each determining compactly supported
inverse-metric test $k^{\mu\nu}$, the metric first variation converges to
\begin{equation}
 \boxed{\chi\int_K\sqrt{-g}\,
       (G_{\mu\nu}+\Lambda g_{\mu\nu})k^{\mu\nu},\qquad \chi\ne0.}
 \label{eq:main-einstein-variation}
\end{equation}
Consequently
\begin{equation}
 \boxed{G_{\mu\nu}+\Lambda g_{\mu\nu}=0}
 \label{eq:main-vacuum-einstein}
\end{equation}
holds distributionally on the determining test core, and on its closure in
the declared test topology.
\end{theorem}

\begin{remark}[Certification is not the definition of the constructive branch]
\label{rem:main-certification-not-tautology}
Conditions \textup{(E1)--(E5)} are sufficient recognition conditions for a
general complete cylinder: they control action stationarity, first-order
gravitational provenance, coupling/torsion, noncollapse/strong coframe
compactness, and curvature propagation/identification.  They do not define
the local writer of \Cref{thm:main-open-3plus1}, and that theorem does not use
them as independent assumptions.  Conversely, the certification theorem is
more general than the constructive population: it applies whenever those
finite and cofinal quantities can be verified, regardless of how the
microscopic writer is generated.
\end{remark}

The nonzero Palatini coefficient is distinct from invertibility of the
Palatini--Holst connection operator.  A nonzero pure Holst coefficient can
eliminate torsion, but its torsion-free metric variation does not supply the
Einstein equation.  Likewise the propagated curvature bound supplies size,
whereas \eqref{eq:main-identified-curvature-hypotheses} identifies the nonlinear
connection curvature.  These distinctions prevent either condition from being
hidden inside a generic ``positive branch.''

\subsection{General vacuum Einstein-or-obstruction alternative}

Call a cofinal regulator a \emph{complete gravitational renewal cylinder} when
the geometric packet, complete accepted determining-field cylinder, generic
protected reward-duration record, and physical endpoint, propagation and
Cartan records occur on compatible restrictions of one operational process.
The records include the physical test lifts, action calibration, actual
curvature interpolation and common-cylinder comparison maps used above.
The word \emph{gravitational} labels the sector being tested; it does not mean
that an Einstein Lagrangian is stored in the entrance record.

\begin{theorem}[General vacuum Einstein-or-obstruction certification]
\label{thm:main-einstein-alternative}
For a complete gravitational renewal cylinder, perform the finite
reconstructions and residual tests of this section at each cutoff.  On a
selected cofinal sequence, the following alternative holds.
\begin{enumerate}[label=\textup{(A\arabic*)},leftmargin=3em]
\item If the certificates \textup{(E1)--(E5)} of
      \Cref{thm:main-renewal-einstein} pass on a cofinal tail, the sequence has
      a subsequence with a nontrivial Lorentzian $3+1$ continuum satisfying
      \begin{equation}
       \boxed{G_{\mu\nu}+\Lambda g_{\mu\nu}=0}
       \label{eq:main-global-vacuum-einstein}
      \end{equation}
      distributionally on the determining test core.  Curvature size is an
      output of the initial-energy propagation certificate, not an additional
      all-time hypothesis.
\item Otherwise a failed certificate is retained on the same physical records:
      a common-action curl or period; a nonvanishing realized metric variation
      or failure of its scaled action-gap bound; a vanishing Palatini
      normalization or failed coupling calibration; a cut/potential witness;
      a face-linearity, resolved-memory or torsion defect; an unbounded initial,
      growth, transfer or differentiated-source budget; an unstable curvature
      interpolation; a Cartan or interface-identification discrepancy; or a
      common-cylinder compactness failure.  A cofinal failure is witnessed by
      a subsequence of residuals bounded away from zero, diverging budgets,
      vanishing positive margins, or escaping compactness tails, as appropriate.
\end{enumerate}
The alternative is unconditional relative to the complete cylinder and the
specified certification procedure.  Its second branch is an obstruction to
that certificate; it is not a theorem that no Einstein subsequence or other
continuum description can exist.
\end{theorem}

% =============================================================================
\section{Vacuum renewal dynamics and explicit Einstein regulators}
\label{sec:vacuum-main}
% =============================================================================

The general certification theorem and the open nonsymmetric population of
Section~\ref{sec:einstein-dynamics} establish the field equation and a
fixed-radius $3+1$ realization, respectively.  We now study complementary
solution-level sectors without inserting an Einstein metric into the
operational entrance.  On the root-symmetric homogeneous branch we leave the
renewal variables arbitrary, derive their full lapse-varied stationarity
equations, and only afterwards identify the flat and de~Sitter solutions.  We
then treat a nonlinear spatially inhomogeneous unpolarized Gowdy sector in
which the finite evolution preserves a discrete momentum constraint and
admits quantitative residual, metric, and action tests.  Finally, the
homogeneous branch is populated by a finite positive acceptance law whose
measured probabilities reconstruct its finite action and whose low-noise
refinement converges to the stationary solution.

\subsection{Homogeneous renewal coordinates and the full lapse constraint}
\label{subsec:homogeneous-renewal-main}

On the $A_3$ root-symmetric branch write
\begin{equation}
 k_{\alpha,h}(t)=\frac{\kappa(t)}{8h^2},
 \qquad
 m_h(t)=\varrho(t)h^3,
 \qquad \kappa(t)>0,\quad\varrho(t)>0.
 \label{eq:main-general-homogeneous-rates}
\end{equation}
The root identity gives $\mathcal B=\kappa I_3$ and reversal symmetry gives
$\beta=0$.  If $g_{ij}=a^2\delta_{ij}$, the ADM inversion
\eqref{eq:adm-inversion-main} is therefore exactly
\begin{equation}
 \boxed{
 a=\varrho^{1/3},\qquad
 N=\varrho^{1/3}\kappa^{1/2},\qquad
 \varrho=a^3,\qquad
 \kappa=\frac{N^2}{a^2}.}
 \label{eq:main-isotropic-renewal-map}
\end{equation}
Thus $\kappa$ is not the scale factor: together with the physical density it
contains the lapse information that must remain independent until after
variation.

On a compact interval and torus of coordinate volume $V_0$, the torsion-free
Palatini--volume branch of \Cref{thm:main-renewal-einstein}, equivalently its
Einstein--Hilbert representative, reduces after the standard
Gibbons--Hawking--York Dirichlet boundary completion
\cite{York1972,GibbonsHawking1977} and division by the nonzero constant
$\chi V_0$, to
\begin{equation}
 S[a,N]=\int\left[-\frac{6a\dot a^2}{N}-2\Lambda Na^3\right]\dd t.
 \label{eq:main-homogeneous-action-an}
\end{equation}
Here the induced spatial metric is fixed at the initial and final slices, so
$\delta a(t_0)=\delta a(t_1)=0$.  The lapse remains an independent bulk
variation and is not gauge-fixed before the Euler equations are taken.  The
four-dimensional reduction and the exact cancellation of the second-derivative
endpoint term by the boundary completion are written out in
Appendix~\ref{supp:homogeneous-dynamics}.

\begin{theorem}[Vacuum Friedmann equations in renewal coordinates]
\label{thm:main-homogeneous-friedmann}
Let $u=\log\kappa$, $v=\log\varrho$, and
$\mathcal H=\dot a/(Na)$, with $D_\tau=N^{-1}\dd/\dd t$.  Under
\eqref{eq:main-isotropic-renewal-map}, the action
\eqref{eq:main-homogeneous-action-an} becomes
\begin{equation}
 \boxed{
 S[\kappa,\varrho]
 =\int\left[-\frac{2\dot\varrho^2}
 {3\varrho^{4/3}\sqrt\kappa}
 -2\Lambda\varrho^{4/3}\sqrt\kappa\right]\dd t.}
 \label{eq:main-renewal-rate-action}
\end{equation}
Its Euler expressions, defined by
$\delta S=\int(E_u\delta u+E_v\delta v)\dd t$, are
\begin{equation}
 \boxed{
 E_u=Na^3(3\mathcal H^2-\Lambda),\qquad
 E_v=Na^3\left[4D_\tau\mathcal H+
 \frac83(3\mathcal H^2-\Lambda)\right].}
 \label{eq:main-renewal-euler}
\end{equation}
Consequently stationarity in the two independent positive renewal variables is
equivalent to
\begin{equation}
 \boxed{3\mathcal H^2=\Lambda,\qquad D_\tau\mathcal H=0,}
 \label{eq:main-vacuum-friedmann}
\end{equation}
which is the complete vacuum Einstein equation within the spatially flat
homogeneous isotropic ansatz, including the lapse constraint.
\end{theorem}

\begin{proof}
The map $(u,v)\mapsto(\log a,\log N)=(v/3,u/2+v/3)$ has nonzero
Jacobian, so no homogeneous variation is lost.  Varying
\eqref{eq:main-homogeneous-action-an} before fixing $N$ gives
\[
 E_N=2a^3(3\mathcal H^2-\Lambda),\qquad
 E_a=6Na^2(2D_\tau\mathcal H+3\mathcal H^2-\Lambda).
\]
The chain rule gives $E_u=(N/2)E_N$ and
$E_v=(a/3)E_a+(N/3)E_N$, proving
\eqref{eq:main-renewal-euler}.  The two equations are therefore equivalent to
\eqref{eq:main-vacuum-friedmann}.  Appendix~\ref{supp:homogeneous-dynamics}
checks directly that these are exactly the $00$ and spatial Einstein equations
for the reconstructed metric, rather than relying on a symmetry-reduction
principle.
\end{proof}

\begin{corollary}[Flat and de Sitter rates are stationary renewal solutions]
\label{cor:main-derived-desitter-rates}
For $\Lambda>0$ put $H_0=\sqrt{\Lambda/3}$ and let
$\tau(t)=\int^tN(s)\dd s$.  Every positive connected stationary solution of
\eqref{eq:main-renewal-rate-action} has
\begin{equation}
 a(\tau)=a_0e^{\sigma H_0(\tau-\tau_0)},\qquad
 \varrho=a^3,\qquad
 \kappa=N^2/a^2,
 \qquad \sigma\in\{-1,1\}.
 \label{eq:main-stationary-homogeneous-solutions}
\end{equation}
After variation, choose proper-time gauge $N=1$.  On the expanding branch,
\begin{equation}
 \boxed{
 k_{\alpha,h}(t)=
 \frac{a_0^{-2}e^{-2H_0(t-t_0)}}{8h^2},\qquad
 m_h(t)=a_0^3e^{3H_0(t-t_0)}h^3.}
 \label{eq:desitter-regulator-main}
\end{equation}
For $\Lambda=0$, $a$ is constant; for $\Lambda<0$ there is no positive real
solution in this spatially flat vacuum sector.
\end{corollary}

The exponential factor in \eqref{eq:desitter-regulator-main} is therefore an
output of the stationary renewal equations on this branch, not an input used to
obtain those equations.  The result does not select $\Lambda$, the expanding
versus contracting sign, or the initial scale $a_0$.

\subsection{A finite lapse-varied renewal action}
\label{subsec:finite-homogeneous-action-main}

The same distinction can be kept at finite cutoff.  The construction is a
midpoint variational discretization in the standard sense of discrete
mechanics~\cite{MarsdenWest2001}, but every equation needed below is derived
explicitly.  Put
$q=a^{3/2}=\sqrt\varrho$ and let
$s_j=N_j\Delta t_j>0$ be independent lapse-weighted intervals.  With
$\bar q_j=(q_j+q_{j+1})/2$, $A_0=8/3$, and $B_0=2\Lambda$, define
\begin{equation}
 \boxed{
 S_{\rm d}(q,s)
 =-\sum_{j=0}^{M-1}\left[
 A_0\frac{(q_{j+1}-q_j)^2}{s_j}
 +B_0s_j\bar q_j^2\right].}
 \label{eq:main-discrete-homogeneous-action}
\end{equation}
This is the midpoint discretization of the already
Gibbons--Hawking--York-completed reduced functional
\eqref{eq:main-homogeneous-action-an}, not of the bare second-derivative
Einstein--Hilbert expression.  The Dirichlet endpoint data are therefore
$q_0,q_M$ fixed (equivalently the induced endpoint three-metrics are fixed),
while every interior $q_j$ and every lapse-weighted interval $s_j$ are varied
independently.  Consequently
\begin{equation}
 \delta S_{\rm d}
 =\sum_{j=1}^{M-1}E_{q_j}\,\delta q_j
  +\sum_{j=0}^{M-1}E_{s_j}\,\delta s_j,
 \qquad \delta q_0=\delta q_M=0,
 \label{eq:main-discrete-dirichlet-variation}
\end{equation}
with no omitted endpoint acceleration term.  If one discretized the uncompleted
Einstein--Hilbert expression instead, a matching discrete gravitational boundary
term would have to be retained explicitly.  Through
\begin{equation}
 \varrho_j=q_j^2,\qquad
 N_{j+1/2}=s_j/\Delta t_j,\qquad
 \kappa_{j+1/2}=\frac{(s_j/\Delta t_j)^2}{\bar q_j^{4/3}},
 \label{eq:main-discrete-renewal-map}
\end{equation}
this is a smooth finite-dimensional action in positive renewal coordinates; no
exponential history appears in its definition.

\begin{theorem}[Finite stationary classification and de Sitter convergence]
\label{thm:main-finite-homogeneous-stationarity}
Let $v_j=(q_{j+1}-q_j)/s_j$.  The full Euler equations of
\eqref{eq:main-discrete-homogeneous-action} are
\begin{align}
 A_0v_j^2-B_0\bar q_j^2&=0,\label{eq:main-discrete-lapse-constraint}\\
 2A_0(v_j-v_{j-1})
 -B_0(s_{j-1}\bar q_{j-1}+s_j\bar q_j)&=0.
 \label{eq:main-discrete-scale-euler}
\end{align}
For $\Lambda>0$, put $\omega=\sqrt{B_0/A_0}=\sqrt{3\Lambda/4}$.
Every positive stationary solution has one common sign
$\sigma\in\{-1,1\}$ and
\begin{equation}
 \boxed{
 q_{j+1}=q_j
 \frac{1+\sigma\omega s_j/2}{1-\sigma\omega s_j/2},
 \qquad 0<\omega s_j/2<1.}
 \label{eq:main-finite-stationary-recurrence}
\end{equation}
Conversely every such recurrence is stationary.  If
$\tau_j=\sum_{i<j}s_i$, $T=\sum_js_j$, $h_t=\max_js_j$, and
$\omega h_t/2\le b<1$, then
\begin{equation}
 \boxed{
 \max_j\left|\log\frac{q_j}{q_0}-\sigma\omega\tau_j\right|
 \le \frac{\omega^3T}{12(1-b^2)}h_t^2.}
 \label{eq:main-finite-profile-convergence}
\end{equation}
Hence $a_j=q_j^{2/3}$, $\varrho_j=q_j^2$, and in proper-time gauge
$\kappa_j=q_j^{-4/3}$ converge with second-order logarithmic error to the
stationary flat/de~Sitter profiles of
\Cref{cor:main-derived-desitter-rates}.
\end{theorem}

The proof, including the lapse variation and the exact rational recurrence, is
in Appendix~\ref{supp:homogeneous-dynamics}.  In particular, fixing $N=1$
before variation would lose \eqref{eq:main-discrete-lapse-constraint} and admit
spurious non-Einstein homogeneous trajectories.

\subsection{Measured acceptance law and action--dynamics matching}
\label{subsec:operational-homogeneous-main}

To relate the finite variational action directly to an operational probability
law, we next exhibit a finite positive class in which the action is recovered
from the trial probabilities themselves.  For $A,B>0$, $C\in\R$, and $s>0$, consider
\begin{equation}
 E_s(x,y)=A\frac{(y-x)^2}{s}
 +Bs\left(\frac{x+y}{2}\right)^2+C(y^2-x^2),
 \qquad
 \begin{pmatrix}A&C\\ C&B\end{pmatrix}\succeq0.
 \label{eq:main-measured-quadratic-cost}
\end{equation}
Within this declared finite parametric class, three measured costs identify
$A,B,C$.  Define the directly measurable rank defect
\begin{equation}
 \boxed{\mathfrak D=AB-C^2\ge0.}
 \label{eq:main-rank-defect}
\end{equation}
For a finite proposal alphabet $F$, let $p_s(y\mid x)>0$ be a normalized
proposal row and accept a proposal with probability
\begin{equation}
 a_s(x,y)=e^{-E_s(x,y)/\varepsilon}.
 \label{eq:main-operational-acceptance}
\end{equation}
Write
\begin{equation}
 Z_s(x)=\sum_{y\in F}p_s(y\mid x)a_s(x,y),\qquad
 Q_s(y\mid x)=\frac{p_s(y\mid x)a_s(x,y)}{Z_s(x)}.
 \label{eq:main-accepted-row-normalizer}
\end{equation}
Both $Q_s$ and the total acceptance probability $Z_s$ are entries of the same
complete trial table.

\begin{theorem}[Measured acceptance reconstructs the finite action]
\label{thm:main-operational-action-matching}
The complete trial table reconstructs the unnormalized cost exactly:
\begin{equation}
 \boxed{
 E_s(x,y)=-\varepsilon\log
 \frac{Z_s(x)Q_s(y\mid x)}{p_s(y\mid x)}.}
 \label{eq:main-observed-cost}
\end{equation}
For a path with fixed endpoints, the observed writer
\begin{equation}
 \mathscr A_{\rm obs}
 =\varepsilon\sum_j\left[
 \log\frac{Q_{s_j}(q_{j+1}\mid q_j)}{p_{s_j}(q_{j+1}\mid q_j)}
 +\log Z_{s_j}(q_j)\right]
 +C(q_M^2-q_0^2)
 \label{eq:main-observed-homogeneous-action}
\end{equation}
is identically
\begin{equation}
 \boxed{
 \mathscr A_{\rm obs}
 =-\sum_j\left[A\frac{(q_{j+1}-q_j)^2}{s_j}
 +Bs_j\bar q_j^2\right]}
 \label{eq:main-offshell-action-match}
\end{equation}
off shell on the parametric domain.  Thus every interior-amplitude and lapse
first variation agrees exactly with the reconstructed finite action.

The one-step minimizing skeleton has proper-time amplitude drift
$c=-C/A$.  The bulk action in
\eqref{eq:main-offshell-action-match} has
$\Lambda_{\rm act}=4B/(3A)$, whereas the limiting constant-Hubble geometry of
the minimizing skeleton has
$\Lambda_{\rm geom}=4C^2/(3A^2)$.  Consequently
\begin{equation}
 \boxed{
 G[g]+\Lambda_{\rm act}g
 =\frac{4\mathfrak D}{3A^2}g.}
 \label{eq:main-measured-einstein-mismatch}
\end{equation}
The local operational update and the vacuum parameter of its measured bulk
action therefore agree exactly if and only if $\mathfrak D=0$.
\end{theorem}

The total acceptance probability in
\eqref{eq:main-observed-cost} is essential: the accepted row $Q_s$ alone is
invariant under addition of a source-dependent cost and therefore does not
identify the lapse action.  The complete trial record removes this ambiguity.
On the rank-one branch with $A=A_0$, $B=A_0\omega^2$, and
$C=-\sigma A_0\omega$, the cost is the square
\begin{equation}
 E_s(x,y)=\frac{A_0}{s}
 \left[y-x-\sigma\omega s\frac{x+y}{2}\right]^2,
 \label{eq:main-rank-one-cost}
\end{equation}
so its zero-cost recurrence is exactly
\eqref{eq:main-finite-stationary-recurrence}.

\begin{theorem}[Finite positive operational family with an Einstein limit]
\label{thm:main-explicit-operational-family}
Fix $T,\omega,q_0>0$ and a sign $\sigma$.  There exists a sequence of genuinely
finite accept/reject cylinders with $M$ time stages, finite positive proposal
alphabets, calibrated step $s=T/M$, rank-one costs
\eqref{eq:main-rank-one-cost}, and a finite retry cap such that, for constants
independent of $M$,
\begin{equation}
 \mathbb E\mathcal E_M\le CM^{-4},\qquad
 \mathbb E\max_j|q_j-q_0e^{\sigma\omega\tau_j}|\le CM^{-2},
 \label{eq:main-stochastic-profile}
\end{equation}
where
$\mathcal E_M=A_0\sum_jR_j^2/s$ and
$R_j=q_{j+1}-q_j-\sigma\omega s\bar q_j$.  The expected absolute first
variation of \eqref{eq:main-discrete-homogeneous-action} against every fixed
smooth scale/lapse test is $O(M^{-2})$, including the lapse constraint.  The
retry-failure probability is at most $M^{-6}$.  Under any common coupling of
cutoffs, failures occur only finitely often almost surely and the successful
piecewise-affine paths converge strongly in $H^1$ and uniformly to the
stationary solution.  Attaching the $A_3$ root race with
\begin{equation}
 \varrho_j=q_j^2,\qquad \kappa_j=q_j^{-4/3},\qquad
 k_{\alpha,h_x,j}=\frac{\kappa_j}{8h_x^2},\qquad
 m_{h_x,j}=\varrho_jh_x^3
 \label{eq:main-stochastic-root-rates}
\end{equation}
therefore gives positive renewal regulators converging to a vacuum metric with
$\Lambda=4\omega^2/3$.
\end{theorem}

Appendix~\ref{supp:homogeneous-dynamics} gives the finite grids, temperature
schedule, entropy--energy estimate, retry bound, and deterministic residual
estimates used in this theorem.  The construction proves nonemptiness of a
matched operational class; it does not assert that an arbitrary Grand Tensor
satisfies the quadratic duration law or the rank-one condition.

\subsection{Conservative blocking preserves the action--law defect}
\label{subsec:homogeneous-blocking-main}

The preceding construction raises a distinct stability question: can ordinary
time coarse graining drive a neighboring quadratic law toward the matched
surface?  For the conservative interval elimination of the measured quadratic
cost, the answer is no.  Assume $0<s\sqrt{B/A}<2$ and set
\begin{equation}
 \mu=\sqrt{AB},\qquad
 \theta=2\operatorname{artanh}\!\left(\frac{s}{2}\sqrt{B/A}\right),\qquad
 \rho=\frac{C}{\mu},\qquad
 \Delta_{\rm rel}=1-\rho^2=\frac{AB-C^2}{AB}.
 \label{eq:main-relative-rank-defect}
\end{equation}
Then the same cost can be written
\begin{equation}
 \mathcal E_{\mu,\theta,C}(x,y)
 =\mu\coth\theta\,(x^2+y^2)-2\mu\operatorname{csch}\theta\,xy
   +C(y^2-x^2).
 \label{eq:main-hyperbolic-cost}
\end{equation}
For two successive intervals define conservative blocking by
\begin{equation}
 (E_1\star_0E_2)(x,y)=\inf_{z>0}\{E_1(x,z)+E_2(z,y)\}.
 \label{eq:main-conservative-blocking}
\end{equation}

\begin{theorem}[No transverse matching attraction under conservative blocking]
\label{thm:main-no-blocking-attractor}
For every $\theta_1,\theta_2>0$,
\begin{equation}
 \boxed{
 \mathcal E_{\mu,\theta_1,C}\star_0
 \mathcal E_{\mu,\theta_2,C}
 =\mathcal E_{\mu,\theta_1+\theta_2,C}.}
 \label{eq:main-hyperbolic-blocking}
\end{equation}
Hence $AB=\mu^2$, $C$, $\mathfrak D=AB-C^2$, and the scale-free mismatch
$\Delta_{\rm rel}$ are invariant under every finite sequence of homogeneous
blockings.  For two identical midpoint intervals of duration $s$, rewritten as
one interval of duration $2s$,
\begin{equation}
 \boxed{
 A'=A+\frac{Bs^2}{4},\qquad
 B'=\frac{AB}{A+Bs^2/4},\qquad C'=C.}
 \label{eq:main-homogeneous-block-map}
\end{equation}
Consequently $\Delta_{\rm rel}>0$ remains strictly positive: the matched surface
$\Delta_{\rm rel}=0$ has no open basin of attraction for this conservative
blocking operation, even after changes of time/amplitude units or common
positive rescaling of the action.  On finite positive amplitude alphabets with a symmetric untwisted trial
kernel, the endpoint twist $C$ is retained exactly under finite blocking, and a controlled
low-temperature/grid limit recovers the same nonzero relative defect.
\end{theorem}

This is a no-go result for a specified natural coarse-graining operation, not a
no-go theorem for every possible renewal renormalization.  In particular, a
driven or conditioned law may define a different physical flow.  The theorem
shows that generic Einstein matching cannot be inferred merely from conservative
elimination of intermediate amplitudes.

\paragraph{Matching is an exact cross-scale consistency condition.}
The absence of an attracting flow does not leave the matched surface without
an intrinsic test.  The same blocking operation characterizes it by asking
whether local minimization commutes with the elimination of an intermediate
amplitude.

\begin{proposition}[Optimizer composition characterizes action--law matching]
\label{prop:main-optimizer-composition}
Let $m_\theta(x)$ minimize
$y\mapsto\mathcal E_{\mu,\theta,C}(x,y)$ over $y>0$ for $x>0$, and use
$\rho=C/\mu$ and $\Delta_{\rm rel}$ from
\eqref{eq:main-relative-rank-defect}.  Then
\begin{equation}
 m_\theta(x)=r_\theta x,\qquad
 r_\theta=(\cosh\theta+\rho\sinh\theta)^{-1},
 \label{eq:main-optimizer-map}
\end{equation}
and, for every $\theta,\phi>0$,
\begin{equation}
 \boxed{r_{\theta+\phi}^{-1}-(r_\theta r_\phi)^{-1}
       =\Delta_{\rm rel}\sinh\theta\sinh\phi.}
 \label{eq:main-optimizer-defect}
\end{equation}
Thus $AB=C^2$ is equivalent to
$m_{\theta+\phi}=m_\phi\circ m_\theta$ for every positive pair, and already
to that equality for one positive pair at one $x>0$.  On this branch
$r_\theta=e^{-\rho\theta}$ with $\rho=\pm1$.
\end{proposition}

The proof is in Appendix~\ref{supp:homogeneous-dynamics}.
Equation~\eqref{eq:main-optimizer-defect} gives a second finite diagnostic of
the relative defect, independent of its reconstruction from three cost
entries.  On compact calibrated ranges away from zero denominators it is
stable under small measurement errors.  It concerns the minimizing maps of
the reconstructed cost; identifying these with observed low-noise updates is
a further test of the accepted law.  It is a characterization of matching,
not a mechanism which makes an arbitrary law approach it.

\subsection{A nonlinear inhomogeneous two-polarization vacuum sector}
\label{subsec:gowdy-main}

The homogeneous calculation leaves open whether the finite evolution and
constraint mechanism survives genuinely spatially varying nonlinear dynamics.
A controlled test is provided by the unpolarized Gowdy $T^3$ vacuum class
\cite{Moncrief1981}.  On a smooth slab
$[t_0,t_1]\times\mathbb T^3$, $t_0>0$, write
\begin{equation}
 g=t^{-1/2}e^{\lambda/2}(-\dd t^2+\dd\theta^2)
 +t\left[e^P(\dd x+Q\dd y)^2+e^{-P}\dd y^2\right].
 \label{eq:main-gowdy-metric}
\end{equation}
Set
\begin{equation}
 a=P_t,\qquad b=P_\theta,\qquad
 c=e^P Q_t,\qquad d=e^P Q_\theta,
 \label{eq:main-gowdy-frame-vars}
\end{equation}
and
\begin{equation}
 e=\frac12(a^2+b^2+c^2+d^2),\qquad f=ab+cd.
 \label{eq:main-gowdy-ef}
\end{equation}
The first-order vacuum system is
\begin{align}
 a_t&=b_\theta-a/t+c^2-d^2, & b_t&=a_\theta,\nonumber\\
 c_t&=d_\theta-c/t-ac+bd, & d_t&=c_\theta+ad-bc,\nonumber\\
 P_t&=a, & Q_t&=e^{-P}c, & \lambda_t&=2te,
 \label{eq:main-gowdy-frame-system}
\end{align}
with coordinate and momentum constraints
\begin{equation}
 P_\theta=b,\qquad e^P Q_\theta=d,\qquad
 \lambda_\theta=2tf.
 \label{eq:main-gowdy-constraints}
\end{equation}
The identities
\begin{equation}
 e_t-f_\theta=-\frac{a^2+c^2}{t},\qquad
 f_t-e_\theta=-\frac ft
 \label{eq:main-gowdy-stress-identities}
\end{equation}
propagate the last constraint and preserve the zero mean of $tf$ required for
periodic $\lambda$.

For the variational test the areal gauge is imposed only after variation.  With
independent fields $(R,\lambda,P,Q,N,\beta)$, $R,N>0$, and
$v_u=\dot u-\beta u'$, use the boundary-marked ADM density
\begin{align}
 \mathscr L_{\rm G}={}&
 -\frac{v_R(v_\lambda-4\beta')}{2N}
 +\frac R{2N}(v_P^2+e^{2P}v_Q^2)\nonumber\\
 &+N\left[-2R''+\frac{R'\lambda'}2
 -\frac R2\{(P')^2+e^{2P}(Q')^2\}\right].
 \label{eq:main-gowdy-adm-density}
\end{align}
Its lapse and shift equations are independent constraints.  The solution of
\eqref{eq:main-gowdy-frame-system}--\eqref{eq:main-gowdy-constraints} lifts to
those equations with $R=t$, $N=1$, $\beta=0$; hence the finite action test below
is not obtained by first gauge-fixing away the lapse and shift variations.

\begin{theorem}[Nonlinear two-polarization renewal regulator]
\label{thm:main-gowdy-regulator}
Let $X_*(t)$ be a smooth solution of
\eqref{eq:main-gowdy-frame-system}--\eqref{eq:main-gowdy-constraints} on a fixed
slab $[t_0,t_1]\times\mathbb T^1$, with the spatial derivatives required below
uniformly bounded and with the reconstructed metric uniformly nondegenerate.
Let $h=\ell$ be the synchronized temporal and spatial mesh.  Then there is a
nearest-neighbour symmetric source/transport/source splitting with the following
properties.
\begin{enumerate}[label=\textup{(G\arabic*)},leftmargin=3em]
\item With
      $U=(u_1,u_2,u_3,u_4)=\sqrt t\,(a,b,c,d)$,
      $\mathcal J(U)=u_1u_2+u_3u_4$,
      $D_+=(S_+-I)/\ell$, and $\mathsf A_+=(I+S_+)/2$, the quantity
      \begin{equation}
       \boxed{\mathcal K_\ell
       =D_+\lambda-2\mathsf A_+\mathcal J(U)}
       \label{eq:main-gowdy-staggered-momentum}
      \end{equation}
      is preserved exactly by every deterministic full step.  The map is
      second-order consistent on the stated smooth chart.
\item If
      $X_{j+1}=\Phi_h(t_j,X_j)+r_j$ and
      \begin{equation}
       \mathcal E_h=\sum_j\frac{\|r_j\|_{2,\ell}^2}{2h},\qquad
       \mathcal F_h=\ell^{-(2r+1)}\mathcal E_h,
       \label{eq:main-gowdy-residual-energy}
      \end{equation}
      then, under the stated smooth-chart bootstrap or an enforced offset
      envelope,
      \begin{equation}
       \boxed{
       \max_j\|X_j-\mathsf S_\ell X_*(t_j)\|_{r,\infty,\ell}
       \le C\bigl(\epsilon_0+h^2+\sqrt{\mathcal F_h}\bigr).}
       \label{eq:main-gowdy-state-error}
      \end{equation}
      For $r\ge1$, the accumulated staggered-constraint error is bounded by its
      initial value plus $C\sqrt{\mathcal F_h}$.
\item On the finite offset alphabet with
      $\|\epsilon_j\|_\infty\le c_0h^4$ and initial $C_h^1$ error $O(h^2)$,
      the shared tensor-product cubic Hermite readout is pathwise $O(h^2)$ in
      $W^{1,\infty}$ and $O(h)$ in its second derivatives.  Consequently, on
      compact subslabs,
      \begin{equation}
       \boxed{
       \|g_h-g_*\|_{W^{1,\infty}}
       +\|D^2g_h-D^2g_*\|_{L^\infty}
       +\|\operatorname{Riem}(g_h)-\operatorname{Riem}(g_*)\|_{L^\infty}
       \le Ch,}
       \label{eq:main-gowdy-full-curvature}
      \end{equation}
      where the first term is in fact $O(h^2)$.  Thus
      $\sup_h\|\operatorname{Riem}(g_h)\|_{L^2(K)}<\infty$ and
      $\|\operatorname{Ric}(g_h)\|_{L^\infty(K)}\le Ch$ pathwise.  The
      Levi--Civita curvature of the limiting metric is therefore identified
      directly.
\item An independent midpoint/centered discretization of
      \eqref{eq:main-gowdy-adm-density}, with $N$ and $\beta$ varied
      independently before restriction to areal gauge, has first-variation
      residual tending to zero against every fixed smooth local test.  A finite
      positive offset/acceptance realization can be chosen with the support cap
      in item~\textup{(G3)} and obeys the exact Gibbs entropy--energy bound; its
      successful histories therefore inherit the pathwise curvature estimate.
\end{enumerate}
\end{theorem}

The construction and proof are given in Appendix~\ref{supp:gowdy}.  It is a
nonlinear and spatially inhomogeneous population result, but its scope remains
symmetry reduced and local to a smooth positive slab.  Unlike the general
certificate of \Cref{thm:main-curvature-propagation}, the curvature bound in
\eqref{eq:main-gowdy-full-curvature} is obtained directly from pathwise
second-jet control and requires no separately assumed all-time curvature budget
on this class.  It does not imply an unrestricted four-dimensional nonlinear
Cauchy theorem.

\paragraph{Measured access to the curvature-controlled histories.}
A finite alphabet with admissible offsets establishes realizability, but does
not by itself give a uniform probability of proposing those offsets.  The
following result separates this operational issue from the deterministic
curvature estimate.

\begin{proposition}[Finite proposal-access certificate]
\label{prop:main-safe-access}
At stage $j$, let $p_j$ be the measured finite proposal row and let
$\mathcal G_j$ contain only proposals satisfying the offset and solver bounds
of \Cref{thm:main-gowdy-regulator}.  Suppose
$p_j(\mathcal G_j)\ge\mu_j>0$ and $0\le E_j\le c_j$ on this set,
uniformly over the supported source histories.  Accept with the actual rule
$\one_{\mathcal G_j}\exp(-E_j/\varepsilon_j)$.  Then
\begin{equation}
 \mu_j e^{-c_j/\varepsilon_j}\le Z_j\le p_j(\mathcal G_j),
 \qquad
 D_{\rm KL}(Q_j\Vert p_j)\ge-\log p_j(\mathcal G_j).
 \label{eq:main-safe-proposal-bounds}
\end{equation}
With a genuine reset to the same source between retries, caps $R_j$ give
\begin{equation}
 \Pr(\hbox{some stage fails})
       \le\sum_j\exp[-R_j\mu_j e^{-c_j/\varepsilon_j}].
 \label{eq:main-safe-retry-bound}
\end{equation}
Every successful history inherits the pathwise curvature and independent
action-test estimates of \Cref{thm:main-gowdy-regulator}.  If the displayed
failure bounds are summable along a countable refinement, these conclusions
hold on every sufficiently fine successful refinement almost surely.
\end{proposition}

The proof in Appendix~\ref{supp:gowdy} also gives a sufficient proposal floor
from finite launch-share and completed-writer records.  A missing safe
proposal mass is a genuine obstruction which changing temperature cannot
remove.  These are bounds on trial counts; an imposed physical deadline must
also accommodate execution of the capped trials.  The selection rule and its
normalizer remain physical records, and its positive deviation cost is not
identified with the signed gravitational action.

\subsection{Flat and de Sitter geometric regulators}

For $\Lambda=0$, choose the constant stationary branch with $a_0=1$.
Then
\begin{equation}
 \mathcal B_h=I_3,
 \qquad \varrho_h=1,
 \qquad g_h=I_3,
 \qquad N_h=1,
 \qquad \beta_h=0,
 \label{eq:flat-adm-main}
\end{equation}
and the graph generators converge to $\tfrac12\Delta_{\mathbb T^3}$.  The
Lorentzian limit is
\begin{equation}
 \boxed{
 (\R\times\mathbb T^3,
 -\dd t^2+\dd x_1^2+\dd x_2^2+\dd x_3^2).}
 \label{eq:flat-metric-main}
\end{equation}
Its curvature and vacuum Einstein tensor vanish.

For $\Lambda>0$, choose the expanding stationary branch of
\Cref{cor:main-derived-desitter-rates} with $a_0=1$, $t_0=0$, and
$H=H_0=\sqrt{\Lambda/3}$.  Then
\begin{equation}
 \mathcal B_h(t)=e^{-2Ht}I_3,
 \qquad
 \varrho_h(t)=e^{3Ht},
 \qquad
 g_h(t)=e^{2Ht}I_3,
 \qquad
 N_h=1,
 \qquad
 \beta_h=0.
 \label{eq:desitter-adm-main}
\end{equation}
The continuum coframe and connection are
\begin{equation}
 \vartheta^0=\dd t,
 \qquad
 \vartheta^a=e^{Ht}\dd x^a,
 \qquad
 \omega^a{}_0=H\vartheta^a,
 \qquad
 \omega^a{}_b=0.
 \label{eq:desitter-cartan-main}
\end{equation}
Cartan's equations give zero torsion and constant sectional curvature $H^2$,
so
\begin{equation}
 \boxed{
 g_H=-\dd t^2+e^{2Ht}
 (\dd x_1^2+\dd x_2^2+\dd x_3^2),
 \qquad
 G_{\mu\nu}+3H^2(g_H)_{\mu\nu}=0.}
 \label{eq:desitter-einstein-main}
\end{equation}
The spatial quotient is globally hyperbolic and is the compact spatial quotient
of the flat de~Sitter slicing.

\begin{theorem}[Stationary flat and de Sitter renewal regulators]
\label{mt:vacuum}
For the operational predictive carrier and geometric branch of
\Cref{mt:grand-tensor,mt:adm}, the stationary homogeneous families above have
the following properties.
\begin{enumerate}[label=\textup{(\roman*)},leftmargin=2.8em]
\item their rate and density profiles solve the lapse-varied homogeneous renewal
      Euler system of \Cref{thm:main-homogeneous-friedmann} and arise as the
      continuum limits of the finite stationary system in
      \Cref{thm:main-finite-homogeneous-stationarity};
\item each spatial cutoff defines a normalized finite $A_3$ renewal law whose
      root moments and speed densities reconstruct the displayed ADM variables
      without a primitive metric;
\item the flat sequence converges to \eqref{eq:flat-metric-main}, with
      $\Lambda=0$, while the curved sequence converges to
      \eqref{eq:desitter-einstein-main}, with $\Lambda=3H^2$;
\item on every compact time slab the cut, Poincar\'e, low-energy counting, and
      compact-screen margins are uniform in $h$;
\item the flat curvature vanishes identically, while on each compact de~Sitter
      slab the reconstructed curvature and observer/transport coefficients are
      uniformly bounded; the initial-energy and integrated propagation/source
      budgets of \Cref{thm:main-curvature-propagation} therefore pass directly,
      and in fact the curvature is uniformly bounded in $L^\infty$;
\item along $h_n=2^{-n^2}$ with time mesh $h_n$ and screen $[-n,n]$, exact
      parallel transport of the reconstructed Levi--Civita connection gives
      metric-unitary links, while the normalized torsion, plaquette-curvature,
      and Palatini first-variation defects are $O_T(h_n)$ and summable.
\end{enumerate}
In addition, \Cref{thm:main-explicit-operational-family} supplies a nontrivial
finite positive accept/reject realization whose measured action and accepted
dynamics converge to the same expanding stationary branch.
\end{theorem}

The detailed homogeneous variational and operational proofs are in
Appendix~\ref{supp:homogeneous-dynamics}; the spatial generator, compact-screen,
Cartan-curvature, and exact-link estimates are in Appendix~\ref{supp:vacuum}.
The result is not an attractor theorem: it does not select $\Lambda$, the time
orientation, or the quadratic operational class from arbitrary microscopic
renewal data.



\section{Outlook}
\label{sec:discussion}

The construction developed here suggests a different order for the relation
between dynamics and spacetime.  In the conventional picture, stochastic or
quantum processes take place within an already defined spacetime.  Here the
logical order is reversed.  The microscopic description is a renewal process
of operational histories; predictive equivalence determines which distinctions
persist; stable relational statistics reconstruct geometry; and only at that
stage does a Lorentzian spacetime description become available.

The key conceptual object is therefore not a microscopic position or a
pre-existing geometric state, but a distinction that remains effective under
future continuation.  Two histories are identified when no admissible future
can distinguish them.  What survives this quotient is the persistent predictive
content of the process.  In this sense, identity is determined by future
consequence rather than by microscopic origin.  The relational structure on
which geometry is eventually built consists precisely of those distinctions
that prediction cannot forget.

This changes the interpretation of the metric.  The metric is not attached to
the microscopic events as a primitive label, nor is the renewal process assumed
to explore an underlying manifold whose geometry already exists.  Instead,
the ADM variables are reconstructed from collective statistics of the renewal
process.  The continuum metric is therefore naturally interpreted as a stable
large-scale organization of those statistics.  Spacetime geometry is not the
arena sampled by the microscopic dynamics; it is an effective structure
acquired by the dynamics when its predictive relations become sufficiently
regular.

The Lorentzian character of this geometry has a correspondingly precise
status.  The positive predictive structure and the root-race bracket determine
a positive spatial geometry, and a positive Gram construction cannot create a
negative direction.  The protected binary orientation is therefore not a
hidden metric assumption but an irreducible signed datum: explicit finite
separating models show that it cannot be reconstructed from the other positive
operational rows, even after common refinement.  Once the spatial module has
dimension three, adjoining exactly this one signed direction gives Lorentzian
inertia $(1,3)$ up to global sign.  The Clifford cell used in the construction
then realizes that already fixed quadratic form; it does not manufacture the
signature by representation design.

A second consequence is that emergent spacetime need not stop at kinematics.
The constructive and certification results play complementary roles.  The
open $3+1$ theorem proves a cutoff-uniform nonsymmetric Einstein population
from a specified finite writer and its local energy estimates; the general
certification theorem tests arbitrary complete cylinders by finite/cofinal
action, noncollapse, and curvature records.  A common action is reconstructed
independently from the renewal record, and the observed evolution can then be
tested for stationarity with respect to that action.  On the gravitational branch, this variational structure converges to
the vacuum Einstein equation.  The significance is that Einstein dynamics is
not introduced merely because a reconstructed metric happens to resemble a
known solution.  Geometry and dynamics arise from different layers of the same
operational structure and are linked only after their consistency has been
tested.  The curvature estimate has the same operational character: it
propagates an initial energy budget and records how action errors are amplified
when they enter the differentiated geometric current.  Failure is attached to
that explicit comparison, rather than hidden in an unspecified regular
continuum limit.

This perspective suggests a relational interpretation of spacetime.  Spacetime
is not treated as a container that exists before physical events and provides
them with positions, distances, and causal relations.  The more primitive
notion is continuation: some distinctions in a history remain capable of
affecting what can happen next, while others become irrelevant.  The stable
organization of the surviving distinctions defines the relational carrier, and
its coarse-grained statistical structure becomes geometric.  Smooth Lorentzian
spacetime is then the macroscopic form taken by a sufficiently regular
predictive renewal process.

The framework should not be read as claiming that every renewal process
generates general relativity, or that the present operational starting point is
necessarily the deepest possible description.  The construction identifies a
particular class in which predictive persistence, Lorentzian geometry, and
vacuum gravitational dynamics can all be reconstructed without assuming a
background spacetime.  Other renewal processes may fail to enter this class,
or may give rise to different effective structures.  Moreover, the conservative
homogeneous blocking tested in \Cref{thm:main-no-blocking-attractor} leaves the
scale-free matching defect invariant, so the matched Einstein surface is not an
attractor of that operation.  Composition of the associated local minimizing
maps nevertheless gives an exact test for matching.  This separates a
consistency principle from a dynamical attraction mechanism: broader
universality would require a genuinely different physical coarse-graining or
conditioning flow.

The dimensional result should likewise be read as a pair of structural
classifications rather than as a universal consequence of positivity.  The
logically weaker endpoint--cycle balance theorem already singles out $K_4$
among connected simple regular relations and uses no alternating-form parity
condition.  The alternating/minimal-source simplex theorem independently
selects the same cell, while the representation-theoretic identity
$\Lambda^2W_4\cong W_4\otimes\mathrm{sgn}$ explains why only at four endpoints
does the complete cycle source become one oriented copy of the spatial source.
The agreement of these criteria is the dimensional statement of the paper; no
claim is made that arbitrary higher-arity pregeometric categories must satisfy
either criterion.

The dynamical evidence is no longer restricted to the homogeneous and Gowdy
sectors or to a single local update.  The writer of
\Cref{thm:main-open-3plus1} gives a cutoff-independent nonlinear Cauchy
interval for all ten metric components and arbitrary three-dimensional
dependence in a fixed Sobolev ball, while
\Cref{thm:main-open-law-basin} places that writer inside a fixed open
$55$-dimensional family with the same whole-sequence Einstein limit.  The same
result gives relatively open neighborhoods of history-dependent finite
transition rows in a complete residual-profile topology and exact finite
action-prepared initial constraint charts from free tensor records.  Thus the
constructive population is robust in differential-law, probability-row and
initial-constraint directions, although not under arbitrary leading-order
forcing or unrestricted rapid law switching.  The homogeneous sector remains
useful because its complete lapse variation and finite acceptance law can be
written in closed form, while the Gowdy sector supplies a nonlinear
large-amplitude symmetry-reduced test with pathwise full-curvature convergence
and exact staggered momentum propagation.

A sharper distinction concerns dynamical selection.  The stable local writer
is an explicitly specified action-compatible realization; it is not identified
with the raw exact finite Euler flow.  The exact-action results
\Cref{thm:main-exact-action-response,prop:main-exact-action-tangency} and
Appendix~\ref{supp:exact-action-audit} now go substantially beyond a static
rank audit.  A transverse source block is exactly regular at every odd cutoff,
the complete signed response separates into positive source reproduction and
a scalar--trace correction, and selected original-action preparations have
zero leading weak target with small literal initial curvature.  On a further
prepared family the active trace resonance is absent and the actual transverse
geometric response decays with refinement.  The remaining obstruction is
dynamical: an exactly prepared branch can fail a higher tangency condition
under its own Euler flow, producing a provable initial layer.  That defect is
not removable by leading multiplier retuning at fixed first jet, but it changes
under admissible canonical first-jet deformation.  The direct exact-action
problem is therefore to find a regular nonlinear root of the complete
transported tangency system and propagate its full multiplier/curvature
estimates on a cutoff-independent physical interval.

Several broader questions now become natural.  The local existence of an open
law-space neighborhood is established here; the remaining universality
question is how far that neighborhood can be enlarged, whether a physical
coarse-graining or conditioning flow selects it from more general microscopic
laws, and whether the independently reconstructed exact action admits a
cofinal regular tangency branch whose own flow enters the same Einstein
universality class.  A second question is how matter and gauge degrees of
freedom should enter the same predictive framework, and whether their
effective actions and couplings can likewise be reconstructed rather than
imposed.  A third concerns residual memory: the microscopic theory need not be
Markovian, while the continuum regimes considered here become effectively
local after predictive compression.  Understanding controlled departures from
this limit may provide a way to characterize non-Markovian corrections to
ordinary local field theory.

There is also a more basic question upstream of the present formalism.  The
construction begins from histories, but its canonical quotient suggests that
histories themselves may not be the most fundamental objects.  What matters is
whether a distinction changes the space of possible futures.  This points
toward a more primitive notion of \emph{effective difference under
continuation}: distinctions are physically retained insofar as they continue
to affect subsequent possibilities, and are identified once their future
effects become identical.  Renewal history is then a way of organizing these
effective differences, while spacetime is the stable geometry they acquire at
large scale.

Within the scope established here, the conceptual chain is therefore


\[
\boxed{
\begin{aligned}
\text{renewal histories}
&\longrightarrow
\text{predictive persistence}
\longrightarrow
\text{relational structure}\\
&\longrightarrow
\text{Lorentzian geometry}
\longrightarrow
\text{gravitational dynamics}.
\end{aligned}
}
\]



The central conclusion is not that spacetime has been eliminated from physics,
but that it need not occupy the most fundamental level of description.
Lorentzian geometry and vacuum Einstein dynamics can instead appear as the
stable large-scale organization of a more primitive predictive process.
Spacetime is then not the container in which renewal occurs, but the form that
renewal acquires when persistent distinctions, orientation, and dynamics become
coherent.

% =============================================================================
% =============================================================================

\section*{Data and formalization availability}
Selected finite algebraic components of the construction have been formalized in
\texttt{Lean}.  Source code is available at
\url{https://github.com/AI-MathPhys/renewal-geometry}.  The analytic continuum
arguments are those given in the present article and its appendices.

\begin{thebibliography}{99}

\bibitem{bombelli-sorkin}
L.~Bombelli, J.~Lee, D.~Meyer, and R.~D.~Sorkin,
\newblock Space-time as a causal set,
\newblock \emph{Physical Review Letters} \textbf{59} (1987), 521--524.

\bibitem{surya}
S.~Surya,
\newblock The causal set approach to quantum gravity,
\newblock \emph{Living Reviews in Relativity} \textbf{22} (2019), 5.

\bibitem{konopkamarkopouloseverini2008}
T.~Konopka, F.~Markopoulou, and S.~Severini,
\newblock Quantum graphity: a model of emergent locality,
\newblock \emph{Physical Review D} \textbf{77} (2008), 104029.

\bibitem{ambjornjurkiewiczloll2009}
J.~Ambj{\o}rn, J.~Jurkiewicz, and R.~Loll,
\newblock Quantum gravity, or the art of building spacetime,
\newblock in \emph{Approaches to Quantum Gravity},
D.~Oriti, ed.,
Cambridge University Press, Cambridge, 2009, pp.~341--359.

\bibitem{oriti2009}
D.~Oriti, ed.,
\newblock \emph{Approaches to Quantum Gravity: Toward a New Understanding of
Space, Time and Matter},
\newblock Cambridge University Press, Cambridge, 2009.

\bibitem{Connes1994}
A.~Connes,
\newblock \emph{Noncommutative Geometry},
\newblock Academic Press, San Diego, 1994.

\bibitem{jacobson1995}
T.~Jacobson,
\newblock Thermodynamics of spacetime: the Einstein equation of state,
\newblock \emph{Physical Review Letters} \textbf{75} (1995), 1260--1263.

\bibitem{verlinde2011}
E.~Verlinde,
\newblock On the origin of gravity and the laws of Newton,
\newblock \emph{Journal of High Energy Physics} \textbf{2011} (2011), no.~4,
029.

\bibitem{vanraamsdonk2010}
M.~Van Raamsdonk,
\newblock Building up spacetime with quantum entanglement,
\newblock \emph{General Relativity and Gravitation} \textbf{42} (2010),
2323--2329.

\bibitem{shalizicrutchfield}
C.~R.~Shalizi and J.~P.~Crutchfield,
\newblock Computational mechanics: pattern and prediction, structure and
simplicity,
\newblock \emph{Journal of Statistical Physics} \textbf{104} (2001),
817--879.

\bibitem{TraversCrutchfield2014}
N.~F.~Travers and J.~P.~Crutchfield,
\newblock Equivalence of history and generator $\epsilon$-machines,
\newblock \emph{Physical Review E} \textbf{89} (2014), 062107.

\bibitem{MarzenCrutchfield2017}
S.~E.~Marzen and J.~P.~Crutchfield,
\newblock Structure and randomness of continuous-time, discrete-event
processes,
\newblock \emph{Journal of Statistical Physics} \textbf{169} (2017),
303--315.

\bibitem{JurgensCrutchfield2021}
A.~M.~Jurgens and J.~P.~Crutchfield,
\newblock Divergent predictive states: the statistical complexity dimension
of stationary, ergodic hidden Markov processes,
\newblock \emph{Chaos} \textbf{31} (2021), 083114.

\bibitem{feller}
W.~Feller,
\newblock \emph{An Introduction to Probability Theory and Its Applications,
Vol.~II},
\newblock Wiley, New York, 1971.

\bibitem{garsialamperti1962}
A.~M.~Garsia and J.~Lamperti,
\newblock A discrete renewal theorem with infinite mean,
\newblock \emph{Commentarii Mathematici Helvetici} \textbf{37} (1962),
221--234.

\bibitem{erickson1970}
K.~B.~Erickson,
\newblock Strong renewal theorems with infinite mean,
\newblock \emph{Transactions of the American Mathematical Society}
\textbf{151} (1970), 263--291.

\bibitem{caravennadoney2019}
F.~Caravenna and R.~Doney,
\newblock Local large deviations and the strong renewal theorem,
\newblock \emph{Electronic Journal of Probability} \textbf{24} (2019),
Paper No.~72, 1--48.

\bibitem{Stinespring1955}
W.~F.~Stinespring,
\newblock Positive functions on $C^*$-algebras,
\newblock \emph{Proceedings of the American Mathematical Society}
\textbf{6} (1955), 211--216.

\bibitem{Choi1975}
M.-D.~Choi,
\newblock Completely positive linear maps on complex matrices,
\newblock \emph{Linear Algebra and its Applications} \textbf{10} (1975),
285--290.

\bibitem{strohmaier}
A.~Strohmaier,
\newblock On noncommutative and pseudo-Riemannian geometry,
\newblock \emph{Journal of Geometry and Physics} \textbf{56} (2006),
175--195.

\bibitem{franco}
N.~Franco,
\newblock Temporal Lorentzian spectral triples,
\newblock \emph{Reviews in Mathematical Physics} \textbf{26} (2014),
1430007.

\bibitem{franco-eckstein}
N.~Franco and M.~Eckstein,
\newblock An algebraic formulation of causality for noncommutative geometry,
\newblock \emph{Classical and Quantum Gravity} \textbf{30} (2013), 135007.

\bibitem{vandendungenpaschkerennie2013}
K.~van den Dungen, M.~Paschke, and A.~Rennie,
\newblock Pseudo-Riemannian spectral triples and the harmonic oscillator,
\newblock \emph{Journal of Geometry and Physics} \textbf{73} (2013),
37--55.

\bibitem{vandendungen2016}
K.~van den Dungen,
\newblock Krein spectral triples and the fermionic action,
\newblock \emph{Mathematical Physics, Analysis and Geometry}
\textbf{19} (2016), 4.

\bibitem{DevastatoFarnsworthLizziMartinetti2018}
A.~Devastato, S.~Farnsworth, F.~Lizzi, and P.~Martinetti,
\newblock Lorentz signature and twisted spectral triples,
\newblock \emph{Journal of High Energy Physics} \textbf{2018} (2018),
no.~3, 089.

\bibitem{DevastatoFilaciMartinettiSingh2020}
A.~Devastato, M.~Filaci, P.~Martinetti, and D.~Singh,
\newblock Actions for twisted spectral triple and the transition from the
Euclidean to the Lorentzian,
\newblock \emph{International Journal of Geometric Methods in Modern Physics}
\textbf{17} (2020), suppl.~1, 2030001.

\bibitem{nieuviarts2026a}
G.~Nieuviarts,
\newblock Emergence of time from a twisted spectral triple in
almost-commutative geometry,
\newblock \emph{Journal of Noncommutative Geometry} (2026), to appear;
arXiv:2512.15450.

\bibitem{nieuviarts2026b}
G.~Nieuviarts,
\newblock Emergence of Lorentz symmetry from an almost-commutative twisted
spectral triple,
\newblock \emph{International Journal of Geometric Methods in Modern Physics}
(2026), to appear; arXiv:2502.18105.

\bibitem{Bregman1967}
L.~M.~Bregman,
\newblock The relaxation method of finding the common point of convex sets and
its application to the solution of problems in convex programming,
\newblock \emph{USSR Computational Mathematics and Mathematical Physics}
\textbf{7} (1967), 200--217.

\bibitem{Chung1997}
F.~R.~K.~Chung,
\newblock \emph{Spectral Graph Theory},
\newblock CBMS Regional Conference Series in Mathematics, vol.~92,
American Mathematical Society, Providence, RI, 1997.

\bibitem{SaloffCoste2002}
L.~Saloff-Coste,
\newblock \emph{Aspects of Sobolev-Type Inequalities},
\newblock London Mathematical Society Lecture Note Series, vol.~289,
Cambridge University Press, Cambridge, 2002.

\bibitem{ACY1998}
A.~Anderson, Y.~Choquet-Bruhat, and J.~W.~York, Jr.,
\newblock Einstein--Bianchi hyperbolic system for general relativity,
\newblock arXiv:gr-qc/9710041, version 2 (1998).

\bibitem{KRS2015}
S.~Klainerman, I.~Rodnianski, and J.~Szeftel,
\newblock The bounded $L^2$ curvature conjecture,
\newblock \emph{Inventiones Mathematicae} \textbf{202} (2015), 91--216.

\bibitem{York1972}
J.~W.~York, Jr.,
\newblock Role of conformal three-geometry in the dynamics of gravitation,
\newblock \emph{Physical Review Letters} \textbf{28} (1972), 1082--1085.

\bibitem{GibbonsHawking1977}
G.~W.~Gibbons and S.~W.~Hawking,
\newblock Action integrals and partition functions in quantum gravity,
\newblock \emph{Physical Review D} \textbf{15} (1977), 2752--2756.

\bibitem{MarsdenWest2001}
J.~E.~Marsden and M.~West,
\newblock Discrete mechanics and variational integrators,
\newblock \emph{Acta Numerica} \textbf{10} (2001), 357--514.

\bibitem{Moncrief1981}
V.~Moncrief,
\newblock Global properties of Gowdy spacetimes with $T^3\times\mathbb R$ topology,
\newblock \emph{Annals of Physics} \textbf{132} (1981), 87--107.

\end{thebibliography}



\clearpage
\appendix

% =============================================================================
\section{Operational reconstruction and positive history algebra}
\label{supp:grand-reconstruction}
% =============================================================================

This section proves \Cref{mt:grand-tensor}.  All operational alphabets and
exposed operator spaces are finite-dimensional.  The construction begins from
histories and conditional future statistics; no generated ledger or persistent
state space is supplied at the entrance.

\subsection{Finite typed operational renewal datum}

\begin{definition}[Finite typed event grammar]
\label{def:supp-event-grammar}
At cutoff $X$, a finite typed event grammar consists of:
\begin{enumerate}[label=\textup{(E\arabic*)},leftmargin=3em]
\item a finite set of cut types $\mathsf O_X$;
\item finite sets $\Sigma_X(x,y)$ of primitive event letters $a:x\to y$;
\item a prefix-closed set $\mathsf H_{X,x}$ of reachable finite typed histories
      ending at each cut type $x$, with typed concatenation and empty words;
\item finite preparation, intervention-outcome, and terminal-Read alphabets,
      together with finite exposed operator ports;
\item a distinguished protected binary orientation record whose two values are
      separated by the declared physical Read policy.
\end{enumerate}
No state label, persistent relation object, metric, scheduler, connection, or
spacetime semantic is assigned to a history.
\end{definition}

For $h\in\mathsf H_{X,x}$ let $\Sigma_{X,h}$ be the event letters admissible
after $h$.

\begin{definition}[Finite conditional operational renewal law]
\label{def:supp-operational-datum}
A finite conditional operational renewal law on
\Cref{def:supp-event-grammar} assigns
\begin{equation}
 p_X(a\mid h)\in[0,1],
 \qquad
 \sum_{a\in\Sigma_{X,h}}p_X(a\mid h)=1,
 \label{eq:supp-one-step-law}
\end{equation}
for every reachable history.  For an admissible future word
$w=a_1\cdots a_k$ define
\begin{equation}
 p_X(w\mid h)
 =\prod_{j=1}^{k}
 p_X(a_j\mid ha_1\cdots a_{j-1}),
 \qquad p_X(\epsilon_x\mid h)=1.
 \label{eq:supp-chain-law}
\end{equation}
The physical intervention letters include, on every exposed Hermitian
Choi-coordinate space, a finite physical $\CP$ frame whose real span is the
full Hermitian space.  The induced multitime frame probabilities are affine
under randomization, additive under coarse graining, normalized on complete
instruments, causal under deterministic discard, compatible on common
prefixes, and satisfy the finite positivity conditions of
\Cref{thm:supp-comb-tomography}.  Actual bounded terminal Reads, marks, and
readout queries may be included as terminal event outcomes; equivalently they
may be appended to the future test bank.
\end{definition}

The distinguished renewal semantics used in the geometric branch is carried by
the admitted completion/return event family.  No ordinary reset state is
needed for the predictive construction below; if a stronger reset law is
imposed, it is an additional property of the same event process rather than a
state-space axiom.

\subsection{Future equivalence and the derived persistent carrier}

\begin{definition}[Complete predictive future signature]
\label{def:supp-future-signature}
For $h\in\mathsf H_{X,x}$ define
\begin{equation}
 \Sigma_X(h)=\bigl(p_X(w\mid h)\bigr)_{w\in\mathsf F_{X,x}},
 \label{eq:supp-future-signature}
\end{equation}
where $\mathsf F_{X,x}$ is the complete finite future-test bank starting at
$x$.  Histories $h,h'\in\mathsf H_{X,x}$ are future equivalent when
\begin{equation}
 \boxed{h\sim_Xh'
 \quad\Longleftrightarrow\quad
 p_X(w\mid h)=p_X(w\mid h')
 \text{ for every }w\in\mathsf F_{X,x}.}
 \label{eq:supp-future-equivalence}
\end{equation}
Histories of different cut type are never identified.
\end{definition}

\begin{theorem}[History-level right congruence]
\label{thm:supp-right-congruence}
Future equivalence is a typed right congruence.  If $h\sim_Xh'$ and
$a:x\to y$ is admissible, then
\begin{equation}
 p_X(a\mid h)=p_X(a\mid h'),
 \label{eq:supp-one-letter-equality}
\end{equation}
and on the positive-probability branch
\begin{equation}
 \boxed{ha\sim_Xh'a.}
 \label{eq:supp-right-congruence}
\end{equation}
\end{theorem}

\begin{proof}
Use the one-letter future $a$ in
\eqref{eq:supp-future-equivalence}.  If the common probability is positive,
then for any continuation $v$ the chain rule gives
\[
 p_X(v\mid ha)
 =\frac{p_X(av\mid h)}{p_X(a\mid h)}
 =\frac{p_X(av\mid h')}{p_X(a\mid h')}
 =p_X(v\mid h'a).
\]
Hence the complete future signatures agree after appending $a$.
\end{proof}

\begin{definition}[Minimum predictive state]
\label{def:supp-predictive-state}
For every cut type $x$ put
\begin{equation}
 \boxed{
 Z_{X,x}^{\min}:=\mathsf H_{X,x}/\!\sim_X,
 \qquad
 q_{X,x}:\mathsf H_{X,x}\twoheadrightarrow Z_{X,x}^{\min}.}
 \label{eq:supp-predictive-state}
\end{equation}
For an admissible event $a:x\to y$ define
\begin{equation}
 \boxed{
 T_{X,a}([h])=[ha],
 \qquad
 \kappa_{X,a}([h])=p_X(a\mid h).}
 \label{eq:supp-predictive-transition}
\end{equation}
\end{definition}

\begin{corollary}[Derived predictive state machine]
\label{cor:supp-derived-state-machine}
The maps in \eqref{eq:supp-predictive-transition} are well defined.  If
$w=a_1\cdots a_k$, $z_0=q_X(h)$, and
$z_j=T_{X,a_j}(z_{j-1})$, then
\begin{equation}
 \boxed{
 p_X(w\mid h)=\prod_{j=1}^{k}\kappa_{X,a_j}(z_{j-1}).}
 \label{eq:supp-state-factorization}
\end{equation}
Thus the predictive state machine is reconstructed from the conditional event
law.
\end{corollary}

\begin{proof}
Well-definedness is \Cref{thm:supp-right-congruence}; the product formula is
\eqref{eq:supp-chain-law} written on future-equivalence classes.
\end{proof}

\begin{definition}[Deterministic predictive realization]
\label{def:supp-predictive-realization}
A reachable deterministic predictive realization consists of finite sets
$S_{X,x}$, history maps $s_x:\mathsf H_{X,x}\to S_{X,x}$, partial updates
$u_a:S_{X,x}\to S_{X,y}$, and branch functions
$r_a:S_{X,x}\to[0,1]$ such that on every reachable branch
\begin{equation}
 s_y(ha)=u_a(s_x(h)),
 \qquad
 p_X(a\mid h)=r_a(s_x(h)).
 \label{eq:supp-realization}
\end{equation}
\end{definition}

\begin{theorem}[Canonical minimum predictive realization]
\label{thm:supp-minimal-record}
The quotient $Z_X^{\min}$ is the unique coarsest reachable deterministic
predictive realization.  For every reachable realization $S_X$ there is a
unique surjection
\begin{equation}
 \boxed{\pi_{X,x}:S_{X,x}\twoheadrightarrow Z_{X,x}^{\min}}
 \label{eq:supp-universal-surjection}
\end{equation}
such that
\begin{equation}
 q_x=\pi_xs_x,
 \qquad
 \pi_yu_a=T_{X,a}\pi_x,
 \qquad
 r_a=\kappa_{X,a}\pi_x.
 \label{eq:supp-universal-intertwining}
\end{equation}
Consequently a generated ledger or hidden-state presentation, when present,
is only an intermediate predictive realization and has a unique quotient onto
$Z_X^{\min}$.
\end{theorem}

\begin{proof}
If $s_x(h)=s_x(h')$, induction on future-word length using
\eqref{eq:supp-realization} gives $p_X(w\mid h)=p_X(w\mid h')$ for every
future word $w$.  Hence $h\sim_Xh'$, so
$\pi_x(s_x(h)):=[h]$ is well defined.  Reachability makes $\pi_x$ unique and
surjective.  The intertwining identities follow from the definitions, and the
quotient itself is a realization by
\Cref{cor:supp-derived-state-machine}.
\end{proof}

\begin{theorem}[Persistence derived from composition]
\label{thm:supp-persistence-derived}
Let $A_1,A_2,\ldots$ be the realized event sequence, let
$H_n=A_1\cdots A_n$, and put $Z_n=q_X(H_n)$.  Then on every realized branch
\begin{equation}
 \boxed{Z_{n+1}=T_{X,A_{n+1}}(Z_n).}
 \label{eq:supp-persistence}
\end{equation}
Consequently
\begin{equation}
 \boxed{\mathfrak p_{X,n}^{\min}=(Z_n,A_{n+1},Z_{n+1})}
 \label{eq:supp-derived-packet}
\end{equation}
is a persistent relational event packet derived from the event process.  Every
other deterministic persistent predictive packet maps uniquely onto
\eqref{eq:supp-derived-packet}.
\end{theorem}

\begin{proof}
By definition,
$Z_{n+1}=q_X(H_nA_{n+1})=T_{X,A_{n+1}}q_X(H_n)$.  For any other predictive
realization, apply the universal map of
\Cref{thm:supp-minimal-record} at consecutive cuts.
\end{proof}

\begin{corollary}[Future-operational transfer]
\label{cor:supp-future-operational-transfer}
Let $\mathcal D$ be any downstream finite readout whose value is unchanged
whenever two candidate states have the same complete declared future
signature.  Then there is a unique $\overline{\mathcal D}$ on
$Z_X^{\min}$ such that
\begin{equation}
 \boxed{\mathcal D=\overline{\mathcal D}\circ\pi_X.}
 \label{eq:supp-compiler-factorization}
\end{equation}
Thus deleting every future-null state coordinate preserves such downstream
outputs exactly.
\end{corollary}

\begin{proof}
The readout is constant on the fibres of the universal surjection
\eqref{eq:supp-universal-surjection}; define $\overline{\mathcal D}$ on a
fibre by the common value.  Surjectivity gives uniqueness.
\end{proof}

\subsection{Conservative readable relational completion}
\label{supp:readable-relational-completion}

The predictive quotient above determines persistence but does not by itself
assert that the quotient class, scheduler path, orientation, and calibration
are jointly readable at one cut.  We now isolate the finite conservative
extension used by the geometric reconstruction.

Let $\Phi_{a,*}$ be the support-minimal Schr\"odinger-picture branch map for an
underlying event letter $a$.  Let $z\in Z_X^{\min}$ be its predictive source
class.  For a finite family of added same-history marks $m=(\gamma,\xi,z')$,
where $\gamma$ is a scheduler-path record, $\xi$ contains orientation and
calibration marks, and $z'$ is the next readable predictive class, let
\begin{equation}
 K_{X,a}(m\mid z)\ge0,
 \qquad
 \sum_m K_{X,a}(m\mid z)=1.
 \label{eq:supp-relational-kernel}
\end{equation}
On block states $\rho=(\rho_z)_z$, define
\begin{equation}
 \boxed{
 [\widetilde\Phi_{a,m,*}(\rho)]_{z'}
 =\sum_z K_{X,a}(m\mid z)\,\Phi_{a,*}(\rho_z).}
 \label{eq:supp-relational-refinement}
\end{equation}
The symbol in \eqref{eq:supp-relational-refinement} denotes a same-history
conditional extension, not a tensor-product decomposition of the state space.

\begin{theorem}[Conservativity and future minimality of the readable packet]
\label{thm:supp-readable-relational-completion}
For the finite extension \eqref{eq:supp-relational-refinement}:
\begin{enumerate}[label=\textup{(C\arabic*)},leftmargin=3em]
\item every refined branch is completely positive and the refined family is
      normalized;
\item if $\mathsf F_X(\rho)=\sum_z\rho_z$ forgets the explicit predictive block
      and all added marks, then
      \begin{equation}
       \boxed{
       \mathsf F_X\!\left(\sum_m\widetilde\Phi_{a,m,*}(\rho)\right)
       =\Phi_{a,*}(\mathsf F_X\rho);}
       \label{eq:supp-relational-forgetful}
      \end{equation}
      hence every underlying cylinder probability and branch map is recovered
      exactly after forgetting the readable enrichment;
\item quotienting raw readable packets by equality of all admitted future
      continuations, readout queries, and terminal Reads terminates at a unique
      coarsest finite right congruence $\mathfrak P_X^{\rm rel,min}$;
\item every other finite readable presentation sufficient for the same future
      readout has a unique surjection onto
      $\mathfrak P_X^{\rm rel,min}$, so unread refinements add no predictive
      content or source dimension.
\end{enumerate}
Consequently the readable relation packet is a conservative physical
realization of the already derived predictive carrier, not an independently
postulated state space.
\end{theorem}

\begin{proof}
Complete positivity follows because every target block in
\eqref{eq:supp-relational-refinement} is a nonnegative scalar combination of
completely positive branch maps.  Summing
\eqref{eq:supp-relational-kernel} over the added marks gives one, and therefore
summing the refined branches proves
\eqref{eq:supp-relational-forgetful}.  Induction on word length gives exact
preservation of every underlying finite cylinder.

For future minimality, define two raw readable packets to be equivalent when
all declared future signatures agree.  Prefixing the same admissible event to
two equal future signatures leaves them equal, so the relation is a typed
right congruence.  Finite partition refinement terminates because every strict
round increases the number of classes.  The usual fibre argument then gives
its universal property: any other sufficient readable presentation identifies
only packets with equal complete future signatures and therefore surjects
uniquely onto the stable quotient.
\end{proof}

\begin{remark}[What the completion adds]
Same-cut readability, scheduler records, protected orientation, and absolute
calibration are actual occurrence rows.  The predictive class appearing in the
packet is nevertheless the quotient already derived from the operational law.
The conservative completion therefore separates physical incidence from
predictive persistence without inserting metric, coframe, connection,
curvature, or spacetime dimension.
\end{remark}

\subsection{The state-conditioned Hankel quotient}

Fix a minimum predictive state sector $(x,z)$.  Let
$\mathsf P_{x,z}$ and $\mathsf F_{x,z}$ denote the finite pasts and futures
meeting at that sector, and set
\begin{equation}
 \mathbb p_{x,z}(f,p)=\Pr(f\circ p).
 \label{eq:supp-Hankel-table}
\end{equation}
For each past $p$ define $h_p(f)=\mathbb p_{x,z}(f,p)$ and
\begin{equation}
 M_{x,z}=\Span_{\R}\{h_p:p\in\mathsf P_{x,z}\}.
 \label{eq:supp-hankel-space}
\end{equation}

\begin{theorem}[Canonical typed Hankel realization]
\label{thm:supp-hankel}
For every primitive event
$a:(x,z)\to(y,T_{X,a}z)$ there is a unique linear map
\begin{equation}
 A_a:M_{x,z}\longrightarrow M_{y,T_{X,a}z},
 \qquad A_ah_p=h_{a\circ p}.
 \label{eq:supp-hankel-transition}
\end{equation}
Every future $f$ defines $\lambda_f(m)=m(f)$ and
\begin{equation}
 \mathbb p(f,p)=\lambda_f(h_p),
 \qquad
 \lambda_fA_a=\lambda_{f\circ a}.
 \label{eq:supp-hankel-evaluation}
\end{equation}
The realization is reachable and future separated.  For any other linear
realization $N_{x,z}$,
\begin{equation}
 N_{x,z}^{\mathrm{reach}}/N_{x,z}^{\mathrm{unobs}}
 \cong M_{x,z};
 \label{eq:supp-hankel-minimality}
\end{equation}
therefore every linear realization has dimension at least $\dim M_{x,z}$,
with equality exactly on the reachable future-separated branch.
\end{theorem}

\begin{proof}
If $\sum_jc_jh_{p_j}=0$, then for every target future $f$,
\[
 \sum_jc_jh_{a\circ p_j}(f)
 =\sum_jc_j\mathbb p(f\circ a,p_j)=0.
\]
Hence appending $a$ respects every column relation and
\eqref{eq:supp-hankel-transition} is well defined.  Reachability is built into
\eqref{eq:supp-hankel-space}; every nonzero response function is detected by
some future.  For another realization, map its reachable state $n$ to the
function $f\mapsto\mu_f(n)$.  The image is $M_{x,z}$ and the kernel is exactly
the future-unobservable subspace, proving
\eqref{eq:supp-hankel-minimality}.
\end{proof}

The maps $A_a$ generate the table-native executable forward algebra.  A
Hilbert adjoint is not introduced at this stage; it is selected only after the
positive process tensor has fixed the support-minimal Hilbert realization.

\subsection{Comb tomography and support-minimal sequential realization}

For each exposed Hermitian Choi space choose a real spanning physical frame
$\{F_\alpha\}$ and a dual frame $\{D^\alpha\}$ satisfying
$\Tr(D^\alpha F_\beta)=\delta^\alpha_\beta$.  For an $N$-step table define
\begin{equation}
 R^{(N)}
 =\sum_{\alpha_1,\ldots,\alpha_N}
 p(\alpha_1,\ldots,\alpha_N)
 D^{\alpha_1}\otimes\cdots\otimes D^{\alpha_N}.
 \label{eq:supp-comb-tomography}
\end{equation}

\begin{theorem}[Finite process-comb tomography]
\label{thm:supp-comb-tomography}
The tensor $R^{(N)}$ in \eqref{eq:supp-comb-tomography} is independent of the
chosen spanning physical frame and is the unique Hermitian tensor representing
all multilinear intervention probabilities.  It is a deterministic finite
quantum process exactly when
\begin{equation}
 R^{(N)}\succeq0,
 \qquad
 \Tr_{2k-1}R^{(k)}=I_{2k-2}\otimes R^{(k-1)},
 \qquad
 R^{(0)}=1.
 \label{eq:supp-comb-conditions}
\end{equation}
On this branch there are memories
$A_k\cong\overline{\supp R^{(k)}}$, $A_0=\C$, and isometries
\begin{equation}
 V_k:H_{2k-2}\otimes A_{k-1}
 \longrightarrow H_{2k-1}\otimes A_k
 \label{eq:supp-comb-isometry}
\end{equation}
whose sequential link realizes $R^{(N)}$.  Every other support-minimal
sequential realization satisfies
\begin{equation}
 V'_k=(I\otimes U_k)V_k(I\otimes U_{k-1}^*)
 \label{eq:supp-comb-unitary}
\end{equation}
for unique memory unitaries $U_k$ compatible with the minimum predictive state.
\end{theorem}

\begin{proof}
Products of the physical frames span the full Hermitian tensor space, so the
dual-frame formula is the unique representative of the multilinear
probability functional.  Positivity and recursive partial-trace identities are
necessary by the Choi theorem and deterministic causal normalization;
conversely they are the finite process-comb conditions.  For support
minimality, vectorize $(R^{(k)})^{1/2}$.  Consecutive causal trace identities
make the two canonical vectors minimal purifications of the same reduced
positive operator, so minimal-purification uniqueness supplies the isometries
and the inter-realization memory unitaries.
\end{proof}

\begin{remark}[Why support reduction is necessary]
An arbitrary nonminimal positive factorization may contain spectator
directions.  Taking ambient Hilbert adjoints before removing them can create
mixed history words absent from the operational process.  The support-minimal
square-root class is therefore the canonical boundary at which a history
$C^*$-algebra becomes well defined.
\end{remark}

\subsection{History algebra, contextual separation, and flat word Grams}

Choose one support-minimal realization.  Let $\mathcal A_X^{\mathrm{raw}}$ be
the finite-dimensional $*$-algebra generated by the represented projections
onto $Z_X^{\min}$, primitive branch maps, and their Hilbert--Schmidt adjoints.
A represented element is \emph{contextually null} when every admitted left and
right history context and every terminal Read has zero matrix coefficient on
it.  Let $\mathcal J_X^{\mathrm{ctx}}$ be the set of all such elements.

\begin{lemma}[Contextual nullity is a two-sided $*$-ideal]
\label{lem:supp-contextual-ideal}
$\mathcal J_X^{\mathrm{ctx}}$ is the largest two-sided $*$-ideal invisible in
every admitted history context.  Hence
\begin{equation}
 \mathcal A_X^{\mathrm{hist}}
 =\mathcal A_X^{\mathrm{raw}}/\mathcal J_X^{\mathrm{ctx}}
 \label{eq:supp-history-algebra}
\end{equation}
is the unique minimum context-preserving separated history envelope.
\end{lemma}

\begin{proof}
If $a$ is null, inserting any represented word on either side only enlarges the
tested context, and reversing the context shows that $a^*$ is null.  Thus the
null set is a two-sided $*$-ideal.  Every context-invisible ideal is contained
in it by definition, which gives the universal property.
\end{proof}

For a finite-dimensional unital $C^*$-algebra $\mathcal A$, define the
normalized regular trace
\begin{equation}
 \widehat\tau_{\mathrm{reg}}^{\mathcal A}(a)
 =\frac{\Tr_{\C}(L_a)}{\dim_{\C}\mathcal A},
 \qquad L_a(b)=ab.
 \label{eq:supp-regular-trace}
\end{equation}
If $\mathcal A\cong\bigoplus_\lambda M_{n_\lambda}(\C)$, then
$\Tr_{\C}(L_a)=\sum_\lambda n_\lambda\Tr(a_\lambda)$, so
$\widehat\tau_{\mathrm{reg}}$ is a faithful tracial state invariant under every
$*$-isomorphism.

Let $\mathcal W_{X,\le r}$ have basis the composable represented history words
of length at most $r$, and let $W_re_w=[w]$ in
$\mathcal A_X^{\mathrm{hist}}$.  Define
\begin{equation}
 \mathbb K_{X,r}(v,w)
 =\widehat\tau_{\mathrm{reg}}([v]^*[w]).
 \label{eq:supp-word-gram}
\end{equation}

\begin{theorem}[Finite flat word reconstruction]
\label{thm:supp-word-flatness}
Every $\mathbb K_{X,r}$ is positive.  There is a least $r_X$ such that
\begin{equation}
 \rank\mathbb K_{X,r_X}=\rank\mathbb K_{X,r_X+1}.
 \label{eq:supp-flat-stop}
\end{equation}
At this depth the represented word span is the whole history algebra, and the
panel through depth $r_X+1$ reconstructs the object units, minimum predictive
state algebra, multiplication, involution, predictive core algebras, and every
represented mixed history.
\end{theorem}

\begin{proof}
Positivity follows from
$\sum_{v,w}\bar c_vc_w\mathbb K(v,w)
=\widehat\tau_{\mathrm{reg}}(a^*a)\ge0$.  Faithfulness makes the Gram kernel
exactly the represented word-relation space.  The nested word spans stabilize
because the history algebra is finite-dimensional.  Equality of two
consecutive ranks therefore means equality of two consecutive spans; the
stabilized span contains the units and is invariant under left multiplication
by every generator, so it equals the full generated algebra.  The one-step
extension reconstructs the generator multiplication matrices and reversed
words reconstruct the involution.
\end{proof}

\subsection{Proof of Theorem~\ref{mt:grand-tensor}}

\begin{proof}[Proof of Theorem~\ref{mt:grand-tensor}]
Item~\textup{(i)} is
\Cref{thm:supp-right-congruence,thm:supp-minimal-record}.  Item~\textup{(ii)}
is \Cref{thm:supp-persistence-derived}.  Item~\textup{(iii)} is
\Cref{thm:supp-hankel}.  Item~\textup{(iv)} is
\Cref{thm:supp-comb-tomography}.  Item~\textup{(v)} is
\Cref{lem:supp-contextual-ideal,thm:supp-word-flatness}.  Item~\textup{(vi)}
follows from \Cref{thm:supp-minimal-record,cor:supp-future-operational-transfer}
and support-minimal unitary uniqueness.  Item~\textup{(vii)} is
\Cref{thm:supp-readable-relational-completion}.  None of these constructions
introduces an independent predictive state space, metric, connection, or
spacetime dimension beyond the actual operational rows explicitly retained.
\end{proof}

\subsection{Relative minimality of the operational entrance}
\label{subsec:supp-relative-minimality}

For the purpose of the foundational audit, keep the typed potential event
grammar fixed and distinguish four semantic rows.  Let $\mathfrak O$ denote
the actual-occurrence interface selecting which potential events, writers, and
Reads are physically present; let $\mathbb P^+$ denote the complete contextual
predictive law, including the finite ancillary/matrix-order positivity and
causal-recursion conditions used in the comb reconstruction; let $\Theta$
denote the protected orientation section; and let $\Lambda$ denote the
absolute physical calibration needed when durations, rates, lengths, or
measures are interpreted in physical units.  This separation is an audit of
information content, not an enlargement of the operational model above.

\begin{theorem}[Relative primitive-basis minimality for the paper-level entrance]
\label{thm:supp-relative-primitive-floor}
Within the fixed category of finite typed operational-predictive models used in
this paper, the semantic entrance
\begin{equation}
 \boxed{
 [\mathfrak G;\mathfrak O,\mathbb P^+,\Theta,\Lambda]}
 \label{eq:supp-relative-floor}
\end{equation}
has the following status.
\begin{enumerate}[label=\textup{(R\arabic*)},leftmargin=3em]
\item The rows already assumed in the paper are sufficient for the foundational
      reconstruction actually used below: $\mathbb P^+$ and composition give
      the minimum future quotient and persistent predictive carrier, the
      contextual positive table gives the support-minimal causal process and
      history algebra, while $\Theta$ and $\Lambda$ supply precisely the signed
      and absolute-scale information required on the Lorentzian ADM branch.
\item Each of $\mathfrak O$, $\mathbb P^+$, $\Theta$, and $\Lambda$ is
      independent of the other three relative to the fixed typed grammar: for
      each forgotten row there are two anchored finite models with identical
      reducts and inequivalent full models.
\item These four independence obstructions survive arbitrary common
      conservative refinement and cofinal reindexing; hence no functorial
      cutoff limit can reconstruct a forgotten row from the other three.
\item Protected orientation can be absorbed into occurrence exactly when the
      occurrence packet contains a named actual Read together with a physical
      decoder selecting one orientation section.  This is a definitional
      bundling, not an informational reduction of the orientation choice.
\end{enumerate}
Consequently the displayed semantic entrance is a minimal known sufficient and
independent basis relative to the declared operational language and target
class.  No claim of absolute minimality among all conceivable foundational
languages is made.
\end{theorem}

\begin{proof}
Sufficiency for the predictive and positive-process part is exactly
\Cref{thm:supp-right-congruence,thm:supp-minimal-record,cor:supp-derived-state-machine,thm:supp-comb-tomography,lem:supp-contextual-ideal,thm:supp-word-flatness}.  On the geometric branch,
\Cref{lem:supp-race-score,thm:supp-adm-frame,lem:supp-fast-clock} show how the
physically calibrated root race supplies the absolute kinetic data used by the
ADM inversion, while
\Cref{lem:supp-positive-krein-distinct,lem:supp-clifford-cell,cor:supp-lorentzian-square} use the protected sign and do not derive it from
positive data.

For row-irredundancy it is enough to give one anchored separating pair for each
row.  For $\mathbb P^+$, fix one complete binary instrument with named outcomes
$a,b$, keep every other row fixed, and compare Bernoulli laws
$\Pr(a)=p$ and $\Pr(a)=q$ with $0<p\ne q<1$.  Forgetting the predictive row
makes the models identical, whereas any full anchored isomorphism fixes the
outcome names and therefore cannot identify the two probabilities.

For $\mathfrak O$, adjoin a potential deterministic writer $W=0$ whose
mathematical extension is unique.  In one model retain only the base event as
actual; in the other declare $W$ to be an actual readable writer as well.  The
structural and probabilistic data agree after the occurrence incidence is
forgotten, but the full models differ on whether $W$ belongs to the physical
interface.  Thus structural realizability, even when unique, does not imply
actual occurrence.

For $\Theta$, take the same anchored real two-plane with named basis
$e_1,e_2$ and positive Gram $I_2$ in both models, but choose the unit
alternating forms $\omega(e_1,e_2)=+1$ and
$\omega(e_1,e_2)=-1$.  Every positive Gram datum agrees.  An anchored
isomorphism fixes the ordered basis and cannot reverse the alternating form,
so the protected orientation is not selected by the positive packet.

For $\Lambda$, fix a projective race with winner law $p_a$ and an internal
waiting variable $S\sim\operatorname{Exp}(1)$ independent of the winner.  Two
models with the same grammar, occurrence, predictive law, and orientation but
physical durations $T=c_0S$ and $T=c_1S$, $c_0\ne c_1$, have identical
uncalibrated reducts and different absolute rates $k_a=p_a/c_i$.  Hence the
common physical scale is not determined by normalized/projective race data.
This is the same scale direction whose loss is visible in
\Cref{lem:supp-race-score} when waiting-time calibration is omitted.

To see cofinal stability, tensor either member of any separating pair with the
same strictly refining finite spectator process and let every cutoff map append
and later discard only spectator records.  The row-forgotten theories are then
naturally isomorphic at every cutoff.  If a full natural isomorphism appeared
on a cofinal tail, composing it with spectator discard would give an anchored
isomorphism of the original finite separating pair, a contradiction.  Thus
common refinement or passage to a cofinal limit cannot create the missing row.

Finally, if occurrence itself contains a named actual orientation Read and a
deterministic physical decoder into the two-element orientation torsor, the
decoded section and a separately named $\Theta$ determine one another.  If the
decoder is omitted, the two opposite sections above remain compatible with the
same positive occurrence statistics.  The former is therefore only a change
of packaging, while the latter retains genuine orientation independence.
\end{proof}

\begin{theorem}[Exact calibration fibre and minimum anchor bank]
\label{thm:supp-calibration-fibre}
Write all positive physical scales and additive origins in log-affine
coordinates $x\in\mathbb R^N$.  Suppose the physically established incidence
relations are
\begin{equation}
 Lx=b.
 \label{eq:supp-calibration-incidence}
\end{equation}
If the solution set is nonempty and $x_0$ is one solution, then:
\begin{enumerate}[label=\textup{(K\arabic*)},leftmargin=3em]
\item the complete calibration fibre is
      \begin{equation}
       \boxed{x_0+\mathcal K_L,
       \qquad \mathcal K_L:=\Ker L,}
       \label{eq:supp-calibration-fibre}
      \end{equation}
      of dimension $N-\rank L$;
\item exactly $\dim\mathcal K_L$ independent scalar anchors are necessary and
      sufficient to determine one calibration in this affine fibre;
\item if $U:\mathbb R^d\to\mathbb R^N$ encodes conventional changes of base
      units with $LU=0$, the intrinsic physical residual is
      \begin{equation}
       \boxed{
       \mathcal K_L^{\rm phys}
       =\mathcal K_L/\operatorname{Im}U,}
       \label{eq:supp-physical-calibration-fibre}
      \end{equation}
      and one anchor per independent quotient direction is necessary and
      sufficient;
\item winner probabilities for a root race instantiate a one-dimensional
      common-rate direction of this type: multiplying all directed rates by the
      same positive constant leaves the winner law unchanged, whereas the
      waiting-time scale separates that direction.
\end{enumerate}
\end{theorem}

\begin{proof}
If $Lx=Lx_0=b$, then $x-x_0\in\Ker L$, and conversely every
$x_0+k$, $k\in\Ker L$, solves \eqref{eq:supp-calibration-incidence}; this gives
\eqref{eq:supp-calibration-fibre} and its dimension by rank--nullity.  A basis
of $\mathcal K_L$ and the corresponding dual linear functionals give a
sufficient scalar anchor bank.  Any smaller bank has nontrivial common kernel
on $\mathcal K_L$ and is therefore insufficient.  Quotienting the subgroup of
pure unit conventions gives
\eqref{eq:supp-physical-calibration-fibre} and the same argument on the quotient
gives the intrinsic anchor count.  For the root race, the winner law is
$p_b=k_b/\sum_ck_c$ and is invariant under $k_b\mapsto ck_b$; the exponential
waiting time has total rate $\sum_ck_c$ and therefore fixes the missing common
scale.
\end{proof}

\section{Minimal relational dimension and ADM reconstruction}
\label{supp:dimension-adm}
% =============================================================================

This section proves \Cref{mt:adm}.  The argument has two logically separate
parts.  The first selects the local relational module.  The second reconstructs
ADM geometry from a scheduler-separated root race on that module.

\subsection{From one serial relation to a complete cell}

\begin{assumption}[Minimal faithful predictive relation branch]
\label{ass:supp-relational-branch}
A local relation event has one future-readable source and one future-readable
target.  Physical reversal retains both orientations.  Future-minimality
identifies parallel relation labels exactly when every continuation identifies
them.  A nonempty relation set is invariant under a group acting transitively
on unordered endpoint pairs.
\end{assumption}

\begin{lemma}[Pair-homogeneous completion]
\label{lem:supp-pair-completion}
Under \Cref{ass:supp-relational-branch}, the relation-colour-free local skeleton
is the complete graph $K_N$ on the endpoint set.
\end{lemma}

\begin{proof}
The selected branch of the derived persistent predictive carrier contains
exactly one future-readable source and one future-readable target, so one
occurring relation gives one unoriented pair after physical reversal.  The
orbit of that pair under a group transitive on $\binom{X}{2}$ is the complete
set of unordered pairs.  Future-minimality removes only unread parallel labels,
not distinct endpoint pairs.  Hence the skeleton is $K_N$.
\end{proof}

Let $X=\{1,\ldots,N\}$, let $u_0=N^{-1/2}\one_N$, and write
$\R^N=\R u_0\oplus W_N$.  Identify the signed edge $[i\to j]$ with
$e_i\wedge e_j\in\Lambda^2\R^N$.  The boundary map is
$\partial(e_i\wedge e_j)=e_j-e_i$.

\begin{theorem}[Cut--cycle decomposition]
\label{thm:supp-cut-cycle}
There is an orthogonal decomposition
\begin{equation}
 C_1(K_N)=u_0\wedge W_N\oplus\Lambda^2W_N,
 \label{eq:supp-cut-cycle}
\end{equation}
with
\begin{equation}
 \partial(u_0\wedge w)=\sqrt N\,w,
 \qquad
 \Ker\partial=\Lambda^2W_N.
 \label{eq:supp-boundary-action}
\end{equation}
Consequently
\begin{equation}
 \dim\Ran\partial=N-1,
 \qquad
 \dim\Ker\partial=\binom{N-1}{2}.
 \label{eq:supp-cut-cycle-dim}
\end{equation}
\end{theorem}

\begin{proof}
The exterior square of
$\R^N=\R u_0\oplus W_N$ is
$u_0\wedge W_N\oplus\Lambda^2W_N$.  For $x,y\in\R^N$,
\[
 \partial(x\wedge y)
 =\left(\sum_ix_i\right)y-
   \left(\sum_iy_i\right)x.
\]
Taking $x=u_0$, $y=w\in W_N$ gives the first identity.  If both vectors lie in
$W_N$, both coordinate sums vanish, so the second summand is in the kernel.
Dimension counting shows that it is the whole kernel and gives
\eqref{eq:supp-cut-cycle-dim}.
\end{proof}

\begin{remark}[Metric tomography does not select $N$]
The root dyadics of a complete simplex define an isomorphism
\[
 \R^{\binom N2}\longrightarrow\Sym(W_N),
 \qquad
 (a_{ij})\longmapsto\sum_{i<j}a_{ij}
 (e_j-e_i)(e_j-e_i)^{\mathsf T}
\]
for every $N$.  Indeed, a vanishing weighted sum is a graph Laplacian whose
off-diagonal entries are $-a_{ij}$.  Thus the equality
$\binom N2=\dim\Sym(W_N)$ is universal and cannot by itself select three
spatial dimensions.  The paper uses two different non-tomographic selectors:
exact endpoint--cycle source-rank balance, and the stronger
alternating/minimal-source simplex criterion.
\end{remark}

\begin{remark}[Conditional character of the selector]
The alternating form in the next theorem is part of the definition of the
selected relational category.  It is not reconstructed from the positive word
Grams and does not, by itself, specify a metric, a Hamiltonian, a symplectic
phase space of continuum fields, or a Lorentzian sign.  Accordingly the theorem
asserts uniqueness only among faithful pair-homogeneous simplex cells satisfying
this alternating/minimal-source query.  It is the second dimensional selector,
not the sole derivation of $K_4$; the regular-graph source-balance result below
uses no alternating form.
\end{remark}

\begin{theorem}[Minimal faithful $K_4$ selector on the alternating simplex branch]
\label{thm:supp-k4-selector}
Assume, in addition to \Cref{ass:supp-relational-branch}, that:
\begin{enumerate}[label=\textup{(D\arabic*)},leftmargin=3em]
\item the temporal--spatial coefficient space is $\R t\oplus W_N$ and carries a
      nondegenerate alternating form;
\item the endpoint and complete signed-edge source Grams are faithful;
\item the cycle source is nonzero;
\item among conservative completions satisfying these queries, the
      future-readable source innovation is minimal.
\end{enumerate}
Then
\begin{equation}
 N=4,
 \qquad
 G_{\mathrm{rel}}=K_4,
 \qquad
 W_{\mathrm{sp}}=W_4\cong\R^3.
 \label{eq:supp-k4-result}
\end{equation}
The faithful endpoint and edge source ranks are
$r_0(N)=N-1$ and $r_1(N)=\binom N2$.
\end{theorem}

\begin{proof}
The alternating coefficient space has dimension
$1+\dim W_N=N$.  A nondegenerate alternating form exists only in even
dimension, so $N$ is even.  A nonzero cycle source requires
$\Lambda^2W_N\ne0$, hence $N\ge3$; together with evenness this gives $N\ge4$.
On the even integers $N\ge4$, both $N-1$ and $\binom N2$ are strictly
increasing.  The unique minimum future-readable source cost is therefore
attained at $N=4$.  The graph conclusion follows from
\Cref{lem:supp-pair-completion}, and $\dim W_4=3$.
\end{proof}

A representation-theoretic strengthening explains why the cycle packet is
particularly economical at $N=4$.  After complexification,
$W_N\cong S^{(N-1,1)}$,
$\Lambda^2W_N\cong S^{(N-2,1,1)}$, and
$W_N\otimes\operatorname{sgn}\cong S^{(2,1^{N-2})}$.  These Young diagrams
coincide exactly for $N=4$.  Hence
\begin{equation}
 \Hom_{S_N}(\Lambda^2W_N,W_N\otimes\operatorname{sgn})
 =\begin{cases}\R,&N=4,\\0,&N\ne4,\end{cases}
 \label{eq:supp-oriented-copy}
\end{equation}
on the nontrivial simplex branch.  Thus only $K_4$ identifies the entire cycle
representation with one oriented spatial copy rather than an additional
inequivalent source type.

\begin{theorem}[Alternating-free endpoint--cycle selector]
\label{cor:supp-regular-graph-k4}
Let $G$ be a connected simple $k$-regular relation graph on $v$ endpoints.  If
its complete endpoint and cycle sources are faithful and have equal rank, then
\begin{equation}
 v-1=\frac{kv}{2}-v+1,
 \qquad\hbox{hence}\qquad
 (4-k)v=4.
 \label{eq:supp-regular-graph-balance}
\end{equation}
Apart from the one-point graph, the unique solution is
\begin{equation}
 \boxed{v=4,\qquad k=3,\qquad G\cong K_4.}
 \label{eq:supp-regular-graph-k4}
\end{equation}
Thus $K_4$ is also the unique nontrivial connected regular simple relation graph
with exact faithful endpoint--cycle rank balance.  This corroborates the
simplex selector but does not extend it to arbitrary higher-arity boundary
systems.
\end{theorem}

\begin{proof}
For a connected graph the endpoint-source rank is $v-1$, while the cycle rank is
$|E|-v+1=kv/2-v+1$.  Equality gives
$2(v-1)=kv/2$, equivalently $(4-k)v=4$.  With $v>1$ and integer $k\ge1$, the
only connected simple regular solution is $v=4$, $k=3$, which is $K_4$.
\end{proof}

\subsection{The \texorpdfstring{$A_3$}{A3} source frame}

In the anchored basis of \eqref{eq:a3-roots-main}, define
$Q_a=r_ar_a^{\mathsf T}$.

\begin{lemma}[Root dyadics span $\Sym_3$]
\label{lem:supp-root-frame}
The six matrices $Q_1,\ldots,Q_6$ form a basis of $\Sym_3$.
\end{lemma}

\begin{proof}
$Q_1,Q_2,Q_3$ are the three diagonal matrix units.  For $i<j$,
\[
 (e_j-e_i)(e_j-e_i)^{\mathsf T}-e_ie_i^{\mathsf T}-e_je_j^{\mathsf T}
 =-(e_ie_j^{\mathsf T}+e_je_i^{\mathsf T}).
\]
The three difference roots therefore recover the three symmetric off-diagonal
matrix units.  The six matrices are linearly independent and span the
six-dimensional space $\Sym_3$.
\end{proof}

Let the twelve oriented roots race with positive rates $k_a^\pm$.  Put
$K=\sum_a(k_a^++k_a^-)$, let $A$ be the winning oriented root, and let $T$ be
the first waiting time.

\begin{lemma}[Winner--waiting score is faithful]
\label{lem:supp-race-score}
For each oriented branch $b$, the log-rate score is
\begin{equation}
 s_b=\one_{\{A=b\}}-k_bT.
 \label{eq:supp-race-score}
\end{equation}
Its Gram matrix is diagonal and positive:
\begin{equation}
 \mathbb E[s_bs_c]=\delta_{bc}\frac{k_b}{K}.
 \label{eq:supp-race-score-gram}
\end{equation}
Winner probabilities alone have covariance
$\diag(p)-pp^{\mathsf T}$ and therefore lose the common rate scale.
\end{lemma}

\begin{proof}
For independent exponential clocks, $T\sim\operatorname{Exp}(K)$, the winner
has probability $p_b=k_b/K$, and $A$ and $T$ are independent.  Therefore
\[
 \begin{aligned}
 \mathbb E[s_bs_c]
 &={}\delta_{bc}p_b-p_bk_c\mathbb ET-k_bp_c\mathbb ET
      +k_bk_c\mathbb ET^2\\
 &={}\delta_{bc}\frac{k_b}{K}
   -2\frac{k_bk_c}{K^2}+2\frac{k_bk_c}{K^2},
 \end{aligned}
\]
which proves \eqref{eq:supp-race-score-gram}.  The covariance of the winner
indicator is singular in the all-ones rate direction because probabilities
are unchanged by common rescaling.
\end{proof}

\subsection{ADM inversion and scheduler scale}

Define $\mathcal B_h$ and $\beta_h$ by
\eqref{eq:bracket-drift-main}.

\begin{theorem}[Complete finite ADM frame]
\label{thm:supp-adm-frame}
The six rate sums map isomorphically to $\mathcal B_h\in\Sym_3$.  The six rate
differences map with rank three onto $\beta_h\in\R^3$ and have a
three-dimensional kernel, the tetrahedral circulation space.  For every
$\mathcal B_h\succ0$ and $\varrho_h>0$, the formula
\eqref{eq:adm-inversion-main} is the unique solution of
\begin{equation}
 \mathcal B_h=N_h^2g_h^{-1},
 \qquad
 \varrho_h=\sqrt{\det g_h}.
 \label{eq:supp-adm-characterization}
\end{equation}
Hence the root rates, waiting-time scale, and speed density reconstruct the six
spatial metric components, three shift components, and lapse.
\end{theorem}

\begin{proof}
The first assertion is \Cref{lem:supp-root-frame}.  The difference map has
image containing $r_1,r_2,r_3$, hence has rank three; rank--nullity gives a
three-dimensional kernel.  For the inversion, take determinants in
$\mathcal B=N^2g^{-1}$:
$\det\mathcal B=N^6/(\det g)=N^6/\varrho^2$.  Thus
$N=\varrho^{1/3}(\det\mathcal B)^{1/6}$.  Substituting back gives
$g=N^2\mathcal B^{-1}$, which is exactly
\eqref{eq:adm-inversion-main}.  Positivity makes the solution unique.
\end{proof}

\begin{lemma}[Fast-clock lower bound]
\label{lem:supp-fast-clock}
Suppose $\norm{r_a}\le R$ and $\mathcal B_h\succeq cI_3$.  Then
\begin{equation}
 \sum_a(k_a^++k_a^-)\ge\frac{3c}{R^2h^2}.
 \label{eq:supp-fast-clock-bound}
\end{equation}
In particular, a total rate $O(h^{-1})$ produces $\mathcal B_h=O(h)$ and a
collapsed continuum bracket.
\end{lemma}

\begin{proof}
Taking traces in \eqref{eq:bracket-drift-main},
\[
 3c\le\tr\mathcal B_h
 \le R^2h^2\sum_a(k_a^++k_a^-).
\]
This gives \eqref{eq:supp-fast-clock-bound}.  The final assertion follows by
substitution.
\end{proof}

\subsection{Finite Krein orientation from the protected sign}

Let $\epsilon_X\in\{+1,-1\}$ be the actual protected orientation Read of
\eqref{eq:orientation-read-main} on the finite support-minimal relational
carrier $\mathcal H_J$.  Functional calculus of this classical two-point Read
gives orthogonal projections $P_\pm$ with $P_++P_-=I$; define
$J=P_+-P_-$.  Thus
\begin{equation}
 \mathcal H_J=\mathcal H_J^+\oplus\mathcal H_J^-,
 \qquad
 \mathcal H_J^\pm=\Ran P_\pm,
 \qquad
 J|_{\mathcal H_J^\pm}=\pm I.
 \label{eq:supp-krein-splitting}
\end{equation}
Assume both summands are nonzero.  Define
\begin{equation}
 [u,v]_J=\langle u,Jv\rangle.
 \label{eq:supp-krein-form}
\end{equation}
Then $[\cdot,\cdot]_J$ is a nondegenerate Hermitian form with one positive and one negative sector in the orientation factor, hence a finite Krein structure.  This statement uses only the protected sign involution and the support-minimal Hilbert realization; no continuum metric is involved.

\begin{lemma}[Positive geometry and Krein signature are distinct layers]
\label{lem:supp-positive-krein-distinct}
The positive word-Gram hierarchy of \Cref{thm:supp-word-flatness} remains positive after adjoining the protected orientation record, while the indefinite sign is carried entirely by $J$.  Consequently no positive combination of the spatial root dyadics can replace the Krein involution, and the Lorentzian signature of the principal symbol requires the signed layer in addition to the $A_3$ bracket.
\end{lemma}

\begin{proof}
The word-Gram hierarchy is a Gram matrix for the faithful positive regular trace and is therefore positive semidefinite.  Likewise every spatial bracket is a positive sum $\sum_a c_a r_ar_a^{\mathsf T}$ with $c_a>0$.  Such a form has no negative direction.  By contrast, $J$ has both $+1$ and $-1$ spectral subspaces by \eqref{eq:supp-krein-splitting}; the form \eqref{eq:supp-krein-form} is therefore indefinite.  The two structures are independent in exactly the required sense: the spatial geometry remains positive while $J$ supplies the sign subsequently represented by the temporal Clifford generator.
\end{proof}

\begin{theorem}[Irreducible sign and Lorentzian inertia before Clifford realization]
\label{thm:supp-signature-minimality}
Let $W_4$ carry the positive spatial form reconstructed by
\Cref{thm:supp-adm-frame}.  The complete positive Grand-Tensor packet cannot
supply a negative direction, while the protected orientation row is independent
of that packet by \Cref{thm:supp-relative-primitive-floor}.  If the geometric
extension retains the positive restriction on $W_4$ and introduces exactly one
independent nondegenerate signed line $\R t$, then, up to overall sign,
\begin{equation}
 \boxed{q(\tau t+w)=-a^2\tau^2+g(w,w),\qquad a>0,
 \qquad \operatorname{inertia}(q)=(1,3).}
 \label{eq:supp-signature-inertia}
\end{equation}
The conclusion depends only on positivity of the spatial form, independence of
the signed row, and one-line minimality; it does not depend on a choice of gamma
matrices.
\end{theorem}
\begin{proof}
The positive word-Gram and root-race forms have no negative eigenvalues by
\Cref{lem:supp-positive-krein-distinct}.  The separating pair in the proof of
\Cref{thm:supp-relative-primitive-floor} shows that reversing the protected
orientation leaves all positive operational data unchanged, so its sign is not
a derived eigenvalue of the positive packet.  A minimal extension introduces
one real line and no additional signed directions.  Nondegeneracy makes its
coefficient nonzero; choosing the relative sign opposite to the positive
spatial restriction gives one negative eigenvalue and the three positive
eigenvalues of $g$.  Sylvester's law of inertia makes this count invariant
under every real change of basis.  Multiplying the entire form by $-1$ reverses
the sign convention but not the Lorentzian type.  No Clifford representation
has been used.
\end{proof}

\subsection{Clifford realization of the signed quadratic form}

The inertia has already been fixed by
\Cref{thm:supp-signature-minimality}.  The following matrices provide a
minimal complex Clifford realization of that form; they are not used to choose
its signature.  Let $X,Y,Z$ denote the Pauli matrices and define the four
matrices in \eqref{eq:clifford-main}.

\begin{lemma}[Depth-two Lorentzian Clifford cell]
\label{lem:supp-clifford-cell}
The matrices $\gamma^\mu$ satisfy
\begin{equation}
 (\gamma^0)^2=-I,
 \qquad
 (\gamma^a)^2=I,
 \qquad
 \{\gamma^\mu,\gamma^\nu\}=0\quad(\mu\ne\nu),
 \label{eq:supp-clifford-relations}
\end{equation}
and generate $M_4(\C)\cong\Cl_{1,3}(\C)$.
\end{lemma}

\begin{proof}
The square relations follow from $X^2=Y^2=Z^2=I$.  Every mixed product contains
exactly one local Pauli anticommutation, so the mixed anticommutators vanish.
The generators contain $X_A\otimes I$ and $Y_A\otimes I$, hence also
$Z_A\otimes I$.  Multiplication by $\gamma^1$ and $\gamma^2$ then gives
$I\otimes Z_B$ and $I\otimes Y_B$, hence $I\otimes X_B$.  Both local matrix
factors are generated, so the total algebra is
$M_2(\C)\otimes M_2(\C)=M_4(\C)$.
\end{proof}

\begin{corollary}[Lorentzian principal square]
\label{cor:supp-lorentzian-square}
Let $E_hE_h^{\mathsf T}=g_h^{-1}$ and define $\sigma_h$ by
\eqref{eq:dirac-symbol-main}.  Then \eqref{eq:dirac-square-main} holds.
If $g_h\succ0$ and $N_h>0$, the characteristic quadratic form realizes the
already determined inertia $(1,3)$ in the convention $(-,+,+,+)$.
\end{corollary}

\begin{proof}
The temporal--spatial cross terms vanish by the Clifford anticommutation
relations.  The temporal square contributes
$-N_h^{-2}(\xi_0-\beta_h\cdot\vec\xi)^2$, while the spatial square is
$\delta^{ab}(E_h)_a{}^i(E_h)_b{}^j\xi_i\xi_j=g_h^{ij}\xi_i\xi_j$.
\end{proof}

\subsection{Cofinal transport and proof of \texorpdfstring{\Cref{mt:adm}}{the ADM reconstruction theorem}}

For adjacent cutoffs, let the complete old/detail positive form have block
matrix
\begin{equation}
 M_Y=\begin{pmatrix}A_{Y/X}&B_{Y/X}\\
                     B_{Y/X}^*&D_{Y/X}\end{pmatrix},
 \label{eq:supp-refinement-block}
\end{equation}
and Schur complement
$S_{Y/X}=A_{Y/X}-B_{Y/X}D_{Y/X}^{\dagger}B_{Y/X}^*$.  Define
\begin{equation}
 \sigma_{Y/X}=\norm{S_{Y/X}-M_X}_{\op},
 \qquad
 \theta_{Y/X}=\norm{B_{Y/X}D_{Y/X}^{\dagger}}_{\op},
 \qquad
 \mu_{Y/X}=\lambda_{\min}^+(D_{Y/X}).
 \label{eq:supp-refinement-defects}
\end{equation}
The harmonic lift
\begin{equation}
 \mathcal H_{Y/X}u
 =J_{Y/X}u-K_{Y/X}D_{Y/X}^{\dagger}B_{Y/X}^*u
 \label{eq:supp-harmonic-lift}
\end{equation}
minimizes the fine energy over detail completions and has energy
$\ip{u}{S_{Y/X}u}$.  This is ordinary completion of the square.

If
\begin{equation}
 \sum_X\bigl(\sigma_X+\theta_X+\mu_X^{-1}\bigr)<\infty,
 \label{eq:supp-summable-refinement}
\end{equation}
then the finite forms and their regular isolated spectral projections are
Cauchy after harmonic transport.  In particular, a separated threshold rank,
a positive pair gap, and a finite-dimensional Clifford source persist.  This
is the finite-dimensional core of the common-screen limit used here.

\begin{proof}[Proof of Theorem~\ref{mt:adm}]
The predictive carrier is supplied by \Cref{thm:supp-persistence-derived}.
For the regular-graph route, the local dimension follows directly from
\Cref{cor:supp-regular-graph-k4}.  For the pair-homogeneous simplex route, the
local graph is complete by \Cref{lem:supp-pair-completion} and the same
$K_4/A_3$ module follows from \Cref{thm:supp-cut-cycle,thm:supp-k4-selector}.
Thus the alternating form is needed only on the second route.  The root sums and
differences, faithful rate coordinates, ADM inversion, and clock lower bound
are \Cref{lem:supp-root-frame,lem:supp-race-score,thm:supp-adm-frame,lem:supp-fast-clock}.
The positive/signed separation is \Cref{lem:supp-positive-krein-distinct};
\Cref{thm:supp-signature-minimality} fixes Lorentzian inertia before the
Clifford matrices are introduced, and
\Cref{lem:supp-clifford-cell,cor:supp-lorentzian-square} realizes that form.
Finally,
\eqref{eq:supp-summable-refinement} makes every corrected finite matrix packet
Cauchy in operator norm on each protected finite screen; spectral separation
therefore preserves the ranks and positive gaps, while the Schur mismatch and
leakage tails vanish.  The canonical limit quotient consequently retains the
$K_4/A_3$ module, the physically occurring temporal calibration, the ADM packet,
and the Lorentzian Clifford cell.
\end{proof}

\begin{remark}[Exact dimension scope]
The result applies to the derived future-minimal, one-source/one-target, reciprocal,
pair-homogeneous, faithful, and minimal faithful branch.  A deliberately
higher-arity boundary, a relation-coloured theory, a nonfaithful cycle quotient,
or an enriched cross-product packet changes the primitive query and is not a
counterexample.  In particular, a seven-dimensional vector cross product
$\Lambda^2W\to W$ has a fourteen-dimensional kernel and requires an additional
stable $G_2$ three-form; it is an enriched theory rather than an equal-cost
realization of the faithful $K_4$ branch.
\end{remark}


% =============================================================================
\section{Operational action reconstruction and the Einstein handoff}
\label{supp:dynamics}
% =============================================================================

This appendix proves the reconstruction and variational statements of
\Cref{sec:einstein-dynamics}.  The order is important: the determining field is
first obtained from the complete accepted cylinder; the common action is then
reconstructed from a protected reward-duration packet and finite exactness
tests; the actual accepted transition is compared with that action; spatial
compactness is derived from a cut margin; and only then is the gravitational
first variation passed to the continuum.

\subsection{Accepted determining-field quotient and kernel}

Let $\widetilde\Theta_X$ be the finite accepted ledger of
\Cref{eq:main-determining-field}.  Two states $x,y$ are future equivalent when
every declared determining coordinate and every admitted future word agrees
from $x$ and $y$, including all branch-domain indicators.  Let
$\Theta_{{\rm fld},X}^{\min}$ be the quotient.

\begin{theorem}[Terminating determining-field and accepted-kernel compiler]
\label{thm:supp-determining-kernel}
The quotient $\Theta_{{\rm fld},X}^{\min}$ is the unique coarsest finite record
through which every selected determining-field future and every primitive update
factors.  Partition refinement beginning with the immediate determining
coordinates and branch domains terminates after finitely many strict rounds.
The kernel $K_X^{\rm acc}$ of \eqref{eq:main-accepted-kernel} is well defined,
table native, and is the unique accepted transition reproducing all selected
post-accepted field futures.
\end{theorem}

\begin{proof}
Future equality is a typed right congruence.  If $x\sim_{{\rm fld},X}y$, then
for every primitive branch $a$ and continuation $w$, every determining read
satisfies
\[
 d(u_wu_a(x))=d(u_{wa}(x))=d(u_{wa}(y))=d(u_wu_a(y)),
\]
so every primitive update descends to the quotient.  At each refinement round,
split two states whenever an immediate coordinate differs or one primitive
update reaches different blocks.  The $k$th partition identifies precisely the
states not separated by a future of length at most $k$.  Every strict round
increases the number of finite blocks, hence the refinement terminates.  Its
stable partition has the universal property by the usual fibre argument.

Accepted branch probabilities and target quotient classes are constant on the
stable fibres.  Summing the resolved accepted branches whose target lies in a
specified quotient class gives \eqref{eq:main-accepted-kernel}, and these sums
are exactly the probabilities of the corresponding post-accepted Reads.
Uniqueness follows because those Reads separate the quotient classes.
\end{proof}

\begin{remark}[Occurrence boundary]
The theorem reconstructs the dynamical state and kernel from the \emph{complete}
accepted cylinder.  It does not reconstruct that cylinder from an accepted
binary clock alone.  Indeed, an identity accepted kernel and a full-support
reset kernel can be coupled to the same state-independent acceptance bit while
producing different terminal field laws.  Thus occurrence of the complete
accepted field record is a genuine physical input row.
\end{remark}

\subsection{Renewal pressure and the common scalar action}

Let $(\Omega_X,p_X)$ be a faithful marked first-return law, let
$\tau_X>0$ be the protected physical duration, and let
$\mathscr G_X:E_X\to L^2(\Omega_X,p_X)$ be the reward synthesis of
\eqref{eq:main-reward-synthesis}.  Put
\[
 \mathcal L_X(q,r)=\mathbb E_X[e^{-g_{q,X}-r\tau_X}],
 \qquad
 \bar\tau_X=\mathbb E_X\tau_X,
 \qquad
 T_Xc=c\tau_X.
\]

\begin{theorem}[Table-derived renewal pressure and action-slope reconstruction]
\label{thm:supp-reward-pressure}
The complete Laplace transform $\mathcal L_X$ and all of its finite derivatives
are reconstructed from the marked-cycle table.  There is a unique analytic
pressure $\mathcal P_X$ near the origin satisfying
\[
 \mathcal P_X(0)=0,
 \qquad
 \mathcal L_X(q,\mathcal P_X(q))=1.
\]
Let $S_Xv=D\log p_{X,q}|_{q=0}[v]$ and $\pi_X=D\mathcal P_X(0)$.  Then
\begin{equation}
 \boxed{
 \pi_X(v)=-\frac{\mathbb E_Xg_{v,X}}{\bar\tau_X},
 \qquad
 S_Xv=-\mathscr G_Xv-\pi_X(v)\tau_X,
 \qquad
 \mathscr G_X=-S_X-T_X\pi_X.}
 \label{eq:supp-pressure-action-slope}
\end{equation}
Thus the physical duration restores the affine action coordinate omitted by the
centered probability score.
\end{theorem}

\begin{proof}
The transform is a finite sum of known probabilities times known protected
writer values.  At the origin, $\mathcal L_X(0,0)=1$ and
$\partial_r\mathcal L_X(0,0)=-\bar\tau_X<0$, so the analytic implicit-function
theorem gives the pressure root.  Differentiating
$\mathcal L_X(q,\mathcal P_X(q))=1$ in direction $v$ gives
\[
 0=-\mathbb E_Xg_{v,X}-\pi_X(v)\bar\tau_X,
\]
which proves the first identity.  Differentiating the normalized tilted law
gives the score formula and hence the synthesis reconstruction.
\end{proof}

\begin{corollary}[Pressure-jet reconstruction of mixed action blocks]
\label{cor:supp-pressure-jet}
With $\bar\tau_X(q)=\mathbb E_{X,q}\tau_X$,
$\delta_X=D\bar\tau_X(0)$, and
$m_{2,X}=\mathbb E_X\tau_X^2$,
\begin{equation}
 \boxed{
 D^2\mathcal P_X(0)=\bar\tau_X^{-1}S_X^*S_X,\qquad
 \delta_X=T_X^*S_X,}
 \label{eq:supp-pressure-jet-a}
\end{equation}
\begin{equation}
 \boxed{
 \mathscr G_X^*\mathscr G_X
 =\bar\tau_XD^2\mathcal P_X(0)
 +\delta_X^*\pi_X+\pi_X^*\delta_X
 +m_{2,X}\pi_X^*\pi_X.}
 \label{eq:supp-pressure-jet-b}
\end{equation}
Consequently one pressure surface reconstructs every marginal and mixed action
block represented on the same marked cycle record.
\end{corollary}

\begin{proof}
Differentiate the pressure normalization twice.  The first factors are the
centered scores $-S_Xu$ and $-S_Xv$, which gives the first identity.
Differentiating the tilted mean duration gives $\delta_X=T_X^*S_X$; expanding
$\mathscr G_X=-S_X-T_X\pi_X$ gives the second formula.
\end{proof}

Let $\mathscr V_X$ be the finite protected variation complex, with coboundaries
$d_0:C^0\to C^1$ and $d_1:C^1\to C^2$, $d_1d_0=0$, and let
$\alpha_X\in C^1(\mathscr V_X;\R)$ be the reconstructed physical action
increment.

\begin{theorem}[Finite common-action exactness]
\label{thm:supp-common-action}
There is a scalar action $\mathscr S_X$ with $\alpha_X=d_0\mathscr S_X$ if and
only if $d_1\alpha_X=0$ on every declared $2$-cell and all periods of
$\alpha_X$ vanish on one integral basis of $H_1(\mathscr V_X;\Z)$.  After fixing
one action anchor, $\mathscr S_X$ is unique and reconstructed by path
integration.  If $P_{\rm har}$ and $P_{\rm coex}$ are the Hodge projections and
$\sigma_{\rm coex}>0$ is the least singular value of $d_1$ on the coexact
range, then
\begin{equation}
 \boxed{
 \operatorname{dist}(\alpha_X,\Ran d_0)^2
 \le\|P_{\rm har}\alpha_X\|^2
 +\sigma_{\rm coex}^{-2}\|d_1\alpha_X\|^2.}
 \label{eq:supp-common-action-distance}
\end{equation}
\end{theorem}

\begin{proof}
If $\alpha_X=d_0\mathscr S_X$, telescoping gives zero face sums and zero periods.
Conversely, fix a base configuration and define $\mathscr S_X(q)$ by summing
$\alpha_X$ along a path from the base to $q$.  Face closure makes the sum
invariant under elementary $2$-cell moves and the period conditions remove the
remaining homology class.  Anchoring removes the additive constant.  For the
estimate use the finite Hodge decomposition
$C^1=\Ran d_0\oplus\mathcal H^1\oplus\Ran d_1^*$ and the singular-value bound
$\|d_1\beta\|\ge\sigma_{\rm coex}\|\beta\|$ on the coexact range.
\end{proof}

\begin{remark}[Provenance is not stationarity]
Reconstructing a common action does not prove that the observed accepted
transition is generated by it.  The accepted/action comparison below is an
independent same-history test.  Conversely, fitting a cost to the observed
kernel after the fact would be tautological and has no stationarity content.
\end{remark}

\subsection{Fisher--Bregman stationarity and physical test lifts}

Let $\Theta$, $\Psi$, $G$, $\mathcal A$, and $\eta$ satisfy the convex-chart
and interiority hypotheses of \Cref{thm:main-action-stationarity}.  Define
$D_\Psi$, $\Phi_\theta$, $y_\theta^*$, $\Delta_\Psi$, and
$\overline\Delta_\Psi$ by
\eqref{eq:main-bregman}--\eqref{eq:main-stationarity-gap}.

\begin{theorem}[Quantitative accepted-action stationarity]
\label{thm:supp-action-stationarity}
Put
\begin{equation}
 \mu_\Phi=(\eta^{-1}-\lambda)m>0,\qquad
 \Lambda_\Phi=(\eta^{-1}+L)M,\qquad
 \ell=\max\{\lambda,L\}.
 \label{eq:supp-action-constants}
\end{equation}
Then
\begin{equation}
 \boxed{\int D_\Psi(y,\theta)\,
       \nu(\dd\theta)K^{\rm acc}(\dd y\mid\theta)
       \le\eta\overline\Delta_\Psi,}
 \label{eq:supp-bregman-bound}
\end{equation}
\begin{equation}
 \boxed{\int\|y-\theta\|^2\,
       \nu(\dd\theta)K^{\rm acc}(\dd y\mid\theta)
       \le\frac{2\eta}{m}\overline\Delta_\Psi,}
 \label{eq:supp-displacement-bound}
\end{equation}
and
\begin{equation}
 \boxed{\int\mathfrak S(\theta)\,\nu(\dd\theta)
       \le C_{\Psi,\mathcal A,\eta}\overline\Delta_\Psi,}
 \label{eq:supp-stationarity-bound}
\end{equation}
where one admissible constant is
\begin{equation}
 C_{\Psi,\mathcal A,\eta}
 =\frac{4\Lambda_\Phi^2}{m\mu_\Phi}
   +\frac{4M^2(1+\eta\ell)^2}{m^2\eta}.
 \label{eq:supp-stationarity-constant}
\end{equation}
\end{theorem}
\begin{proof}
State independence of $\vartheta>0$ and stationarity of $\nu$ imply
$\nu K^{\rm acc}=\nu$.  Since $y_\theta^*$ minimizes $\Phi_\theta$ and
$\Phi_\theta(\theta)=\mathcal A(\theta)$,
\[
 \mathcal A(y)-\mathcal A(\theta)+\eta^{-1}D_\Psi(y,\theta)
       \le\Delta_\Psi(\theta,y).
\]
The action difference cancels on averaging, proving
\eqref{eq:supp-bregman-bound}.  The lower Hessian bound gives
$D_\Psi(y,\theta)\ge m\|y-\theta\|^2/2$, hence
\eqref{eq:supp-displacement-bound}.

The Hessian of $\Phi_\theta$ lies between $\mu_\Phi I$ and
$\Lambda_\Phi I$ on the convex chart.  Interiority of $y_\theta^*$ gives
$\nabla\Phi_\theta(y_\theta^*)=0$.  Strong convexity and the upper Hessian
bound therefore imply
\[
 \|y-y_\theta^*\|^2\le\frac{2\Delta_\Psi(\theta,y)}{\mu_\Phi},
 \qquad
 \|\nabla\Phi_\theta(y)\|^2
       \le\frac{2\Lambda_\Phi^2}{\mu_\Phi}\Delta_\Psi(\theta,y).
\]
For
\[
 r_\Psi(\theta,y)=\eta\nabla\Phi_\theta(y)
       =\nabla\Psi(y)-\nabla\Psi(\theta)+\eta\nabla\mathcal A(y)
\]
we have
\[
 \eta\nabla\mathcal A(\theta)
 =r_\Psi-[\nabla\Psi(y)-\nabla\Psi(\theta)]
       +\eta[\nabla\mathcal A(\theta)-\nabla\mathcal A(y)].
\]
The last two terms together have norm at most
$M(1+\eta\ell)\|y-\theta\|$.  Squaring with
$\|a+b\|^2\le2\|a\|^2+2\|b\|^2$, using $G^{-1}\preceq m^{-1}I$, and
averaging yields \eqref{eq:supp-stationarity-bound} with
\eqref{eq:supp-stationarity-constant}.
\end{proof}

\begin{remark}[Boundary minima]
The first two inequalities do not require interiority, but the ordinary
gradient conclusion does.  For example, on $\Theta=[0,1]$ with
$\Psi(x)=x^2/2$ and $\mathcal A(x)=x$, the source $\theta=0$ has proximal
minimum $y_\theta^*=0$.  The stationary identity transition supported at zero
has zero gap although $\mathcal A'(0)=1$.  A boundary formulation must retain
the normal-cone or projected variational residual.  It cannot be substituted
for unrestricted determining metric variations without checking their
admissibility.
\end{remark}

\begin{proposition}[Cutoff-aware stationarity on the determining test core]
\label{prop:supp-physical-test-stationarity}
Let $k$ be a fixed physical test and let its represented lift satisfy
\eqref{eq:main-test-lift-budget}.  Then
\eqref{eq:main-physical-stationarity} holds.  If
\begin{equation}
 B_{X,K}(k)^2C_X\overline\Delta_{\Psi,X}\longrightarrow0
 \label{eq:supp-scaled-test-gap}
\end{equation}
for each test in a countable determining core, the corresponding first
variations converge to zero in mean square.  Under any common coupling there
is a deterministic subsequence of cutoffs on which all these first variations
converge to zero almost surely.
\end{proposition}
\begin{proof}
At each represented field value, duality for the positive $G_X$ pairing gives
\[
 |D\mathcal A_X[v]|^2
 \le\|D\mathcal A_X\|_{G_X^{-1}}^2\|v\|_{G_X}^2.
\]
Take the supported supremum of the second factor and apply
\Cref{thm:supp-action-stationarity}.  This proves the mean-square estimate.
Enumerate the determining core as $k_1,k_2,\ldots$.  Select increasing cutoffs
$X_j$ so that the right-hand side of
\eqref{eq:main-physical-stationarity} is at most $2^{-j}$ for each of
$k_1,\ldots,k_j$.  For a fixed $k_m$, the sum over $j\ge m$ of the expected
squared residuals is finite.  Tonelli's theorem implies that the corresponding
sum of squared residuals is finite almost surely, so the residual tends to
zero.  A countable intersection gives the assertion for the whole core.
No independence between cutoffs is used.  Uniform continuity of the limiting
variational functional on the declared test topology then extends zero from
the core to all determining tests.
\end{proof}

\begin{corollary}[Finite Gibbs/KL certificate]
\label{cor:supp-gibbs-gap}
For a finite reference row $p_\theta(y)>0$ and reconstructed same-history cost
$c_\theta(y)$, the row $q_\theta^*$ of \eqref{eq:main-gibbs-row} is the unique
entropy-proximal row and
\begin{equation}
 \boxed{\Delta_{{\rm Gibbs},\theta}
       =\eta^{-1}D_{\rm KL}
       \bigl(K^{\rm acc}(\cdot\mid\theta)\Vert q_\theta^*\bigr)}
 \label{eq:supp-gibbs-gap}
\end{equation}
is exactly its row variational gap.
\end{corollary}
\begin{proof}
For any probability row $q$,
\[
 \sum_yq(y)c_\theta(y)+\eta^{-1}D_{\rm KL}(q\Vert p_\theta)
 =\eta^{-1}D_{\rm KL}(q\Vert q_\theta^*)
       -\eta^{-1}\log\sum_zp_\theta(z)e^{-\eta c_\theta(z)}.
\]
Relative entropy is nonnegative and vanishes only at equality.  This proves a
statement about probability rows; it does not identify their sampled points
with interior critical points of the field action.
\end{proof}

\subsection{Global cut margin and compact-screen reconstruction}

Let a finite physical endpoint graph have vertex masses $m_{v,X}$, measure
$\mu_X$, symmetric conductances $c_{uv,X}$, and efficient edge capacities
$u_{e,X}$.  We use the two cut constants $I_X^{\rm eff}$ and
$I_X^{\rm sp}$ of \eqref{eq:main-cut-constant}; on the positive spatial branch
they are related by the finite capacity--conductance comparison
\eqref{eq:main-capacity-conductance}.

\begin{definition}[Spatial Dirichlet normalization]
\label{def:supp-spatial-dirichlet}
On the finite weighted graph $G_X=(V_X,E_X)$ set
\begin{equation}
 \mathcal E_X^{\rm sp}(f)
 =\sum_{\{u,v\}\in E_X}c_{uv,X}|f(u)-f(v)|^2,
 \qquad
 \langle f,\Delta_X^{\rm sp}f\rangle_{L^2(\mu_X)}
 =\mathcal E_X^{\rm sp}(f).
 \label{eq:supp-spatial-dirichlet}
\end{equation}
All scalar Sobolev estimates below are first stated for real-valued functions;
the real weighted Laplacian has the same spectrum after complexification.  Let
$P_{0,X}$ be the orthogonal projection onto the mean-zero subspace and put
$T_X(t)=e^{-t\Delta_X^{\rm sp}}P_{0,X}$.  Since the conductances are symmetric
and nonnegative, $e^{-t\Delta_X^{\rm sp}}$ is a finite Markov semigroup and is
contractive on every $L^q(\mu_X)$, $1\le q\le\infty$.
\end{definition}

\begin{lemma}[Cut inequality implies the half-volume $L^1$ Sobolev bound]
\label{lem:supp-half-volume-sobolev}
Assume
\begin{equation}
 h_Xc_X(\partial A)\ge I_*^{\rm sp}\,\mu_X(A)^{2/3}
 \label{eq:supp-cut-isoperimetry}
\end{equation}
whenever $0<\mu_X(A)\le\mu_X(V_X)/2$.  If $g\ge0$ and
$\mu_X(\supp g)\le\mu_X(V_X)/2$, then
\begin{equation}
 \boxed{
 \|g\|_{L^{3/2}(\mu_X)}
 \le \frac{h_X}{I_*^{\rm sp}}
 \sum_{\{u,v\}\in E_X}c_{uv,X}|g(u)-g(v)|.}
 \label{eq:supp-l1-sobolev}
\end{equation}
\end{lemma}

\begin{proof}
For $t\ge0$ put $A_t=\{v:g(v)>t\}$.  Since $A_t\subseteq\supp g$, every
nonempty $A_t$ lies on the half-volume side of
\eqref{eq:supp-cut-isoperimetry}.  The layer-cake identity gives
$g=\int_0^\infty\one_{A_t}\,\dd t$.  Minkowski's integral inequality in
$L^{3/2}(\mu_X)$ therefore yields
\begin{equation}
 \|g\|_{3/2}
 \le\int_0^\infty\|\one_{A_t}\|_{3/2}\,\dd t
 =\int_0^\infty\mu_X(A_t)^{2/3}\,\dd t.
 \label{eq:supp-layer-minkowski}
\end{equation}
Applying the cut inequality to every level set and then using the finite
coarea identity gives
\begin{align}
 \|g\|_{3/2}
 &\le \frac{h_X}{I_*^{\rm sp}}
       \int_0^\infty c_X(\partial A_t)\,\dd t \\
 &=\frac{h_X}{I_*^{\rm sp}}
   \sum_{\{u,v\}\in E_X}c_{uv,X}|g(u)-g(v)|,
\end{align}
which is \eqref{eq:supp-l1-sobolev}.
\end{proof}

\begin{lemma}[Discrete $L^6$ Sobolev and Poincar\'e inequalities]
\label{lem:supp-l6-poincare}
Assume \eqref{eq:supp-cut-isoperimetry},
$\mu_X(V_X)\le V_*$, and
\begin{equation}
 \sum_uc_{uv,X}\le D_*h_X^{-2}m_{v,X}
 \qquad(v\in V_X).
 \label{eq:supp-weighted-degree}
\end{equation}
Then every mean-zero $f$ satisfies
\begin{equation}
 \boxed{
 \|f\|_6^2\le C_{\rm S}\mathcal E_X^{\rm sp}(f),
 \qquad
 C_{\rm S}=\frac{128D_*}{(I_*^{\rm sp})^2},}
 \label{eq:supp-sobolev}
\end{equation}
and
\begin{equation}
 \boxed{
 \|f\|_2^2\le C_{\rm P}\mathcal E_X^{\rm sp}(f),
 \qquad C_{\rm P}=V_*^{2/3}C_{\rm S}.}
 \label{eq:supp-poincare}
\end{equation}
In particular
$\lambda_1(\Delta_X^{\rm sp}|_{\one^\perp})\ge C_{\rm P}^{-1}$.
\end{lemma}

\begin{proof}
First suppose $f\ge0$ and $\mu_X(\supp f)\le\mu_X(V_X)/2$.  Apply
\Cref{lem:supp-half-volume-sobolev} to $g=f^4$.  Since
$\|f^4\|_{3/2}=\|f\|_6^4$ and
$|a^4-b^4|\le4(a^3+b^3)|a-b|$ for $a,b\ge0$,
\begin{align}
 \|f\|_6^4
 &\le \frac{4h_X}{I_*^{\rm sp}}
 \sum_{\{u,v\}}c_{uv,X}(f(u)^3+f(v)^3)|f(u)-f(v)| \\
 &\le \frac{4h_X}{I_*^{\rm sp}}
 \left[\sum_{\{u,v\}}c_{uv,X}(f(u)^3+f(v)^3)^2\right]^{1/2}
 \mathcal E_X^{\rm sp}(f)^{1/2}.
 \label{eq:supp-sobolev-cs}
\end{align}
Using $(a^3+b^3)^2\le2(a^6+b^6)$ and
\eqref{eq:supp-weighted-degree},
\begin{align}
 \sum_{\{u,v\}}c_{uv,X}(f(u)^3+f(v)^3)^2
 &\le2\sum_v f(v)^6\sum_uc_{uv,X} \\
 &\le2D_*h_X^{-2}\|f\|_6^6.
 \label{eq:supp-degree-six}
\end{align}
Substitution into \eqref{eq:supp-sobolev-cs} and division by
$\|f\|_6^3$ gives
\begin{equation}
 \|f\|_6^2
 \le \frac{32D_*}{(I_*^{\rm sp})^2}\mathcal E_X^{\rm sp}(f)
 \label{eq:supp-half-sobolev-six}
\end{equation}
for every half-volume-supported nonnegative $f$.

For a general real $f$, choose a weighted median $m$ so that
$\mu_X(f>m)$ and $\mu_X(f<m)$ are each at most half the total volume.  Put
$f_+=(f-m)_+$ and $f_-=(m-f)_+$.  Truncation is $1$-Lipschitz and, edgewise,
\begin{equation}
 |f_+(u)-f_+(v)|^2+|f_-(u)-f_-(v)|^2
 \le |f(u)-f(v)|^2,
 \label{eq:supp-truncation-energy}
\end{equation}
so
$\mathcal E(f_+)+\mathcal E(f_-)\le\mathcal E(f)$.  Since the two functions
have half-volume support, \eqref{eq:supp-half-sobolev-six} gives
\begin{align}
 \|f-m\|_6^2
 &=\bigl(\|f_+\|_6^6+\|f_-\|_6^6\bigr)^{1/3} \\
 &\le \|f_+\|_6^2+\|f_-\|_6^2
 \le \frac{32D_*}{(I_*^{\rm sp})^2}\mathcal E(f).
 \label{eq:supp-median-six}
\end{align}
If $f$ has mean zero, then
\begin{equation}
 |m|\,\mu_X(V_X)^{1/6}
 \le \|f-m\|_6,
 \label{eq:supp-median-mean}
\end{equation}
because
$|m|=|\mu_X(V_X)^{-1}\int(f-m)\,\dd\mu_X|$ and H\"older gives
$\|f-m\|_1\le\mu_X(V_X)^{5/6}\|f-m\|_6$.  Therefore
$\|f\|_6\le2\|f-m\|_6$, which yields \eqref{eq:supp-sobolev} with the
stated factor four.  Finally,
\begin{equation}
 \|f\|_2^2\le\mu_X(V_X)^{2/3}\|f\|_6^2
 \le V_*^{2/3}C_{\rm S}\mathcal E(f),
\end{equation}
which is \eqref{eq:supp-poincare}.  Rayleigh--Ritz gives the spectral-gap
bound.
\end{proof}

\begin{lemma}[Nash inequality, ultracontractivity, and spectral counting]
\label{lem:supp-nash-counting}
Under the hypotheses of \Cref{lem:supp-l6-poincare}, every mean-zero $f$ obeys
\begin{equation}
 \boxed{
 \|f\|_2^{10/3}
 \le C_{\rm S}\mathcal E_X^{\rm sp}(f)\|f\|_1^{4/3}.}
 \label{eq:supp-nash}
\end{equation}
Moreover, for every $t>0$,
\begin{equation}
 \boxed{
 \|T_X(t)\|_{L^1\to L^\infty}
 \le4\left(\frac{3C_{\rm S}}{2t}\right)^{3/2}.}
 \label{eq:supp-ultracontractive}
\end{equation}
Consequently, if
$N_X(R)=\rank\one_{[0,R]}(\Delta_X^{\rm sp})$, then for $R>0$
\begin{equation}
 \boxed{
 N_X(R)\le1+4e^{3/2}V_*(C_{\rm S}R)^{3/2}.}
 \label{eq:supp-global-weyl-count}
\end{equation}
\end{lemma}

\begin{proof}
Interpolation between $L^1$ and $L^6$ gives
$\|f\|_2\le\|f\|_1^{2/5}\|f\|_6^{3/5}$.  Raising this inequality to the
power $10/3$ and using \eqref{eq:supp-sobolev} yields
\eqref{eq:supp-nash}.

Let $F_t=T_X(t)f$ and $u(t)=\|F_t\|_2^2$.  Since $F_t$ has mean zero and the
finite Markov semigroup is $L^1$-contractive,
\begin{equation}
 \|F_t\|_1
 \le\|P_{0,X}f\|_1
 \le2\|f\|_1.
 \label{eq:supp-p0-l1}
\end{equation}
Self-adjointness and finite dimensionality permit differentiation under the
spectral decomposition:
\begin{equation}
 u'(t)=-2\mathcal E_X^{\rm sp}(F_t).
 \label{eq:supp-energy-decay}
\end{equation}
Combining \eqref{eq:supp-nash}, \eqref{eq:supp-p0-l1}, and
\eqref{eq:supp-energy-decay} gives
\begin{equation}
 u'(t)
 \le-\frac{2}{C_{\rm S}2^{4/3}\|f\|_1^{4/3}}u(t)^{5/3}.
 \label{eq:supp-nash-ode}
\end{equation}
Integrating the differential inequality for $u^{-2/3}$ yields
\begin{equation}
 \|T_X(t)f\|_2^2
 \le \frac{3^{3/2}}{2}C_{\rm S}^{3/2}t^{-3/2}\|f\|_1^2.
 \label{eq:supp-l1-l2}
\end{equation}
Because $T_X(t/2)$ is self-adjoint and
$T_X(t)=T_X(t/2)T_X(t/2)$ on the mean-zero subspace, duality and the semigroup
property give
\begin{align}
 \|T_X(t)\|_{1\to\infty}
 &\le\|T_X(t/2)\|_{1\to2}\|T_X(t/2)\|_{2\to\infty} \\
 &=\|T_X(t/2)\|_{1\to2}^2
 \le4\left(\frac{3C_{\rm S}}{2t}\right)^{3/2},
\end{align}
which is \eqref{eq:supp-ultracontractive}.  This is the standard
Sobolev--Nash ultracontractivity route in finite weighted form; compare
\cite[Chs.~10--11]{Chung1997} and \cite{SaloffCoste2002}.

Let $K_t^0(v,w)$ be the integral kernel of $T_X(t)$ with respect to $\mu_X$.
The spectral expansion gives
$K_t^0(v,v)=\sum_{j\ge1}e^{-t\lambda_j}|\varphi_j(v)|^2\ge0$, and therefore
\begin{align}
 \Tr(T_X(t))
 &=\sum_v m_{v,X}K_t^0(v,v) \\
 &\le \mu_X(V_X)\|T_X(t)\|_{1\to\infty} \\
 &\le4V_*\left(\frac{3C_{\rm S}}{2t}\right)^{3/2}.
 \label{eq:supp-heat-trace}
\end{align}
Since the positive cut margin implies connectedness, the zero eigenvalue is
simple.  Hence
\begin{equation}
 (N_X(R)-1)e^{-tR}
 \le\sum_{0<\lambda_j\le R}e^{-t\lambda_j}
 \le\Tr(T_X(t)).
 \label{eq:supp-counting-trace}
\end{equation}
Choosing $t=3/(2R)$ gives \eqref{eq:supp-global-weyl-count}.
\end{proof}

\begin{theorem}[Operational Sobolev, Poincar\'e, and compact-screen bounds]
\label{thm:supp-spatial-screen}
Assume uniformly that $\mu_X(V_X)\le V_*$,
\eqref{eq:supp-weighted-degree} holds, the capacity--conductance comparison
\eqref{eq:main-capacity-conductance} holds, and
$I_X^{\rm eff}\ge I_*^{\rm eff}>0$.  Then
$I_X^{\rm sp}\ge I_*^{\rm eff}/a_+=I_*^{\rm sp}>0$ and the estimates
\eqref{eq:supp-sobolev}, \eqref{eq:supp-poincare}, and
\eqref{eq:supp-global-weyl-count} hold.  In addition,
\begin{equation}
 \boxed{
 \|(I-P_{X,R})f\|_2^2
 \le R^{-1}\mathcal E_X^{\rm sp}(f),
 \qquad P_{X,R}=\one_{[0,R]}(\Delta_X^{\rm sp}).}
 \label{eq:supp-screen-count-tail}
\end{equation}
Thus $P_{X,R}$ is a uniformly exhausted finite-rank compact-screen family.
A sequence of positive cut constants tending to zero gives the normalized
cut/potential obstruction \eqref{eq:main-neck-obstruction}; a zero-capacity
cut is retained directly as disconnected support.
\end{theorem}

\begin{proof}
The comparison \eqref{eq:main-cut-comparison} converts the efficient-capacity
margin into \eqref{eq:supp-cut-isoperimetry}.  The Sobolev and Poincar\'e
estimates are \Cref{lem:supp-l6-poincare}; the counting law is
\Cref{lem:supp-nash-counting}.  By the spectral theorem,
$\Delta_X^{\rm sp}\succeq R(I-P_{X,R})$, so
\eqref{eq:supp-screen-count-tail} follows immediately.  For a positive minimizing cut $A$, set
$\phi=\one_A/u_X(\partial A)$.  Then
$\sum_eu_{e,X}|\dd\phi(e)|=1$, and \eqref{eq:main-cut-demand} gives
$|\sum_vb_{A,X}(v)\phi(v)|=\mu_X(A)^{2/3}/u_X(\partial A)
=(I_X^{\rm eff})^{-1}$.  This is an explicit feasible dual witness.
A zero-capacity cut is recorded directly as disconnected support.
\end{proof}

\begin{proposition}[Compact positive-screen upgrade]
\label{prop:supp-compact-upgrade}
Let $Y_X\rightharpoonup Y$ in a Hilbert space and let positive operators
$\mathbb A_X$ satisfy
$mI\preceq\mathbb A_X\preceq MI$ and
$\mathbb A_X^{-1/2}\to\mathbb A^{-1/2}$ strongly.  Put
$Z_X=\mathbb A_X^{1/2}Y_X$ and suppose $Z_X\rightharpoonup Z$.  If for every
$R$ there are finite-rank screens $S_{X,R}\to S_R$ in operator norm, with
$S_R$ compact, and
\[
 \lim_{R\to\infty}\sup_X\|(I-S_{X,R})Z_X\|=0,
 \qquad
 \lim_{R\to\infty}\|(I-S_R)Z\|=0,
\]
then $Y_X\to Y$ strongly.  Consequently every bounded quadratic insertion in
$Y_X$ converges in $L^1$ on the corresponding compact spacetime cylinder.
\end{proposition}

\begin{proof}
For fixed $R$, compactness of $S_R$ converts weak convergence into
$S_RZ_X\to S_RZ$ strongly, and operator-norm convergence gives
$S_{X,R}Z_X-S_RZ_X\to0$.  Decompose
\[
 Z_X-Z=(I-S_{X,R})Z_X+(S_{X,R}Z_X-S_RZ)+(S_R-I)Z.
\]
Taking the limsup and then $R\to\infty$ gives $Z_X\to Z$ strongly.  Uniform
coercivity and strong convergence of the inverse square roots then give
$Y_X\to Y$.  The quadratic statement follows from
$\|B(Y_X,Y_X)-B(Y,Y)\|_1\le C\|Y_X-Y\|_2(\|Y_X\|_2+\|Y\|_2)$.
\end{proof}

\subsection{Derivative-free face linearization and complete memory shorting}

The first-order gravitational score is not assumed as a derivative at a zero
face.  It is obtained from renewal subdivision after all resolved cut-memory
coordinates have been retained.

\begin{theorem}[Derivative-free renewal face linearization]
\label{thm:supp-face-linearization}
Let $E$ be a finite-dimensional real normed space, $U\subset E$ a balanced
convex neighborhood of zero, and $A:U\to\R$ with $A(0)=0$.  If
\begin{equation}
 \boxed{A(z)=2A(z/2)}
 \label{eq:supp-face-divisibility}
\end{equation}
for every $z\in U$, $A$ is bounded near zero, and
\[
 |A(x+y)-A(x)-A(y)|\le H\|x\|\,\|y\|
\]
near zero, then $A$ is the restriction of a unique real linear functional.
More generally, if the renormalized scores
$L_m(z)=2^mA_m(2^{-m}z)$ have summable adjacent-cutoff defects and their
additivity defect tends to zero, then $L_m$ converges locally uniformly to one
continuous linear functional.
\end{theorem}

\begin{proof}
For the exact branch iterate \eqref{eq:supp-face-divisibility}:
\[
 A(x+y)-A(x)-A(y)
 =2^m[A(2^{-m}(x+y))-A(2^{-m}x)-A(2^{-m}y)].
\]
The right-hand side is bounded by $H2^{-m}\|x\|\|y\|$ and therefore vanishes as
$m\to\infty$.  Thus $A$ is additive; boundedness near zero gives continuity and
hence real homogeneity.  On the stable branch the adjacent-cutoff estimate makes
$L_m$ locally uniformly Cauchy; passing to the limit in the additivity defect
gives a continuous linear limit.
\end{proof}

\begin{theorem}[Complete resolved-memory subtraction]
\label{thm:supp-memory-short}
Let the resolved cut-memory state be
$\rho_{\rm cq}=\bigoplus_\omega p_\omega\rho_\omega$ with faithful blocks and
linear perturbation
$K_{\rm cq}(z)=\bigoplus_\omega K_\omega(z)$.  Put
\begin{equation}
 \boxed{
 M_{\rm cq}(z)
 =\log\Tr\exp(\log\rho_{\rm cq}+K_{\rm cq}(z))-\bar m(z),}
 \label{eq:supp-memory-logpartition}
\end{equation}
where $\bar m$ is the average linear slope.  Then
$M_{\rm cq}(0)=DM_{\rm cq}(0)=0$, its Hessian is the sum of a classical branch
covariance and a positive Kubo--Mori form, and
\begin{equation}
 \boxed{A_{\rm sh}=A_{\rm coarse}-M_{\rm cq}}
 \label{eq:supp-memory-short}
\end{equation}
is the canonical fully memory-shorted face score.  Thus unresolved memory is
removed at the complete log-partition level rather than only to second order.
\end{theorem}

\begin{proof}
The block exponential and trace split over the classical label.  Duhamel's
formula gives the first derivative in each faithful block and the Kubo--Mori
quadratic form as the second logarithmic derivative.  Differentiating the outer
finite log-sum-exp adds the covariance of the branch means.  Both terms are
positive.  Subtracting the complete log-partition therefore removes exactly the
resolved cut-memory contribution and leaves the displayed shorted score.
\end{proof}

\subsection{Connection stationarity and torsion}

Let $B^{IJ}=e^I\wedge e^J$, $T^I=De^I$, and
$P_{\alpha,\beta}=\alpha I+\beta\star$, with the same
spacetime-constant coefficients as in \eqref{eq:main-phv-coefficients}.

\begin{theorem}[Connection stationarity and the torsion alternative]
\label{thm:supp-palatini-torsion}
If $(\alpha,\beta)\ne(0,0)$, then
\begin{equation}
 \boxed{
 P_{\alpha,\beta}^{-1}
 =\frac{\alpha I-\beta\star}{\alpha^2+\beta^2}.}
 \label{eq:supp-palatini-inverse}
\end{equation}
For the Palatini--Holst action the spinless connection Euler equation is
$D(P_{\alpha,\beta}B)=0$.  For a nondegenerate coframe, the Cartan map
\[
 \mathcal K_e(T)^{IJ}=T^I\wedge e^J-e^I\wedge T^J
\]
is injective.  Consequently connection stationarity implies $T=0$; at finite
cutoff, a uniform lower singular-value bound on $\mathcal K_{e,X}$ controls the
torsion norm by the connection Euler residual.
\end{theorem}

\begin{proof}
The inverse formula follows from $\star^2=-I$ on Lorentzian bivectors.  Varying
the connection and integrating by parts gives the Euler source.  To prove
injectivity, expand $T^I=\frac12T^I{}_{KL}e^K\wedge e^L$ in a coframe basis.
Choosing a fourth index in
$T^I\wedge e^J-e^I\wedge T^J=0$ first eliminates all components with
$K,L\ne I$; the remaining coefficients satisfy pairwise sign relations that,
in four dimensions, force them all to vanish.  Inverting the two maps gives the
finite estimate.
\end{proof}

\subsection{Lorentz-natural classification, curvature propagation, and Palatini handoff}

Define, up to overall normalization,
\begin{align}
 L_{\rm Palatini}
 &=\frac12\epsilon_{IJKL}e^I\wedge e^J\wedge R^{KL},\\
 L_{\rm Holst}
 &=e^I\wedge e^J\wedge R_{IJ},\\
 L_{\rm vol}
 &=\frac1{4!}\epsilon_{IJKL}
   e^I\wedge e^J\wedge e^K\wedge e^L.
 \label{eq:supp-phv-densities}
\end{align}

\begin{lemma}[Strong interpolation for the reconstructed coframes]
\label{lem:supp-coframe-interpolation}
Let $K$ be a compact spacetime cylinder with finite reconstructed volume measure.
Suppose the common-cylinder interpolants of the finite coframes satisfy
\begin{equation}
 e_X\to e\quad\hbox{strongly in }L^2(K),
 \qquad
 \sup_X\|e_X\|_{L^6(K)}\le M_6<\infty.
 \label{eq:supp-coframe-l2-l6}
\end{equation}
Then $e\in L^6(K)$ and, for every $2\le q<6$,
\begin{equation}
 \boxed{e_X\to e\quad\hbox{strongly in }L^q(K).}
 \label{eq:supp-strong-lq}
\end{equation}
More precisely, if
$\theta_q=3/q-1/2\in(0,1]$, then
\begin{equation}
 \|e_X-e\|_{L^q}
 \le \|e_X-e\|_{L^2}^{\theta_q}
      \|e_X-e\|_{L^6}^{1-\theta_q}.
 \label{eq:supp-interpolation-explicit}
\end{equation}
If $p>3/2$ and $p'=p/(p-1)$, then $2p'<6$ and
\begin{equation}
 \boxed{
 e_X\wedge e_X\longrightarrow e\wedge e
 \quad\hbox{strongly in }L^{p'}(K).}
 \label{eq:supp-bivector-strong}
\end{equation}
\end{lemma}

\begin{proof}
The uniform $L^6$ bound gives a weakly convergent subsequence in $L^6(K)$.
Because $K$ has finite measure, weak $L^6$ convergence implies distributional
convergence, while the strong $L^2$ convergence has the unique distributional
limit $e$.  Hence every weak $L^6$ cluster point is $e$, so $e\in L^6(K)$ and
$\|e\|_6\le M_6$.  The difference $e_X-e$ is therefore uniformly bounded in
$L^6$.  The Riesz--Thorin/H\"older interpolation identity
$1/q=\theta_q/2+(1-\theta_q)/6$ gives
\eqref{eq:supp-interpolation-explicit}; its right-hand side tends to zero for
$q<6$, proving \eqref{eq:supp-strong-lq}.

For $p>3/2$, $p'<3$ and hence $q=2p'<6$.  The exterior product is a bounded
bilinear map on the finite-dimensional coframe fibres, so pointwise
$|a\wedge b|\le C|a||b|$.  Therefore H\"older's inequality gives
\begin{align}
 \|e_X\wedge e_X-e\wedge e\|_{L^{p'}}
 &\le C\|(e_X-e)\wedge e_X+e\wedge(e_X-e)\|_{L^{p'}} \\
 &\le C(\|e_X\|_{L^{2p'}}+\|e\|_{L^{2p'}})
       \|e_X-e\|_{L^{2p'}}\longrightarrow0,
\end{align}
which proves \eqref{eq:supp-bivector-strong}.
\end{proof}

\subsubsection{Propagation and differentiated source control}

\begin{theorem}[Finite curvature-energy propagation from initial data]
\label{thm:supp-curvature-propagation}
At a fixed cutoff let $y_X=(E_X,B_X)$, $M_X\succ0$, $K_X$, $L_X$, and $f_X$
be as in \eqref{eq:main-curvature-propagation-writer}, and set
$z_X^2=y_X^TM_Xy_X$.  Then
\begin{equation}
 \frac12\frac{\dd}{\dd t}z_X^2
 \le a_Xz_X^2+z_X\|f_X\|_{M_X},
 \label{eq:supp-curvature-energy-differential}
\end{equation}
and therefore \eqref{eq:main-curvature-energy-propagation} holds.  Under the
stable incidence and uniform budgets
\eqref{eq:main-curvature-incidence}--\eqref{eq:main-curvature-budgets},
\eqref{eq:main-derived-l2-curvature} follows.
\end{theorem}

\begin{proof}
Differentiate $z_X^2=y_X^TM_Xy_X$ and substitute
\eqref{eq:main-curvature-propagation-writer}.  The principal term cancels
exactly because $M_XK_X+K_X^TM_X=0$.  By
\eqref{eq:main-curvature-growth}, the remaining quadratic term is at most
$2a_Xz_X^2$, while Cauchy--Schwarz in the positive $M_X$-Gram bounds the source
term by $2z_X\|f_X\|_{M_X}$.  Applying the integrating-factor argument to
$(z_X^2+\varepsilon^2)^{1/2}$ and sending $\varepsilon\downarrow0$ gives
\eqref{eq:main-curvature-energy-propagation}.  The interpolation estimate in
\eqref{eq:main-curvature-incidence}, integration against the physical lapse and
slice measure, and the uniform budgets give
\eqref{eq:main-derived-l2-curvature}.
\end{proof}

\begin{theorem}[Midpoint propagation with the physical Gram transfer]
\label{thm:supp-discrete-curvature}
For the actual normalized recurrence
\eqref{eq:main-discrete-curvature-update}, assume
$\|T_j\|_{\op}\le e^{d_j}$, $K_j^T=-K_j$,
$(L_j+L_j^T)/2\preceq a_jI$, and $s_ja_j/2\le b<1$.
Then \eqref{eq:main-discrete-curvature-bound} holds.  With the stable
curvature incidence of \eqref{eq:main-curvature-incidence} and stable interval
interpolation, it gives the fully discrete curvature certificate of
\Cref{thm:main-discrete-curvature}.
\end{theorem}
\begin{proof}
Write $A_j=K_j+L_j$.  For any real vector $v$,
\[
 \langle v,(I-s_jA_j/2)v\rangle
       \ge(1-s_ja_j/2)\|v\|^2.
\]
Cauchy--Schwarz gives
$\|(I-s_jA_j/2)v\|\ge(1-s_ja_j/2)\|v\|$.  Thus the finite matrix is
invertible and its inverse norm is at most $(1-s_ja_j/2)^{-1}$.

Let $v=T_jY_j$ and first let $w$ solve the step with zero residual.  Pair
$w-v=s_jA_j(w+v)/2$ with $(w+v)/2$ to obtain
\[
 \|w\|^2-\|v\|^2
 \le\frac{s_ja_j}{2}(\|w\|+\|v\|)^2.
\]
Dividing by $\|w\|+\|v\|$ when this number is nonzero, and treating the
zero case directly, gives
\[
 \|w\|\le\frac{1+s_ja_j/2}{1-s_ja_j/2}\|v\|.
\]
The inhomogeneous solution adds
$(I-s_jA_j/2)^{-1}\varepsilon_j$.  Consequently
\[
 \|Y_{j+1}\|\le e^{d_j}
       \frac{1+s_ja_j/2}{1-s_ja_j/2}\|Y_j\|
       +\frac{\|\varepsilon_j\|}{1-b}.
\]
For $0\le x\le b$,
$\log((1+x)/(1-x))\le2x/(1-b)$.  Every product of step factors is
therefore bounded by $\exp(D_X+A_X^{\rm d}/(1-b))$.
Iteration and weighted Cauchy--Schwarz,
\[
 \sum_j\|\varepsilon_j\|
       \le\sqrt{T_X\sum_j\|\varepsilon_j\|^2/s_j},
\]
prove \eqref{eq:main-discrete-curvature-bound}.  The size of the skew matrix
$K_j$ has entered neither estimate.

For raw curvature records $y_j$ with Gram $M_j$, the transfer is
$T_j=M_{j+1}^{1/2}\mathcal T_jM_j^{-1/2}$, where $\mathcal T_j$ is the
actual recorded frame/transport map.  Its norm excess is retained in $d_j$.
Stable interval interpolation bounds the reconstructed slice norm by a fixed
multiple of the largest endpoint norm.  Integrating with the physical lapse
and adding the bounded reconstruction error yields, for an appropriate
interpolation constant $C_I$,
\[
 \|R_X\|_{L^2(K)}
 \le C_I\sqrt{N_*T_X}\,e^{D_X+A_X^{\rm d}/(1-b)}
       \left(\|Y_0\|+\frac{\sqrt{T_X\mathcal D_X}}{1-b}\right)+C_\zeta.
\]
All interface and incoming boundary contributions must already be included
in the actual transfer, residual or reconstruction comparison.  This proof
does not assert that arbitrary nonlinear coordinate changes preserve the
midpoint recurrence: any such discrepancy is part of $\varepsilon_j$.
\end{proof}

\paragraph{Differentiated gravitational source.}
The source norm is not the norm of an undifferentiated Einstein residual.
On the actual shared-link branch the forcing contains the differentiated,
volume-normalized Palatini metric row, together with the differentiated
Cartan/torsion defect and the differential Bianchi defect.  Lower-order
transport terms are retained in $a_X$.  To state a quantitative criterion,
use the physical incidence \eqref{eq:main-current-incidence}, with all
untransferred contributions collected in $r_X^{\rm src}$.  Its verification
is an assertion about the actual action, connection and observation records;
it is not obtained simply by introducing an auxiliary source variable.

\begin{proposition}[Exact amplification of the differentiated action current]
\label{prop:supp-sharp-current-transfer}
Let $G_X\succ0$ define the action dual norm and let $H_X\succ0$ be the target
source Gram.  For the actual linear transfer $\mathsf T_X$, the least constant
in
\[
 \|\mathsf T_X\mathfrak a\|_{H_X}^2
       \le\Gamma_X\|\mathfrak a\|_{G_X^{-1}}^2
\]
is \eqref{eq:main-sharp-current-transfer}.  If
$\|\mathfrak a_X\|_{G_X^{-1}}^2\le C_X\Delta_X$ pointwise, then
\begin{equation}
 \|f_X\|_{H_X}^2
       \le2\Gamma_XC_X\Delta_X+2\|r_X^{\rm src}\|_{H_X}^2.
 \label{eq:supp-curvature-source-budget}
\end{equation}
This proves \Cref{prop:main-action-current}.  With deterministic supported
upper bounds on $\Gamma_X$, the same inequality holds in expectation when
the action-gradient bound holds only in expectation.
\end{proposition}
\begin{proof}
Write $\mathfrak a=G_X^{1/2}v$.  Then
$\|\mathfrak a\|_{G_X^{-1}}=\|v\|$ and
\[
 \|\mathsf T_X\mathfrak a\|_{H_X}
       =\|H_X^{1/2}\mathsf T_XG_X^{1/2}v\|.
\]
The maximum squared ratio is the largest squared singular value.  A unit
maximizing right singular vector attains it, so the constant is sharp and
also supplies a finite witness to large source amplification.
Apply $\|u+v\|^2\le2\|u\|^2+2\|v\|^2$ to
\eqref{eq:main-current-incidence}.  Time Cauchy--Schwarz proves
\eqref{eq:main-action-source-budget}; summing the same pointwise inequality
with weights $s_j$ proves its discrete counterpart.  For varying field
values the argument is pointwise before averaging.  Without a deterministic
supported bound, the expected estimate must retain
$\mathbb E[\Gamma_X\|\mathfrak a_X\|_{G_X^{-1}}^2]$, not a product of
separately averaged factors.
\end{proof}

\begin{proposition}[Sharp grid-scale source amplification]
\label{prop:supp-current-derivative-loss}
On a periodic grid of $m$ sites, with $m$ divisible by four, spacing $h$ and
inner product $h\sum_i u_iv_i$, let
$(D_hu)_i=(u_{i+1}-u_{i-1})/(2h)$.  With the same positive source and target
norms,
\begin{equation}
 \|D_h\|_{2\to2}^2=h^{-2}.
 \label{eq:supp-current-derivative-loss}
\end{equation}
In particular, a vanishing unscaled action-gradient norm does not imply a
bounded differentiated current.
\end{proposition}
\begin{proof}
The Fourier multiplier is $i\sin(2\pi k/m)/h$, whose maximum modulus is
$h^{-1}$ and is attained at $k=m/4$.  The same singular value occurs on the
real sine/cosine subspace.  Choose a unit vector $v_h$ in that subspace and
set $\mathfrak a_h=h^{1/2}v_h$.  Then
$\|\mathfrak a_h\|_2^2=h\to0$, whereas
$\|D_h\mathfrak a_h\|_2^2=h^{-1}\to\infty$.
\end{proof}

Thus, for this derivative transfer and uniformly controlled remaining
factors, $C_X\Delta_X=O(h_X^2)$ suffices for bounded source and
$o(h_X^2)$ suffices for vanishing source.  These rates are not inferred from
$\Delta_X\to0$.  Cancellations in the actual Bianchi transfer may improve
$\Gamma_X$, but they must be retained in that map and its physical Gram.

\begin{corollary}[Expected curvature energy and a sufficient pathwise upgrade]
\label{cor:supp-stochastic-curvature}
Suppose the finite propagation equation holds pathwise on $[0,T_*]$, with
$\int a_X\le A_*$ deterministically.  Suppose its forcing satisfies
\eqref{eq:main-current-incidence}, with the target norm already transferred
to the curvature-energy norm, and
$\mathbb E\|\mathfrak a_X(t)\|_{G_X^{-1}}^2\le C_X(t)\Delta_X(t)$.
Using deterministic supported upper bounds on $\Gamma_X$, one obtains
\begin{equation}
 \begin{split}
 \mathbb E\sup_{t\le T_*}z_X(t)^2
 \le 2e^{2A_*}\bigg[&\mathbb Ez_X(0)^2
       +2T_*\int_0^{T_*}\Gamma_XC_X\Delta_X\,\dd t\\
       &+2T_*\mathbb E\int_0^{T_*}\|r_X^{\rm src}\|_{H_X}^2\,\dd t\bigg].
 \end{split}
 \label{eq:supp-expected-curvature}
\end{equation}
Together with the stable interpolation and lapse bounds, uniformity of these
budgets gives a uniform expected squared spacetime curvature norm.  For a
countable cofinal sequence, a sufficient pathwise strengthening is
\begin{equation}
 \sup_Xz_X(0)<\infty\ \hbox{almost surely},\qquad
 \sum_X\mathbb E\int_0^{T_*}\|f_X(t)\|_{H_X}^2\,\dd t<\infty,
 \label{eq:supp-pathwise-source-budget}
\end{equation}
with pathwise uniform growth, interpolation, lapse and boundary budgets.
Under these conditions $\sup_X\|R_X\|_{L^2(K)}<\infty$ almost surely.
\end{corollary}
\begin{proof}
The pathwise estimate gives
$\sup_tz_X(t)\le e^{A_*}(z_X(0)+\int\|f_X\|)$.
Square and apply Cauchy--Schwarz in time to obtain
\[
 \mathbb E\sup_tz_X(t)^2
 \le2e^{2A_*}\left(\mathbb Ez_X(0)^2
       +T_*\mathbb E\int_0^{T_*}\|f_X\|_{H_X}^2\,\dd t\right).
\]
Apply \Cref{prop:supp-sharp-current-transfer} and the interpolation estimate.
For the last assertion, Tonelli makes the sum of the nonnegative integrated
source energies finite almost surely.  Their supremum is then finite, and
the deterministic propagation estimate applies to the whole realized
sequence.  A uniform expected norm alone gives boundedness in probability
by Markov's inequality; it is not substituted for this pathwise argument.
\end{proof}

\subsubsection{Curvature identification and interface lifting}

\begin{theorem}[Weak curvature compactness with an identified connection limit]
\label{thm:supp-curvature-compactness}
Let $K$ be a compact cylinder and $1<p<\infty$.  A bound
\begin{equation}
 \sup_X\|R_X\|_{L^p(K)}\le C_{K,p}<\infty
 \label{eq:supp-curvature-lp-bound}
\end{equation}
gives a weakly convergent subsequence
\begin{equation}
 \boxed{R_X\rightharpoonup R\quad\hbox{in }L^p(K).}
 \label{eq:supp-curvature-reflexive-limit}
\end{equation}
If, in addition, $\omega_X\to\omega$ strongly in $L^2_{\rm loc}$ and
\begin{equation}
 R_X-[\dd_{\rm dist}\omega_X+\omega_X\wedge\omega_X]
       \longrightarrow0\quad\hbox{in distributions},
 \label{eq:supp-curvature-identification}
\end{equation}
then $R=\dd\omega+\omega\wedge\omega$ distributionally.
The finite-writer conditions
\eqref{eq:main-cartan-curvature-residual}--\eqref{eq:main-cartan-incidence},
with exhaustive determining tests and their interpolation control, give
\eqref{eq:supp-curvature-identification}.  In particular the propagated
$p=2$ bound proves \Cref{thm:main-curvature-compactness}.
\end{theorem}
\begin{proof}
For $1<p<\infty$, reflexivity gives weak sequential compactness of bounded
balls of $L^p(K)$.  Strong $L^2$ convergence of the connection gives
convergence of its distributional derivative, and
\[
 \|\omega_X\wedge\omega_X-\omega\wedge\omega\|_{L^1(K)}
 \le C(\|\omega_X\|_2+\|\omega\|_2)\|\omega_X-\omega\|_2\to0.
\]
Thus \eqref{eq:supp-curvature-identification} identifies the distributional
limit of $R_X$.  Weak $L^p$ convergence has the same distributional limit,
proving the assertion.  For the finite-writer version, split the tested
residual into $R_X-\mathcal R_X(\omega_X)$ and the discrepancy
\eqref{eq:main-cartan-incidence}.  The first tends to zero on the exhaustive
core by \eqref{eq:main-cartan-curvature-residual}, and the second by the
stated consistency.  The declared test-interpolation control extends the
identity to the determining test topology.  A diagonal subsequence works on
a countable compact exhaustion.

If the curvature norms are instead unbounded, choose a subsequence on which
they diverge and retain the unit fields $R_X/\|R_X\|_{L^p(K)}$ together
with their diverging scales.  This witnesses failure of the specified
curvature-size certificate, not nonexistence of every possible weaker limit.
\end{proof}

For a piecewise smooth connection, its full distributional derivative includes
face measures.  A bound on cell interiors alone therefore does not control the
curvature entering the continuum variation.  The following construction
specifies a sufficient volume realization without discarding these measures.

Work on a periodic rectangular four-dimensional domain with a uniform cubical
mesh of side $h$, or on a compact interior subdomain together with its
neighbouring cells.  Polynomial degrees are fixed, positive Grams and physical
measures are uniformly comparable to the reference ones, and all face
coefficients are compared in a common local trivialization.  For an oriented
interior face $f$, set $J_f=n_f^\flat\wedge[\omega_h]_f$, with the sign
chosen so that
\begin{equation}
 \dd_{\rm dist}\omega_h=\dd_b\omega_h+\sum_fJ_f\delta_f.
 \label{eq:supp-connection-jumps}
\end{equation}
Let $\eta\ge0$ be a fixed smooth function supported in $(-1/4,1/4)$ with
integral one.  On the normal prism of $f$, define
\begin{equation}
 (\mathcal L_fJ_f)(y+rn_f)=h^{-1}\eta(r/h)J_f(y),\qquad
 \mathcal L_hJ=\sum_f\mathcal L_fJ_f,
 \label{eq:supp-interface-lifting}
\end{equation}
and extend each term by zero outside its prism.

\begin{lemma}[Interface lifting and its distributional error]
\label{lem:supp-interface-lifting}
With $\mathcal J_h$ defined in \eqref{eq:main-interface-budget},
\begin{equation}
 \|\mathcal L_hJ\|_2\le C\mathcal J_h,\qquad
 \left|\left\langle\mathcal L_hJ-\sum_fJ_f\delta_f,\Phi\right\rangle\right|
       \le C_Kh\mathcal J_h\|\nabla\Phi\|_\infty
 \label{eq:supp-interface-lifting-bounds}
\end{equation}
for each fixed smooth compactly supported or periodic test $\Phi$.
\end{lemma}
\begin{proof}
Integration in the normal coordinate gives
\[
 \|\mathcal L_fJ_f\|_2^2
       =h^{-1}\|\eta\|_{L^2(\R)}^2\|J_f\|_{L^2(f)}^2.
\]
The prisms have bounded overlap depending only on dimension.  Pointwise
Cauchy--Schwarz over the overlapping prisms proves the first bound.
For the second, subtract $\Phi(y)$ from $\Phi(y+rn_f)$ and use
$|r|\le h/4$.  The total error is at most
$Ch\|\nabla\Phi\|_\infty\sum_f\|J_f\|_{L^1(f)}$.  Moreover
\[
 \sum_fh\|J_f\|_{L^1(f)}
 \le\left(\sum_fh^{-1}\|J_f\|_{L^2(f)}^2\right)^{1/2}
       \left(\sum_fh^3|f|\right)^{1/2}
 \le C_Kh\mathcal J_h.
\]
There are $O_K(h^{-4})$ faces of area $h^3$, which proves the last
inequality and the result.
\end{proof}

\begin{theorem}[Constructive connection-and-interface curvature criterion]
\label{thm:supp-interface-curvature}
Suppose
\begin{equation}
 \|\dd_b\omega_h\|_{L^2(K)}+\|\omega_h\|_{L^4(K)}^2+
       \mathcal J_h\le C_K.
 \label{eq:supp-connection-interface-budget}
\end{equation}
Then the volume curvature $R_h^{\rm lift}$ in
\eqref{eq:main-lifted-curvature} is uniformly bounded in $L^2(K)$ and has
the displayed distributional error.  If also
$\omega_h\to\omega$ strongly in $L^2_{\rm loc}$, every weak $L^2$ cluster
point equals $\dd\omega+\omega\wedge\omega$ distributionally.
\end{theorem}
\begin{proof}
The matrix exterior product satisfies
$|\omega_h\wedge\omega_h|\le C|\omega_h|^2$.  H\"older's inequality and
\Cref{lem:supp-interface-lifting} give
\eqref{eq:main-lifted-curvature}.  The discrepancy from
$\dd_{\rm dist}\omega_h+\omega_h\wedge\omega_h$ is exactly the lifting
error in \eqref{eq:supp-interface-lifting-bounds}, which tends to zero on
each fixed test.  Apply \Cref{thm:supp-curvature-compactness} with $p=2$.
\end{proof}

The theorem is a sufficient identification and size criterion with its own
explicit connection and jump budgets.  A stationary connection estimate in
$L^2$ does not imply \eqref{eq:supp-connection-interface-budget}.
Nor is $R_h^{\rm lift}$ silently substituted for the action's curvature
writer: that writer must agree with the lift on the determining tests up to
a vanishing recorded discrepancy.  On a conforming connection reconstruction
$J_f=0$, and the interface part is absent.

\subsubsection{The weak--strong Palatini passage}

\begin{lemma}[Weak--strong curvature pairing and volume convergence]
\label{lem:supp-weak-strong-palatini}
Assume \eqref{eq:supp-coframe-l2-l6}, let $p>3/2$, and suppose
\begin{equation}
 R_X\rightharpoonup R\quad\hbox{weakly in }L^p(K).
 \label{eq:supp-curvature-weak}
\end{equation}
Then, for every bounded compactly supported coefficient tensor $\Phi$,
\begin{equation}
 \boxed{
 \int_K\!\langle\Phi(e_X\wedge e_X),R_X\rangle
 \longrightarrow
 \int_K\!\langle\Phi(e\wedge e),R\rangle.}
 \label{eq:supp-weak-strong-pairing}
\end{equation}
The same conclusion holds with an internal Hodge star applied to either
bivector factor.  Moreover
\begin{equation}
 \boxed{
 e_X\wedge e_X\wedge e_X\wedge e_X
 \longrightarrow
 e\wedge e\wedge e\wedge e
 \quad\hbox{strongly in }L^1(K).}
 \label{eq:supp-volume-strong}
\end{equation}
\end{lemma}

\begin{proof}
Weak convergence in $L^p$ implies $\sup_X\|R_X\|_{L^p}<\infty$.  Put
$B_X=\Phi(e_X\wedge e_X)$ and $B=\Phi(e\wedge e)$.  By
\Cref{lem:supp-coframe-interpolation}, $B_X\to B$ strongly in $L^{p'}$.
Split the difference as
\begin{align}
 \int_K\langle B_X,R_X\rangle-\int_K\langle B,R\rangle
 &=\int_K\langle B_X-B,R_X\rangle
   +\int_K\langle B,R_X-R\rangle.
 \label{eq:supp-pairing-split}
\end{align}
The first term tends to zero by H\"older and the uniform $L^p$ bound on
$R_X$; the second tends to zero by the definition of weak $L^p$ convergence,
because $B\in L^{p'}$.  The internal Hodge star is a fixed bounded linear map
on the bivector fibre and changes none of these estimates.

For the volume term, \Cref{lem:supp-coframe-interpolation} with $q=4$ gives
$e_X\to e$ strongly in $L^4(K)$ and uniform $L^4$ boundedness.  Expanding the
fourfold difference one factor at a time and applying H\"older with four
$L^4$ factors gives
\begin{align}
 &\|e_X^{\wedge4}-e^{\wedge4}\|_{L^1} \\
 &\qquad\le C\|e_X-e\|_{L^4}
 (\|e_X\|_{L^4}+\|e\|_{L^4})^3\longrightarrow0,
\end{align}
which is \eqref{eq:supp-volume-strong}.
\end{proof}

\begin{remark}[Weak--strong passage without compensated compactness]
\label{rem:supp-no-divcurl}
The product passage above does not use a div--curl or compensated-compactness
lemma.  Such a device would be needed if both the coframe bivector and curvature
were only weakly convergent.  Here the compact-screen hypothesis gives strong
$L^2$ convergence of the coframe together with a uniform $L^6$ bound; the
explicit interpolation in \Cref{lem:supp-coframe-interpolation} upgrades this
to strong $L^{2p'}$, so ordinary weak--strong duality suffices.  At this stage
all fields are already represented by their canonical common-cylinder
interpolants; equivalently, the same estimates may be read cellwise with the
physical cell measures because the displayed $L^q$ hypotheses are precisely
the cutoff-uniform interpolation bounds used in the handoff.
\end{remark}

\begin{theorem}[Localized renewal--Palatini handoff]
\label{thm:supp-renewal-palatini}
Assume the face-linearization and complete memory-short branches, real Lorentz
naturality, scalar residual typed commutant on the bivector carrier, and
represented determinant volume.  Let the constant calibrated coefficients
converge to $(\alpha,\beta,\lambda)$ with $\beta\ne0$, and let the spinless
connection Euler residual vanish with the uniform Cartan inverse of
\Cref{thm:supp-palatini-torsion}.  On every compact $K$, assume
\eqref{eq:supp-coframe-l2-l6}, nondegeneracy of the limiting coframe, and either
the continuous or the discrete initial-energy curvature certificate.
Assume also the connection identification of
\Cref{thm:supp-curvature-compactness} and the bounded, convergent physical test
lifts of \textup{(E1)} in \Cref{thm:main-renewal-einstein}.
Then the limiting score has the classified form
\begin{equation}
 \boxed{\alpha L_{\rm Holst}+\beta L_{\rm Palatini}+\lambda L_{\rm vol},}
 \label{eq:supp-phv-classification}
\end{equation}
limiting torsion vanishes, and the metric first variations converge to
\begin{equation}
 \boxed{\chi\int_K\sqrt{-g}\,
       (G_{\mu\nu}+\Lambda g_{\mu\nu})k^{\mu\nu},\qquad\chi\ne0.}
 \label{eq:supp-einstein-insertion}
\end{equation}
Here $\chi$ is the nonzero Palatini normalization and $\Lambda$ is the
corresponding calibrated volume-to-Palatini coefficient ratio, in the density
conventions of \eqref{eq:supp-phv-densities}.
\end{theorem}
\begin{proof}
Both coframe bivectors and Lorentz curvature lie in
$\Lambda^2\R^{1,3}$.  After complexification the bivector module splits into
nonisomorphic self-dual and anti-self-dual irreducibles.  The real commutant is
therefore
\begin{equation}
 \operatorname{End}_{SO^+(1,3)}(\Lambda^2\R^{1,3})
       =\Span_\R\{I,\star\}.
 \label{eq:supp-lorentz-commutant}
\end{equation}
The residual typed commutant supplies no additional internal tensor, and the
only alternating four-coframe scalar at this order is determinant volume.
Thus \eqref{eq:supp-phv-classification} follows in the stated first-order
class.  The calibrated coefficient limits ensure that these are the
coefficients of the limiting variational functional, rather than unrelated
values chosen after taking the limit.

By \Cref{thm:supp-palatini-torsion}, the vanishing connection Euler residual
and the uniform inverse bounds give vanishing torsion.  Strong local $L^2$
convergence of $e_X$ and $\omega_X$ passes their product in $L^1$; the
distributional derivative of $e_X$ passes against every fixed smooth test.
The limit is therefore a metric-compatible torsion-free represented
connection.  Its curvature is identified by
\Cref{thm:supp-curvature-compactness}, not merely by matching finite
curvature labels.

The continuous or discrete propagation theorem supplies
$\sup_X\|R_X\|_{L^2(K)}<\infty$.  After passage to the identified weak
subsequence, \Cref{lem:supp-coframe-interpolation} gives
$e_X\wedge e_X\to e\wedge e$ strongly in $L^2$.  The curvature terms
therefore converge by \Cref{lem:supp-weak-strong-palatini} with $p=p'=2$.
The same argument applies to their physical coframe variations: their
coefficient tensors are uniformly bounded and converge, so multiplying the
strong $L^2$ bivectors by these tensors preserves strong convergence.
The four-coframe volume term and its test insertions converge in $L^1$ by
the same lemma.  No product of two uncontrolled weak limits is used.

For the torsion-free connection, the first Bianchi identity removes the
Holst coframe variation.  The Palatini and volume variations give the
Einstein and cosmological insertions in
\eqref{eq:supp-einstein-insertion}.  The coefficient of the first is nonzero
because $\beta\ne0$.  A nonzero coefficient pair alone would only justify
the torsion step, and would not justify division by this Einstein coefficient.
\end{proof}

\begin{corollary}[Vacuum equation from certified accepted-action variations]
\label{cor:supp-renewal-einstein}
Under the preceding hypotheses, vanishing realized common-action variations
on the determining metric-test core imply
\[
 \boxed{G_{\mu\nu}+\Lambda g_{\mu\nu}=0}
\]
distributionally on every compact cylinder.  The scaled hypothesis
\eqref{eq:supp-scaled-test-gap} is a sufficient stochastic route on the
almost-sure subsequence of \Cref{prop:supp-physical-test-stationarity},
provided the remaining compactness and identification budgets hold on the
same realized sequence.
\end{corollary}
\begin{proof}
The handoff identifies the limit of each certified first variation with
\eqref{eq:supp-einstein-insertion}.  Its value is zero on the determining
core, and $\chi\ne0$ permits division.  Continuity in the declared test
topology extends the identity to the closure of that core.  In the stochastic
case, first take the realized vanishing subsequence supplied by
\Cref{prop:supp-physical-test-stationarity}; the assumed pathwise compactness
budgets then permit the curvature subsequence used in the handoff.
A mean-square bound is not interpreted as stationarity of every finite
sample.
\end{proof}

\subsection{Proof of the certification-or-obstruction alternative}

\begin{proof}[Proof of Theorem~\ref{thm:main-einstein-alternative}]
For a complete cylinder, \Cref{thm:supp-determining-kernel} reconstructs the
finite determining field and accepted kernel.  The reward-pressure and
common-action theorems either reconstruct the common action or expose its
curl/period failure.  The accepted-action theorem and the physical test-lift
estimate give a quantitative stationarity certificate with every cutoff
constant retained.  The spatial cut theorem gives the stated spatial screen
or a cut witness, while the common-cylinder comparison packet separately
records the temporal and transported-screen conditions needed for strong
coframe convergence.

Face linearization and memory subtraction give the classified first-order
score on their positive branches.  The calibrated coupling and Cartan tests
supply the nonzero Einstein coefficient and torsion control.  Continuous or
fully discrete propagation bounds the actual curvature from the finite
initial energy and growth/transfer/source records; the sharp current transfer
may supply the last budget from scaled action errors.  Finally the connection
and interface comparisons identify that curvature.  If all certificates pass,
\Cref{thm:supp-renewal-palatini,cor:supp-renewal-einstein} prove branch
\textup{(A1)}, using a diagonal subsequence on the compact exhaustion.

For completeness, the cofinal witnesses in branch \textup{(A2)} are obtained
from the quantifiers of the failed certificates.  A nonnegative residual
which does not tend to zero has a subsequence bounded below by some
$\epsilon>0$.  An unbounded budget has a subsequence tending to infinity.
A failed strictly positive lower margin has a subsequence tending to zero;
for the cut margin the finite graph supplies the corresponding potential
witness.  Failure of a uniform compactness tail supplies an $\epsilon>0$
and normalized components outside arbitrarily large declared screens.
Failure of a required calibrated Cauchy comparison retains the separated
pairs of finite records themselves.  A finite algebraic failure retains its
actual curl, period or operator residual.  Discarding a finite prefix is
allowed when checking the cofinal tail.  These records identify which
sufficient certificate failed; they do not establish that every alternative
compactness argument or every other subsequence must fail.
\end{proof}

\begin{remark}[What is discharged and what remains dynamical]
Neither a continuum curvature bound nor an unidentified nonlinear connection
product is inserted into the conclusion.  The former follows from the
continuous or discrete propagation certificate; the latter is handled by
strong connection convergence and the full distributional Cartan comparison,
including interfaces.  The differentiated source budget can follow from
\eqref{eq:main-action-source-budget}, whereas bounded symmetric growth,
physical occurrence and noncollapse remain genuine conditions of the
microscopic evolution.  A bounded mean curvature norm does not by itself
supply the pathwise version.  The fixed-amplitude nonsymmetric Cauchy theorem
is supplied independently by Appendix~\ref{supp:open-3plus1} for its declared
local writer; the Gowdy sector below remains a complementary nonlinear
symmetry-reduced class with pathwise full-curvature control.
\end{remark}

% =============================================================================
\section{Open nonsymmetric \texorpdfstring{$3+1$}{3+1} population}
\label{supp:open-3plus1}
% =============================================================================

This section proves \Cref{thm:main-open-3plus1} directly from the declared
local writer.  The purpose is to supply a fixed-radius, genuinely
three-dimensional nonlinear population theorem on the same reconstructed
Lorentzian coefficient chart used by the general certification theorem,
without invoking \textup{(E1)--(E5)} as independent assumptions.  The
evolution law is declared explicitly and is kept logically distinct from the
exact finite Euler flow studied in Appendix~\ref{supp:exact-action-audit}.
No symmetry reduction or analytic Fourier tail is used.

\subsection{Execution graph, exact difference calculus, and the local writer}

Let $N=2m+1$, $h=N^{-1}$ and
$\Lambda_h=(h\mathbb Z/\mathbb Z)^3$.  Let $\mathcal I_h$ denote the real
trigonometric interpolation on the signed frequency cube
$\{-m,\ldots,m\}^3$, and let $\mathcal S_h$ denote grid sampling.  The
tetrahedral $A_3\cong D_3$ root
execution directions may be represented as
\begin{equation}
 \mathcal R_F=\{\pm(e_i+e_j),\ \pm(e_i-e_j):1\le i<j\le3\}.
 \label{eq:supp-open-fcc-roots}
\end{equation}
All these integer vectors have even coordinate sum.  This harmless parity at
fixed smooth frequency becomes relevant for a cutoff-uniform Cartesian
Sobolev realization.

\begin{lemma}[Root-parity obstruction and a minimal execution connector]
\label{lem:supp-open-parity}
On the odd periodic grid, the root Dirichlet form
\begin{equation}
 \mathcal D_F(u)=\sum_{r\in\mathcal R_F/\{\pm1\}}
       \|(S_r-I)u/h\|_h^2
 \label{eq:supp-open-root-dirichlet}
\end{equation}
is not uniformly coercive for
$\|\nabla\mathcal I_hu\|_{L^2}^2$.  Moreover the root-graph distance from
$0$ to $e_1$ is $N-1$.  If the oriented execution edges $\pm e_1$ are
adjoined, then with
\begin{equation}
 \mathcal D_+(u)=\mathcal D_F(u)+\|D_{e_1}^+u\|_h^2,
 \qquad \mathcal D_C(u)=\sum_i\|D_i^+u\|_h^2,
 \label{eq:supp-open-connector-energy}
\end{equation}
one has
\begin{equation}
 \boxed{\frac15\mathcal D_C(u)\le\mathcal D_+(u)\le9\mathcal D_C(u).}
 \label{eq:supp-open-connector-coercivity}
\end{equation}
At least one added fixed direction of odd coordinate sum is necessary within
this fixed-direction class.  The connector is an execution incidence only;
it does not enter the $A_3$ moment reconstruction.
\end{lemma}
\begin{proof}
Take
$u_h(x)=\exp(2\pi i m(x_1+x_2+x_3))$.  Every difference in a minus root
vanishes and every plus-root multiplier remains bounded, while
$\|\nabla\mathcal I_hu_h\|_2\asymp m$ diverges.  The same argument applies to
every fixed-radius word in even-parity directions.  A root path ending at
$e_1$ has integer displacement $e_1+Nz$ with even coordinate sum; since $N$
is odd, $\sum z_i$ is odd and one coordinate has modulus at least $N-1$.
The explicit alternating pair $-e_1+e_2,-e_1-e_2$, repeated $m$ times,
attains that distance.

For the connector, the identity
\[
 S_{e_2}-I=S_{e_1}^{-1}\bigl[(S_{e_1+e_2}-I)-(S_{e_1}-I)\bigr]
\]
and its $e_3$ analogue give
$\mathcal D_C\le5\mathcal D_+$.  Conversely each root difference is bounded
by the sum of the two participating coordinate differences, so
$\mathcal D_F\le8\mathcal D_C$.  Adding the connector gives the upper bound.
If every added direction had even coordinate sum, the high-frequency test
above would remain unchanged.
\end{proof}

Define $S_i,D_i^\pm,D_i^0$ by \eqref{eq:main-open-differences}, use
$\langle u,v\rangle_h=h^3\sum_xu(x)\cdot v(x)$, and for a multi-index put
$D^\alpha=\prod_i(D_i^+)^{\alpha_i}$ and
$S^\alpha=\prod_iS_i^{\alpha_i}$.  Set
\begin{equation}
 \|u\|_{r,h}^2=\sum_{|\alpha|\le r}\|D^\alpha u\|_h^2,
 \qquad
 \|X\|_{\mathcal X_h^r}^2=\|q\|_{r+1,h}^2+\|v\|_{r,h}^2.
 \label{eq:supp-open-norms}
\end{equation}

\begin{lemma}[Exact shifted product rule and uniform Sobolev calculus]
\label{lem:supp-open-calculus}
On the periodic array space,
\begin{equation}
 (D_i^+)^*=-D_i^-,\qquad (D_i^0)^*=-D_i^0,
 \label{eq:supp-open-adjoints}
\end{equation}
and
\begin{equation}
 \boxed{D^\alpha(fw)=\sum_{\beta\le\alpha}\binom\alpha\beta
 (S^{\alpha-\beta}D^\beta f)D^{\alpha-\beta}w.}
 \label{eq:supp-open-shifted-product}
\end{equation}
For $r\ge3$, the norms \eqref{eq:supp-open-norms} are uniformly equivalent
to the trigonometric $H^r$ norms, form a uniform product algebra, and the
shifted commutator
\begin{equation}
 \mathcal C_\alpha(f,w)=D^\alpha(fw)-(S^\alpha f)D^\alpha w
 \label{eq:supp-open-commutator}
\end{equation}
satisfies, for a constant coefficient $f_*$,
\begin{align}
 \|\mathcal C_\alpha(f,w)\|_h
 &\le C_r\|f-f_*\|_{r,h}\|w\|_{r-1,h},\\
 \|D_i^-\mathcal C_\alpha(f,w)\|_h
 +\|D_i^0\mathcal C_\alpha(f,w)\|_h
 &\le C_r\|f-f_*\|_{r+1,h}\|w\|_{r,h}.
 \label{eq:supp-open-commutator-bounds}
\end{align}
The constants are independent of $h$.
\end{lemma}
\begin{proof}
Translation of the finite sum gives the adjoints, and
$D_i^+(fw)=(S_if)D_i^+w+(D_i^+f)w$ gives
\eqref{eq:supp-open-shifted-product} by iteration.  The forward symbol
$d_i(k)=(e^{2\pi ihk_i}-1)/h$ satisfies
$4|k_i|\le|d_i(k)|\le2\pi|k_i|$ on $|k_i|\le m$.  Parseval therefore gives
uniform norm equivalence.  Multiplication is convolution on the finite
frequency group and the exact identity
$d_i([k+p])=d_i(k)+e^{2\pi ihk_i}d_i(p)$ gives the usual uniform weighted
convolution estimate.  In every commutator term at least one difference falls
on $f$; one additional output difference gives the second line.  Since
$r>3/2$, the required Fourier $\ell^1$ weights are summable uniformly in the
finite cube.  Analytic coefficient compositions follow from their convergent
power series on a smaller ball.
\end{proof}

For the harmonic reduced vacuum equation use the coefficient functions in
\eqref{eq:main-harmonic-normal-row}.  The skew transport is
\eqref{eq:main-open-skew-transport} and the first-jet compensator is
\eqref{eq:main-open-compensator}.  Pointwise multiplication operators are
denoted $M_f$.

\begin{proposition}[Locality and continuum consistency of the writer]
\label{prop:supp-open-locality}
The acceleration in \eqref{eq:main-open-writer} at a site depends only on the
site and the twelve coordinate-stencil neighbours
$\{\pm e_i,\pm(e_i-e_j):i<j\}$.  These lie within two hops of the execution
graph obtained in \Cref{lem:supp-open-parity}.  On every smooth metric jet,
the continuum limit of the second row is exactly
\eqref{eq:main-harmonic-normal-row}.
\end{proposition}
\begin{proof}
The flux at $x$ is
\[
 h^{-2}\sum_{i,j}\{c^{ij}(x)[q(x+he_j)-q(x)]
 -c^{ij}(x-he_i)[q(x-he_i+he_j)-q(x-he_i)]\}.
\]
The skew transport uses only $x\pm he_i$, and the compensator only first
forward differences.  In the smooth limit the divergence contributes
$(\partial_ic^{ij})\partial_jq$ and the skew transport contributes
$(\partial_ib^i)v$ in addition to the desired second- and mixed-derivative
terms.  The two coefficient-differential terms in
\eqref{eq:main-open-compensator} cancel exactly these additions by the
ordinary chain rule.  No discrete Leibniz rule is invoked.
\end{proof}

\subsection{Shifted nonlinear energy and a cutoff-independent lifespan}

On a sufficiently small fixed $\mathcal X_h^s$ ball, $a\ge a_*>0$ and
$c\succeq c_*I$.  Put
$a_\alpha=S^\alpha a$, $c_\alpha=S^\alpha c$,
$q_\alpha=D^\alpha q$, and $v_\alpha=D^\alpha v$, and define
\begin{equation}
 \boxed{\mathscr E_{s,h}(X)=\frac12\sum_{|\alpha|\le s}
 \left[\langle v_\alpha,a_\alpha v_\alpha\rangle_h
 +\sum_{i,j}\langle D_i^+q_\alpha,c_\alpha^{ij}D_j^+q_\alpha\rangle_h
 +\|q_\alpha\|_h^2\right].}
 \label{eq:supp-open-energy}
\end{equation}
It is uniformly equivalent to $\|X\|_{\mathcal X_h^s}^2$.

\begin{theorem}[All-component shifted nonlinear energy]
\label{thm:supp-open-energy}
For $s\ge11$, every solution of \eqref{eq:main-open-writer} remaining in the
fixed chart satisfies
\begin{equation}
 \boxed{\frac{\dd}{\dd t}\mathscr E_{s,h}
 \le C_s\mathscr E_{s,h}+C_s\mathscr E_{s,h}^{3/2},}
 \label{eq:supp-open-energy-inequality}
\end{equation}
with $C_s$ independent of $h$.
\end{theorem}
\begin{proof}
Apply $D^\alpha$ to the mass row.  The exact shifted product rule gives
\begin{equation}
 a_\alpha\dot v_\alpha
 =\sum_{i,j}D_i^-(c_\alpha^{ij}D_j^+q_\alpha)
 -\mathsf K_{b_\alpha}v_\alpha+\mathsf R_\alpha,
 \label{eq:supp-open-commuted-row}
\end{equation}
where every term in $\mathsf R_\alpha$ contains either a shifted commutator
or the quadratic first-jet source.  In particular the outer-divergence
commutator is retained, rather than treated as an undifferentiated error.
\Cref{lem:supp-open-calculus} and the analytic coefficient bounds imply
\begin{equation}
 \left(\sum_{|\alpha|\le s}\|\mathsf R_\alpha\|_h^2\right)^{1/2}
 \le C_s\|X\|_{\mathcal X_h^s}^2.
 \label{eq:supp-open-remainder}
\end{equation}
Pair \eqref{eq:supp-open-commuted-row} with $v_\alpha$.
Summation by parts cancels the divergence against the derivative of the
spatial energy, while
$\langle v_\alpha,\mathsf K_{b_\alpha}v_\alpha\rangle_h=0$.  Thus
\begin{equation}
 \begin{split}
 \dot{\mathscr E}_{s,h}=\sum_{|\alpha|\le s}\bigg[&
 \tfrac12\langle v_\alpha,\dot a_\alpha v_\alpha\rangle_h
 +\tfrac12\sum_{i,j}\langle D_i^+q_\alpha,
          \dot c_\alpha^{ij}D_j^+q_\alpha\rangle_h\\
 &+\langle q_\alpha,v_\alpha\rangle_h
 +\langle v_\alpha,\mathsf R_\alpha\rangle_h\bigg].
 \end{split}
 \label{eq:supp-open-energy-identity}
\end{equation}
The coefficient time derivatives are analytic multiples of $v$ and have
$\ell^\infty$ size $O(\|X\|_{\mathcal X_h^s})$.  Energy equivalence,
Cauchy--Schwarz and \eqref{eq:supp-open-remainder} prove the result.  The
translated coefficient in each differentiated row is essential: replacing
it by an unshifted continuum Leibniz rule would leave uncancelled top-order
terms.
\end{proof}

\begin{corollary}[Common nonlinear time and forced stability]
\label{cor:supp-open-lifespan}
There are $\varepsilon_s,T_s,C_s>0$, independent of $h$, such that every
$\|X_h(0)\|_{\mathcal X_h^s}\le\varepsilon_s$ has a unique solution through
$T_s$ satisfying \eqref{eq:main-open-uniform-time}.  If an acceleration defect
$f_h$ is added, then on the same smaller chart
\begin{equation}
 (\mathscr E_{s,h}^{1/2})'
 \le C_s\mathscr E_{s,h}^{1/2}+C_s\mathscr E_{s,h}
       +C_s\|f_h\|_{s,h}.
 \label{eq:supp-open-forced-energy}
\end{equation}
Two top-bounded records satisfy the corresponding lower-order difference
estimate
\begin{equation}
 \sup_{t\le T_s}\|X_h-Y_h\|_{\mathcal X_h^r}
 \le C_s\left(\|X_h(0)-Y_h(0)\|_{\mathcal X_h^r}
 +\int_0^{T_s}\|f_h-\widetilde f_h\|_{r,h}\,\dd t\right)
 \label{eq:supp-open-difference}
\end{equation}
for $3\le r\le s-1$.
\end{corollary}
\begin{proof}
The finite ODE is analytic on the coefficient chart.  A scalar comparison for
\eqref{eq:supp-open-energy-inequality} keeps its energy strictly inside a
smaller fixed ball on a time chosen before the cutoff; finite-dimensional
continuation then gives the common lifespan.  Adding the differentiated force
to the remainder proves \eqref{eq:supp-open-forced-energy}.  For differences,
put one record's coefficients in the principal row and every coefficient
difference in the source.  The product and commutator estimates at one lower
order, followed by Gronwall, give \eqref{eq:supp-open-difference} without an
$e^{CT/h}$ loss.
\end{proof}

\subsection{Compatible initial data, subsidiary propagation, and the Einstein limit}

For $(\gamma,K)\in\mathcal D_s(\varepsilon)$ define the harmonic Cauchy
normalization
\begin{equation}
 \begin{gathered}
 g_{00}^0=-1,\qquad g_{0i}^0=0,\qquad g_{ij}^0=\gamma_{ij},\\
 v_{ij}^0=-2K_{ij},\qquad v_{00}^0=2\tr_\gamma K,\qquad
 v_{0i}^0=\gamma^{jk}(\partial_j\gamma_{ki}
                         -\tfrac12\partial_i\gamma_{jk}).
 \end{gathered}
 \label{eq:supp-open-harmonic-initial}
\end{equation}
Let $\mathcal S_h$ be grid sampling and initialize
$q_h(0)=\mathcal S_h(g^0-\eta)$, $v_h(0)=\mathcal S_hv^0$ with the harmless
finite projection needed for the canonical interpolation.  Sampling a
future Einstein history is not part of this preparation.

Let $c^\mu(g)=g^{\alpha\beta}\Gamma^\mu_{\alpha\beta}$ be the harmonic gauge
vector.  For a smooth interpolant the reduced equation may be written
\begin{equation}
 G(g_h)-\mathfrak H(g_h,c(g_h))=r_h,
 \label{eq:supp-open-reduced-residual}
\end{equation}
where $\mathfrak H$ is the standard linear first-derivative gauge insertion.
The consistency estimates following from
\Cref{lem:supp-open-calculus,prop:supp-open-locality} give
\begin{equation}
 \|r_h\|_{C_tH^{s-2}}+
 \|\partial_tr_h\|_{C_tH^{s-3}}
 \le C_sh\varepsilon
 \label{eq:supp-open-reduced-rate}
\end{equation}
on the common chart.

\begin{lemma}[Subsidiary control on the common interval]
\label{lem:supp-open-subsidiary}
For the harmonic initialization \eqref{eq:supp-open-harmonic-initial},
\begin{equation}
 \sup_{t\le T_s}\bigl(
 \|c(g_h)\|_{H^{s-2}}+\|\partial_tc(g_h)\|_{H^{s-3}}
 +\|G(g_h)\|_{H^{s-3}}\bigr)\le C_sh\varepsilon.
 \label{eq:supp-open-subsidiary}
\end{equation}
\end{lemma}
\begin{proof}
The continuum constraints in \eqref{eq:main-open-data-class} and the harmonic
normalization make the unprojected initial gauge field and its first time
derivative zero.  Sampling and the actual finite acceleration perturb these
quantities by $O(h\varepsilon)$ at the displayed orders.  The contracted
Bianchi identity applied to \eqref{eq:supp-open-reduced-residual} gives the
inhomogeneous subsidiary system
\[
 \Box_{g_h}c_b+\operatorname{Ric}_b{}^d(g_h)c_d
       =-2\nabla^ar_{h,ab}+\mathcal O(h\varepsilon)
\]
on the smooth canonical interpolant.  The coefficients are uniformly
hyperbolic on the common Cauchy ball.  The standard coordinate wave energy,
commuted through $s-3$ spatial derivatives, and
\eqref{eq:supp-open-reduced-rate} give the first two terms by Gronwall.
Equation \eqref{eq:supp-open-reduced-residual} then gives the Einstein
residual.  The argument is performed after the independent nonlinear energy
has supplied the common interval; no stopped-time assumption is used.
\end{proof}

\begin{theorem}[Whole-sequence nonsymmetric vacuum limit]
\label{thm:supp-open-einstein}
For sufficiently small $(\gamma,K)\in\mathcal D_s(\varepsilon)$, the complete
sequence of interpolated solutions of \eqref{eq:main-open-writer} converges on
$[0,T_s]\times\mathbb T^3$ to a unique $g$ with those initial data,
\begin{equation}
 \boxed{G(g)=0,}
 \label{eq:supp-open-vacuum}
\end{equation}
and the quantitative estimate
\begin{equation}
 \sup_{t\le T_s}\left(
 \|g_h-g\|_{H^{s-1}}+\|\partial_tg_h-\partial_tg\|_{H^{s-2}}
 +\|\partial_t^2g_h-\partial_t^2g\|_{H^{s-3}}
 \right)\le C_sh\varepsilon
 \label{eq:supp-open-metric-rate}
\end{equation}
holds.  In particular the corresponding Levi--Civita curvatures converge one
further derivative below, and are uniformly bounded on the slab.
\end{theorem}
\begin{proof}
The uniform $H^{s+1}\times H^s$ bound and the actual time-jet bounds obtained
from the equation give compactness in every slightly weaker
$C_tH^{s+1-\delta}\times C_tH^{s-\delta}$ topology.  Pass to a cluster point.
The exact difference calculus and consistency
\eqref{eq:supp-open-reduced-rate} pass the local writer to the harmonic
reduced equation, while \Cref{lem:supp-open-subsidiary} gives $c(g)=0$.
Hence $G(g)=0$ and the initial data are the prescribed pair.

If two cluster points have the same data, subtract their harmonic reduced
wave systems and apply the ordinary Sobolev energy at index $s-2$; the
coefficient differences are lower-order sources controlled by the bounded top
norms, so Gronwall gives equality.  Thus every cluster point is the same and
the entire sequence converges.  To obtain the rate, subtract the sampled
limit equation from the finite writer.  First-difference and product
consistency give an $O(h\varepsilon)$ source in $H^{s-2}$ and the sampled
initial tail is $O(h\varepsilon)$ in the lower Cauchy norm.  Apply
\eqref{eq:supp-open-difference}; solving the mass-row difference for the
acceleration gives the third term.  Curvature is an analytic function of a
nondegenerate metric and its first two jets at these Sobolev orders, proving
the final assertion.
\end{proof}

\begin{corollary}[Direct constructive Einstein limit]
\label{cor:supp-open-direct-route}
The Einstein conclusion of \Cref{thm:supp-open-einstein} follows from the
cutoff-uniform nonlinear energy, consistency of the declared local writer,
harmonic subsidiary estimate, and uniqueness of the limiting reduced Cauchy
problem.  No item \textup{(E1)--(E5)} of
\Cref{thm:main-renewal-einstein} is invoked as an independent hypothesis.
Accordingly those five items function as sufficient certification conditions
for the broader class of complete renewal cylinders, while
\Cref{thm:main-open-3plus1} supplies an independently constructed
non-symmetry-reduced positive population.
\end{corollary}

\subsection{Positive rate lift and finite renewal words}

The small metric chart has an explicit positive $A_3$ rate lift.  In the
$D_3$ realization
\begin{equation}
 r_a\in\{(1,1,0),(1,-1,0),(1,0,1),(1,0,-1),(0,1,1),(0,1,-1)\},
 \label{eq:supp-open-d3-roots}
\end{equation}
put $B=N_{\rm ADM}^2\gamma^{-1}$, $\varrho=\sqrt{\det\gamma}$ and, for a
coordinate pair $ij$ with remaining index $k$,
\begin{equation}
 c_{ij,\pm}=\frac{B_{ii}+B_{jj}-B_{kk}}4\pm\frac{B_{ij}}2,
 \qquad d_a=\frac14r_a\cdot\beta,\qquad
 k_a^\pm=\frac{c_a}{2h^2}\pm\frac{d_a}{2h}.
 \label{eq:supp-open-rate-lift}
\end{equation}
At the flat point all $c_a=1/4$, so a fixed small metric chart makes all
rates positive and comparable to $h^{-2}$.  Direct calculation gives
\begin{equation}
 h^2\sum_a(k_a^++k_a^-)r_ar_a^T=B,
 \qquad h\sum_a(k_a^+-k_a^-)r_a=\beta,
 \label{eq:supp-open-rate-inversion}
\end{equation}
so the ADM reconstruction returns the same metric.  The resulting weighted
periodic graph has a cutoff-uniform cut margin by perturbation of the flat
$D_3$ regulator and \Cref{thm:main-spatial-screen}.

To obtain finite histories rather than a continuous-time finite ODE, fix a
segment duration $\tau_h\asymp h^2$.  At a source $x=(q,v)$ use finite
registers $z=(\alpha,\beta)$ and the degree-seven two-endpoint Hermite curve
whose endpoint acceleration and jerk are those of the local vector field.
Evaluate the normalized equation residual at two fixed interior times and
apply the inverse of the resulting universal two-slot Hermite matrix.  Denote
the resulting local residual map by $\mathscr D_x(z)$.

\begin{theorem}[Finite local successor words]
\label{thm:supp-open-finite-words}
On the fixed small Cauchy ball there are constants independent of $h$ such
that, for $\tau_h\le c_sh^2$,
\begin{equation}
 \|D_z\mathscr D_x-I\|_{\op}+\|D_z^2\mathscr D_x\|_{\op}
 \le C_s\tau_h/h.
 \label{eq:supp-open-selector-inverse}
\end{equation}
With zero initial register, six parallel residual corrections and absolute
local arithmetic error at most $c_s\varepsilon h^6$ suffice, for all
sufficiently small $h$, to commit a successor with fixed execution-graph
radius.  Sharing endpoint acceleration and jerk caches between adjacent
segments gives a finite piecewise-Hermite history satisfying
\begin{equation}
 \boxed{
 \sup_t\|(Q_h,\dot Q_h)\|_{\mathcal X_h^s}\le C_s\varepsilon,
 \quad \|f_h^{\rm num}\|_{C_tH_h^s}\le C_s\varepsilon h^4,
 \quad \|\partial_tf_h^{\rm num}\|_{C_tH_h^s}\le C_s\varepsilon h^2.}
 \label{eq:supp-open-finite-defects}
\end{equation}
If the initial finite record is $O(\varepsilon h^4)$ from the compatible
preparation, every such history has the same Einstein limit as
\Cref{thm:supp-open-einstein}.
\end{theorem}
\begin{proof}
The local acceleration loses one discrete spatial order and its jerk loses
one further order.  Differentiating the fixed Hermite residual therefore gives
\eqref{eq:supp-open-selector-inverse}; every residual evaluation has spatial
radius three in the coordinate stencil.  For $C_sh\le1/2$, the correction
map satisfies, with local error $\nu_h$,
\[
 d_{n+1}\le C_shd_n+\nu_h,
 \qquad d_0\le C_s\varepsilon h.
\]
Thus
$d_6\le C_s\varepsilon h^7+2\nu_h=O(\varepsilon h^6)$ for the stated
arithmetic floor.  Fixed-order Hermite remainder estimates with
$\tau_h\asymp h^2$ then give the two defects in
\eqref{eq:supp-open-finite-defects}.  The forced energy
\eqref{eq:supp-open-forced-energy} prevents exit from the common chart and the
lower-order difference estimate compares the finite word to the exact local
trajectory.  The defects are smaller than the $O(h\varepsilon)$ consistency
budget in \Cref{thm:supp-open-einstein}, which proves the same limit.  The
number of communication hops is fixed because both the residual radius and
the number of corrections are fixed.  This proof concerns information range;
it does not provide a hardware latency bound for executing those operations
inside the physical duration $\tau_h$.
\end{proof}

The results above prove the population statement required in the main text.
They do not identify the local writer with the exact finite Euler equations of
the reconstructed action.  That stronger question has a different finite
constraint and response structure, to which we now turn.


% =============================================================================
\section{Open law-space robustness and finite action-prepared initial data}
\label{supp:open-law-basin}
% =============================================================================

This section proves \Cref{thm:main-open-law-basin}.  The first part enlarges
the central local writer to a fixed open family whose members may differ
substantially at the ultraviolet edge while sharing one nonlinear energy and
one Einstein limit.  The second part places finite chronological probability
rows in relatively open residual-profile neighborhoods.  The third constructs
exact initial zeroes of the original finite lapse--shift constraint map from
free tensor records.  These are three distinct notions of openness: openness
of differential coefficients, openness of finite transition rows, and
relative openness of the original finite constraint surface.

\subsection{A fixed open family of local differential writers}

Put
\begin{equation}
 \Lambda_h=-\sum_{i=1}^3D_i^-D_i^+,
 \qquad
 \ell_h(k)=4h^{-2}\sum_i\sin^2(\pi hk_i),
 \qquad0\le h^2\ell_h(k)\le12.
 \label{eq:supp-law-laplacian}
\end{equation}
On the metric chart of Appendix~\ref{supp:open-3plus1}, let
$c\succeq c_*I$.  Choose $0<b_0<c_*/24$ and set
\begin{equation}
 \mathcal B=\{B=B^T\in\mathbb R^{10\times10}:\|B\|_{\op}<b_0\}.
 \label{eq:supp-law-ball}
\end{equation}
This is an open ball of dimension $55$.  For $B\in\mathcal B$ define
\begin{equation}
 q_t=v,\qquad
 a(g)v_t=\sum_{i,j}D_i^-\!\bigl(c^{ij}(g)D_j^+q\bigr)
       -\mathsf K_bv+\mathsf G_h(q,v)-h^2B\Lambda_h^2q.
 \label{eq:supp-law-family}
\end{equation}
The coefficient $B$ is constant during one physical history.  It may vary
with the cutoff or be sampled once as a finite initial mark.

\begin{lemma}[Uniform fourth-order multiplier bounds]
\label{lem:supp-law-multiplier}
For every array $w$,
\begin{equation}
 h^2\|\Lambda_hw\|_h^2\le12\sum_i\|D_i^+w\|_h^2.
 \label{eq:supp-law-domination}
\end{equation}
At every fixed index used below,
\begin{align}
 h^2\|\Lambda_h^2q\|_{s-1,h}&\le C_s\|q\|_{s+1,h},&
 h^2\|\Lambda_h^2q\|_{s-2,h}&\le C_sh\|q\|_{s+1,h},\\
 h^2\|\Lambda_h^2q\|_{s-3,h}&\le C_sh^2\|q\|_{s+1,h},&
 h^2\|\Lambda_h^2v\|_{s-3,h}&\le C_sh\|v\|_{s,h}.
 \label{eq:supp-law-orders}
\end{align}
\end{lemma}
\begin{proof}
Parseval reduces the first claim to
$h^2\ell_h(k)^2\le12\ell_h(k)$.  Uniform equivalence of the forward-difference
Sobolev norm with the Fourier weight $(1+\ell_h)^{r/2}$ reduces the four
remaining claims to the bounded ratios
\[
 \frac{h^2\ell_h^2}{1+\ell_h},\quad
 \frac{h^2\ell_h^2}{(1+\ell_h)^{3/2}},\quad
 \frac{h^2\ell_h^2}{(1+\ell_h)^2},
\]
which are bounded by $12$, $\sqrt{12}h$, and $h^2$, respectively.
\end{proof}

The stencil $\Lambda_h^2$ has fixed coordinate radius two.  Its contribution
is the negative $q$-gradient of
$h^2\langle\Lambda_hq,B\Lambda_hq\rangle_h/2$; this identity concerns the
added term only and does not identify the complete local writer with the raw
finite Euler flow of the reconstructed signed action.

For $q_\alpha=D^\alpha q$ define
\begin{equation}
 \boxed{\mathscr E^B_{s,h}=\mathscr E_{s,h}
   +\frac{h^2}{2}\sum_{|\alpha|\le s}
       \langle\Lambda_hq_\alpha,B\Lambda_hq_\alpha\rangle_h,}
 \label{eq:supp-law-energy}
\end{equation}
where $\mathscr E_{s,h}$ is \eqref{eq:supp-open-energy}.

\begin{theorem}[Fixed-radius propagation on an open writer ball]
\label{thm:supp-law-stability}
For $s\ge11$, uniformly for $B\in\overline{\mathcal B}$ on a sufficiently
small common Cauchy chart,
\begin{equation}
 c_s\|X\|_{\mathcal X_h^s}^2\le\mathscr E^B_{s,h}
 \le C_s\|X\|_{\mathcal X_h^s}^2,
 \qquad
 (\mathscr E^B_{s,h})'\le C_s\mathscr E^B_{s,h}
                     +C_s(\mathscr E^B_{s,h})^{3/2}.
 \label{eq:supp-law-energy-bound}
\end{equation}
Consequently there are $\varepsilon_s,T_s>0$, independent of $h$ and $B$,
such that every initial record with
$\|X_h(0)\|_{\mathcal X_h^s}\le\varepsilon_s$ has a unique solution through
$T_s$ and
\begin{equation}
 \sup_{t\le T_s}\|X_h(t)\|_{\mathcal X_h^s}\le C_s\|X_h(0)\|_{\mathcal X_h^s}.
 \label{eq:supp-law-lifespan}
\end{equation}
With an additional acceleration defect $f$,
\begin{equation}
 (\sqrt{\mathscr E^B_{s,h}})'
 \le C_s\sqrt{\mathscr E^B_{s,h}}+C_s\mathscr E^B_{s,h}+C_s\|f\|_{s,h}.
 \label{eq:supp-law-forced}
\end{equation}
\end{theorem}
\begin{proof}
By \Cref{lem:supp-law-multiplier}, the absolute value of the added quadratic
energy at each commuted level is at most
$6b_0\sum_i\|D_i^+q_\alpha\|_h^2$.  The pre-existing spatial energy is at
least $c_*\sum_i\|D_i^+q_\alpha\|_h^2/2$, so
$b_0<c_*/24$ leaves a cutoff-independent positive margin.

Commute \eqref{eq:supp-law-family} with $D^\alpha$.  Since $B$ and
$\Lambda_h$ commute with every lattice shift, the new principal term is
$-h^2B\Lambda_h^2q_\alpha$.  Pairing with $v_\alpha$ gives exactly
\[
 -h^2\langle v_\alpha,B\Lambda_h^2q_\alpha\rangle_h
 =-\frac{\dd}{\dd t}\frac{h^2}{2}
       \langle\Lambda_hq_\alpha,B\Lambda_hq_\alpha\rangle_h.
\]
Thus the additional high-frequency term cancels at principal order rather
than being estimated as a small top-order force.  Every remaining commutator
is one of those controlled in \Cref{thm:supp-open-energy}; the only change is
that the actual acceleration appears in analytic mass-coefficient terms, and
\Cref{lem:supp-law-multiplier} bounds it uniformly in $H_h^{s-1}$.  This gives
\eqref{eq:supp-law-energy-bound}.  The finite-dimensional vector field is
analytic on the common chart, and the same first-exit argument used for
\eqref{eq:supp-open-forced-energy} proves the lifespan and the forced estimate.
\end{proof}

\begin{theorem}[Uniform Einstein limit for the whole coefficient family]
\label{thm:supp-law-einstein}
Let $X_{B,h}$ and $X_{0,h}$ have the same sufficiently small compatible
initial data.  On their common interval,
\begin{equation}
 \boxed{\sum_{j=0}^3
 \|q_{B,h}^{(j)}-q_{0,h}^{(j)}\|_{C_tH_h^{s-1-j}}
       \le C_sh\|B\|_{\op}\varepsilon.}
 \label{eq:supp-law-jet-comparison}
\end{equation}
For any sequence $B_h\in\overline{\mathcal B}$, with no convergence of $B_h$
required, the whole interpolated sequence converges to the same harmonic
vacuum development as the central writer.  In particular the metric and its
first two time derivatives satisfy the $O(h\varepsilon)$ bounds of
\Cref{thm:supp-open-einstein}, and the subsidiary field, rate positivity,
cut margin and noncollapse estimates remain uniform.
\end{theorem}
\begin{proof}
For comparison only, regard the added term on the $B$-trajectory as an
external force
\[
 f_{B,h}=-h^2a(g_{B,h})^{-1}B\Lambda_h^2q_{B,h}.
\]
Both trajectories already exist on the common top chart by
\Cref{thm:supp-law-stability}.  The lower-order bounds in
\Cref{lem:supp-law-multiplier} give
\[
 \|f_{B,h}\|_{C_tH_h^{s-2}}
 +\|\partial_tf_{B,h}\|_{C_tH_h^{s-3}}
 \le C_sh\|B\|\varepsilon,
\]
while the top $H_h^{s-1}$ norm need not be small.  Apply the forced difference
estimate \eqref{eq:supp-open-difference} at the reserved lower orders to get
the first two jet slots.  Subtracting the acceleration rows gives the second
time derivative; differentiating once and using the velocity differential of
the local vector field gives the third.  This proves
\eqref{eq:supp-law-jet-comparison}.  Add the central-writer convergence of
\Cref{thm:supp-open-einstein}.  All constants are uniform on
$\overline{\mathcal B}$, so no convergence of $B_h$ is needed.
\end{proof}

\begin{proposition}[Rapid switching is not included]
\label{prop:supp-law-switching}
For two distinct scalar marks $B_i=\sigma_iI$ in the admissible interval, one
resolved corner mode has frequency
\begin{equation}
 \Omega_{\sigma,h}(k)^2=\ell_h(k)(1+\sigma h^2\ell_h(k)).
 \label{eq:supp-law-dispersion}
\end{equation}
Holding each law for a quarter period gives transfer matrices whose two-step
product has eigenvalues
$-\Omega_{\sigma_1,h}/\Omega_{\sigma_2,h}$ and its reciprocal.  Their modulus
ratio stays separated from one at the ultraviolet corner when
$\sigma_1\ne\sigma_2$.  Repeated $O(h)$-time switching can therefore amplify
that mode like $\exp(cT/h)$ on a fixed slab.
\end{proposition}
\begin{proof}
The flat scalar mode is an oscillator.  A quarter-period transfer is
$\bigl(\begin{smallmatrix}0&\Omega^{-1}\\-\Omega&0\end{smallmatrix}\bigr)$.
Multiplying the two transfers gives the stated diagonal matrix.  At the
corner $h^2\ell_h(k)\to12$, so the frequency ratio approaches
$\sqrt{(1+12\sigma_1)/(1+12\sigma_2)}\ne1$.
\end{proof}

\subsection{Relatively open neighborhoods of chronological transition rows}

Fix $B\in\overline{\mathcal B}$ and $T<T_s$.  Partition $[0,T]$ into cells
$I_j$ with $c_-h^2\le\tau_j\le c_+h^2$.  A chronological letter is a finite
physical polynomial metric segment $Q_e$ with shared endpoint value,
velocity, acceleration and jerk records.  Before a possible guard exit define
\begin{equation}
 f_e(t)=Q_e''(t)-V_{B,h}(Q_e(t),Q_e'(t))
 \label{eq:supp-law-residual}
\end{equation}
and the cost \eqref{eq:main-law-cost}; a grammar or execution failure has cost
one and terminates the word.  The common-energy margin is chosen so that any
word with total cost below one is evaluated entirely inside the coefficient
chart.

\begin{lemma}[A stopped physical cost controls the whole word]
\label{lem:supp-law-cost-control}
If $C_h\le d_h^2<1$ and $d_h$ is sufficiently small, then no failure or guard
exit occurs and
\begin{equation}
 \|f_h\|_{L_t^2H_h^s}+\|\partial_tf_h\|_{L_t^2H_h^{s-3}}
       \le2\varepsilon d_h.
 \label{eq:supp-law-cost-force}
\end{equation}
Relative to the exact unforced $B$-history with the same initial record,
\begin{equation}
 \|Q_h-q_{B,h}\|_{C_tH_h^s}
 +\|Q_h'-q_{B,h}'\|_{C_tH_h^{s-1}}
 +\|Q_h''-q_{B,h}''\|_{C_tH_h^{s-3}}
 \le C_s\varepsilon d_h.
 \label{eq:supp-law-cost-shadow}
\end{equation}
\end{lemma}
\begin{proof}
Until a first guard hit, Cauchy--Schwarz gives
$\int_0^{t_*}\|f_h\|_{s,h}\dd t\le\sqrt T\,\varepsilon d_h$.
For small $d_h$ the forced energy \eqref{eq:supp-law-forced} contradicts a
first exit through the strict guard margin; a failure is excluded because it
would already contribute unit cost.  Hence the integrals are over the whole
interval and \eqref{eq:supp-law-cost-force} follows.  Subtract the exact and
forced mass equations at the reserved lower order, use
\eqref{eq:supp-open-difference} and Gronwall, and then solve the actual
acceleration difference for the last displayed slot.
\end{proof}

\begin{theorem}[A chronological law neighborhood with the same Einstein limit]
\label{thm:supp-law-row-neighbourhood}
Let $K_{h,j}(e\mid\pi)$ be an arbitrary finite transition row at each live
prefix, possibly depending on the entire past.  If
\begin{equation}
 \sum_eK_{h,j}(e\mid\pi)c_{h,j}(\pi,e)<M\tau_jh^4
 \label{eq:supp-law-row-moment}
\end{equation}
for every live prefix, then
\begin{equation}
 \boxed{\mathbb EC_h\le MTh^4,
 \qquad\mathbb P\{C_h>h^2\}\le MTh^2.}
 \label{eq:supp-law-row-probability}
\end{equation}
On $\{C_h\le h^2\}$ the history obeys
\Cref{lem:supp-law-cost-control} with $d_h=h$ and therefore has the same
Einstein limit as \Cref{thm:supp-law-einstein}.  At fixed cutoff and finite
grammar the strict inequalities \eqref{eq:supp-law-row-moment} define
nonempty relatively open convex subsets of the product of row simplexes when
a uniformly accurate completion is available at every live source.  They
contain strictly positive rows whenever all candidate letters have finite
cost.  Along the odd sequence $h=(2m+1)^{-1}$, the exceptional probabilities
are summable, so the conclusion holds eventually almost surely under any
common coupling.
\end{theorem}
\begin{proof}
Assign zero future cost to stopped prefixes.  Conditional on any live prefix,
\eqref{eq:supp-law-row-moment} bounds the expected stage cost by
$M\tau_jh^4$, independently of the probability of reaching that prefix.
Summation gives the expectation bound; Markov at threshold $h^2$ gives the
probability bound.  The deterministic cost lemma applies on the complement.
The series $\sum_m(2m+1)^{-2}$ converges, so the first Borel--Cantelli lemma
gives the eventual almost-sure assertion without independence between
cutoffs.  At fixed finite grammar each strict conditional mean bound is an
open affine half-space in one row simplex; finite products preserve relative
openness and convexity.  Small positive mixing preserves a strict margin.
\end{proof}

For an observable neighborhood description, attach to each ordinary letter
the profile
\begin{equation}
 \Phi_e=(f_e/\varepsilon,\partial_tf_e/\varepsilon,0),
 \label{eq:supp-law-profile}
\end{equation}
and to a failure the vector $(0,0,1)$ in
$L^2(I_j;H_h^s)\oplus L^2(I_j;H_h^{s-3})\oplus\mathbb R$.  Let
$d(e,e')=\|\Phi_e-\Phi_{e'}\|$ and let $W_{2,d}$ be the finite optimal
transport distance induced by $d$.

\begin{proposition}[Physical profile transport certifies a row neighborhood]
\label{prop:supp-law-transport}
For finite rows $K,P$,
\begin{equation}
 \left(\sum_eK(e)c_e\right)^{1/2}
 \le\left(\sum_eP(e)c_e\right)^{1/2}+W_{2,d}(K,P).
 \label{eq:supp-law-transport-bound}
\end{equation}
If $\sum_eP(e)c_e\le M_0\tau_jh^4$, then every row satisfying
\begin{equation}
 W_{2,d}(K,P)<\gamma\sqrt{\tau_j}\,h^2
 \label{eq:supp-law-transport-ball}
\end{equation}
obeys \eqref{eq:supp-law-row-moment} for any
$M>(\sqrt{M_0}+\gamma)^2$.  These strict balls are relatively open in the
finite probability simplex.
\end{proposition}
\begin{proof}
For any coupling of $K$ and $P$, Minkowski in the profile Hilbert space gives
$\|\Phi_e\|_{L^2}\le\|\Phi_{e'}\|_{L^2}+\|\Phi_e-\Phi_{e'}\|_{L^2}$.
Average under the coupling and minimize.  The finite coupling polytope is
compact, so a minimizer exists.  Continuity in the row probabilities follows
from finite-dimensional optimal transport.
\end{proof}

The scale in \eqref{eq:supp-law-transport-ball} is intentional.  A fixed
unscaled total-variation ball is not stable under refinement: if there are
$J_h\asymp h^{-2}$ stages and an absorbing failure of probability $p_h$ is
added independently at every live stage, then the failed-word probability is
$1-(1-p_h)^{J_h}$, which tends to zero exactly when $J_hp_h\to0$.  Thus even
$p_h\to0$ can be too large.  The theorem controls the accumulated physical
defect rather than hiding this refined-word effect.

\subsection{Exact finite initial constraints from free records}

We now remove the need to start the constructive population from a continuum
constraint solution.  This subsection uses the original reconstructed finite
action only at the initial cut.  Let the odd phase derivative be
\begin{equation}
 \widehat{\delta_i u}(k)=i\kappa_i(k)\widehat u(k),
 \qquad \kappa_i(k)=2h^{-1}\sin(\pi hk_i),
 \qquad \Omega_h^2=-\sum_i\delta_i^2.
 \label{eq:supp-initial-phase-derivative}
\end{equation}
For $r\ge3$ put
\begin{equation}
 \mathcal P_h^r=H_h^{r+2}(\operatorname{Sym}_3)
                 \times H_h^{r+1}(\operatorname{Sym}_3),
 \qquad \mathcal Y_h^r=H_h^r(\mathbb R\oplus\mathbb R^3),
 \label{eq:supp-initial-spaces}
\end{equation}
and write $X=(u,\pi)$, $\gamma=I+u$.  Let
$\mathscr C_h:\mathcal P_h^r\to\mathcal Y_h^r$ be
\eqref{eq:main-action-initial-constraints}.

\begin{lemma}[Uniform original-action constraint calculus]
\label{lem:supp-initial-calculus}
On a common ball $\mathscr C_h$ is analytic, $\mathscr C_h(0)=0$, and
\begin{equation}
 D\mathscr C_h(0)(u,\pi)=
 \left(\delta_i\delta_j u_{ij}-\sum_i\delta_i^2\tr u,
                         -2\delta_j\pi^{ij}\right)=:L_h(u,\pi).
 \label{eq:supp-initial-linear-map}
\end{equation}
Writing $\mathscr C_h=L_h+\mathscr N_h$,
\begin{equation}
 \|\mathscr N_h(X)\|_{\mathcal Y_h^r}\le C_r\|X\|^2,
 \qquad \|D\mathscr N_h(X)\|\le C_r\|X\|,
 \label{eq:supp-initial-nonlinear-bounds}
\end{equation}
with cutoff-independent constants.  No lapse or shift time derivative and no
metric acceleration appears in $\mathscr C_h$.
\end{lemma}
\begin{proof}
Stationary connection and Legendre elimination are analytic on the retained
small chart.  In a lapse or shift derivative, every spatial load has one phase
derivative; exact skew-adjointness transfers it to the stationary connection,
so the output loses at most the declared Sobolev orders.  At flat data the
quadratic phase-compatible action has the Fierz--Pauli rows displayed in
\eqref{eq:supp-initial-linear-map}; the full link remainder starts at cubic
order in the flat connection and contributes no linear initial row.  Taylor's
theorem gives \eqref{eq:supp-initial-nonlinear-bounds}.
\end{proof}

Let $P_0$ be spatial averaging in $\mathcal Y_h^r$ and
$P_\perp=I-P_0$.  On mean-zero targets define
\begin{equation}
 R_h(f,m)=\left(\tfrac12(\Omega_h^2)^{-1}f\,I,
  \tfrac12\mathcal L_h^{\rm vec}\mathcal D_h^{-1}m\right),
 \label{eq:supp-initial-right-inverse}
\end{equation}
where
$(\mathcal L_h^{\rm vec}Y)_{ij}=\delta_iY_j+\delta_jY_i
 -\frac23\delta_{ij}\delta_kY_k$ and
$\mathcal D_hY=-\delta_j(\mathcal L_h^{\rm vec}Y)_{ij}$.

\begin{lemma}[Range inverse and four missing means]
\label{lem:supp-initial-range}
One has $L_hR_h=I$ on the mean-zero target and
$\|R_h\|\le C_r$ uniformly.  For $d=N^3$,
\begin{equation}
 \operatorname{rank}L_h=4d-4,
 \qquad\dim\ker L_h=8d+4.
 \label{eq:supp-initial-rank}
\end{equation}
For each sufficiently small $z\in\ker L_h$ there is a unique mean-zero
$y_h(z)=O(\|z\|^2)$ such that
\begin{equation}
 P_\perp\mathscr C_h(z+R_hy_h(z))=0.
 \label{eq:supp-initial-range-zero}
\end{equation}
The remaining mean map
$\Theta_h(z)=P_0\mathscr C_h(z+R_hy_h(z))$ is analytic and has zero constant
and linear jets.
\end{lemma}
\begin{proof}
For $u=\phi I$ the scalar row equals $2\Omega_h^2\phi$.  For
$\pi=\mathcal L_h^{\rm vec}Y$ the vector row equals $2\mathcal D_hY$, whose
nonzero-mode symbol is $|\kappa|^2I+\kappa\kappa^T/3$ and is uniformly
invertible on the unit odd torus.  This proves the right inverse and rank.
The nonlinear range equation is
$y=-P_\perp\mathscr N_h(z+R_hy)$; the bounds of
\Cref{lem:supp-initial-calculus} make it a uniform contraction on a ball of
radius $C\|z\|^2$.  The expansion of $\Theta_h$ follows because
$P_0L_h=0$.
\end{proof}

To resolve the four means, let $T_i\in\operatorname{Sym}_3$ have its two
transverse off-diagonal entries equal to $1/\sqrt2$ and all other entries
zero, and define
\begin{equation}
 u_*=\sum_{i=1}^3T_i\cos(2\pi x_i),\qquad
 p(\vartheta)=-\tfrac23\kappa I
       +\sum_{i=1}^3b_iT_i\sin(2\pi x_i),
 \qquad \vartheta=(\kappa,b_1,b_2,b_3).
 \label{eq:supp-initial-seed}
\end{equation}
These are balancing directions only; arbitrary complementary modes remain
free.

\begin{lemma}[Exact four-mode mean calibration]
\label{lem:supp-initial-balance}
With $\omega_h=2h^{-1}\sin(\pi h)$, the quadratic mean jet of the full
retained action on the seed \eqref{eq:supp-initial-seed} is
\begin{equation}
 Q_h(\vartheta)=
 \begin{pmatrix}
 \frac23\kappa^2-\frac12\sum_i b_i^2-\frac38\omega_h^2\\
 -\frac12\omega_hb_1\\-\frac12\omega_hb_2\\-\frac12\omega_hb_3
 \end{pmatrix}.
 \label{eq:supp-initial-quadratic-mean}
\end{equation}
It has the exact root
\begin{equation}
 \vartheta_h^*=(3\omega_h/4,0,0,0),
 \qquad
 DQ_h(\vartheta_h^*)=
 \operatorname{diag}(\omega_h,-\omega_h/2,-\omega_h/2,-\omega_h/2),
 \label{eq:supp-initial-balance-root}
\end{equation}
whose inverse is uniformly bounded.  The opposite sign of the first component
gives the second branch.
\end{lemma}
\begin{proof}
Keep spatially constant lapse and shift independent and expand the canonical
Hamiltonian to quadratic order about flat data.  For the transverse trace-free
configuration seed and momentum equal to a homogeneous trace plus transverse
trace-free modes, the coefficient is
\[
 \mathscr H_h^{[2]}(u,\pi;N,\beta)=
 N\left(\|\pi\|_h^2-\tfrac12\|\tr\pi\|_h^2
          +\tfrac14\sum_i\|\delta_i u\|_h^2\right)
 +\sum_i\beta^i\langle\pi,\delta_i u\rangle_h.
\]
The complete link remainder starts at cubic order in the stationary connection
at flat data, so it does not modify this quadratic mean coefficient.  Taking
minus the lapse derivative and the three shift derivatives, then using
orthogonality of the constant and unit Fourier modes on every odd grid, gives
\eqref{eq:supp-initial-quadratic-mean}.  Direct differentiation yields
\eqref{eq:supp-initial-balance-root}; $\omega_h$ is uniformly separated from
zero on the declared grids.
\end{proof}

Let $\mathcal Z_h=\ker L_h$ and remove the four balancing seed directions
from an arbitrary small free record, leaving $W$ in a fixed complementary
subspace.  Write $Z_h(\vartheta)$ for the seed and, for $t\ne0$, define
\begin{equation}
 \mathcal Q_h(t,\vartheta,W)=t^{-2}\Theta_h(t[Z_h(\vartheta)+W]).
 \label{eq:supp-initial-normalized-mean}
\end{equation}
Since $\Theta_h$ has zero constant and linear jets, the integral Taylor formula
extends this map analytically through $t=0$.

\begin{theorem}[Exact finite preparation from free initial records]
\label{thm:supp-action-prepared-chart}
There are $t_0,\rho,d_0,h_0>0$, independent of $h$, such that for
$0<|t|<t_0$, $h<h_0$, and every complementary free record $W$ with
$\|W\|<\rho$, there is a unique $\vartheta_h(t,W)$ near
$\vartheta_h^*$ for which \eqref{eq:supp-initial-normalized-mean} vanishes.
The resulting record
\begin{equation}
 \boxed{\mathfrak I_h(t,W)
 =t[Z_h(\vartheta_h(t,W))+W]
   +R_hy_h(t[Z_h(\vartheta_h(t,W))+W])}
 \label{eq:supp-initial-prepared-record}
\end{equation}
satisfies
\begin{equation}
 \mathscr C_h(\mathfrak I_h(t,W))=0,
 \qquad \|\mathfrak I_h(t,W)\|_{\mathcal P_h^r}\le C_r|t|.
 \label{eq:supp-initial-exact-zero}
\end{equation}
At every such record $D\mathscr C_h$ is onto and admits a right inverse with
\begin{equation}
 \boxed{\|\mathscr R_{h,t}f\|_{\mathcal P_h^r}
 \le C_r\bigl(\|P_\perp f\|_{\mathcal Y_h^r}+|t|^{-1}|P_0f|\bigr).}
 \label{eq:supp-initial-full-inverse}
\end{equation}
For fixed $t\ne0$, the image contains a relatively open neighborhood in the
complete finite zero set; that zero set has codimension $4N^3$ in the
$12N^3$-dimensional canonical phase.  No symmetry, CMC, conformal-flatness,
frequency-support restriction, or continuum constraint solution is imposed
on $W$.
\end{theorem}
\begin{proof}
The analytic extension of \eqref{eq:supp-initial-normalized-mean} converges in
$C^1$ uniformly to the quadratic map of
\Cref{lem:supp-initial-balance} as $|t|+\|W\|\to0$.  Its root has a uniform
inverse margin.  Therefore
$\vartheta\mapsto\vartheta-[DQ_h(\vartheta_h^*)]^{-1}
\mathcal Q_h(t,\vartheta,W)$ is a contraction on one cutoff-independent ball
after choosing its radius first and then $t_0,\rho$.  The range equation was
already solved in \Cref{lem:supp-initial-range}, giving the complete zero.

For the derivative, eliminate the mean-zero block with the uniformly bounded
range inverse.  Along a balancing direction the remaining mean derivative is
$t^2D_\vartheta\mathcal Q_h$, while a range direction contributes only
$O(t)$ to the mean.  Solving for the balancing coefficient and multiplying by
the physical factor $t$ gives exactly the $|t|^{-1}$ term in
\eqref{eq:supp-initial-full-inverse}.  The inverse function theorem on the
zero set then gives relative openness and the codimension count.  The same
estimate shows why uniformity through the symmetric tip $t=0$ is neither
claimed nor true in general.
\end{proof}

The harmonic completion of a prepared record sets
$g_{00}=-1$, $g_{0i}=0$, $g_{ij}=\gamma_{ij}$ and uses its actual canonical
metric velocity $w$ for the spatial slot, with
\begin{equation}
 v_{ij}=w_{ij},\qquad
 v_{00}=-\gamma^{ij}w_{ij},\qquad
 v_{0i}=\gamma^{jk}\left(\delta_j\gamma_{ki}
             -\tfrac12\delta_i\gamma_{jk}\right).
 \label{eq:supp-initial-harmonic-completion}
\end{equation}
These lapse--shift velocities do not enter the four initial rows
\eqref{eq:main-action-initial-constraints}.

\begin{corollary}[Action-prepared inputs populate the open law family]
\label{cor:supp-action-prepared-einstein}
Fix sufficiently small nonzero $t$ and take a convergent sequence of free
records in \Cref{thm:supp-action-prepared-chart}.  For any sequence
$B_h\in\overline{\mathcal B}$, with no convergence of $B_h$ required, the
harmonically completed histories exist on the common interval and converge to
the vacuum Einstein development of the limiting prepared data.  At the exact
initial cut all four original lapse--shift rows vanish identically,
independently of $B_h$.  The chronological transition-row theorem
\Cref{thm:supp-law-row-neighbourhood} applies to these inputs unchanged.
\end{corollary}
\begin{proof}
The preparation has size $O(|t|)$ in Sobolev spaces strictly above the
evolution Cauchy norm, so \eqref{eq:supp-initial-harmonic-completion} lies in
the fixed ball of \Cref{thm:supp-law-stability}.  Consistency of the phase
derivatives and the stationary/Legendre maps gives convergence of the
prepared finite records to their continuum vacuum-constraint limit; the
harmonic initial gauge error tends to zero.  Apply
\Cref{thm:supp-law-einstein}.  The initial action rows vanish before the
acceleration law is chosen and therefore remain independent of the coefficient
mark.  The probabilistic statement depends only on the physical residual
costs and the common Cauchy estimates, so no new probability hypothesis is
needed.
\end{proof}

The three results of this section give an open population in a precise sense,
but not universal attraction.  The coefficient ball is fixed and genuinely
open; the transition-row neighborhoods are open in the scaled complete-profile
topology appropriate to $O(h^{-2})$ stages; and the initial records are
relatively open in the exact finite constraint surface.  The raw exact-action
evolution itself remains the separate question addressed next.

% =============================================================================
\section{Exact-action response geometry, preparation, and dynamical tangency}
\label{supp:exact-action-audit}
% =============================================================================

This section studies the original phase-compatible finite action without
identifying its Euler flow with the stable harmonic writer.  The organization
is mathematical rather than chronological.  We first determine the exact
finite constraint and response geometry, then exhibit nontrivial prepared
branches and quantitative response estimates, and finally isolate the higher
tangency condition that obstructs a direct cutoff-uniform flow theorem.

\subsection{Canonical constraints and the finite characteristic quotient}

Let $X=(u,\pi)\in\mathbb R^{12V}$ and let
$\lambda=(N,\beta)\in\mathbb R^{4V}$.  For the canonical Hamiltonian
$H_h(X,\lambda)$ put
\begin{equation}
 P_h=\nabla_\lambda H_h,\qquad
 F_h^{\rm can}=\mathbb J\nabla_XH_h,\qquad
 B_h=D_XP_hF_h^{\rm can},\qquad M_h=D_\lambda P_h.
 \label{eq:supp-exact-canonical-objects}
\end{equation}
The exact Euler system is
\begin{equation}
 X_t=F_h^{\rm can}(X,\lambda),\qquad P_h(X,\lambda)=0,
 \label{eq:supp-exact-euler}
\end{equation}
and differentiating the constraint gives
\begin{equation}
 B_h+M_h\dot\lambda=0.
 \label{eq:supp-exact-first-consistency}
\end{equation}
Phase-compatible multiplier homogeneity gives
\begin{equation}
 H_h(X,c\lambda)=cH_h(X,\lambda),\qquad
 M_h\lambda=0,
 \qquad \langle\lambda,B_h\rangle_h=0.
 \label{eq:supp-exact-homogeneity}
\end{equation}
These are global scaling identities, not local discrete diffeomorphism
identities.

Introduce multiplier momenta and the extended constraint bracket
\begin{equation}
 \mathcal A_h=D_XP_h\,\mathbb J(D_XP_h)^*,\qquad
 \mathfrak D_h=\begin{pmatrix}0&-M_h\\M_h&\mathcal A_h\end{pmatrix}.
 \label{eq:supp-exact-dirac-matrix}
\end{equation}

\begin{theorem}[Finite characteristic quotient on a regular one-null branch]
\label{thm:supp-exact-constraint-rank}
Suppose $D_XP_h$ is onto, $\ker M_h=\operatorname{span}\{\lambda\}$ and
\eqref{eq:supp-exact-first-consistency} is solvable.  Then
\begin{equation}
 \boxed{
 \ker\mathfrak D_h=\operatorname{span}\left\{
 \binom{\lambda}{0},\binom{-M_h^\dagger B_h}{\lambda}\right\},
 \qquad \rank\mathfrak D_h=8V-2,}
 \label{eq:supp-exact-dirac-rank}
\end{equation}
and the local characteristic quotient has dimension $12V-2$.
\end{theorem}
\begin{proof}
If $(r,s)\in\ker\mathfrak D_h$, the first row gives $s=c\lambda$.
Differentiating the homogeneity identity in $X$ and pairing with the
Hamiltonian vector field gives $\mathcal A_h\lambda=B_h$.  The second row
therefore gives $r=-cM_h^\dagger B_h+d\lambda$.  The extended phase has
dimension $20V$ and the $8V$ independent primary and secondary constraints
leave dimension $12V$; quotienting the two characteristic directions gives
$12V-2$.
\end{proof}

\subsection{An exact regular transverse block and lossless source elimination}

Use the odd-grid phase derivatives
\begin{equation}
 \widehat{\delta_i u}(k)=i\kappa_i(k)\widehat u(k),
 \qquad \kappa_i(k)=2h^{-1}\sin(\pi hk_i),
 \qquad \Omega_h^2=-\sum_i\delta_i^2.
 \label{eq:supp-exact-phase-derivative}
\end{equation}
Let
\begin{equation}
 \mathcal T_h=\{\beta:\widehat\beta(0)=0,
 \ \kappa(k)^T\widehat\beta(k)=0\ (k\ne0)\}
 \label{eq:supp-exact-transverse-space}
\end{equation}
be the transverse-shift space, of dimension $2(V-1)$.  On the scalar-momentum
family the first source restricted to $\mathcal T_h$ has the form
\begin{equation}
 A_h\beta=-\frac{h\varkappa_h}{2}\,\mathcal L_h
          \left(-\tfrac12\operatorname{curl}_\delta\beta\right),
 \label{eq:supp-exact-transverse-source}
\end{equation}
where $\varkappa_h$ is uniformly bounded above and below and
$\mathcal L_hw$ is the symmetric matrix with zero diagonal and
$(\mathcal L_hw)_{ij}=S_iw_j+S_jw_i$.

\begin{theorem}[Exact all-cutoff transverse rank and natural-order bounds]
\label{thm:supp-exact-transverse}
At every odd cutoff,
\begin{equation}
 \rank A_h=2(V-1),\qquad \mathsf P_{\rm tr}A_h=0,
 \label{eq:supp-exact-transverse-rank}
\end{equation}
and
\begin{equation}
 \boxed{c h^2\|\Omega_h\beta\|_h^2
 \le\|A_h\beta\|_h^2
 \le C h^2\|\Omega_h\beta\|_h^2,}
 \qquad \beta\in\mathcal T_h.
 \label{eq:supp-exact-transverse-bounds}
\end{equation}
The same estimates hold after multiplication by a scalar Fourier Sobolev
weight, and
\begin{equation}
 A_h^*K_0^{-1}A_h=\tfrac12A_h^*A_h\succ0.
 \label{eq:supp-exact-transverse-positive}
\end{equation}
\end{theorem}
\begin{proof}
At one Fourier mode the off-diagonal map is unitarily equivalent to
$w\mapsto(w_1+w_2,w_1+w_3,w_2+w_3)$, whose squared singular values are
$4,1,1$.  For transverse $\beta$,
$\|\operatorname{curl}_\delta\beta\|_h=\|\Omega_h\beta\|_h$.
Substitution in \eqref{eq:supp-exact-transverse-source} proves the two-sided
estimate and injectivity.  The output diagonal vanishes, so the trace
projection vanishes and $K_0^{-1}$ acts by $1/2$ on this range.
\end{proof}

Write the complete leading graded source as $F_h=(A_h,G_h)$ with target
$(0,d_h)$, and put $\overline G_h=(I-\Pi_{A,h})G_h$.

\begin{corollary}[Lossless elimination of the transverse source columns]
\label{cor:supp-exact-transverse-elimination}
One has
\begin{equation}
 \rank F_h=(2V-2)+\rank\overline G_h,
 \label{eq:supp-exact-rank-split}
\end{equation}
and
\begin{equation}
 F_h^*K_0^{-1}F_h\sim
 \operatorname{diag}\left(\tfrac12A_h^*A_h,
          \overline G_h^*K_0^{-1}\overline G_h\right).
 \label{eq:supp-exact-Gram-split}
\end{equation}
The triangular column transformation producing this congruence leaves the
actual target unchanged and preserves every remaining source equation.
\end{corollary}
\begin{proof}
Use the coefficient
$C_h=(A_h^*A_h)^{-1}A_h^*G_h$ and right-multiply by
$\left(\begin{smallmatrix}I&-C_h\\0&I\end{smallmatrix}\right)$.
The transformed source is $(A_h,\overline G_h)$ with orthogonal column
ranges.  Since $K_0^{-1}A_h=A_h/2$, the signed cross term vanishes as well.
\end{proof}

\subsection{Exact scalar--trace reconstruction of the signed response}

Work after any regular positive block elimination and let $F$ denote the
remaining full-rank source with actual target $b$.  Define
\begin{equation}
 \Pi=F(F^*F)^{-1}F^*,\qquad
 w=F(F^*F)^{-1}b,
 \qquad E\tau=\tau I/\sqrt3,
 \label{eq:supp-exact-Pi-w}
\end{equation}
then
\begin{equation}
 C=E^*\Pi E,\qquad H=2I-3C,\qquad \zeta=E^*w.
 \label{eq:supp-exact-trace-data}
\end{equation}

\begin{theorem}[Exact scalar--trace response reconstruction]
\label{thm:supp-exact-scalar-reconstruction}
The signed Gram $F^*K_0^{-1}F$ is invertible if and only if $H$ is
invertible.  If $\chi=H^{-1}\zeta$, then the original physical response is
\begin{equation}
 \boxed{v=w-3(I-\Pi)E\chi,\qquad E^*v=-\chi,}
 \label{eq:supp-exact-scalar-v}
\end{equation}
and
\begin{equation}
 \boxed{\|v\|^2=\|w\|^2+9\langle\chi,(I-C)\chi\rangle.}
 \label{eq:supp-exact-scalar-norm}
\end{equation}
Consequently
\begin{equation}
 v_h\to0\quad\Longleftrightarrow\quad
 w_h\to0\ \hbox{and}\ H_h^{-1}E^*w_h\to0
 \label{eq:supp-exact-response-equivalence}
\end{equation}
on any regular cofinal family.
\end{theorem}
\begin{proof}
Write $F=QR$ with $Q^*Q=I$.  The signed Gram is congruent to
$(2I-3Q^*EE^*Q)/4$.  The nonzero eigenvalues of
$Q^*EE^*Q$ and $E^*QQ^*E=C$ agree, proving the invertibility equivalence.
The response is characterized by $F^*v=b$ and $K_0v\in\Ran F$.  The first
identity gives $\Pi v=w$ and the second gives
$(I-\Pi)v=3(I-\Pi)EE^*v$.  Applying $E^*$ yields
$H(E^*v)=-\zeta$, hence $E^*v=-\chi$ and
\eqref{eq:supp-exact-scalar-v}.  Orthogonality of the two displayed terms
gives \eqref{eq:supp-exact-scalar-norm}.
\end{proof}

Diagonalizing the compression $C$ gives the trace-angle form.  Equivalently,
with the positive polar decomposition $F=QR$ and
$T=Q^*\mathsf P_{\rm tr}Q$, $\eta=R^{-1}b$, one obtains
\begin{equation}
 v=(2I-3\mathsf P_{\rm tr})Q(2I-3T)^{-1}\eta.
 \label{eq:supp-exact-response-polar}
\end{equation}
If $Tu_j=\mu_ju_j$ and $\eta_j=\langle u_j,\eta\rangle$, then
\begin{equation}
 \boxed{\|\mathcal E^0\|^2
 =\sum_j\frac{4-3\mu_j}{(2-3\mu_j)^2}|\eta_j|^2,}
 \label{eq:supp-exact-physical-norm}
\end{equation}
so $\mu=2/3$ is the unique signed trace resonance.  This is the spectral
form of \Cref{thm:supp-exact-scalar-reconstruction}, not a different response
law.

For a cofinal regular family define
$\nu_h=\sum_j|\eta_{h,j}|^2\delta_{\mu_{h,j}}$.  Then
$\|\mathcal E_h^0\|\to0$ if and only if
$\nu_h([0,1])\to0$ and
\begin{equation}
 \lim_{\delta\downarrow0}\limsup_{h\to0}
 \int_{|2-3\mu|<\delta}\frac{\dd\nu_h(\mu)}{|2-3\mu|^2}=0.
 \label{eq:supp-exact-resonance-ui}
\end{equation}
The proof is the direct split of \eqref{eq:supp-exact-physical-norm} into a
region separated from $2/3$ and a shrinking neighborhood of the resonance.

\subsection{Time-dependent response and the necessary clock test}

Static response control is insufficient for a spacetime comparison.  Let
$F(t),b(t)$ be $C^r$ with full column rank and regular signed Gram in a fixed
physical tangent space, and let $\Pi,w,C,H,\zeta,\chi,v$ be the preceding
objects.

\begin{theorem}[Derivative hierarchy in scalar--trace coordinates]
\label{thm:supp-exact-time-hierarchy}
Put $\chi_0=\chi$ and, for $1\le j\le r$,
\begin{equation}
 \boxed{\chi_j=H^{-1}\left[\zeta^{(j)}+
 3\sum_{\ell=1}^j\binom j\ell C^{(\ell)}\chi_{j-\ell}\right].}
 \label{eq:supp-exact-time-sources}
\end{equation}
Then $\chi_j=\chi^{(j)}=-E^*v^{(j)}$.  In particular
\begin{equation}
 \chi'=H^{-1}(\zeta'+3C'\chi),\qquad
 v'=w'+3\Pi'E\chi-3(I-\Pi)E\chi'.
 \label{eq:supp-exact-first-time-source}
\end{equation}
Thus a small static response does not control the differentiated response
without bounds on the transported source geometry.
\end{theorem}
\begin{proof}
Differentiate $H\chi=\zeta$ repeatedly and use $H^{(\ell)}=-3C^{(\ell)}$.
The identity $E^*v=-\chi$ holds at every time; differentiating
\eqref{eq:supp-exact-scalar-v} gives the second formula.
\end{proof}

The normalized exact continuation has an active/weak amplitude grading
$S_a=\operatorname{diag}(aI_A,a^2I_W)$.  If $u_a$ is the normalized source
coefficient entering the multiplier solve, then
\begin{equation}
 \|\lambda_t\|_h^2=a^{-2}\|(u_a)_A\|_h^2
                    +a^{-4}\|(u_a)_W\|_h^2.
 \label{eq:supp-exact-rate-norm}
\end{equation}

\begin{proposition}[Necessary lapse--shift clock condition for harmonic shadowing]
\label{prop:supp-exact-clock-test}
Suppose $a(h)\to0$ and an exact-action metric is to shadow a small harmonic
metric in a topology controlling the initial first time derivative, with the
same lapse, shift and spatial metric at the initial cut.  Then necessarily
\begin{equation}
 \boxed{a(h)^{-1}\|(u_{a(h)})_A\|_h\to0,
 \qquad a(h)^{-2}\|(u_{a(h)})_W\|_h\to0.}
 \label{eq:supp-exact-rate-necessary}
\end{equation}
This requirement is independent of smallness of the instantaneous electric
coefficient.
\end{proposition}
\begin{proof}
At $N=1$, $\beta=0$ in ADM coordinates,
$\partial_tg_{00}=-2\partial_tN$ and
$\partial_tg_{0i}=\gamma_{ij}\partial_t\beta^j$.  First-jet shadowing
therefore forces $\|\lambda_t\|\to0$.  Equation
\eqref{eq:supp-exact-rate-norm} gives the two conditions.
\end{proof}

\subsection{Prepared exact-action families}

The generic probe-source calculation remains useful for comparison.  For the
smooth first jet used in the original response audit, the actual weak source
obeys
\begin{equation}
 \|\mathcal I_hd_{h,\rho}-h^2d_\rho^{(2)}\|_{H^r}\le C_rh^4,
 \qquad c_rh^2\le\|d_{h,\rho}\|_{H_h^r}\le C_rh^2.
 \label{eq:supp-exact-source-expansion}
\end{equation}
Thus the generic weak Ward defect is genuinely second order, but
\Cref{thm:supp-exact-scalar-reconstruction} shows why its norm alone does not
control the physical response.

We next record two exact preparations with different purposes.

\paragraph{A rationally detuned all-cutoff active family.}
Let $N=2m+1$, $h=N^{-1}$ and define the one-dimensional product-defect operator
\begin{equation}
 D_hu=\delta(fu)-f\delta u-(\delta f)u,
 \qquad f(x)=\cos(2\pi x).
 \label{eq:supp-exact-cosine-defect}
\end{equation}
It is self-adjoint, has kernel the constants, and its nonconstant part is an
irreducible zero-diagonal Jacobi matrix.  For positive amplitudes
$\alpha=(\alpha_1,\alpha_2,\alpha_3)$ put
\begin{equation}
 U_\alpha=\sum_{i=1}^3\alpha_iT_i\cos(2\pi x_i),\qquad
 p_{h,\alpha}=-s_{h,\alpha}I,
 \qquad
 s_{h,\alpha}=\omega_h\sqrt{\frac{\alpha_1^2+\alpha_2^2+\alpha_3^2}{6}},
 \label{eq:supp-exact-detuned-seed}
\end{equation}
where $T_i$ is the transverse off-diagonal tensor used in
\eqref{eq:supp-initial-seed}.  The diagonal lapse part of the complete first
source is
\begin{equation}
 \mathscr D_{h,\alpha}n=
 \begin{pmatrix}
 (\alpha_2D_{2,h}-\alpha_3D_{3,h})n\\
 (\alpha_3D_{3,h}-\alpha_1D_{1,h})n\\
 (\alpha_1D_{1,h}-\alpha_2D_{2,h})n
 \end{pmatrix}.
 \label{eq:supp-exact-diagonal-source}
\end{equation}

\begin{lemma}[Finite rational detuning]
\label{lem:supp-exact-rational-detuning}
For every sufficiently fine odd cutoff there is a rational
$\theta_h$ with $1<\theta_h<1+h$ such that, for
$\alpha_h=(1,\theta_h,1)$,
\begin{equation}
 \|\mathscr D_{h,\alpha_h}n\|_h\ge ch^6\|n\|_h
 \qquad(\overline n=0).
 \label{eq:supp-exact-diagonal-margin}
\end{equation}
The selection can be made by a finite algebraic search at each cutoff.
\end{lemma}
\begin{proof}
The Jacobi representation of $D_h$ has simple algebraic spectrum
$\Sigma_h$.  Its reciprocal-eigenvalue sum is $O(h^{-3})$.  The union of
ratios $\theta$ for which
$|\lambda_a-\theta\lambda_b|<4h^6$ has measure $O(h^2)$, smaller than a
fixed fraction of the interval $(1,1+h)$ for small $h$.  Choose a dyadic
rational outside that union and sufficiently close to the selected point.
In the tensor product eigenbasis, the three squared components of
\eqref{eq:supp-exact-diagonal-source} vanish simultaneously only for the
constant mode, and the spectral separation gives the stated margin.
\end{proof}

\begin{theorem}[All-cutoff regular active source]
\label{thm:supp-exact-active-regular}
The exact original-action initial constraints have analytic preparations
\begin{equation}
 \mathscr C_h(X_{h,\alpha_h}(a))=0,
 \qquad X_{h,\alpha_h}(a)=a(U_{\alpha_h},p_{h,\alpha_h})+O(a^2),
 \label{eq:supp-exact-detuned-preparation}
\end{equation}
uniformly for all sufficiently fine odd cutoffs.  On their first jets the
complete leading active source has
\begin{equation}
 \ker\mathcal A_h^{\rm act}=\mathbb R(1,0)
 \oplus\{(0,\beta):\operatorname{curl}_\delta\beta=0\},
 \qquad \rank\mathcal A_h^{\rm act}=3V-3,
 \label{eq:supp-exact-active-rank}
\end{equation}
its range is trace-free, and
\begin{equation}
 \boxed{(\mathcal A_h^{\rm act})^*K_0^{-1}\mathcal A_h^{\rm act}
 =\tfrac12(\mathcal A_h^{\rm act})^*\mathcal A_h^{\rm act}
 \succeq ch^{16}I.}
 \label{eq:supp-exact-active-Gram}
\end{equation}
Hence the complete remaining signed rank problem is confined to the weak
columns.
\end{theorem}
\begin{proof}
The exact preparation follows from the finite initial-constraint chart and
the nondegenerate four-mean calibration; the $O(h)$ detuning changes the
prepared data by $O(|a|h)$ and not their continuum first jet.  The diagonal
lapse margin is \Cref{lem:supp-exact-rational-detuning}.  The transverse
shift block is controlled by \Cref{thm:supp-exact-transverse}.  Together they
leave only constant lapse and curl-free shifts in the first-source kernel.
The entire first source is trace-free on this three-amplitude family, so
$K_0^{-1}$ acts by $1/2$ on its range.  The lapse estimate costs $h^{-6}$ and
the transverse reconstruction an additional two powers, giving a sufficient
smallest singular value $ch^8$ and therefore the displayed Gram bound.
\end{proof}

\paragraph{A source-neutral nontrivial family.}
For each axis let $T_i$ be the transverse off-diagonal tensor and let $P_i$
be the transverse diagonal trace-free tensor with entries $+1,-1$ on the two
coordinates orthogonal to axis $i$.  Put
\begin{equation}
 U=\sum_{i=1}^3r_i\left[T_i\cos(2\pi x_i+\phi_i)
          +\sigma_iP_i\sin(2\pi x_i+\phi_i)\right],
 \qquad r_i>0,\quad\sigma_i\in\{-1,1\}.
 \label{eq:supp-exact-balanced-family}
\end{equation}
The two polarizations on each axis are orthogonal and have equal Frobenius
norm.

\begin{theorem}[Exact source-neutral finite preparation]
\label{thm:supp-exact-neutral-preparation}
For the family \eqref{eq:supp-exact-balanced-family}, the complete weak
quadratic target vanishes at every odd cutoff.  On compact nonzero parameter
sets the original finite constraint map admits exact preparations
\begin{equation}
 \boxed{\mathscr C_h(X_h(a))=0,
 \qquad X_h(a)=a(U,p_h)+O(a^2)}
 \label{eq:supp-exact-neutral-prepared}
\end{equation}
uniformly for sufficiently fine odd cutoffs.  At a concrete three-site seed
the complete leading signed Gram is nonsingular; hence a relatively open
balanced neighborhood has an exact local original-action continuation for
small nonzero $a$.  On that branch the literal initial temporal and spatial
link curvatures obey
\begin{equation}
 \boxed{\|E_h(a,0)\|_h+\|F_h^{\rm sp}(a,0)\|_h\le C_h|a|.}
 \label{eq:supp-exact-initial-curvature}
\end{equation}
The constant in this last estimate is fixed-cutoff; no common physical time
is inferred.
\end{theorem}
\begin{proof}
The weak quadratic drift on one axis is a sum of two polarization invariants:
the norm difference and the Frobenius cross product.  Equal norms and
orthogonality make both vanish, independently of $h$.  The linear constraints
hold because the configurations and oscillatory momenta are transverse and
trace-free.  The nonconstant rows are solved by the uniform range inverse of
the initial-constraint chart.  The four means are solved by a homogeneous
trace momentum and three oscillatory balancing parameters; the derivative of
the quadratic mean map is uniformly invertible on compact nonzero
polarization sets.  This gives \eqref{eq:supp-exact-neutral-prepared}.  At the
explicit three-site seed, finite exact elimination gives a nonzero determinant
for the complete signed Gram, and analyticity preserves it on a neighborhood.
Because the normalized source target starts one amplitude order later on the
balanced branch, its true finite-amplitude response is $O_h(a)$.  Substitution
in the literal temporal-link derivative and spatial plaquette formulas gives
\eqref{eq:supp-exact-initial-curvature}.
\end{proof}

\subsection{Direct transverse response and the positive scalar flux equation}

A further weak-neutral preparation can be chosen so that the leading active
range is trace-free at every sufficiently fine odd cutoff.  On this family
the transverse source admits an exact least-squares factorization.  Let
$T_h$ be the transverse source, $B_h=T_h^*T_h$, and let the normalized
positive normal matrix $\mathsf H_h$ be the finite-band operator obtained by
factoring the phase curl and coefficient multiplication.  It satisfies
\begin{equation}
 \frac9{16}I\preceq\mathsf H_h\preceq324I.
 \label{eq:supp-exact-H-bounds}
\end{equation}
The actual pseudoinverse has the factorization
\begin{equation}
 T_h^\dagger=\frac{2\sqrt2}{h\omega_h}
 \Omega_h^{-1}\mathsf H_h^{-1}R_h\widehat{\mathcal L}_h^*\iota_o^*,
 \label{eq:supp-exact-T-dagger}
\end{equation}
with the operators defined by the physical Hodge and coefficient
factorization of the source; no inverse-then-project substitution is made.

\begin{theorem}[Actual direct transverse geometric suppression]
\label{thm:supp-exact-direct-transverse}
For the actual prepared transverse target $b_{\beta,h}$ and direct solution
$\beta_{0,h}=-2B_h^{-1}b_{\beta,h}$,
\begin{equation}
 \boxed{\|\beta_{0,h}\|_{H_h^{s+3}}\le C_sh^{-1},
 \qquad \|\mathfrak E_h(\beta_{0,h})\|_{H_h^s}\le C_sh,}
 \label{eq:supp-exact-direct-result}
\end{equation}
while the folded-face tail is exponentially small.  These are estimates of
the actual direct source response, not of a prescribed low-frequency
projection.
\end{theorem}
\begin{proof}
The factorization \eqref{eq:supp-exact-T-dagger} and
\eqref{eq:supp-exact-H-bounds} give two derivatives of Sobolev gain for
$B_h^{-1}$ at cost $h^{-2}$.  The prepared transverse target is $O(h)$ on a
fixed Fourier bank.  Applying the geometric observation operator, whose
symbol contributes the compensating $h^2$ factor at the required order,
gives the second estimate.  Exponential off-diagonal decay of
$\mathsf H_h^{-1}$ in cyclic Fourier distance gives the folded-face tail.
\end{proof}

After eliminating the transverse variable, the actual leading lapse solves
\begin{equation}
 \mathscr S_hn_h=g_h.
 \label{eq:supp-exact-scalar-equation}
\end{equation}
Let $Q_h$ project onto the physical off-diagonal complement and let
$K_{Q,h}$ be the induced positive coefficient metric.

\begin{theorem}[Exact positive scalar energy]
\label{thm:supp-exact-scalar-flux}
At every odd cutoff,
\begin{equation}
 \frac1{324}I\preceq K_{Q,h}\preceq\frac{16}{9}I,
 \label{eq:supp-exact-KQ-bounds}
\end{equation}
and the unchanged scalar operator is
\begin{equation}
 \boxed{\mathscr S_h
 =\tfrac12\mathscr D_h^*\mathscr D_h
  +\tfrac{h^2}{2}\sigma_h^*K_{Q,h}^{-1}\sigma_h.}
 \label{eq:supp-exact-S-flux}
\end{equation}
Hence
\begin{equation}
 2\langle n,\mathscr S_hn\rangle_h
 =\|\mathscr D_hn\|_h^2
  +h^2\langle\sigma_h(n),K_{Q,h}^{-1}\sigma_h(n)\rangle_h.
 \label{eq:supp-exact-scalar-energy}
\end{equation}
On the bare diagonal soft subspace the normalized full energy diverges under
refinement, so a pole predicted by the diagonal compression alone is not a
pole of the full physical scalar energy.  This compression result does not by
itself prove a uniform inverse for \eqref{eq:supp-exact-scalar-equation}.
\end{theorem}
\begin{proof}
The coefficient map giving the off-diagonal source is uniformly invertible,
so compression of its positive inverse metric gives
\eqref{eq:supp-exact-KQ-bounds}.  Orthogonal projection of the physical
source yields the exact residual factor
$hM_h^{-*}Q_hK_{Q,h}^{-1}\sigma_h(n)$; squaring and adding the diagonal
source gives \eqref{eq:supp-exact-S-flux}.  On the finite-dimensional soft
subspace the limiting flux is nonzero on every nonzero vector, while the
bare diagonal contribution is higher order; compactness therefore forces the
normalized full energy to diverge.
\end{proof}

\subsection{Dynamical tangency: an exact obstruction and an upstream escape direction}

Exact preparation is not the same as invariance of the prepared set under the
original action vector field.  Let $\mathcal W$ be the weak curl-free
multiplier space and let $\mathbb K_1(Y)$ be the first Poisson matrix at a
canonical first jet $Y$.  For a leading multiplier tilt $\ell$ define the
weak null space
\begin{equation}
 \mathcal N_W(Y)=\{\nu\in\mathcal W:
 \mathbb K_1(Y)[\nu,\mu]=0\ \hbox{for every multiplier }\mu\},
 \label{eq:supp-exact-weak-null}
\end{equation}
and the first tangency functional
\begin{equation}
 \boxed{\Delta_\nu(Y,\ell)=
 Dd_{2,W}(Y)[F_1Y+G\ell,\nu]
 +\mathbb K_1(F_1Y+G\ell)[\nu,\ell].}
 \label{eq:supp-exact-tangency-functional}
\end{equation}
It is determined entirely by the canonical first jet and leading multiplier
tilt; higher preparation coefficients do not enter this leading null test.

\begin{theorem}[Initial-layer obstruction on an exactly prepared branch]
\label{thm:supp-exact-initial-layer}
There is a nontrivial exactly prepared three-site original-action family
$X_a$, $\lambda_a$ with
\begin{equation}
 \mathsf B_W(X_a,\lambda_a)=0,
 \qquad (\lambda_t)_W(a,0)=0,
 \qquad \|\operatorname{Riem}(g_h)(a,0)\|=O_h(a),
 \label{eq:supp-exact-prepared-zero-rate}
\end{equation}
but with a weak null test $\nu_*$ satisfying
$\Delta_{\nu_*}(Y,\ell)\ne0$.  Along the exact continuation,
\begin{equation}
 \left.\frac{\dd}{\dd t}\mathsf B_W[\nu_*]\right|_{t=0}
 =a^2\Delta_{\nu_*}(Y,\ell)+O_h(a^3),
 \label{eq:supp-exact-tangency-growth}
\end{equation}
which forces a weak multiplier acceleration of order $a^{-2}$ in the
normalized grading.  For every fixed $S$ the exact histories exist on
$0\le t\le a^3S$ for sufficiently small $a$, and the unscaled rates and
curvature remain $O_h(a)$ there.  Nevertheless no constants $T,C,a_0>0$ can
make
\begin{equation}
 \sup_{0\le t\le T}\|R_{0i0j}(g_h)(a,t)\|\le Ca
 \label{eq:supp-exact-no-linear-fixed-time}
\end{equation}
hold for every $0<a<a_0$ on this family.
\end{theorem}
\begin{proof}
Taylor expansion of the exact weak drift gives
$D_X\mathsf B_W=a\mathcal C_\ell+O_h(a^2)$ and
$D_\lambda\mathsf B_W=a\mathcal K_W+O_h(a^2)$.  Pairing with a Poisson-null
$\nu_*$ removes the unknown active multiplier acceleration at leading order
and yields \eqref{eq:supp-exact-tangency-growth}.  Rescale time by
$t=a^3s$; the finite-dimensional analytic vector field converges on every
fixed $s$-interval to a linear displacement system, giving the shrinking
layer and the nonzero limiting curvature slope.  If a uniform $Ca$ bound
held on a fixed physical interval, choose a fixed $s$ beyond the corresponding
normalized bound and then take $a$ sufficiently small so that $a^3s<T$,
contradicting that limiting slope.
\end{proof}

The leading multiplier tilt has a finite affine freedom after the lower
quadratic, normalization and cubic-mean conditions are imposed.  On the
preceding canonical first jet the tangency functional is constant and nonzero
on that entire admissible tilt family; hence changing only the leading
multiplier tilt cannot repair the obstruction.

The canonical first jet itself has additional admissible directions.  Write
$\mathcal L(Y,\ell)$ for the complete lower incidence consisting of the
quadratic drift, multiplier normalizations and cubic mean rows.  At the
selected base point its multiplier derivative has constant rank, and one
canonical direction $Z_2$ admits an exact multiplier tangent $m_2$ with
\begin{equation}
 D\mathcal L(Y,\ell)[Z_2,m_2]=0.
 \label{eq:supp-exact-lower-tangent}
\end{equation}

\begin{theorem}[The fixed-first-jet tangency obstruction is deformable]
\label{thm:supp-exact-upstream-deformation}
For the exact lower-incidence tangent $(Z_2,m_2)$ above, one weak null test
$\nu_\diamond$ remains Poisson-null to first order and
\begin{equation}
 \boxed{D\Delta_{\nu_\diamond}(Y,\ell)[Z_2,m_2]\ne0.}
 \label{eq:supp-exact-upstream-nonzero}
\end{equation}
Thus the nonzero tangency obstruction on the preceding prepared first jet is
not invariant under all admissible canonical first-jet deformations of the
unchanged finite action.
\end{theorem}
\begin{proof}
Differentiate every row of the finite lower-incidence map in the selected
canonical direction and solve the resulting characteristic-zero linear system
for $m_2$.  The same weak test has zero first variation of its Poisson-null
condition in each coordinate direction of the local first-jet chart.  The
tangency functional is quadratic in the simultaneous first-jet/multiplier
increment, so its derivative is obtained exactly by polarization.  Direct
finite algebra gives a nonzero rational coefficient; hence the derivative is
nonzero over characteristic zero.  The calculation establishes a tangent,
not an exact nonlinear root.
\end{proof}

\subsection{Exact scope of the raw-action route}

The preceding results change the status of the direct exact-action problem in
three ways.  First, a large part of the multiplier response is analytically
regular and exactly eliminable at every cutoff.  Second, source-neutral and
active-regular finite preparations exist and give genuine original-action
curvature information.  Third, the first unresolved dynamical condition is
now an explicit tangency hierarchy rather than an unspecified failure of a
finite determinant.

A direct raw-action Einstein theorem would still require one common family on
which all of the following hold:
\begin{enumerate}[label=\textup{(R\arabic*)},leftmargin=3em]
\item an exact nonlinear root of the complete transported tangency system in a
      regular canonical first-jet chart;
\item persistence of the active and higher weak signed ranks and of the
      physical source/observation bounds along a cofinal refinement;
\item a source-specific uniform estimate for the full scalar solve
      $\mathscr S_hn_h=g_h$, not merely for a diagonal compression;
\item differentiated multiplier and curvature-source estimates strong enough
      to control the literal connection/metric observations; and
\item a cutoff-independent physical interval on which the exact constrained
      Euler history remains in the common noncollapsed chart.
\end{enumerate}
These are not assumed elsewhere in the paper.  Conversely, the finite
characteristic-rank difference is not a no-go theorem for singular continuum
restoration of gauge symmetry.  The open-law Einstein population and the
exact-action hierarchy are therefore complementary: the first supplies a
proved nonsymmetric $3+1$ continuum class, while the second identifies how
much of that conclusion can already be recovered from the unchanged finite
action and isolates the remaining dynamical bridge.

% =============================================================================
% =============================================================================
\section{Homogeneous renewal dynamics and operational action matching}
\label{supp:homogeneous-dynamics}
% =============================================================================

This appendix proves the homogeneous dynamical statements of
Section~\ref{sec:vacuum-main}.  Arbitrary positive homogeneous renewal variables
are first varied with the full lapse constraint retained; the flat/de~Sitter
profiles are solved only after the stationary equations have been obtained.  A second layer constructs a finite positive
trial law whose measured probabilities reconstruct the corresponding finite
action off shell.

\subsection{Homogeneous ADM map and reduced vacuum action}

On the root-symmetric branch
$k_{\alpha,h}=\kappa/(8h^2)$ and $m_h=\varrho h^3$.  Since
$\sum_{\alpha\in\Phi_{A_3}}\alpha\alpha^{\mathsf T}=8I$, one has
$\mathcal B=\kappa I$ and $\beta=0$.  Substitution in
\eqref{eq:adm-inversion-main} gives
\eqref{eq:main-isotropic-renewal-map}.

For completeness, let
\[
 \dd s^2=-N(t)^2\dd t^2+a(t)^2\delta_{ij}\dd x^i\dd x^j.
\]
The nonzero Christoffel symbols are
\[
 \Gamma^0_{00}=\dot N/N,\qquad
 \Gamma^0_{ij}=a\dot a\,\delta_{ij}/N^2,\qquad
 \Gamma^i_{0j}=\dot a\,\delta^i_j/a.
\]
Direct contraction gives
\begin{align}
 R_{00}&=-3\left(\ddot a/a-\dot a\dot N/(aN)\right),\\
 R_{ij}&=\left(a\ddot a/N^2+2\dot a^2/N^2
              -a\dot a\dot N/N^3\right)\delta_{ij},\\
 R&=6\left(\ddot a/(aN^2)+\dot a^2/(a^2N^2)
              -\dot a\dot N/(aN^3)\right).
 \label{eq:supp-homogeneous-curvature}
\end{align}
Thus
\begin{equation}
 Na^3R=6\frac{\dd}{\dd t}\left(\frac{a^2\dot a}{N}\right)
       -\frac{6a\dot a^2}{N}.
 \label{eq:supp-homogeneous-total-derivative}
\end{equation}
With the orientation/sign convention used here, the bare torsion-free
Einstein--Hilbert representative on $[t_0,t_1]\times\mathbb T^3$ is therefore
\begin{equation}
 \frac{S_{\rm EH}}{\chi V_0}
 =6\left[\frac{a^2\dot a}{N}\right]_{t_0}^{t_1}
 +\int_{t_0}^{t_1}\left[-\frac{6a\dot a^2}{N}
                         -2\Lambda Na^3\right]\dd t.
 \label{eq:supp-homogeneous-eh-split}
\end{equation}
The Gibbons--Hawking--York Dirichlet boundary term
\cite{York1972,GibbonsHawking1977} contributes
\begin{equation}
 \frac{S_{\rm GHY}}{\chi V_0}
 =-6\left[\frac{a^2\dot a}{N}\right]_{t_0}^{t_1},
 \label{eq:supp-homogeneous-ghy}
\end{equation}
so $S_{\rm EH}+S_{\rm GHY}=\chi V_0S[a,N]$ with $S[a,N]$ exactly
\eqref{eq:main-homogeneous-action-an}.  The corresponding Dirichlet problem
fixes the induced spatial metric at the two endpoint slices,
$\delta a(t_0)=\delta a(t_1)=0$, while the lapse remains independently variable
in the interior.  The finite action \eqref{eq:main-discrete-homogeneous-action}
is a midpoint variational discretization of this already boundary-completed
first-order functional; hence fixing $q_0,q_M$ is precisely the reduced discrete
counterpart of the same induced-metric boundary condition.  Substituting
$a=\varrho^{1/3}$ and $N=\varrho^{1/3}\sqrt\kappa$ gives
\eqref{eq:main-renewal-rate-action}.

\begin{proof}[Detailed proof of \Cref{thm:main-homogeneous-friedmann}]
Varying the lapse and scale factor independently gives
\begin{equation}
 E_N=2a^3(3\mathcal H^2-\Lambda),\qquad
 E_a=6Na^2(2D_\tau\mathcal H+3\mathcal H^2-\Lambda).
 \label{eq:supp-homogeneous-EaEN}
\end{equation}
Because
$a=e^{v/3}$ and $N=e^{u/2+v/3}$,
\[
 E_u=\frac N2E_N,\qquad
 E_v=\frac a3E_a+\frac N3E_N,
\]
which proves \eqref{eq:main-renewal-euler}.  The Jacobian of
$(u,v)\mapsto(\log a,\log N)$ is $-1/6\ne0$, so this change of
variables loses no variation.  The reconstructed Einstein residual is
\begin{equation}
 G_{00}+\Lambda g_{00}=N^2(3\mathcal H^2-\Lambda),
 \qquad
 G_{ij}+\Lambda g_{ij}
 =-a^2(2D_\tau\mathcal H+3\mathcal H^2-\Lambda)\delta_{ij},
 \label{eq:supp-homogeneous-Einstein-residual}
\end{equation}
with $G_{0i}=0$.  Hence the two renewal Euler equations are equivalent to every
Einstein component within this ansatz.
\end{proof}

\begin{proof}[Proof of \Cref{cor:main-derived-desitter-rates}]
For $\Lambda>0$, the lapse constraint gives
$\mathcal H^2=H_0^2$ with $H_0=\sqrt{\Lambda/3}$.  Continuity on a
connected interval fixes one sign $\sigma$, while
$D_\tau\mathcal H=0$ makes it constant.  Integrating
$D_\tau\log a=\sigma H_0$ gives
\eqref{eq:main-stationary-homogeneous-solutions}; the renewal variables then
follow from \eqref{eq:main-isotropic-renewal-map}.  The cases $\Lambda=0$ and
$\Lambda<0$ follow from the same squared constraint.
\end{proof}

\begin{remark}[Why the lapse is varied before gauge fixing]
If $q=a^{3/2}$ and one sets $N=1$ before variation, the remaining equation is
$q''=\omega^2q$ with $\omega=\sqrt{3\Lambda/4}$.  Its general solution
$q=Ae^{\omega t}+Be^{-\omega t}$ satisfies the omitted lapse constraint
$(q')^2=\omega^2q^2$ only when $AB=0$.  Thus a positive cosh-type solution is a
spurious critical point of the prematurely gauge-fixed action.
\end{remark}

\subsection{Finite lapse-varied action and stationary recurrence}

Put $q_j>0$, $s_j>0$, $\bar q_j=(q_j+q_{j+1})/2$, and define
\eqref{eq:main-discrete-homogeneous-action}.  Because this discretizes the
boundary-completed first-order action, the discrete Dirichlet problem fixes
$q_0,q_M$ and varies $(q_1,\ldots,q_{M-1},s_0,\ldots,s_{M-1})$ independently.
A direct summation-by-parts calculation gives exactly
\eqref{eq:main-discrete-dirichlet-variation}; the endpoint momenta multiply
$\delta q_0$ and $\delta q_M$ and therefore vanish under the fixed induced
endpoint geometry.  No condition is imposed on $\sum_js_j$, since that would
modify the lapse variation and hence the Hamiltonian constraint.  The change of
variables \eqref{eq:main-discrete-renewal-map}, together with
$q_j=\sqrt{\varrho_j}$, is a smooth bijection on the positive finite domain.
Hence critical points in $(q,s)$ and in the corresponding positive renewal
coordinates are equivalent.

\begin{proof}[Proof of the Euler equations and stationary classification in
\Cref{thm:main-finite-homogeneous-stationarity}]
Differentiating an edge term with respect to $s_j$ gives
\[
 A_0\frac{(q_{j+1}-q_j)^2}{s_j^2}-B_0\bar q_j^2,
\]
which is \eqref{eq:main-discrete-lapse-constraint}.  At an interior node the
right derivative of the preceding edge and left derivative of the following
edge give \eqref{eq:main-discrete-scale-euler}.

For $\Lambda>0$, the lapse equation says
$v_j=\sigma_j\omega\bar q_j$ with $\sigma_j\in\{\pm1\}$.  Solving for the
terminal amplitude gives
\[
 q_{j+1}=q_j\frac{1+\sigma_j\omega s_j/2}
                      {1-\sigma_j\omega s_j/2}.
\]
Positivity requires $\omega s_j/2<1$.  Substitution in the interior momentum
matching condition shows
$2A_0\omega q_j(\sigma_j-\sigma_{j-1})=0$, so every positive stationary path
has one common sign.  Conversely a common-sign recurrence satisfies every
finite Euler equation.  For $\Lambda=0$ the lapse equation forces constant
$q$; for $\Lambda<0$ it has no positive real solution.

To prove \eqref{eq:main-finite-profile-convergence}, set
$x_j=\omega s_j/2$.  For $0\le x\le b$,
\[
 0\le2\operatorname{arctanh}x-2x
 =2\int_0^x\frac{t^2}{1-t^2}\dd t
 \le\frac{2x^3}{3(1-b^2)}.
\]
Taking logarithms of the recurrence and summing gives
\[
 \left|\log(q_j/q_0)-\sigma\omega\tau_j\right|
 \le\frac{\omega^3}{12(1-b^2)}\sum_i s_i^3
 \le\frac{\omega^3T}{12(1-b^2)}h_t^2.
\]
The power relations between $q,a,\varrho,\kappa$ then give the stated
convergence of renewal variables.
\end{proof}

For later residual estimates it is useful to define
\begin{equation}
 R_j^\sigma=q_{j+1}-q_j-\sigma\omega s_j\bar q_j.
 \label{eq:supp-homogeneous-residual}
\end{equation}
A direct expansion gives the exact factorization
\begin{equation}
 \boxed{
 S_{\rm d}=-A_0\sum_j\frac{(R_j^\sigma)^2}{s_j}
 -A_0\sigma\omega(q_M^2-q_0^2).}
 \label{eq:supp-homogeneous-square-factor}
\end{equation}
Thus at fixed endpoints the finite stationary branch is exactly the zero set of
the nonnegative residual energy
$\mathcal E=A_0\sum_j(R_j^\sigma)^2/s_j$.

\begin{lemma}[Full finite variation controlled by the residual energy]
\label{lem:supp-homogeneous-euler-control}
For endpoint-vanishing nodal variations $w_j$ and lapse variations
$\delta s_j=s_jn_j$, put
\begin{equation}
 \mathcal B_s(w,n;q)=\sum_js_j\left(
 \frac{w_{j+1}-w_j}{s_j}-\sigma\omega\bar w_j
 -\sigma\omega\bar q_jn_j\right)^2.
 \label{eq:supp-homogeneous-variation-norm}
\end{equation}
Then
\begin{equation}
 \boxed{
 |\delta S_{\rm d}[w,sn]|
 \le2\sqrt{A_0\mathcal E}\sqrt{\mathcal B_s(w,n;q)}
 +\mathcal E\|n\|_\infty.}
 \label{eq:supp-homogeneous-euler-bound}
\end{equation}
If $q_j^*$ is the exact common-sign stationary recurrence with the same
$q_0,s_j$, and $\omega\max s_j/2\le b<1$, then on a fixed positive amplitude
range
\begin{equation}
 \max_j|q_j-q_j^*|\le C\sqrt{\mathcal E}.
 \label{eq:supp-homogeneous-path-control}
\end{equation}
\end{lemma}

\begin{proof}
Differentiate \eqref{eq:supp-homogeneous-square-factor}; weighted
Cauchy--Schwarz gives \eqref{eq:supp-homogeneous-euler-bound}.  For the path
estimate, write the forced recurrence as
$q_{j+1}=r_jq_j+R_j/d_j$ with
$d_j=1-\sigma\omega s_j/2$ and
$r_j=(1+\sigma\omega s_j/2)/d_j$.  Iteration, $d_j\ge1-b$, and
$\sum_j|R_j|\le\sqrt{T\sum_jR_j^2/s_j}$ give the result.
\end{proof}

\subsection{Measured quadratic costs and exact off-shell action matching}

We now prove that the finite action can be read from an actual normalized trial
law rather than appended as an unrelated deterministic writer.  Let
$E_s(x,y)$ be \eqref{eq:main-measured-quadratic-cost}.

\begin{proposition}[Three-cost identification and rank defect]
\label{prop:supp-three-costs}
Fix $s>0$ and distinct $u,v>0$.  Within the declared quadratic class,
$E_s(u,u)$, $E_s(u,v)$, and $E_s(v,u)$ determine
\begin{align}
 B&=\frac{E_s(u,u)}{su^2},\\
 C&=\frac{E_s(u,v)-E_s(v,u)}{2(v^2-u^2)},\\
 A&=\frac{s}{2(v-u)^2}\left[
 E_s(u,v)+E_s(v,u)
 -2Bs\left(\frac{u+v}{2}\right)^2\right].
 \label{eq:supp-three-cost-identification}
\end{align}
Moreover, with $\mathfrak D=AB-C^2$, $c=-C/A$, and
$\nu=\mathfrak D/A$,
\begin{equation}
 \boxed{
 E_s(x,y)=\frac A s\left[y-x-cs\frac{x+y}{2}\right]^2
 +\nu s\left(\frac{x+y}{2}\right)^2.}
 \label{eq:supp-cost-completion}
\end{equation}
Thus $\mathfrak D=0$ is exactly the rank-one compatibility condition.
\end{proposition}

\begin{proof}
The diagonal measurement determines $B$.  Reversal exchanges $x$ and $y$,
leaving the first two terms invariant and changing the sign of the $C$ term,
which determines $C$; the reversal sum then determines $A$.  Completing the
square gives \eqref{eq:supp-cost-completion}.
\end{proof}

\begin{proof}[Proof of \Cref{thm:main-operational-action-matching}]
From \eqref{eq:main-accepted-row-normalizer},
$Z_sQ_s=p_sa_s$.  Taking logarithms and using
$a_s=e^{-E_s/\varepsilon}$ proves \eqref{eq:main-observed-cost}.  Summing over
one path gives
\[
 \varepsilon\sum_j\left[\log(Q_j/p_j)+\log Z_j\right]
 =-\sum_jE_{s_j}(q_j,q_{j+1}).
\]
The $C$ term telescopes to $C(q_M^2-q_0^2)$, which is exactly cancelled by the
endpoint term in \eqref{eq:main-observed-homogeneous-action}; this proves the
off-shell identity \eqref{eq:main-offshell-action-match}.  Differentiating that
identity on the smooth positive parametric domain proves equality of all
interior and lapse first variations.

For the mismatch statement, minimizing $y\mapsto E_s(x,y)$ gives, for small
$s$, a ratio with logarithmic drift $c=-C/A$, so the limiting scale factor has
Hubble rate $2c/3$.  Hence its geometric cosmological parameter is
$\Lambda_{\rm geom}=4C^2/(3A^2)$.  Dividing the bulk action by $A/A_0$ shows
that its parameter is $\Lambda_{\rm act}=4B/(3A)$.  Therefore
\[
 G[g]+\Lambda_{\rm act}g
 =(\Lambda_{\rm act}-\Lambda_{\rm geom})g
 =\frac{4(AB-C^2)}{3A^2}g,
\]
which is \eqref{eq:main-measured-einstein-mismatch}.
\end{proof}

\begin{remark}[Why the acceptance normalizer is physical]
Retaining only $p_s$ and $Q_s$ would not determine the action: replacing
$E_s(x,y)$ by $E_s(x,y)+c_s(x)$ leaves $Q_s$ unchanged but multiplies $Z_s(x)$
by $e^{-c_s(x)/\varepsilon}$.  Such a source-dependent term can change the
lapse variation.  The total acceptance frequency is therefore a load-bearing
part of the complete physical trial table.
\end{remark}

\subsection{An explicit finite stochastic cylinder}

We give the construction underlying
\Cref{thm:main-explicit-operational-family}.  Set $s=T/M$ and, for large $M$,
\begin{equation}
 d=1-\sigma\omega s/2,\qquad
 r=\frac{1+\sigma\omega s/2}{1-\sigma\omega s/2},\qquad
 \delta_M=M^{-3},\qquad
 \varepsilon_M=\frac{M^{-5}}{1+\log M}.
 \label{eq:supp-explicit-scales}
\end{equation}
Starting from $F_0=\{q_0\}$, define support intervals recursively by
\[
 L_0=q_0/2,\quad U_0=3q_0/2,\qquad
 L_{j+1}=rL_j-\delta_M,\quad U_{j+1}=rU_j+\delta_M,
\]
and let $F_j$ be finite grids of mesh at most $\delta_M$ in these intervals.
Each source proposes uniformly from $F_{j+1}$ and accepts with
\begin{equation}
 a_j(x,y)=\exp\left[-\frac{A_0d^2(y-rx)^2}{s\varepsilon_M}\right].
 \label{eq:supp-explicit-acceptance}
\end{equation}
Because
$d(y-rx)=y-x-\sigma\omega s(x+y)/2$, this is exactly the rank-one cost
\eqref{eq:main-rank-one-cost}.

\begin{lemma}[History-level entropy--energy estimate]
\label{lem:supp-entropy-energy}
At stage $j$, let a finite reference row $p_j$ and nonnegative cost $E_j$ define
$Q_j=p_je^{-E_j/\varepsilon}/Z_j$.  Let $\mathbb Q$ be the product of these
accepted rows and $\mathbb P\ll\mathbb Q$ any complete path law with the same
initial value, possibly with memory.  If each relevant row contains a proposal
of cost at most $e_j$ and every proposal weight is at least $p_{*,j}$, then
\begin{equation}
 \boxed{
 \mathbb E_{\mathbb P}\sum_jE_j
 \le\sum_je_j+\varepsilon\left[
 \sum_j\log(1/p_{*,j})+D_{\rm KL}(\mathbb P\Vert\mathbb Q)\right].}
 \label{eq:supp-entropy-energy}
\end{equation}
\end{lemma}

\begin{proof}
For a fixed history prefix, direct substitution gives
\[
 \mathbb E_P E_j
 =\varepsilon D_{\rm KL}(P_j\Vert Q_j)
 -\varepsilon D_{\rm KL}(P_j\Vert p_j)-\varepsilon\log Z_j.
\]
The candidate bound gives
$Z_j\ge p_{*,j}e^{-e_j/\varepsilon}$.  Integrating over prefixes, summing, and
using the chain rule for path relative entropy proves
\eqref{eq:supp-entropy-energy}.
\end{proof}

\begin{proof}[Proof of \Cref{thm:main-explicit-operational-family}]
The rational stationary multiplier and its inverse are uniformly bounded on
$0\le j\le M$ when $\omega s/2\le b<1$.  Iterating the support recursion shows
that all $F_j$ lie in one fixed interval $0<q_-\le q\le q_+<\infty$ and have
cardinality $O(M^3)$.  For every source $x$, the exact target $rx$ lies inside
the next support and has a grid point within $\delta_M$.  The corresponding
cost is at most
$C\delta_M^2/s=CM^{-5}$.  Hence the proposal floor is $p_*\ge cM^{-3}$ and the
total acceptance probability has a positive explicit lower bound.

Apply \Cref{lem:supp-entropy-energy} to the ideal accepted product law
$\mathbb Q_M$ itself.  With the scales
\eqref{eq:supp-explicit-scales},
\begin{equation}
 \mathbb E_{\mathbb Q_M}\mathcal E_M\le CM^{-4}.
 \label{eq:supp-explicit-energy}
\end{equation}
The deterministic estimates
\eqref{eq:supp-homogeneous-euler-bound} and
\eqref{eq:supp-homogeneous-path-control}, Jensen's inequality, and the
second-order stationary-profile estimate then give the two bounds in
\eqref{eq:main-stochastic-profile} and an $O(M^{-2})$ expected first variation
against every fixed smooth scale/lapse test.

To make every cutoff history genuinely finite, permit at most
\[
 R_M=\left\lceil7\underline Z_M^{-1}\log M\right\rceil
\]
independent proposals at each stage, where $\underline Z_M$ is the explicit
row-wise acceptance lower bound.  A union bound gives total failure probability
at most $M^{-6}$.  Conditional on complete success, the path-law correction has
relative entropy at most $-\log(1-M^{-6})$, which changes
\eqref{eq:supp-explicit-energy} only by a summable term.  The expected residual
energies are summable in $M$, so Tonelli and the first Borel--Cantelli lemma give
$\mathcal E_M\to0$ almost surely under any common coupling.  The deterministic
path estimates then imply uniform and strong $H^1$ convergence to
$q_0e^{\sigma\omega\tau}$.

Finally, \eqref{eq:main-stochastic-root-rates} is the positive ADM inverse on
the proper-time nodal branch.  The $A_3$ root identity gives exactly
$\mathcal B_j=\kappa_jI$, and ordinary Taylor expansion of the symmetric root
generator yields the spatial continuum coefficient
$\tfrac12q_*^{-4/3}\Delta$.  Together with the homogeneous Einstein identity,
this proves the vacuum limit with $\Lambda=4\omega^2/3$.
\end{proof}

\begin{remark}[Exact scope of the operational construction]
The quadratic duration law, the rank-one condition $\mathfrak D=0$, the
nonzero action normalization, and the value of $\Lambda$ are not deduced from
the bare predictive quotient.  They are finite operational conditions that can
be measured and can fail.  The construction proves that this matched class is
nonempty and robust under vanishing measured action/likelihood errors; it does
not prove generic microscopic attraction to that class.
\end{remark}

\subsection{Exact conservative blocking and the relative defect}
\label{subsec:supp-homogeneous-blocking}

Put $u=A/s+Bs/4$, $v=A/s-Bs/4$, $\mu=\sqrt{AB}$ and
$\theta=2\operatorname{artanh}((s/2)\sqrt{B/A})$.  Then
$u^2-v^2=\mu^2$, $u=\mu\coth\theta$ and
$v=\mu\operatorname{csch}\theta$, giving
\eqref{eq:main-hyperbolic-cost}.

\begin{lemma}[Two-edge Schur identity]
\label{lem:supp-two-edge-schur}
Let
$E_i(x,y)=p_ix^2-2q_ixy+r_iy^2$ with $p_i,q_i,r_i>0$ and
$d_i=p_ir_i-q_i^2\ge0$.  Then the minimizer of
$E_1(x,z)+E_2(z,y)$ is positive on positive endpoints and the blocked quadratic
has
\begin{equation}
 p'=p_1-\frac{q_1^2}{r_1+p_2},\qquad
 q'=\frac{q_1q_2}{r_1+p_2},\qquad
 r'=r_2-\frac{q_2^2}{r_1+p_2},
 \label{eq:supp-schur-coefficients}
\end{equation}
with
\begin{equation}
 \boxed{p'r'-(q')^2=\frac{r_2d_1+p_1d_2}{r_1+p_2}.}
 \label{eq:supp-schur-determinant}
\end{equation}
\end{lemma}
\begin{proof}
Complete the square in $z$.  The minimizer is
$z_*=(q_1x+q_2y)/(r_1+p_2)>0$.  The three coefficients in
\eqref{eq:supp-schur-coefficients} are the remaining quadratic coefficients;
direct expansion gives \eqref{eq:supp-schur-determinant}.
\end{proof}

\begin{proof}[Proof of \Cref{thm:main-no-blocking-attractor}]
For $\mathcal E_{\mu,\theta_i,C}$ put
$u_i=\mu\coth\theta_i$ and $v_i=\mu\operatorname{csch}\theta_i$.
The two intermediate $Cz^2$ terms cancel.  Applying
\Cref{lem:supp-two-edge-schur} and $u_i^2-v_i^2=\mu^2$ gives
\begin{equation}
 u'=\frac{u_1u_2+\mu^2}{u_1+u_2},\qquad
 v'=\frac{v_1v_2}{u_1+u_2},\qquad C'=C.
\end{equation}
The hyperbolic addition formulas identify these with
$\mu\coth(\theta_1+\theta_2)$ and
$\mu\operatorname{csch}(\theta_1+\theta_2)$, proving
\eqref{eq:main-hyperbolic-blocking}.  Converting the two-identical-edge result
back to midpoint coefficients gives \eqref{eq:main-homogeneous-block-map}.
Hence $AB$, $C$, $\mathfrak D$, and
$\Delta_{\rm rel}=\mathfrak D/(AB)$ are invariant.

For a change of amplitude and time units, $x=r\widetilde x$ and
$s=\ell\widetilde s$, followed by common action rescaling by $d>0$, one has
$\widetilde A=r^2A/(d\ell)$,
$\widetilde B=r^2\ell B/d$, and $\widetilde C=r^2C/d$; therefore
$C^2/(AB)$ is unchanged.

For the finite positive statement, let $F\subset(0,\infty)$ be finite and
write the subprobability kernel as
\begin{equation}
 L_C(x,y)=e^{Cx^2/\varepsilon}L_0(x,y)e^{-Cy^2/\varepsilon},
 \qquad L_0(x,y)=L_0(y,x)>0.
 \label{eq:supp-finite-twist-similarity}
\end{equation}
Then $L_C^n=D_CL_0^nD_C^{-1}$ and hence, for $x\ne y$,
\begin{equation}
 \boxed{\frac{\varepsilon}{2(x^2-y^2)}
 \log\frac{L_C^n(x,y)}{L_C^n(y,x)}=C.}
 \label{eq:supp-finite-twist-read}
\end{equation}
Thus finite blocking preserves the measured endpoint twist exactly.  If the
intermediate positive amplitudes are integrated over a grid of covering error
$d$ and cardinality $m$, completion of the positive interior quadratic form
around its exact bridge minimizer gives
\begin{equation}
 0\le E_{F,n}^{\varepsilon}(x,y)-
 \mathcal E_{\mu,n\theta,C}(x,y)
 \le4u(n-1)d^2+\varepsilon(n-1)\log m.
 \label{eq:supp-finite-block-error}
\end{equation}
At fixed total duration one may take $d_n=n^{-2}$,
$m_n=O(n^2)$, and $\varepsilon_n=n^{-3}/(1+\log n)$, making the right side
$O(n^{-2})$.  The three-probe coefficient reconstruction is continuous on a
compact calibrated range with $AB>0$, so every initial
$\Delta_{\rm rel}>0$ remains bounded away from zero for sufficiently fine
finite blocking.  This proves the operational statement.
\end{proof}

% =============================================================================

% =============================================================================
\subsubsection{The cross-scale optimizer test}

\begin{proof}[Proof of \Cref{prop:main-optimizer-composition}]
The coefficient of $y^2$ in
$\mathcal E_{\mu,\theta,C}(x,y)$ is $\mu\coth\theta+C>0$, since
$|C|\le\mu$ and $\theta>0$.  Its unique minimizer is therefore
\[
 y=\frac{\mu\operatorname{csch}\theta}{\mu\coth\theta+C}\,x
       =\frac{x}{\cosh\theta+\rho\sinh\theta}>0.
\]
Set $d_\theta=\cosh\theta+\rho\sinh\theta$.  Hyperbolic addition gives
\[
 d_{\theta+\phi}-d_\theta d_\phi
       =(1-\rho^2)\sinh\theta\sinh\phi,
\]
which proves \eqref{eq:main-optimizer-defect}.  All denominators and both
hyperbolic sines are positive.  Equality of the maps for even one positive
pair and one nonzero amplitude therefore forces $1-\rho^2=0$, equivalently
$AB=C^2$.  Conversely, $\rho=\pm1$ gives $d_\theta=e^{\rho\theta}$ and
exact composition for every pair.  For $\Delta_{\rm rel}>0$ the formula
also gives $r_{\theta+\phi}<r_\theta r_\phi$.

Solving the same identity for the defect gives
\begin{equation}
 \Delta_{\rm rel}
 =\frac{r_{\theta+\phi}^{-1}-(r_\theta r_\phi)^{-1}}
              {\sinh\theta\sinh\phi}.
 \label{eq:supp-optimizer-defect-read}
\end{equation}
On a compact range where the observed ratios and
$\sinh\theta\sinh\phi$ are bounded away from zero, the right-hand side
has bounded first derivatives in its measured arguments.  It is consequently
locally Lipschitz there.  This is stability of a finite diagnostic, not
transverse attraction under the conservative blocking operation.
\end{proof}

\section{Nonlinear two-polarization vacuum regulator}
\label{supp:gowdy}
% =============================================================================

This appendix proves \Cref{thm:main-gowdy-regulator}.  The purpose is not to
establish unrestricted four-dimensional global regularity, but to show that the
finite renewal evolution, constraint, readout, and independent action tests can
be carried through a genuinely nonlinear spatially inhomogeneous vacuum sector.
We use the unpolarized Gowdy $T^3$ class in areal coordinates
\cite{Moncrief1981}; the lapse and shift are nevertheless retained as
independent variables until after the action variation.

\subsection{First-order frame system and constraints}

On $[t_0,t_1]\times\mathbb T^3$, $t_0>0$, consider
\begin{equation}
 g=t^{-1/2}e^{\lambda/2}(-\dd t^2+\dd\theta^2)
 +t\left[e^P(\dd x+Q\dd y)^2+e^{-P}\dd y^2\right].
 \label{eq:supp-gowdy-metric}
\end{equation}
With
$a=P_t$, $b=P_\theta$, $c=e^PQ_t$, $d=e^PQ_\theta$ and
\begin{equation}
 e=\frac12(a^2+b^2+c^2+d^2),\qquad f=ab+cd,
 \label{eq:supp-gowdy-ef}
\end{equation}
the vacuum equations take the first-order form
\begin{align}
 a_t&=b_\theta-a/t+c^2-d^2,&b_t&=a_\theta,\nonumber\\
 c_t&=d_\theta-c/t-ac+bd,&d_t&=c_\theta+ad-bc,\nonumber\\
 P_t&=a,&Q_t&=e^{-P}c,&\lambda_t&=2te,
 \label{eq:supp-gowdy-frame-system}
\end{align}
with
\begin{equation}
 P_\theta=b,\qquad e^PQ_\theta=d,\qquad
 \lambda_\theta=2tf.
 \label{eq:supp-gowdy-coordinate-constraints}
\end{equation}
A direct differentiation using \eqref{eq:supp-gowdy-frame-system} gives
\begin{equation}
 \boxed{
 e_t-f_\theta=-\frac{a^2+c^2}{t},\qquad
 f_t-e_\theta=-\frac ft.}
 \label{eq:supp-gowdy-stress-identities}
\end{equation}
Consequently
$\partial_t(\lambda_\theta-2tf)=0$ whenever the displayed equations hold, and
$\partial_t\int_{\mathbb T^1}tf\,\dd\theta=0$.  Thus the momentum constraint
and the periodic zero-mean condition for $\lambda$ propagate.

\subsection{Independent lapse and shift variation before areal gauge}

Let $(R,\lambda,P,Q,N,\beta)$ be independent fields with $R,N>0$, put
$v_u=\dot u-\beta u'$, and define
\begin{align}
 \mathscr L_{\rm G}={}&
 -\frac{v_R(v_\lambda-4\beta')}{2N}
 +\frac R{2N}(v_P^2+e^{2P}v_Q^2)\nonumber\\
 &+N\left[-2R''+\frac{R'\lambda'}2
 -\frac R2\{(P')^2+e^{2P}(Q')^2\}\right].
 \label{eq:supp-gowdy-adm-density}
\end{align}
The canonical momenta are
\begin{equation}
 p_R=-\frac{v_\lambda-4\beta'}{2N},\qquad
 p_\lambda=-\frac{v_R}{2N},\qquad
 \pi=\frac{Rv_P}{N},\qquad
 \varpi=\frac{Re^{2P}v_Q}{N}.
 \label{eq:supp-gowdy-momenta}
\end{equation}
Lapse and shift variation give respectively
\begin{align}
 \mathcal C={}&-2p_Rp_\lambda
 +\frac{\pi^2+e^{-2P}\varpi^2}{2R}
 +\frac R2\{(P')^2+e^{2P}(Q')^2\}
 +2R''-\frac12R'\lambda',\label{eq:supp-gowdy-H}\\
 \mathcal D={}&p_RR'+p_\lambda\lambda'+\pi P'+\varpi Q'-4p_\lambda'.
 \label{eq:supp-gowdy-D}
\end{align}
On a solution of \eqref{eq:supp-gowdy-frame-system}, set
\begin{equation}
 R=t,\quad N=1,\quad\beta=0,\quad
 p_\lambda=-\frac12,\quad p_R=-te,\quad
 \pi=ta,\quad\varpi=te^Pc.
 \label{eq:supp-gowdy-areal-lift}
\end{equation}
Substitution of \eqref{eq:supp-gowdy-coordinate-constraints} shows
$\mathcal C=0$ and $\mathcal D=0$; the remaining area equation follows from
\eqref{eq:supp-gowdy-stress-identities}.  Hence the areal-gauge solution is the
restriction of a critical point of the unreduced symmetry-class action rather
than a solution obtained after deleting the lapse and shift equations.

\subsection{Local splitting and exact staggered momentum}

Set
\begin{equation}
 U=(u_1,u_2,u_3,u_4)=\sqrt t\,(a,b,c,d),\qquad
 \mathcal J(U)=u_1u_2+u_3u_4,
 \qquad n(U)=|U|^2.
 \label{eq:supp-gowdy-rescaled-frame}
\end{equation}
The local source flow, with the evolution parameter denoted by $s$, is
\begin{align}
 \frac{\dd u_1}{\dd s}&=-\frac{u_1}{2t}
 +\frac{u_3^2-u_4^2}{\sqrt t},&
 \frac{\dd u_2}{\dd s}&=\frac{u_2}{2t},\nonumber\\
 \frac{\dd u_3}{\dd s}&=-\frac{u_3}{2t}
 +\frac{-u_1u_3+u_2u_4}{\sqrt t},&
 \frac{\dd u_4}{\dd s}&=\frac{u_4}{2t}
 +\frac{u_1u_4-u_2u_3}{\sqrt t},\nonumber\\
 \frac{\dd P}{\dd s}&=\frac{u_1}{\sqrt t},&
 \frac{\dd Q}{\dd s}&=e^{-P}\frac{u_3}{\sqrt t},&
 \frac{\dd t}{\dd s}&=1,
 \label{eq:supp-gowdy-source-flow}
\end{align}
with $\dd\lambda/\dd s=0$.  Let $\mathsf S_\sigma$ be its near-identity
implicit-midpoint map over source duration $\sigma$.

On a periodic grid of spacing $\ell$, let $S_\pm$ be the two shifts,
$D_+=(S_+-I)/\ell$, and $\mathsf A_+=(I+S_+)/2$.  Put
\begin{equation}
 w^+=\frac1{\sqrt2}(u_1+u_2,u_3+u_4),\qquad
 w^-=\frac1{\sqrt2}(u_1-u_2,u_3-u_4),\qquad
 g^\pm=|w^\pm|^2.
 \label{eq:supp-gowdy-characteristics}
\end{equation}
At frozen physical time define the transport map by
\begin{equation}
 (w^+)^+=S_+w^+,
 \qquad (w^-)^+=S_-w^-,
 \qquad
 \lambda^+=\lambda+\frac\ell2
 [(I+S_+)g^++(I+S_-)g^-],
 \label{eq:supp-gowdy-transport}
\end{equation}
leaving $P,Q$ fixed.

\begin{theorem}[Exact staggered momentum under symmetric splitting]
\label{thm:supp-gowdy-momentum}
With synchronized full step $h=\ell$, the source-half/transport/source-half
composition preserves exactly
\begin{equation}
 \boxed{\mathcal K_\ell
 =D_+\lambda-2\mathsf A_+\mathcal J(U).}
 \label{eq:supp-gowdy-staggered-momentum}
\end{equation}
The update is nearest neighbour, its background increment is a nonnegative sum
of squared characteristic amplitudes, and it is second-order consistent on
bounded smooth reference charts with $t\ge t_0>0$.
\end{theorem}

\begin{proof}
For the source flow, direct substitution in
\eqref{eq:supp-gowdy-source-flow} gives
$\nabla\mathcal J(U)\cdot\dot U=0$.  The difference of a quadratic function at
two endpoints equals its derivative at the arithmetic midpoint paired with the
increment, so implicit midpoint preserves $\mathcal J$ exactly while leaving
$\lambda$ fixed.  For the transport step,
\begin{equation}
 \Delta\mathcal J
 =\frac12[(S_+-I)g^+-(S_--I)g^-],\qquad
 D_+\Delta\lambda=2\mathsf A_+\Delta\mathcal J,
 \label{eq:supp-gowdy-transport-identity}
\end{equation}
which proves exact preservation without invoking an assumed discrete product
rule.  The $\lambda$ increment in \eqref{eq:supp-gowdy-transport} is a positive
sum of squared norms.  Characteristic translation is exact for the transport
fields, while the background increment is a trapezoidal approximation to the
integral of the translated squares with local error $O(\ell^3)$.  Implicit
midpoint has the same local order for the smooth source, and symmetric
composition therefore gives second-order consistency.
\end{proof}

\subsection{Residual, constraint, metric, and action budgets}

Let
$X_j=(U_j,P_j,Q_j,\lambda_j)$ satisfy
\begin{equation}
 X_{j+1}=\Phi_h(t_j,X_j)+r_j,
 \qquad h=\ell,
 \label{eq:supp-gowdy-residual-update}
\end{equation}
and define
\begin{equation}
 \mathcal E_h=\sum_j\frac{\|r_j\|_{2,\ell}^2}{2h},\qquad
 \mathcal F_h=\ell^{-(2r+1)}\mathcal E_h.
 \label{eq:supp-gowdy-residual-energy}
\end{equation}
Here $\mathsf S_\ell$ denotes grid sampling.

\begin{proposition}[Residual and constraint control]
\label{prop:supp-gowdy-residual-control}
On a fixed smooth reference chart, with initial $C^r_\ell$ error $\epsilon_0$
and either a smallness bootstrap or an enforced offset envelope,
\begin{equation}
 \boxed{
 \max_j\|X_j-\mathsf S_\ell X_*(t_j)\|_{r,\infty,\ell}
 \le C(\epsilon_0+h^2+\sqrt{\mathcal F_h}).}
 \label{eq:supp-gowdy-state-error}
\end{equation}
Writing $\widehat U_j$ for the deterministic predicted frame, the exact
constraint increment is
\begin{align}
 \mathcal K_\ell(X_{j+1})-\mathcal K_\ell(X_j)
 ={}&D_+r_j^\lambda\nonumber\\
 &-2\mathsf A_+
 [\nabla\mathcal J(\widehat U_j)\cdot r_j^U
   +\mathcal J(r_j^U)].
 \label{eq:supp-gowdy-constraint-increment}
\end{align}
For $r\ge1$, its accumulated norm is bounded by the initial residual plus
$C\sqrt{\mathcal F_h}$.
\end{proposition}

\begin{proof}
The deterministic split map has smooth-chart Lipschitz constant $1+Ch$ and
consistency error $Ch^3$.  Discrete Gronwall therefore bounds the global error
by
$C(\epsilon_0+h^2+\sum_j\|r_j\|_{r,\infty,\ell})$.
The one-dimensional inverse grid estimate and Cauchy--Schwarz bound the last
sum by $C\sqrt{\mathcal F_h}$, proving
\eqref{eq:supp-gowdy-state-error}.  Expanding the quadratic $\mathcal J$ about
$\widehat U_j$ and using exact preservation
\eqref{eq:supp-gowdy-staggered-momentum} gives
\eqref{eq:supp-gowdy-constraint-increment}; the same inverse estimate controls
its accumulated norm.
\end{proof}

For the pathwise curvature statement specialize the finite offset alphabet so
that every accepted offset obeys
\begin{equation}
 \|\epsilon_j\|_\infty\le c_0h^4,
 \label{eq:supp-gowdy-offset-cap}
\end{equation}
and take initial $C_h^1$ error $O(h^2)$.  The retained source implementation may
add a solver error $O(h^5)$ in $C_h^1$.  Since
$\|D_hv\|_\infty\le2h^{-1}\|v\|_\infty$, the offset has $C_h^1$ size
$O(h^3)$.  The one-step Lipschitz factor $1+Lh$, second-order consistency, and
discrete Gronwall then give, pathwise,
\begin{equation}
 \boxed{
 \max_j\|X_j-\mathsf S_hX_*(t_j)\|_{1,\infty,h}\le Ch^2.}
 \label{eq:supp-gowdy-pathwise-c1}
\end{equation}
For $f=P,Q,\lambda$, subtraction of the reference step and division by $h$
also gives an $O(h^2)$ bound for the nodal time differences of the errors.

Use the unchanged derivative records
\begin{align}
 u&=\left(u_1/\sqrt t,\ e^{-P}u_3/\sqrt t,\ |U|^2\right),\nonumber\\
 v&=\left(u_2/\sqrt t,\ e^{-P}u_4/\sqrt t,\ 2\mathcal J(U)\right),
 \qquad w=D_0u,
 \label{eq:supp-gowdy-hermite-jets}
\end{align}
where $D_0=(S_+-S_-)/(2h)$.  These are respectively the time, space, and mixed
jet records for $(P,Q,\lambda)$.  On each spacetime square let
$\mathcal H_h(f,u,v,w)$ be the tensor-product cubic Hermite interpolant using
\begin{align*}
 H_0(q)&=2q^3-3q^2+1,&H_1(q)&=-2q^3+3q^2,\\
 G_0(q)&=q^3-2q^2+q,&G_1(q)&=q^3-q^2.
\end{align*}
The identities $H_0+H_1=1$, $H_0'+H_1'=0$, and
$H_0''+H_1''=0$ imply the stability estimate
\begin{equation}
 \|F\|_\infty+\|DF\|_\infty\le C\eta,
 \qquad \|D^2F\|_\infty\le C\eta/h,
 \label{eq:supp-gowdy-hermite-second-jet}
\end{equation}
whenever the two nodal jet banks and the first nodal differences of their value
records differ by at most $\eta$.  Shared corner jets make the interpolant
$C^1$ across cell interfaces, so its weak second derivative is the cellwise
second derivative and has no interface measure.

\begin{theorem}[Pathwise full-curvature control on the Gowdy alphabet]
\label{thm:supp-gowdy-full-curvature}
Assume the smooth reference has bounded derivatives through the order required
by the Hermite remainder, \eqref{eq:supp-gowdy-offset-cap}, and
\eqref{eq:supp-gowdy-pathwise-c1}.  For the shared Hermite reconstruction of
$f=(P,Q,\lambda)$,
\begin{equation}
 \boxed{
 \|f_h-f_*\|_{W^{1,\infty}}\le Ch^2,
 \qquad
 \|D^2f_h-D^2f_*\|_{L^\infty}\le Ch.}
 \label{eq:supp-gowdy-second-jet}
\end{equation}
For the reconstructed metric \eqref{eq:supp-gowdy-metric},
\begin{align}
 \|g_h-g_*\|_{W^{1,\infty}}&\le Ch^2,\nonumber\\
 \|D^2g_h-D^2g_*\|_{L^\infty}
 +\|\operatorname{Riem}(g_h)-\operatorname{Riem}(g_*)\|_{L^\infty}
 &\le Ch.
 \label{eq:supp-gowdy-full-curvature}
\end{align}
Hence on every compact positive subslab
\begin{equation}
 \boxed{\sup_h\|\operatorname{Riem}(g_h)\|_{L^2}<\infty,
 \qquad \|\operatorname{Ric}(g_h)\|_{L^\infty}\le Ch.}
 \label{eq:supp-gowdy-curvature-certificate}
\end{equation}
All estimates are pathwise and the curvature limit is the Levi--Civita
curvature of $g_*$.  No all-time curvature norm is assumed.
\end{theorem}
\begin{proof}
The right sides of \eqref{eq:supp-gowdy-hermite-jets} are smooth functions on
the standing chart.  Equation \eqref{eq:supp-gowdy-pathwise-c1}, the nodal time
difference estimate, and centered-difference consistency therefore make every
corner value/first/mixed jet error $O(h^2)$.  Apply
\eqref{eq:supp-gowdy-hermite-second-jet}; the exact-reference cubic Hermite
remainder is $O(h^{4-k})$ through derivative order $k=2$.  This proves
\eqref{eq:supp-gowdy-second-jet}.

On $t\ge t_0>0$ the Gowdy metric map is smooth on the bounded chart and has
determinant $-te^\lambda$, so its inverse stays uniformly bounded.  The metric
estimates follow by differentiating this smooth composition twice.  In
coordinates, $\Gamma=g^{-1}*Dg$, hence
$\|\Gamma_h-\Gamma_*\|_\infty\le Ch^2$ and
$\|D\Gamma_h-D\Gamma_*\|_\infty\le Ch$.  Because $g_h$ is $C^1$ across
interfaces there is no distributional surface term.  Substitution in
$\operatorname{Riem}=D\Gamma+\Gamma*\Gamma$ gives
\eqref{eq:supp-gowdy-full-curvature}.  The reference is vacuum, so contraction
with the uniformly bounded metric and inverse gives the Ricci estimate; finite
volume gives the $L^2$ curvature certificate.
\end{proof}

A finite positive realization is obtained by giving each predicted local state
a finite offset list containing zero and contained in
\eqref{eq:supp-gowdy-offset-cap}, a strictly positive reference row, and
acceptance probability $\exp[-E_j/\varepsilon]$ for a nonnegative deviation
cost.  If $p_{*,j}$ is the reference mass of the zero offset, the exact Gibbs
identity gives
\begin{equation}
 \boxed{
 \mathbb E\sum_jE_j
 \le\varepsilon\left[
 \mathcal K_{\rm path}+\sum_j\log(1/p_{*,j})\right],}
 \label{eq:supp-gowdy-gibbs-bound}
\end{equation}
where $\mathcal K_{\rm path}$ is the chain-rule relative entropy of the observed
history law against the declared accepted law.  Product reference masses count
all local trials; solver residuals, finite retry failures, and iteration errors
remain separate budgets.  The positive selection cost is not identified with
the signed ADM action.  Since the curvature estimates above are deterministic
and uniform over the permitted alphabet, every successful history inherits
\eqref{eq:supp-gowdy-curvature-certificate}.

\subsection{Measured proposal access and finite execution}

The deterministic second-jet estimate is uniform over the stated offset
alphabet.  It therefore survives changes of the actual proposal weights,
provided the acceptance rule retains that support and successful proposals
really occur.  At a fixed stage and source history write
$\mu=p(\mathcal G)$ for the measured mass of the admissible set, and use
\[
 a(u)=\one_{\mathcal G}(u)e^{-E(u)/\varepsilon},\qquad
 Z=\sum_up(u)a(u),\qquad Q(u)=p(u)a(u)/Z.
\]
This is a declared selection rule, not an unrecorded replacement of the
original microscopic transition.

\begin{proof}[Proof of \Cref{prop:main-safe-access}]
Since $0\le E\le c$ on $\mathcal G$,
$\mu e^{-c/\varepsilon}\le Z\le\mu$.
For $\mu>0$, let $p_{\mathcal G}=p/\mu$ on $\mathcal G$.  Every
$Q\ll p$ supported on $\mathcal G$ obeys the exact identity
\[
 D_{\rm KL}(Q\Vert p)=D_{\rm KL}(Q\Vert p_{\mathcal G})-\log\mu
       \ge-\log\mu.
\]
In particular $\mu=0$ precludes such an accepted row altogether.
No choice of temperature can correct this support failure.

On the accepted support the complete trial table still reconstructs its cost:
$E(u)=-\varepsilon\log(ZQ(u)/p(u))$.  Averaging also gives
$\mathbb E_QE=-\varepsilon\log Z-
\varepsilon D_{\rm KL}(Q\Vert p)$.  These identities concern the recorded
positive deviation cost; the independently varied signed action is unchanged.

With independent reset retries conditional on the same source, the probability
that $R_j$ attempts fail is $(1-Z_j)^{R_j}\le e^{-R_jZ_j}$.
The same bound follows for history-dependent retries when the displayed lower
success bound holds conditionally after every failed attempt.  Taking the
supported lower bound for each stage and using a union bound proves
\eqref{eq:main-safe-retry-bound}.  For $J$ stages and total failure target
$0<\delta<1$, one sufficient finite choice is
\begin{equation}
 R_j=\left\lceil\mu_j^{-1}e^{c_j/\varepsilon_j}
                         \log(J/\delta)\right\rceil.
 \label{eq:supp-safe-retry-cap}
\end{equation}
Conditioning the entire path on success may reweight its intermediate source
states.  No equality with the unconditioned product of accepted rows is needed:
every successful path has the same deterministic admissible support, and hence
the same curvature and action bounds.  Summability of the complete failure
probabilities and the first Borel--Cantelli lemma give the final assertion,
without independence between refinements.
\end{proof}

\begin{lemma}[A finite execution certificate for the proposal floor]
\label{lem:supp-executed-proposal-floor}
At a retained source snapshot let a launch race have finitely many bins.
Conditional on no earlier launch, let $a_\ell$ be the probability of any
launch and $f_\ell$ the probability of a predictor-marked launch in bin
$\ell$.  Set $S_0=1$, $S_\ell=\prod_{i\le\ell}(1-a_i)$.
Suppose
\begin{equation}
 f_\ell\ge\rho_*a_\ell,\qquad
 -\log S_J\ge A_*>0,\qquad 0<\rho_*\le1.
 \label{eq:supp-launch-share}
\end{equation}
If, conditional on every predictor launch, the actually executed writer
returns a certified member of $\mathcal G$ before the attempt deadline with
probability at least $q_*>0$, then
\begin{equation}
 \boxed{p(\mathcal G)\ge q_*\rho_*(1-e^{-A_*}).}
 \label{eq:supp-executed-proposal-floor}
\end{equation}
\end{lemma}
\begin{proof}
The disjoint first-launch events give predictor launch probability
$\sum_{\ell=1}^J S_{\ell-1}f_\ell$.  Since
$S_{\ell-1}a_\ell=S_{\ell-1}-S_\ell$, this is at least
$\rho_*(1-S_J)\ge\rho_*(1-e^{-A_*})$.  On each such event the conditional
completed-safe-output probability is at least $q_*$.  Multiplication of
conditional probabilities and summation prove the claim.  Launch,
interruption, completion and output certificates belong to the actual
same-history record; a possible but unexecuted numerical predictor supplies
none of these probabilities.
\end{proof}

The execution probability and deadline in this lemma are physical inputs,
not consequences of a fast geometric root clock.  For example, one root clock
competing with $D$ clocks of equal rate has share $1/(D+1)$, regardless of
their common rate scale.  Moreover, a cap on the number of retries bounds their
physical duration only when each attempt has a certified duration bound.  With
attempt durations at most $d_j$, the allowed stage duration must accommodate
$R_jd_j$.  Thus the proposition removes the need to assume a separate
continuum proposal profile on a class with the displayed finite execution
budgets, without asserting a uniform computation or sampling cost for every
microscopic law.

For the general residual/action estimate below put
\begin{equation}
 \kappa_h:=\epsilon_0+h^2+\sqrt{\mathcal F_h}.
 \label{eq:supp-gowdy-kappa}
\end{equation}

Finally discretize \eqref{eq:supp-gowdy-adm-density} by midpoint time sampling
and centered spatial differences, keep $N$ and $\beta$ independent, and replace
$-2NR''$ by $2N'R'$ under the periodic boundary convention.  Evaluation at the
recorded areal-gauge fields gives, for every fixed smooth compactly supported
independent test $(\delta R,\delta\lambda,\delta P,\delta Q,\delta N,\delta\beta)$,
\begin{equation}
 \boxed{|\delta S_h-\delta S_*|\le C\kappa_h.}
 \label{eq:supp-gowdy-action-test}
\end{equation}
The estimate follows from the same first-jet comparison used above, and the
variation is taken before the areal-gauge restriction.  Hence
$\mathbb E\mathcal F_h=O(h^4)$ and $\epsilon_0=O(h^2)$ yield second-order mean
metric and action-test errors.  Exact preservation of
$\mathcal K_\ell$ is one constraint statement and is not substituted for the
independent lapse/shift action tests.

\begin{proof}[Proof of \Cref{thm:main-gowdy-regulator}]
Item~\textup{(G1)} is \Cref{thm:supp-gowdy-momentum}.  Item~\textup{(G2)} is
\Cref{prop:supp-gowdy-residual-control}.  The pathwise $C_h^1$ estimate and
\Cref{thm:supp-gowdy-full-curvature} give item~\textup{(G3)}.  The unreduced
variational calculation
\eqref{eq:supp-gowdy-adm-density}--\eqref{eq:supp-gowdy-areal-lift}, the action
test \eqref{eq:supp-gowdy-action-test}, and the Gibbs identity
\eqref{eq:supp-gowdy-gibbs-bound} give item~\textup{(G4)} while retaining the
lapse and shift variations before gauge restriction.
\end{proof}

\section{Flat and de Sitter regulator limits}
\label{supp:vacuum}
% =============================================================================

This section proves the spatial and Cartan-continuum parts of
\Cref{mt:vacuum}.  The temporal profiles are no longer posited here: they are
the stationary renewal solutions derived in
Appendix~\ref{supp:homogeneous-dynamics}.  The present appendix verifies that
the associated $A_3$ spatial generators have the required normalization and
compact screens, and that the reconstructed coframes and exact link transports
converge to the corresponding flat and de~Sitter Cartan geometries.

\subsection{The \texorpdfstring{$A_3$}{A3} identity and periodic regulators}

Let
\begin{equation}
 W_0=\left\{x\in\R^4:\sum_{i=1}^4x_i=0\right\},
 \qquad
 \Lambda_{A_3}=W_0\cap\Z^4,
 \label{eq:supp-a3-lattice}
\end{equation}
and let
\begin{equation}
 \Phi_{A_3}=\{e_i-e_j:1\le i\ne j\le4\}
 \label{eq:supp-a3-oriented-roots}
\end{equation}
be the twelve oriented roots.  The scalar product on $W_0$ is the restriction
of the standard Euclidean scalar product on $\R^4$.

\begin{lemma}[$A_3$ second-moment identity]
\label{lem:supp-a3-second-moment}
On $W_0$,
\begin{equation}
 \boxed{\sum_{\alpha\in\Phi_{A_3}}
        \alpha\alpha^{\mathsf T}=8I_{W_0}.}
 \label{eq:supp-a3-second-moment}
\end{equation}
Equivalently, for every $v\in W_0$,
\begin{equation}
 \sum_{i\ne j}(v_i-v_j)^2=8\norm{v}^2.
 \label{eq:supp-a3-quadratic-identity}
\end{equation}
\end{lemma}

\begin{proof}
For $v\in W_0$, expand
\[
 \sum_{i\ne j}(v_i-v_j)^2
 =\sum_{i\ne j}(v_i^2+v_j^2-2v_iv_j).
\]
Each $v_i^2$ occurs six times in the first two sums, while
\[
 \sum_{i\ne j}v_iv_j
 =\left(\sum_i v_i\right)^2-\sum_i v_i^2
 =-\sum_i v_i^2.
\]
Hence the displayed expression is
$6\sum_i v_i^2+2\sum_i v_i^2=8\norm{v}^2$.
Polarization gives \eqref{eq:supp-a3-second-moment}.
\end{proof}

For an integer $d\ge2$, put $h=d^{-1}$ and
\begin{equation}
 V_h=\Lambda_{A_3}/d\Lambda_{A_3}.
 \label{eq:supp-periodic-vertices}
\end{equation}
We regard $hx$, $x\in V_h$, as a periodic mesh of the flat torus
$\mathbb T^3_{A_3}=W_0/\Lambda_{A_3}$.  An oriented root $\alpha$ connects
$x$ to $x+\alpha$ in $V_h$.  Let $m_h>0$ be the mass of each vertex and let
$k_{\alpha,h}>0$ be the jump rate along the oriented root $\alpha$.  The
Markov generator on functions $f:V_h\to\C$ is
\begin{equation}
 (L_hf)(x)
 =\sum_{\alpha\in\Phi_{A_3}}k_{\alpha,h}
  \bigl[f(x+\alpha)-f(x)\bigr].
 \label{eq:supp-root-generator}
\end{equation}
Its positive Dirichlet form on
$L^2(V_h,m_h)$ is
\begin{equation}
 \mathcal E_h(f)
 =-\ip{f}{L_hf}_{L^2(m_h)}
 =\frac{m_h}{2}\sum_{x\in V_h}
   \sum_{\alpha\in\Phi_{A_3}}k_{\alpha,h}
   \abs{f(x+\alpha)-f(x)}^2,
 \label{eq:supp-root-form}
\end{equation}
whenever the rates are reversal symmetric.

\begin{lemma}[Fourier symbol and flat generator limit]
\label{lem:supp-flat-symbol}
Suppose
\begin{equation}
 k_{\alpha,h}=\frac1{8h^2}
 \qquad(\alpha\in\Phi_{A_3}).
 \label{eq:supp-flat-rates}
\end{equation}
For a dual-lattice mode $e_k(x)=e^{2\pi i h\langle k,x\rangle}$,
$-L_h$ has eigenvalue
\begin{equation}
 \lambda_h(k)
 =\frac1{4h^2}\sum_{a=1}^{6}
  \left[1-\cos\bigl(2\pi h\langle r_a,k\rangle\bigr)\right],
 \label{eq:supp-flat-symbol}
\end{equation}
where $r_1,\ldots,r_6$ is one orientation of the six root pairs.  For every
fixed $k$,
\begin{equation}
 \lambda_h(k)\longrightarrow2\pi^2\norm{k}^2,
 \label{eq:supp-flat-symbol-limit}
\end{equation}
and, under the natural piecewise-constant embeddings,
\begin{equation}
 L_h\longrightarrow\frac12\Delta_{\mathbb T^3_{A_3}}
 \label{eq:supp-flat-generator-limit}
\end{equation}
in strong resolvent sense and strongly at the semigroup level.
\end{lemma}

\begin{proof}
Pair the roots $\alpha$ and $-\alpha$ in
\eqref{eq:supp-root-generator}.  This gives
\eqref{eq:supp-flat-symbol}.  Taylor expansion yields
\[
 \lambda_h(k)
 =\frac{\pi^2}{2}\sum_{a=1}^6\langle r_a,k\rangle^2+O_k(h^2).
\]
Because the sum over all twelve roots is twice the sum over the six unoriented
roots, \Cref{lem:supp-a3-second-moment} gives
$\sum_{a=1}^6r_ar_a^{\mathsf T}=4I$, and therefore the leading term is
$2\pi^2\norm{k}^2$.

The trigonometric symbol is nonnegative and converges pointwise on the common
Fourier core.  On each fixed finite Fourier band the convergence is uniform.
The corresponding quadratic forms converge on trigonometric polynomials and
have the common lower bound zero.  Their closure gives Mosco convergence to
the form of $\tfrac12\Delta$; the resolvent and semigroup conclusions follow
from the Mosco--resolvent equivalence.  Equivalently, one may apply dominated
convergence to the Fourier multipliers and then use density.
\end{proof}

\begin{lemma}[Uniform spatial screens for the periodic $A_3$ graph]
\label{lem:supp-periodic-screens}
Let the vertex mass be $m_h=a^3h^3$ and the symmetric root rate be
$k_{\alpha,h}=a^{-2}/(8h^2)$, where $a$ ranges in a compact interval
$0<a_-\le a\le a_+<\infty$.  Then there are constants
$c_{\rm P},C_{\rm W},c_{\rm cut}>0$, depending only on
$a_-,a_+$ and the fixed torus, such that for all sufficiently small $h$:
\begin{enumerate}[label=\textup{(\roman*)},leftmargin=2.8em]
\item the first nonzero eigenvalue of $-L_h$ is at least $c_{\rm P}$;
\item the low-energy counting function satisfies
      \begin{equation}
       N_h(R):=\#\{\lambda\in\spec(-L_h):\lambda\le R\}
       \le C_{\rm W}(1+R^{3/2});
       \label{eq:supp-weyl-count}
      \end{equation}
\item with edge conductance
      $c_{e,h}=m_hk_{\alpha,h}=ah/8$, every vertex set $A$ with
      $0<\mu_h(A)\le\mu_h(V_h)/2$ obeys
      \begin{equation}
       h\,c_h(\partial A)
       \ge c_{\rm cut}\,\mu_h(A)^{2/3}.
       \label{eq:supp-periodic-cut}
      \end{equation}
\end{enumerate}
Consequently the spectral projections
$P_{h,R}=\one_{[0,R]}(-L_h)$ have uniformly bounded rank and
\begin{equation}
 \norm{(I-P_{h,R})f}_{L^2(m_h)}^2
 \le R^{-1}\mathcal E_h(f).
 \label{eq:supp-screen-tail}
\end{equation}
Thus they form a uniformly exhausted compact-screen family.
\end{lemma}

\begin{proof}
The scalar $a^{-2}$ multiplies the flat symbol
\eqref{eq:supp-flat-symbol}, and the compact interval for $a$ changes all
spectral comparisons only by fixed constants.  For $\abs{2\pi h\langle
r_a,k\rangle}\le\pi$, the inequalities
$2s^2/\pi^2\le1-\cos s\le s^2/2$, together with the root-frame identity, give
\begin{equation}
 c_1\norm{k}^2\le\lambda_h(k)\le c_2\norm{k}^2
 \label{eq:supp-symbol-comparison}
\end{equation}
for modes in a fundamental dual cell, with constants independent of $h$.
The first nonzero dual vector has fixed positive length, proving (i).
Counting dual-lattice points in a three-dimensional ball gives (ii).

For (iii), retain only the three basis-root directions
$\pm e_1,\pm e_2,\pm e_3$ from \eqref{eq:a3-roots-main}.  Their Cayley
subgraph is a periodic three-dimensional grid and each retained edge has
conductance $ah/8\ge a_-h/8$.  The discrete edge-isoperimetric inequality on
the $d^3$ periodic grid follows by slicing $A$ along the three coordinate
directions and applying the Loomis--Whitney inequality to the three coordinate
projections; in physical units it reads
$h^2\#\partial A\ge c_0(h^3\#A)^{2/3}$.  Since
$\mu_h(A)=a^3h^3\#A$ and $c_{e,h}=ah/8$, it follows that
$h c_h(\partial A)\ge[c_0/(8a_+)]\mu_h(A)^{2/3}$.
Adding the other $A_3$ edges only increases this boundary capacity, proving
\eqref{eq:supp-periodic-cut}.  Finally,
$-L_h\succeq R(I-P_{h,R})$ gives
\eqref{eq:supp-screen-tail} by the spectral theorem.
\end{proof}

\subsection{Flat renewal vacuum}

\begin{construction}[Flat periodic renewal regulator]
\label{con:supp-flat-regulator}
At mesh $h=d^{-1}$, take the persistent spatial state $Z_n\in V_h$, one point
in every nongeometric fibre, the twelve root branches
\eqref{eq:supp-flat-rates}, vertex mass
\begin{equation}
 m_h=h^3,
 \label{eq:supp-flat-mass}
\end{equation}
the identity root router, the identity slow row, and zero nonvacuum reward.
The event alphabet consists of the root-race outcomes together with the
protected orientation and one-point spectator fibres.  Its conditional law is
normalized directly, so the predictive carrier and scheduler records are
literal parts of the operational regulator.
\end{construction}

\begin{theorem}[Flat ADM and vacuum limit]
\label{thm:supp-flat-vacuum}
The regulator of \Cref{con:supp-flat-regulator} satisfies
\begin{equation}
 \boxed{
 \mathcal B_h=I_3,\qquad
 \varrho_h=1,\qquad
 g_h=I_3,\qquad N_h=1,\qquad\beta_h=0.}
 \label{eq:supp-flat-adm}
\end{equation}
Every undirected root edge has conductance $h/8$.  The spatial generators
converge to $\tfrac12\Delta_{\mathbb T^3_{A_3}}$, and the Lorentzian continuum
is
\begin{equation}
 \boxed{
 (M_0,g_0)=
 \left(\R\times\mathbb T^3_{A_3},
 -\dd t^2+\dd x_1^2+\dd x_2^2+\dd x_3^2\right).}
 \label{eq:supp-flat-lorentzian}
\end{equation}
Its Levi--Civita connection and curvature vanish, so
$G(g_0)=0$.  The finite torsion, curvature, and vacuum Palatini first-variation
residuals vanish identically.
\end{theorem}

\begin{proof}
The root identity and rates give
\[
 \mathcal B_h
 =h^2\sum_{\alpha\in\Phi_{A_3}}
   \frac1{8h^2}\alpha\alpha^{\mathsf T}=I_3.
\]
Reversal symmetry gives $\beta_h=0$.  The mass density per physical volume is
one, so $\varrho_h=1$.  Substitution in the ADM inversion theorem
\Cref{thm:supp-adm-frame} gives \eqref{eq:supp-flat-adm}.  Multiplication of
one directed rate by one vertex mass gives $h/8$.

Generator convergence is \Cref{lem:supp-flat-symbol}.  The coframe
$\vartheta^0=\dd t$, $\vartheta^a=\dd x^a$ is constant, and the exact identity
router gives the zero spin connection.  Hence
$T^A=\dd\vartheta^A+\omega^A{}_B\wedge\vartheta^B=0$ and
$\Omega^A{}_B=\dd\omega^A{}_B+\omega^A{}_C\wedge\omega^C{}_B=0$.
Therefore $\Ric=0$, $R=0$, and $G=0$.  Any face-based finite Cartan and
Palatini writer constructed from these exact constant data is identically zero,
as is its first variation.  Normalization and predictive persistence follow directly from the operational event law and \Cref{thm:supp-persistence-derived}.
\end{proof}

\subsection{Homogeneous de Sitter renewal vacuum}

Let $\Lambda>0$, $H=\sqrt{\Lambda/3}$, and choose the expanding proper-time
stationary branch of \Cref{cor:main-derived-desitter-rates} with
$a_0=1$ and $t_0=0$.  On a compact time slab the resulting renewal rates and
masses are
\begin{equation}
 k_{\alpha,h}(t)=\frac{e^{-2Ht}}{8h^2},
 \qquad
 m_h(t)=e^{3Ht}h^3,
 \label{eq:supp-desitter-rates-mass}
\end{equation}
with physical speed density $\varrho_h(t)=e^{3Ht}$.  These profiles are
stationary outputs of the homogeneous renewal action, not assumptions made in
this appendix.

\begin{theorem}[Reconstructed de Sitter ADM packet]
\label{thm:supp-desitter-adm}
For every $h$ and $t$,
\begin{equation}
 \boxed{
 \mathcal B_h(t)=e^{-2Ht}I_3,
 \quad
 \varrho_h(t)=e^{3Ht},
 \quad
 g_h(t)=e^{2Ht}I_3,
 \quad
 N_h(t)=1,
 \quad
 \beta_h(t)=0.}
 \label{eq:supp-desitter-adm}
\end{equation}
Each undirected edge has conductance
\begin{equation}
 \boxed{c_{e,h}(t)=\frac{e^{Ht}}8h.}
 \label{eq:supp-desitter-conductance}
\end{equation}
On every compact time slab, the cut, Poincar\'e, counting, and compact-screen
constants are uniform in $h$ and $t$.
\end{theorem}

\begin{proof}
By \Cref{lem:supp-a3-second-moment},
\[
 \mathcal B_h(t)
 =h^2\frac{e^{-2Ht}}{8h^2}
  \sum_{\alpha\in\Phi_{A_3}}\alpha\alpha^{\mathsf T}
 =e^{-2Ht}I_3.
\]
The rates are reversal symmetric, so $\beta_h=0$.  Formula
\eqref{eq:adm-inversion-main} gives
\[
 g_h=e^{2Ht}I_3,
 \qquad
 N_h=(e^{3Ht})^{1/3}(e^{-6Ht})^{1/6}=1.
\]
The conductance is the product of vertex mass and directed rate:
$e^{3Ht}h^3e^{-2Ht}/(8h^2)=e^{Ht}h/8$.  On $[-T,T]$, the scalar
$a(t)=e^{Ht}$ is bounded above and below by positive constants, so
\Cref{lem:supp-periodic-screens} applies uniformly.
\end{proof}

Let internal indices $A,B\in\{0,1,2,3\}$ be raised and lowered with
$\eta=\diag(-1,1,1,1)$.  Define
\begin{equation}
 \vartheta^0=\dd t,
 \qquad
 \vartheta^a=e^{Ht}\dd x^a,
 \qquad
 \omega^a{}_0=H\vartheta^a,
 \qquad
 \omega^a{}_b=0,
 \label{eq:supp-desitter-cartan}
\end{equation}
with $\omega_{AB}=-\omega_{BA}$.

\begin{theorem}[Cartan curvature and the vacuum Einstein equation]
\label{thm:supp-desitter-einstein}
The connection \eqref{eq:supp-desitter-cartan} is the torsion-free
Levi--Civita spin connection of
\begin{equation}
 g_H=-\dd t^2+e^{2Ht}\sum_{a=1}^3(\dd x^a)^2.
 \label{eq:supp-desitter-metric}
\end{equation}
Its curvature two-form is
\begin{equation}
 \boxed{\Omega^A{}_B=H^2\vartheta^A\wedge\vartheta_B.}
 \label{eq:supp-desitter-curvature}
\end{equation}
Consequently
\begin{equation}
 \Ric_{\mu\nu}=3H^2g_{\mu\nu},
 \qquad
 R=12H^2,
 \qquad
 \boxed{G_{\mu\nu}+3H^2g_{\mu\nu}=0.}
 \label{eq:supp-desitter-einstein}
\end{equation}
The slices $t=\mathrm{const}$ are compact Cauchy hypersurfaces, so the spatial
quotient is globally hyperbolic.
\end{theorem}

\begin{proof}
Since
$\dd\vartheta^a=H\vartheta^0\wedge\vartheta^a$,
\[
 T^a=\dd\vartheta^a+\omega^a{}_0\wedge\vartheta^0
 =H\vartheta^0\wedge\vartheta^a
  +H\vartheta^a\wedge\vartheta^0=0,
\]
and $T^0=0$.  The coframe is nondegenerate, so this metric-compatible
connection is Levi--Civita.

Cartan's second equation gives
\[
 \Omega^a{}_0=\dd(H\vartheta^a)=H^2\vartheta^0\wedge\vartheta^a,
 \qquad
 \Omega^a{}_b=\omega^a{}_0\wedge\omega^0{}_b
 =H^2\vartheta^a\wedge\vartheta^b,
\]
which is \eqref{eq:supp-desitter-curvature} in invariant notation.  Hence
$R_{ABCD}=H^2(\eta_{AC}\eta_{BD}-\eta_{AD}\eta_{BC})$,
so contraction gives $\Ric_{AB}=3H^2\eta_{AB}$ and $R=12H^2$.  Therefore
$G=\Ric-\tfrac12Rg=-3H^2g$, proving
\eqref{eq:supp-desitter-einstein}.

The function $t$ has timelike gradient and every inextendible causal curve
meets each compact slice $\{t\}\times\mathbb T^3_{A_3}$ exactly once; this is
the standard global-hyperbolicity criterion for a product metric with positive
scale factor.
\end{proof}

\subsection{Exact links and summable finite Cartan--Palatini defects}

Let $K_h$ be the spacetime cell complex obtained by combining the spatial root
mesh with a time mesh of size $h$.  For an oriented edge $e$, let
\begin{equation}
 U_{e,h}=\mathcal P\exp\left(-\int_e\omega\right)
 \label{eq:supp-exact-link}
\end{equation}
be exact parallel transport of the reconstructed continuum connection along
that edge.  Then
$U_{e,h}^{\mathsf T}\eta U_{e,h}=\eta$ exactly.  For an oriented face $p$ with
area $|p|$, define the normalized curvature defect
\begin{equation}
 \mathfrak r_{p,h}
 =|p|^{-1}\log U_{\partial p}+\Omega(x_p),
 \label{eq:supp-curvature-defect}
\end{equation}
where $x_p$ is the face midpoint and the logarithm is taken on the principal
chart for sufficiently small $h$.  Define the normalized torsion defect by the
covariant developed-edge closure
\begin{equation}
 \mathfrak t_{p,h}
 =|p|^{-1}\sum_{e\subset\partial p}
  U_{x_p\leftarrow s(e),h}
  \int_e\vartheta.
 \label{eq:supp-torsion-defect}
\end{equation}
Finally, let $\delta S_h(e_h,U_h)$ be the midpoint quadrature of the
Palatini--cosmological first variation and let
$\delta S(e,\omega)$ be its continuum value.

\begin{theorem}[First-order consistency of the exact-link regulator]
\label{thm:supp-finite-defects}
For every compact time slab $[-T,T]$ there is $C_T<\infty$, independent of
$h$, such that
\begin{equation}
 \sup_p\bigl(\norm{\mathfrak r_{p,h}}+
\norm{\mathfrak t_{p,h}}\bigr)
 +\sup_{\norm{v_h}_{C^1}\le1}
  \abs{\delta S_h[v_h]-\delta S[v_h]}
 \le C_Th.
 \label{eq:supp-finite-defect-bound}
\end{equation}
For the flat regulator the left-hand side is zero.  For the de Sitter
regulator, take
\begin{equation}
 h_n=2^{-n^2},
 \qquad
 T_n=n.
 \label{eq:supp-diagonal-exhaustion}
\end{equation}
Then every normalized defect is summable along the diagonal exhaustion:
\begin{equation}
 \sum_{n=1}^{\infty}C_{T_n}h_n<\infty.
 \label{eq:supp-defect-summability}
\end{equation}
\end{theorem}

\begin{proof}
On a compact slab, the coframe, connection, curvature, and the finitely many
derivatives needed below are bounded.  The nonabelian Stokes expansion for the
exact link gives
\[
 \log U_{\partial p}
 =-|p|\Omega(x_p)+O_T(h^3),
\]
because $|p|=O(h^2)$ and the first variation of the curvature across the face
is $O_T(h)$.  Division by $|p|$ gives the curvature part of
\eqref{eq:supp-finite-defect-bound}.  The covariant sum of developed edge
coframes is the integral of the torsion over the face plus an $O_T(h^3)$
quadrature remainder.  Since the continuum torsion vanishes, division by area
gives the torsion bound.

The Palatini--cosmological first variation is a finite sum of products of
coframe, curvature, and test-field coefficients.  Midpoint quadrature on each
cell has local error $O_T(h^5)$ for a four-dimensional cell of volume
$O(h^4)$; summing $O(h^{-4})$ cells on a fixed slab gives $O_T(h)$.
Equivalent estimates follow directly from Taylor expansion of the finite
wedge and holonomy writers.  In the flat case all coefficients are constant
and the exact links are identities, so the defects vanish.

For de Sitter, derivatives of order bounded independently of $n$ grow at most
exponentially on $[-n,n]$: $C_{T_n}\le C e^{cn}$ for fixed constants $C,c$.
Hence
$C_{T_n}h_n\le C\exp(cn-(\log2)n^2)$, whose series converges.
\end{proof}

\subsection{Proof of Theorem~\ref{mt:vacuum}}

\begin{proof}[Proof of Theorem~\ref{mt:vacuum}]
The predictive carrier at every cutoff is reconstructed by
\Cref{thm:supp-persistence-derived}.  The homogeneous temporal profiles and
their finite stationary origin are
\Cref{thm:main-homogeneous-friedmann,cor:main-derived-desitter-rates,thm:main-finite-homogeneous-stationarity};
the nonempty operational matching family is
\Cref{thm:main-explicit-operational-family}.  The spatial ADM reconstruction is
\Cref{thm:supp-flat-vacuum,thm:supp-desitter-adm}.  The strong-resolvent and
semigroup spatial limits follow from \Cref{lem:supp-flat-symbol}; the scalar
time-dependent factor is uniformly controlled on compact slabs.  Uniform cut,
Poincar\'e, counting, and compact-screen bounds are
\Cref{lem:supp-periodic-screens}.  The flat Einstein equation follows from the
zero connection in \Cref{thm:supp-flat-vacuum}.  The de Sitter Cartan and
Einstein calculation is \Cref{thm:supp-desitter-einstein}.  Exact metric-unitary
links and the $O_T(h)$, diagonally summable consistency estimates are
\Cref{thm:supp-finite-defects}.  These statements establish the six numbered clauses and the operational
realization asserted in the main theorem.
\end{proof}

\begin{remark}[Meaning of constructive vacuum emergence]
The temporal rate profiles used here are not chosen in order to reproduce a
target metric: on the homogeneous branch they solve the lapse-varied renewal
Euler system of Appendix~\ref{supp:homogeneous-dynamics}, and a nonempty finite
accept/reject class reconstructs the same reduced action from its measured
trial probabilities.  The present appendix verifies the downstream spatial and
Cartan geometry of those stationary profiles.  The result still does not claim
that an arbitrary renewal table selects the quadratic operational class, the
rank-one compatibility condition, the value of $\Lambda$, or the expanding
sign, nor that de~Sitter is a universal attractor under coarse graining.
\end{remark}



\end{document}
